\documentclass[11pt]{article}

\usepackage[T1]{fontenc}
\usepackage[utf8]{inputenc}
\usepackage[english]{babel}
\usepackage[margin=1in]{geometry}
\usepackage{lmodern}
\usepackage{microtype}
\usepackage{mathtools}
\usepackage{amsmath,amssymb,amsthm}
\usepackage{array}
\usepackage{booktabs}
\usepackage{tabularx}
\usepackage{longtable}
\usepackage{float}
\usepackage{graphicx}
\usepackage{pdflscape}
\usepackage{caption}
\usepackage{enumitem}
\usepackage{fancyhdr}
\usepackage[numbers]{natbib}
\usepackage{doi}
\usepackage{xurl}
\usepackage{hyperref}
\usepackage{cleveref}
\usepackage{orcidlink}
\usepackage{titling}

\usepackage{ifluatex}
\ifluatex
  \usepackage{fontspec}
  \setmonofont{DejaVu Sans Mono}[Scale=MatchLowercase]
\fi



\providecommand{\BirdInt}{\mathbf{BirdInt}^{\mathrm{aud}}_{\mathrm{fin}}}

\providecommand{\factor}{\mathrel{\!\vdash\!}}

\providecommand{\Pdrive}{\mathrm{P6}_{\mathrm{drive}}}
\providecommand{\PdriveH}[1]{\mathrm{P6}_{\mathrm{drive}}^{#1}}

\providecommand{\tightlist}{}
\providecommand{\pandocbounded}[1]{#1}
\providecommand{\real}[1]{#1}
 
\newtheorem{theorem}{Theorem}[section]
\newtheorem{lemma}[theorem]{Lemma}
\newtheorem{proposition}[theorem]{Proposition}
\newtheorem{corollary}[theorem]{Corollary}
\theoremstyle{definition}
\newtheorem{definition}[theorem]{Definition}
\newtheorem{remark}[theorem]{Remark}
\floatstyle{ruled}
\newfloat{algorithm}{tbp}{loa}
\floatname{algorithm}{Algorithm}

\urlstyle{same}
\captionsetup{
  font=small,
  labelfont=bf,
  labelsep=period
}
\captionsetup[table]{skip=6pt,position=top}
\captionsetup[figure]{skip=6pt}
\captionsetup[algorithm]{skip=6pt,position=top}
\crefname{algorithm}{algorithm}{algorithms}
\Crefname{algorithm}{Algorithm}{Algorithms}
\crefname{theorem}{theorem}{theorems}
\Crefname{theorem}{Theorem}{Theorems}
\crefname{lemma}{lemma}{lemmas}
\Crefname{lemma}{Lemma}{Lemmas}
\crefname{proposition}{proposition}{propositions}
\Crefname{proposition}{Proposition}{Propositions}
\crefname{corollary}{corollary}{corollaries}
\Crefname{corollary}{Corollary}{Corollaries}
\crefname{definition}{definition}{definitions}
\Crefname{definition}{Definition}{Definitions}
\crefname{remark}{remark}{remarks}
\Crefname{remark}{Remark}{Remarks}
\setlist[enumerate]{leftmargin=*, itemsep=2pt, topsep=4pt}
\setlength{\droptitle}{-3.2em}
\pretitle{\begin{center}\LARGE\bfseries}
\posttitle{\par\end{center}\vspace{0.5em}}
\preauthor{\begin{center}\normalsize}
\postauthor{\par\end{center}\vspace{-0.4em}}
\predate{\begin{center}\small}
\postdate{\par\end{center}\vspace{0.6em}}
\setlength{\emergencystretch}{2em}
\clubpenalty=10000
\widowpenalty=10000
\displaywidowpenalty=10000
\interfootnotelinepenalty=10000
\allowdisplaybreaks[2]
\fancypagestyle{firstpage}{\fancyhf{}
  \fancyfoot[C]{\scriptsize Automorph Inc., Wilmington, DE, USA \quad \texttt{ioannis@automorph.io} \quad \textcopyright\ 2026 \quad CC-BY 4.0}
  \renewcommand{\headrulewidth}{0pt}
  \renewcommand{\footrulewidth}{0pt}
}

\title{Six Birds Foundations III:\\
A Finite Audited Interaction Calculus for SBT}
\author{Ioannis Tsiokos\,\orcidlink{0009-0009-7659-5964}}
\date{\today}
\hypersetup{
  pdftitle={Six Birds Foundations III: A Finite Audited Interaction Calculus for SBT},
  pdfauthor={Ioannis Tsiokos},
  hidelinks
}

\begin{document}

\maketitle
\thispagestyle{firstpage}

\begin{abstract}
Earlier work in the Six-Bird program supplies six primitive role families ($P_1,\ldots,P_6$ --- descent, representability, route mismatch, refinement, packaging, audit), an admissibility and exact-six scope discipline over $\mathsf{FATCD}$, inherited status, threshold, visibility, and audit machinery, and concrete model fragments such as PICA and Cantor. It does not, however, supply a paper-level finite interaction calculus that records, in one finite typed framework, how the six primitive roles interact, how typed directed cells are judged, how defects induce statuses, how promotion is audited, and how model realizations are scoped. This paper constructs $\BirdInt$, a finite audited typed calculus for six-bird interaction. The calculus carries finite typed schemas of theory packages, instruments, visibility and suppression records, witness and update families, role-channel, directed-cell, pair-observable, promotion, dependency, and no-go judgments, defect records, statuses, gates, audit records, and model-realization maps. Every component is finite, admissibility is checked on finite records, and every claim is recorded under an explicit host or instrument together with a nonclaim register. The main theorem groups are: no total six-symbol algebra; typed non-collapse and status-family separation; finite well-formedness and covered-cell status uniqueness; descent, closure-deficit, completion, saturation, strict-extension, and definability results; graph and cohomology no-drive results; promotion-gate soundness; and no-overreading and finite rotating-audit theorems. The numbered stack comprises theorem-grade Theorems 1--28 and 30--32, model-grade Model Theorems 29, 33, and 34, and audit-grade Proposition 35. The repository includes a Lean 4 validation track, pinned to \texttt{leanprover/lean4:v4.28.0}, whose manifest-driven strict validator audits the paper-to-Lean declaration coverage. The mathematical content of these theorems is standard --- descent through a finite quotient, idempotent saturation, the Kullback--Leibler decomposition, finite cycle/cochain duality, the universal property of the image factorization, and so on --- and is not claimed as novel; what the paper contributes is the calculus, its scope discipline, the countermodel atlas, and the explicit nonclaim register within which these results are restated and used. Twelve countermodels $\mathsf{CM}_1$--$\mathsf{CM}_{12}$ each block a specific overclaim, and four scoped model-realization statements cover finite stochastic PICA, audited Cantor shells, finite abstract-interpretation hosts, and finite graph and cohomology supports. Each model realization is host-bound, fragment-scoped, and certified at claim-strength \texttt{model\_realization} rather than theorem-grade \texttt{accepted}. The paper does not prove a universal exact-six theorem, a universal six-symbol algebra, a unique decomposition theorem, an empirical realization claim, or complete same-level self-certification.
\end{abstract}

\medskip
\noindent\textbf{Keywords.} finite typed calculus; audited interaction; emergence calculus; status semantics; defect calculus; theory packages; promotion bridge; model realization; rotating audit; nonfactorization.

\section{Introduction}

This paper constructs \(\mathbf{BirdInt}^{\mathrm{aud}}_{\mathrm{fin}}\), a finite audited typed interaction calculus for six-bird theory, and records the theorem-grade, model-grade, and audit-grade results that it admits. The introduction states the gap that the calculus closes, lists what the paper proves and what it does not prove, and gives a reader's map of the remaining sections. Subsection 1.1 fixes the missing mathematical object that earlier work in the program leaves open. Subsection 1.2 lists the paper's positive results across five layers --- the core calculus, descent/packaging/completion/strictness, no-go and countermodel results, promotion and instrument results, and model-realization results. Subsection 1.3 records the negative scope discipline. Subsection 1.4 gives a section-by-section reading guide.

\subsection{The Missing Mathematical Object}

This subsection records the prior work in the six-bird program, the gap that the present paper closes, and the constructed object that closes it.

\subsubsection{Starting Point}

Earlier work in the program supplies the following:

\begin{itemize}
\tightlist
\item
  \emph{Six primitive role families.} The primitives \(P_{1},\ldots,P_{6}\) --- descent, representability, route mismatch, refinement, packaging, audit --- are recorded in Foundations I and Foundations II as role labels indexed by stable concepts (closure, gating, holonomy, refinement, packaging, audit) rather than as elements of a total algebra.
\item
  \emph{Admissibility and exact-six scope discipline.} The scoped exact-six discipline of Foundations II distinguishes admissible finite audited typed closure descriptions in \(\mathsf{FATCD}\) from naive total decompositions, and records role channels with status inside that inherited domain. The directed-cell and profile-relative form is the refinement introduced by the present paper.
\item
  \emph{Status, threshold, visibility, and audit machinery.} Earlier work supplies inherited channel statuses, threshold records, audit records, and instrument-relative visibility discipline; the present paper sharpens these ingredients into the multiple status families and classifiers of Sections 4 and 5.
\item
  \emph{Concrete model examples.} The PICA and Cantor model families, together with abstract interpretation and graph/cohomology fragments, record host-specific fragments of the six-bird structure on finite stochastic, audited shell, lattice-based, and graph-based hosts.
\end{itemize}

These four kinds of input record the program's content at the level of primitives, admissibility predicates, status and audit machinery, and host-specific examples, but they do not, by themselves, supply a paper-level finite interaction calculus over the six primitives.

\subsubsection{Gap}

The gap addressed by the present paper is the absence of a paper-level finite interaction calculus that records, in one finite typed framework:

\begin{itemize}
\tightlist
\item
  how the six primitive roles interact under directed-cell judgments \(P_{i}\leftarrow P_{j}\);
\item
  how a finite typed cell record is judged for admissibility under a host instrument;
\item
  how defects induce statuses across the role, cell, pair, promotion, claim, gate, and square families;
\item
  how promotion bridges between theory packages are audited and classified;
\item
  how model realizations are scoped, host-bound, and certified at \(\texttt{model\_realization}\) rather than theorem-grade verdicts;
\item
  how claims under one instrument are constrained by visibility and audit, and how self-reference is handled by finite rotating audits without claiming complete self-certification.
\end{itemize}

Prior work supplies fragments of each of these pieces; the present paper supplies the finite calculus that records them in a single typed framework with explicit admissibility, audit, status, and nonclaim discipline.

\subsubsection{This Paper's Answer}

The constructed object is

\[
\mathbf{BirdInt}^{\mathrm{aud}}_{\mathrm{fin}},
\]

the finite audited interaction calculus for typed six-bird interaction. It is a finite typed object --- defined formally in Sections 3 and 4, with status and defect semantics in Section 5 --- whose components include:

\begin{enumerate}
\def\labelenumi{\arabic{enumi}.}
\tightlist
\item
  a finite primitive set \(\mathbb P=\{P_{1},\ldots,P_{6}\}\) with stable role glosses;
\item
  a finite typed cell record \(\mathrm{Cell}\) over actor, informant, witness, update, threshold, audit, visibility, and defect components, indexed by a per-cell profile \(\lambda\);
\item
  seven finite status families \(\mathrm{RoleStatus}\), \(\mathrm{CellStatus}\), \(\mathrm{PairStatus}\), \(\mathrm{PromotionStatus}\), \(\mathrm{ClaimStatus}\), \(\mathrm{GateStatus}\), \(\mathrm{SqStatus}\), with frozen priority orders within each family;
\item
  a finite typed schema of theory packages, instruments, visibility/suppression records, source-of-truth records, audit records, promotion bridges, and visibility bridges;
\item
  a finite admissibility predicate on each component;
\item
  a finite typed family of defect records that route claims to non-\(\texttt{accepted}\) statuses through the classifier of Section 5;
\item
  a finite nonclaim register \(\mathcal N\) that records the negative boundary of every admissible claim.
\end{enumerate}

The calculus admits numbered theorem-grade results, scoped model-grade realization results, twelve countermodels, and one Lean-validation audit report, recorded in Sections 6--15 and in the appendices. The introduction's role is to record what the calculus is, what it admits, and what it does not admit, before the formal development begins in Section 2.

\subsection{What the Paper Proves}

This subsection records the paper's positive results across five layers: the core finite-calculus theorems, the descent/packaging/completion/strictness theorems, the no-go and countermodel results, the promotion and instrument results, and the model-realization results. The lists below are summary pointers; full statements and proofs are in the body of the paper.

\subsubsection{Core Calculus Results}

The paper records five core finite-calculus theorems in Section 6:

\begin{itemize}
\tightlist
\item
  \emph{No total six-symbol algebra} (Theorem 1, Section 6): cell status does not factor through the primitive pair \((P_{i},P_{j})\) alone, blocking any total binary algebra on \(\mathbb P\) as a status-decoding map.
\item
  \emph{Typed non-collapse} (Theorem 2, Section 6): an admissible directed-cell judgment \(P_i\leftarrow P_j\), even with status \(\texttt{action}\), does not identify the actor and informant primitives.
\item
  \emph{Cell well-formedness decidability} (Theorem 3, Section 6): the well-formedness predicate on a finite typed cell record is decidable on finite records.
\item
  \emph{Covered-cell status uniqueness} (Theorem 4, Section 6): on a covered finite cell, the classifier of Subsection 5.6 returns a unique cell status.
\item
  \emph{Status-family separation} (Theorem 5, Section 6): the five judgment-attached status families remain typed-separately judgment-indexed, with gate and square status families likewise kept separate by their own inputs.
\end{itemize}

\subsubsection{Descent, Packaging, Completion, and Strictness Results}

The paper records nine theorems on descent, packaging, completion, and strictness in Section 7:

\begin{itemize}
\tightlist
\item
  \emph{Descent through quotient} (Theorem 6) and \emph{descent defect equivalence} (Theorem 7): the well-definedness equivalence \(q(x)=q(x')\Rightarrow qF(x)=qF(x')\) characterizes when \(F\) descends along \(q\), with a defect record characterizing failure.
\item
  \emph{Markov closure-deficit theorem} (Theorem 8): a stochastic closure deficit obstructs macro-dynamics closure on a finite Markov host, in the sense fixed in Section 7.
\item
  \emph{Noncommuting completions} (Theorem 9) and \emph{completion-pasting defect} (Theorem 10): two finite idempotent completions \(E_{1},E_{2}:C\to C\) may fail to commute, and the failure is a finite typed pasting defect.
\item
  \emph{Idempotent saturation} (Theorem 11): a finite extensive monotone idempotent operator is a closure operator with finite saturation.
\item
  \emph{Strict-extension nonfactorization} (Theorem 12) and \emph{factorization defect equivalence} (Theorem 13): \(\pi_{1}\not\factor\pi_{0}\) is equivalent to a nonempty fiber-mismatch witness on the carrier.
\item
  \emph{Fixed-interface definability bound} (Theorem 14): for a fixed finite interface \(f\), the full finite-lens definability family has cardinality \(|\mathrm{Def}(f)|=2^{|\mathrm{im}(f)|}\).
\end{itemize}

\subsubsection{No-Go and Countermodel Results}

The paper records the cohomological no-drive theorems of Section 8 and the countermodel atlas of Section 9:

\begin{itemize}
\tightlist
\item
  \emph{Forest no-drive} (Theorem 15): \(\beta_{1}(G)=0\Rightarrow H^{1}(G;k)=0\).
\item
  \emph{Exact-form null-drive} (Theorem 16): \(a=d\phi\Rightarrow\oint_{\gamma}a=0\) for every cycle \(\gamma\).
\item
  \emph{Nonzero-affinity equivalence} (Theorem 17): \([a]\neq 0\iff\exists\gamma:\oint_{\gamma}a\neq 0\).
\item
  \emph{Gating-affinity suppression} (Theorem 18): gating to a forest kills cycle-supported drive.
\item
  \emph{Countermodels CM\_1--CM\_12} (Section 9 and Appendix C): finite countermodels blocking package-without-closure, holonomy-without-drive, local-without-global, cell-without-pair-realness, refinement-worsens-closure, same-level self-audit circularity, same-pair-different-status, idempotent-noncommuting completions, strictness-without-closure, strictness-without-drive, fallback-as-real-source, and gating-destroys-cycle-supported-drive overclaims.
\end{itemize}

\subsubsection{Promotion and Instrument Results}

The paper records the promotion-and-stacking theorems of Section 10 and the self-reference theorems of Section 11:

\begin{itemize}
\tightlist
\item
  \emph{Promotion gate soundness} (Theorem 19) and \emph{strict promotion implies nonfactorization} (Theorem 20): any accepted-family promotion verdict requires every required core gate to exist and pass, and a strict promotion verdict implies a nonfactorization witness on the carrier.
\item
  \emph{Non-strict promotion implies factorization} (Theorem 21): a non-strict promotion verdict on an admissible bridge implies a factorization filler.
\item
  \emph{Toy strict bridge} (Example Theorem 22): a finite two-layer bridge witnesses the strict-promotion pattern without proving any general stacking-composition theorem.
\item
  \emph{No-overreading suppression} (Theorem 23, Section 11): no \(I\)-accepted claim depends on unbridged suppressed content.
\item
  \emph{Same-level self-audit failure} (Theorem 24, Section 11): an instrument \(I_{0}\) does not certify its own complete soundness; the verdict is \(\texttt{outside\_scope}\) or \(\texttt{undefined\_circular}\).
\item
  \emph{Finite rotating-audit theorem} (Theorem 25, Section 11): a finite stack \(\mathbb I_{\le N}\) admits a level-\(N+1\) extension instrument that audits the lower stack and accepts compliance iff the stack-audit defect is empty, while remaining non-final on its own soundness.
\end{itemize}

\subsubsection{Model-Realization Results}

The paper records four scoped model-realization theorems plus secondary high-structure recovery results:

\begin{itemize}
\tightlist
\item
  \emph{PICA model-realization} (Model Theorem 33, Subsection 13.2): \(\mathsf{PICAReal}_{\mathrm{fin}}\models\mathbf{BirdInt}^{\mathrm{aud}}_{\mathrm{fin\text{-}stoch}}{}^{\,\mathrm{cell}/\mathrm{pair}/\mathrm{prov}}:\texttt{model\_realization}\).
\item
  \emph{Cantor strict-bridge realization} (Model Theorem 34, Subsection 13.3): \(\mathsf{CantorShell}_{\mathrm{aud}}\models\mathbf{StrictBridge}_{0\to 1}^{\mathrm{audited\ shell}}:\texttt{model\_realization}\).
\item
  \emph{Abstract-interpretation model-family theorem} (Model Theorem 29, Section 12): \(\mathsf{AIFam}_{\mathrm{fin}}\models\mathbf{BirdInt}^{\mathrm{aud}}_{\mathrm{fin}}{}^{\,P_{1}/P_{2}/P_{4}/P_{5}/P_{6}}:\texttt{model\_realization}\), supported by Theorems 26--28 (closure as packaging, descent as sound transformer, representability as fixedness).
\item
  \emph{Graph/cohomology realization} (Subsection 13.4): \(\mathsf{CohGraph}_{\mathrm{fin}}\models\mathbf{BirdInt}^{\mathrm{aud}}_{\mathrm{fin}}{}^{\,P_{2}/\mathrm{P6}_{\mathrm{drive}}^{H^{1}}/P_{3}\text{-separation}}:\texttt{model\_realization}\).
\item
  \emph{High-structure recovery} (Theorems 30--32, Section 14): descent, completion-pasting, and strict-extension data are recovered as decorated squares in a finite double-category fragment, as secondary organization rather than a replacement for the core calculus.
\item
  \emph{Lean validation audit} (Proposition 35, Section 15): the repository contains a Lean 4 track whose manifest-driven strict validator confirms that the required finite definitions and numbered claims are represented by Lean declarations, with the audit-grade verdict \(\texttt{audit\_report}\).
\end{itemize}

\subsection{What the Paper Does Not Prove}

This subsection records the negative scope discipline of the paper, placed early in the introduction so that the reader has the boundary in hand before encountering the formal development. Each item below is a scoped negative statement; the full nonclaim register is recorded in Section 16.

\subsubsection{No Universal Algebra Claim}

The paper does not construct a total binary algebra on the six primitive symbols. There is no faithful total operation \(\ast:\mathbb P^{2}\to\mathbb P\) whose composition with a status-decoding map recovers typed, statused, profile-relative directed-cell status across all profiles. Theorem 1 of Section 6 records the formal counterexample: cell status does not factor through the primitive pair, so no such algebra is admissible as a status decoder.

\subsubsection{No Universal Exact-Six Claim}

The paper does not assert that every phenomenon, every theory, or all of mathematics reduces to six primitives. Foundations II's scoped exact-six theorem is a statement over \(\mathsf{FATCD}\), not an unrestricted decomposition theorem, and the present paper's directed-cell schema is likewise admissibility-bound. The calculus accepts a directed-cell judgment \(P_{i}\leftarrow P_{j}\) only under the per-cell admissibility predicate and per-cell profile of Subsection 4.3, and does not assert that every mathematical structure is decomposable into the six primitives.

\subsubsection{No Activation Equivalence}

The paper distinguishes between:

\begin{itemize}
\tightlist
\item
  the six role channels (or status slots) in an admissible description, which are typed-separately judgment-indexed and remain in scope under the host instrument; and
\item
  the active projections or active cells in a particular record, which require the relevant witness, threshold, audit, and classifier checks in their own judgment family.
\end{itemize}

For inherited role-channel records, statuses such as \(\texttt{below\_threshold}\), \(\texttt{collapsed\_boundary\_case}\), \(\texttt{absent\_with\_record}\), or \(\texttt{inapplicable\_outside\_scope}\) do not count as active projections. For Foundations III directed cells, statuses such as \(\texttt{blocked}\), \(\texttt{absent}\), \(\texttt{outside\_scope}\), \(\texttt{undefined\_circular}\), or \(\texttt{below\_threshold}\) do not count as \(\texttt{action}\) cells. The paper records this distinction throughout, in keeping with the status-discipline of Section 5.

\subsubsection{No Model Universality}

The paper does not assert that the four model families --- PICA, Cantor, abstract interpretation, and graph/cohomology --- realize the entire calculus or supply universal theorem sources. Each realization is host-bound, fragment-scoped, and certified at \(\texttt{model\_realization}\) on a declared substructure of \(\mathbf{BirdInt}^{\mathrm{aud}}_{\mathrm{fin}}\), in keeping with the model-realization convention of Subsection 13.1.2. A \(\texttt{model\_realization}\) verdict is not a theorem-grade \(\texttt{accepted}\) verdict, and the paper does not promote it to one without an admissible bridge.

\subsubsection{No Complete Self-Certification}

The paper does not assert that any instrument certifies its own complete soundness at the same level, nor that the calculus admits a final self-certifying instrument on a finite tower. Theorem 24 of Section 11 records the same-level self-audit failure (an instrument \(I_{0}\) does not certify its own complete soundness under \(I_0\); the verdict is \(\texttt{outside\_scope}\) or \(\texttt{undefined\_circular}\)), and Theorem 25 records the finite rotating-audit theorem together with the non-finality clause: a finite rotating-audit chain audits the lower stack but does not certify its top.

The paper records what an admissible audit chain delivers and what it does not. It does not assert complete instrument-family completeness, an infinite reflection tower, or a final instrument that audits itself.

\subsection{Reader's Map}

The paper is organized into sixteen sections plus a conclusion (Section 17) and eight appendices (A--H). The reader's map is:

\begin{itemize}
\tightlist
\item
  \emph{Section 2} records preliminaries and scope: finite host mathematics, primitive labels and levels, claim strengths, status discipline, and the global nonclaim conventions.
\item
  \emph{Sections 3 and 4} define the finite audited calculus: the finite typed domain, the central object \(\mathbf{BirdInt}^{\mathrm{aud}}_{\mathrm{fin}}\) as a finite typed tuple, theory packages, instruments, source-of-truth records, audit records, visibility bridges, directed-cell records, pair-observable records, promotion bridges, and the per-cell profile schema.
\item
  \emph{Section 5} records the status and defect semantics: the seven status families, the per-component defect schemas, the visibility-defect family, and the classifier priority orders.
\item
  \emph{Sections 6--8} prove the finite theorem spine: the no-total-algebra theorem and the four core finite-interaction theorems (Section 6); the nine descent/packaging/completion/strictness theorems (Section 7); and the four cohomological no-drive theorems on graph-supported drive (Section 8).
\item
  \emph{Section 9} records the countermodel atlas in main-text form, with full sheets for CM\_1--CM\_6 and named pointers to CM\_7--CM\_12. The full twelve-countermodel atlas is recorded in Appendix C.
\item
  \emph{Section 10} records promotion and stacking: the promotion-bridge schema, the eight-gate machinery, the promotion-status consequences of the classifier, Theorems 19--21, the Toy22 strict-bridge example, and the stacking claims not proved.
\item
  \emph{Section 11} records instruments, visibility, and rotating audit: the instrument-indexed claim semantics, the classifier priority, the visibility/suppression discipline, the no-overreading theorem (Theorem 23), the same-level self-audit failure theorem (Theorem 24), the \(P_{6}\leftarrow P_{6}\) status separation, the rotating-audit chain and bridge, the finite rotating-audit theorem (Theorem 25), and the instrument-transfer rule.
\item
  \emph{Section 12} records the abstract-interpretation model family: the host data, the realized fragment \(P_{1}/P_{2}/P_{4}/P_{5}/P_{6}\), Theorems 26--28 (closure as packaging, descent as sound transformer, representability as fixedness), and Model Theorem 29.
\item
  \emph{Section 13} records the four scoped model realizations (PICA, Cantor strict-bridge, abstract interpretation, graph/cohomology), Model Theorems 33 and 34, the combined model-realization matrix, and the cross-model role-coverage observation.
\item
  \emph{Section 14} records the secondary high-structure semantics: the finite decorated double-category fragment and the recovery theorems 30--32 (descent-square recovery, completion-pasting recovery, strict extension as no-filler).
\item
  \emph{Section 15} records the Lean 4 validation artifact, the manifest-driven validator, and Proposition 35.
\item
  \emph{Section 16} records limitations, nonclaims, and future work, with the final paper-inclusion nonclaim block of thirty-six entries.
\item
  \emph{Section 17} records the conclusion.
\item
  \emph{Appendices A--H} record the long-form supplementary material: full definitions and tables (A), long proofs (B), the full countermodel atlas (C), the four model-realization sheets (D), the high-structure semantics extension (E), the Lean mechanization appendix (F), the dependency graph (G), and the deferred and dropped claims register (H).
\end{itemize}

Sections 2--5 are the formal-development prerequisites for every subsequent section. Sections 6--11 are the theorem spine: they build on Sections 2--5 and, where stated, on earlier theorem sections. Sections 12 and 13 depend on Sections 2--5 and Section 11 (model-realization claims are recorded under instruments). Section 14 depends on Sections 2--5 and Section 7. Section 15 depends on the Lean manifest, the full \texttt{SixBirdsIII} import graph, and the validation track described in Appendix F. Section 16 collects scope discipline from every preceding section. The reader interested only in the theorem spine can read Sections 2--8, 10, and 11 sequentially; the reader interested in model realizations can additionally read Sections 12, 13, and Appendix D; the reader interested in mechanization can read Section 15 and Appendix F.
 \section{Preliminaries and Scope}

This section fixes the host mathematics, the primitive labels, the level set, the claim-strength discipline, and the scope statement that the rest of the paper will assume. Every later definition, theorem, and model realization is interpreted inside the domain declared here, or inside an explicitly named finite extension host. No object, judgment, or proof in this paper depends on hosts that fall outside the inventory below.

\subsection{Finite Host Mathematics}

\subsubsection{Base Host}

The base host of the calculus is

\[
\mathsf H_{\mathrm{fin}} \;=\; \text{finite sets, finite maps, finite relations, finite records, finite diagrams, finite graphs, finite audit records}.
\]

All carriers, lenses, packages, instruments, profiles, witnesses, updates, defects, gates, audits, statuses, and judgment instances introduced in Sections 3 and 4 are objects of \(\mathsf H_{\mathrm{fin}}\) unless an extension host is declared. In particular, every set used in a theorem statement is finite, every map is a finite function on finite sets, every diagram has finitely many vertices and edges, and every audit record is a finite tuple of finite fields.

The base host carries the following structures, each used by a specific judgment family later in the paper:

\begin{longtable}[]{@{}ll@{}}
\toprule
\begin{minipage}[b]{0.47\columnwidth}\raggedright
Structure\strut
\end{minipage} & \begin{minipage}[b]{0.47\columnwidth}\raggedright
Use in the calculus\strut
\end{minipage}\tabularnewline
\midrule
\endhead
\begin{minipage}[t]{0.47\columnwidth}\raggedright
finite sets\strut
\end{minipage} & \begin{minipage}[t]{0.47\columnwidth}\raggedright
carriers \(Z\), primitive labels \(\mathbb P\), witness/update domains, status sets\strut
\end{minipage}\tabularnewline
\begin{minipage}[t]{0.47\columnwidth}\raggedright
finite maps\strut
\end{minipage} & \begin{minipage}[t]{0.47\columnwidth}\raggedright
lenses \(f:Z\to X\), quotients \(q\), completion endomaps \(E\), object maps \(\pi_j\), update maps \(U_i\)\strut
\end{minipage}\tabularnewline
\begin{minipage}[t]{0.47\columnwidth}\raggedright
finite relations\strut
\end{minipage} & \begin{minipage}[t]{0.47\columnwidth}\raggedright
equivalence on a carrier, dependency between primitives, gluing compatibility on a cover\strut
\end{minipage}\tabularnewline
\begin{minipage}[t]{0.47\columnwidth}\raggedright
finite records\strut
\end{minipage} & \begin{minipage}[t]{0.47\columnwidth}\raggedright
theory packages \(\mathcal T\), instruments \(I\), directed cells, promotion bridges, claim records, defect records, gate records\strut
\end{minipage}\tabularnewline
\begin{minipage}[t]{0.47\columnwidth}\raggedright
finite diagrams\strut
\end{minipage} & \begin{minipage}[t]{0.47\columnwidth}\raggedright
finite decorated squares, descent squares, refinement and gluing diagrams\strut
\end{minipage}\tabularnewline
\begin{minipage}[t]{0.47\columnwidth}\raggedright
finite graphs\strut
\end{minipage} & \begin{minipage}[t]{0.47\columnwidth}\raggedright
support graphs, finite cycle/cochain complexes for the drive module\strut
\end{minipage}\tabularnewline
\begin{minipage}[t]{0.47\columnwidth}\raggedright
finite audit records\strut
\end{minipage} & \begin{minipage}[t]{0.47\columnwidth}\raggedright
provenance, replay traces, source identifiers, threshold checks, nonclaims\strut
\end{minipage}\tabularnewline
\bottomrule
\end{longtable}

\subsubsection{Permitted Extension Hosts}

The following finite or scoped extension hosts are permitted in the paper. Each is named, finite under its own rules, and used only for the specific theorem groups indicated. No theorem in the paper uses an extension host implicitly; every appearance is gated by an explicit declaration on the surrounding judgment.

\[
\mathsf{Host}\;=\;\bigl\{\mathsf H_{\mathrm{fin}},\,\mathsf H_{\mathrm{prob}},\,\mathsf H_{\mathrm{graph}},\,\mathsf H_{\mathrm{AI}},\,\mathsf H_{\mathrm{PICA}},\,\mathsf H_{\mathrm{CantorShell}},\,\mathsf H_{\mathrm{DC}},\,\mathsf H_{\mathrm{instr}}\bigr\}.
\]

\begin{longtable}[]{@{}lll@{}}
\toprule
\begin{minipage}[b]{0.30\columnwidth}\raggedright
Host\strut
\end{minipage} & \begin{minipage}[b]{0.30\columnwidth}\raggedright
Content\strut
\end{minipage} & \begin{minipage}[b]{0.30\columnwidth}\raggedright
Scope\strut
\end{minipage}\tabularnewline
\midrule
\endhead
\begin{minipage}[t]{0.30\columnwidth}\raggedright
\(\mathsf H_{\mathrm{fin}}\)\strut
\end{minipage} & \begin{minipage}[t]{0.30\columnwidth}\raggedright
finite sets, maps, relations, records, diagrams, graphs, audit records\strut
\end{minipage} & \begin{minipage}[t]{0.30\columnwidth}\raggedright
base host for typed-calculus theorems\strut
\end{minipage}\tabularnewline
\begin{minipage}[t]{0.30\columnwidth}\raggedright
\(\mathsf H_{\mathrm{prob}}\)\strut
\end{minipage} & \begin{minipage}[t]{0.30\columnwidth}\raggedright
finite stochastic kernels on finite state spaces, finite information quantities\strut
\end{minipage} & \begin{minipage}[t]{0.30\columnwidth}\raggedright
Markov closure-deficit theorem and PICA realization, under explicit kernel/partition assumptions\strut
\end{minipage}\tabularnewline
\begin{minipage}[t]{0.30\columnwidth}\raggedright
\(\mathsf H_{\mathrm{graph}}\)\strut
\end{minipage} & \begin{minipage}[t]{0.30\columnwidth}\raggedright
finite graphs, finite cycle bases, finite cochain/coboundary complexes\strut
\end{minipage} & \begin{minipage}[t]{0.30\columnwidth}\raggedright
graph/cohomology drive module of Section 8\strut
\end{minipage}\tabularnewline
\begin{minipage}[t]{0.30\columnwidth}\raggedright
\(\mathsf H_{\mathrm{AI}}\)\strut
\end{minipage} & \begin{minipage}[t]{0.30\columnwidth}\raggedright
finite concrete domains, finite abstract domains, finite abstraction/concretization pairs, finite sound transformers\strut
\end{minipage} & \begin{minipage}[t]{0.30\columnwidth}\raggedright
abstract-interpretation model family of Section 12\strut
\end{minipage}\tabularnewline
\begin{minipage}[t]{0.30\columnwidth}\raggedright
\(\mathsf H_{\mathrm{PICA}}\)\strut
\end{minipage} & \begin{minipage}[t]{0.30\columnwidth}\raggedright
finite PICA realization records (rows, columns, statuses, source records, audit restrictions)\strut
\end{minipage} & \begin{minipage}[t]{0.30\columnwidth}\raggedright
PICA model-realization theorem of Section 13\strut
\end{minipage}\tabularnewline
\begin{minipage}[t]{0.30\columnwidth}\raggedright
\(\mathsf H_{\mathrm{CantorShell}}\)\strut
\end{minipage} & \begin{minipage}[t]{0.30\columnwidth}\raggedright
audited Cantor shell records (base/target packages, interfaces, audit)\strut
\end{minipage} & \begin{minipage}[t]{0.30\columnwidth}\raggedright
Cantor strict-bridge realization of Section 13\strut
\end{minipage}\tabularnewline
\begin{minipage}[t]{0.30\columnwidth}\raggedright
\(\mathsf H_{\mathrm{DC}}\)\strut
\end{minipage} & \begin{minipage}[t]{0.30\columnwidth}\raggedright
finite decorated double-category fragment (objects, horizontal and vertical arrows, decorated squares, finite strict pasting)\strut
\end{minipage} & \begin{minipage}[t]{0.30\columnwidth}\raggedright
secondary high-structure semantics of Section 14\strut
\end{minipage}\tabularnewline
\begin{minipage}[t]{0.30\columnwidth}\raggedright
\(\mathsf H_{\mathrm{instr}}\)\strut
\end{minipage} & \begin{minipage}[t]{0.30\columnwidth}\raggedright
instrument records, visibility maps, rotating-audit chains\strut
\end{minipage} & \begin{minipage}[t]{0.30\columnwidth}\raggedright
instrument and self-reference theorems of Section 11\strut
\end{minipage}\tabularnewline
\bottomrule
\end{longtable}

A host appearing in this list is finite under its own definition: \(\mathsf H_{\mathrm{prob}}\) uses finite state spaces and finite kernels, \(\mathsf H_{\mathrm{graph}}\) uses finite vertex and edge sets, \(\mathsf H_{\mathrm{AI}}\) uses finite abstract domains and finite sound transformers, and so on. An extension host never licenses an infinitary construction.

\subsubsection{Excluded Hosts}

The following hosts are excluded. A claim that depends on any of them, however informally stated, falls outside the formal scope of this paper.

\begin{itemize}
\tightlist
\item
  \textbf{Unrestricted infinite reflection towers.} No theorem in the paper requires an infinite tower of theories, languages, or instruments closed under reflection. Every audit chain used in Section 11 is finite, and the rotating-audit theorem is stated with explicit finite chain length \(N\).
\item
  \textbf{Unrestricted semantic universes.} No theorem assumes a model class closed under arbitrary semantic operations. The model families of Section 13 are scoped finite hosts and are accepted only as model realizations of named fragments of the calculus.
\item
  \textbf{Empirical systems without bridge records.} No theorem appeals to a physical, biological, or computational system unless its content has been packaged into a finite record with explicit bridge data, source records, audit, and nonclaims. An empirical-bridge assertion in this paper is admissible only as a finite record on a named host, never as a bare claim about an external system.
\item
  \textbf{Untyped total six-symbol products.} No definition or theorem assumes a total binary operation \(*:\mathbb P^2\to\mathbb P\) on the primitive labels, and no later object is interpreted as a primitive product. The later no-total-algebra theorem is stated against the typed witness, update, profile, and status requirements of the calculus, so subsequent uses of pair notation \((P_i,P_j)\) or directed-cell notation \(P_i\leftarrow P_j\) are typed bridge judgments rather than algebraic products.
\end{itemize}

The host inventory above, together with these exclusions, fixes what counts as an admissible mathematical environment for the paper. Section 3 builds the calculus over this inventory; Section 4 defines its central judgment families; later sections invoke individual extension hosts only where named.

\subsection{Primitive Labels and Levels}

\subsubsection{Primitive Set}

The primitive set of the calculus is the fixed finite set

\[
\mathbb P \;=\; \{P_1,P_2,P_3,P_4,P_5,P_6\}.
\]

The labels \(P_1,\ldots,P_6\) are role names. They are not elements of an interaction algebra: there is no total binary operation on \(\mathbb P\) in this paper, and the later no-total-algebra theorem makes that exclusion formal. Pair notation \((P_i,P_j)\) names a primitive pair only; it does not determine status, effect, claim strength, visibility, or audit. Directed-cell notation \(P_i\leftarrow P_j\) refers to a typed judgment whose admissibility depends on the witness, update, profile, level, visibility, threshold, audit, and defect data attached to the judgment, not on the primitive labels alone.

The set \(\mathbb P\) is reserved for primitive labels for the duration of the paper. Where a probability kernel and a primitive label appear in the same local context, the kernel is written with an explicit subscript or as a named transition operator so that no expression \(P\cdot P'\) or \(Px\) is ambiguous between a primitive label and a kernel application.

\subsubsection{Informal Role Gloss}

Each primitive carries a stable role description, fixed for the entire paper. The descriptions are role labels, not propositions: they identify the kind of typed witness, update, defect, and audit data that a primitive participates in, and they are not theorems about what a primitive does in any particular host.

\begin{itemize}
\tightlist
\item
  \(P_1\): descent, closure of updates, rewrite compatibility, transport of content. Witness data record descent successes, descent obstructions, and rewrite/closure deficits; updates install descended maps or rewrites.
\item
  \(P_2\): representability, gating, admissibility, support control. Witness data record representability or its failure, support, capacity, and gate decisions; updates mark representable or non-representable status and adjust support.
\item
  \(P_3\): route mismatch, holonomy, critical pairs, pasting. Witness data record commutator failures, holonomy values, critical-pair mismatches, and pasting defects; updates record these as typed defect entries.
\item
  \(P_4\): refinement, staging, gluing, local-to-global structure. Witness data record stages, refinements, covers, and local/global obstructions; updates refine the lens, advance the stage, or record a gluing defect.
\item
  \(P_5\): packaging, completion, objecthood, idempotence. Witness data record package maps, quotients, completion outcomes, and fixed points; updates install a package, apply a completion, or record a fixed-point residual.
\item
  \(P_6\): audit, provenance, replay, source-of-truth, drive certification. Witness data record audit checks, provenance, source identity, threshold checks, and (where present) drive certificates; updates add an audit record, commit a source, or certify drive or no-drive.
\end{itemize}

The audit face \(P_6\) and the drive face \(\mathrm{P6}_{\mathrm{drive}}\) are distinguished: a generic \(P_6\) audit witness is not a drive certificate. Section 8 isolates the drive face inside the graph/cohomology module, and the no-drive theorems there are stated for \(\mathrm{P6}_{\mathrm{drive}}\) specifically. The judgment discipline of the calculus preserves this separation throughout.

\subsubsection{Levels}

The finite level set used by the paper is

\[
\mathsf{Lev} \;=\; \{\mathsf{Beh},\,\mathsf{Ver},\,\mathsf{Str}\},
\]

with the following roles.

\begin{longtable}[]{@{}ll@{}}
\toprule
\begin{minipage}[b]{0.47\columnwidth}\raggedright
Level\strut
\end{minipage} & \begin{minipage}[b]{0.47\columnwidth}\raggedright
Role\strut
\end{minipage}\tabularnewline
\midrule
\endhead
\begin{minipage}[t]{0.47\columnwidth}\raggedright
\(\mathsf{Beh}\)\strut
\end{minipage} & \begin{minipage}[t]{0.47\columnwidth}\raggedright
behavioral level: model traces, runs, observed dynamics, raw data\strut
\end{minipage}\tabularnewline
\begin{minipage}[t]{0.47\columnwidth}\raggedright
\(\mathsf{Ver}\)\strut
\end{minipage} & \begin{minipage}[t]{0.47\columnwidth}\raggedright
verification level: instruments, audits, checks, replay, certification\strut
\end{minipage}\tabularnewline
\begin{minipage}[t]{0.47\columnwidth}\raggedright
\(\mathsf{Str}\)\strut
\end{minipage} & \begin{minipage}[t]{0.47\columnwidth}\raggedright
structural level: theorems, objects, packages, interfaces, definitional content\strut
\end{minipage}\tabularnewline
\bottomrule
\end{longtable}

Every directed-cell, pair-observable, promotion, and claim judgment in Sections 4 and beyond carries an explicit level annotation drawn from \(\mathsf{Lev}\). A judgment whose data spans more than one level is admissible only when an explicit level bridge is recorded. The admissible level bridges are

\[
\texttt{translation},\quad\texttt{lift},\quad\texttt{forgetting},\quad\texttt{reflection},
\]

each interpreted as a finite typed map between level-tagged records. A claim that crosses levels without an admissible bridge has status \(\texttt{blocked}\) or \(\texttt{outside\_scope}\) in the relevant judgment family of Section 4.

Levels are not a substitute for the channel/activation distinction: a primitive may be associated with a level annotation without producing an active witness at that level, and no theorem in this paper infers an active witness from a level tag alone. Levels are not a substitute for instruments either: a level fixes which kind of content a judgment refers to, while an instrument fixes which content is visible, suppressed, or outside scope and what claim strengths are admissible. Both are required.

\subsection{Claim Strengths and Status Discipline}

\subsubsection{Claim Strengths}

Every numbered statement in this paper carries a claim strength, drawn from the following finite taxonomy. Claim strengths apply to theorems, propositions, definitions, classifier rules, model-realization statements, countermodels, and the contents of the limitations section. These names are paper-facing labels for the viability categories used in the source packet: theorem-grade corresponds to \texttt{proved}, proposition-grade to \texttt{provable} once the proof is supplied in this paper, schema-grade to \texttt{framework-only}, model-only to \texttt{model-only}, future-work to \texttt{conjectural}, and dropped to \texttt{drop}. Countermodel-grade is the theorem-grade subcase in which the proof is an explicit finite countermodel.

\begin{longtable}[]{@{}ll@{}}
\toprule
\begin{minipage}[b]{0.47\columnwidth}\raggedright
Strength\strut
\end{minipage} & \begin{minipage}[b]{0.47\columnwidth}\raggedright
Meaning\strut
\end{minipage}\tabularnewline
\midrule
\endhead
\begin{minipage}[t]{0.47\columnwidth}\raggedright
theorem-grade\strut
\end{minipage} & \begin{minipage}[t]{0.47\columnwidth}\raggedright
A statement with a proof in the calculus, stated under explicit host, instrument, package, profile, visibility, threshold, audit, and nonclaim assumptions.\strut
\end{minipage}\tabularnewline
\begin{minipage}[t]{0.47\columnwidth}\raggedright
proposition-grade\strut
\end{minipage} & \begin{minipage}[t]{0.47\columnwidth}\raggedright
A statement that follows directly from definitions and prior theorems and is proved here without introducing new mathematical content.\strut
\end{minipage}\tabularnewline
\begin{minipage}[t]{0.47\columnwidth}\raggedright
schema-grade\strut
\end{minipage} & \begin{minipage}[t]{0.47\columnwidth}\raggedright
A definition, judgment form, classifier specification, gate list, or discipline rule. Not a standalone theorem; later theorems may quote it.\strut
\end{minipage}\tabularnewline
\begin{minipage}[t]{0.47\columnwidth}\raggedright
model-only\strut
\end{minipage} & \begin{minipage}[t]{0.47\columnwidth}\raggedright
A statement valid only as a scoped realization or application in a named model host. Not a universal theorem of the calculus.\strut
\end{minipage}\tabularnewline
\begin{minipage}[t]{0.47\columnwidth}\raggedright
countermodel-grade\strut
\end{minipage} & \begin{minipage}[t]{0.47\columnwidth}\raggedright
A finite construction together with the no-go statement it supports. The construction and the no-go are both required.\strut
\end{minipage}\tabularnewline
\begin{minipage}[t]{0.47\columnwidth}\raggedright
future-work\strut
\end{minipage} & \begin{minipage}[t]{0.47\columnwidth}\raggedright
A statement that is plausible from the calculus but not proved in this paper. Stated only in the limitations and future-work section.\strut
\end{minipage}\tabularnewline
\begin{minipage}[t]{0.47\columnwidth}\raggedright
dropped\strut
\end{minipage} & \begin{minipage}[t]{0.47\columnwidth}\raggedright
A statement that earlier exploration considered but that this paper does not assert as a theorem, model realization, schema, or future-work claim. Recorded in the dropped-claims appendix or limitations discussion for traceability.\strut
\end{minipage}\tabularnewline
\bottomrule
\end{longtable}

Each later theorem environment in the paper will be tagged with one of these strengths in its statement. A theorem-grade statement does not become more general by being labelled theorem-grade; it remains scoped by the host, instrument, package, profile, visibility, threshold, audit, and nonclaim assumptions on its judgment line. A model-only statement does not gain universal force by being placed beside a theorem-grade statement; the surrounding model-realization discipline of Section 13 keeps the two distinct.

The countermodel-grade label is reserved for statements of the form ``\(X\not\Rightarrow Y\)'' supported by a finite explicit construction that exhibits an instance of \(X\) together with a failure of \(Y\). Sections 6, 7, 8, 9, and 10 use this label for their no-go and separation theorems, and the full atlas is recorded in Appendix C.

\subsubsection{Status Families}

The calculus uses several distinct status families, one per judgment family. They are introduced informally here and defined fully in Sections 4 and 5. The families are:

\begin{itemize}
\tightlist
\item
  role-channel statuses, classifying the standing of a primitive role channel under an instrument, with values such as \(\texttt{active\_projection}\), \(\texttt{below\_threshold}\), \(\texttt{collapsed\_boundary\_case}\), \(\texttt{absent\_with\_record}\), and \(\texttt{inapplicable\_outside\_scope}\) inherited from the channel discipline of Foundations II;
\item
  directed-cell statuses, classifying typed directed-cell judgments \(P_i\leftarrow P_j\) with values such as \(\texttt{action}\), \(\texttt{implicit}\), \(\texttt{trivial}\), \(\texttt{blocked}\), \(\texttt{undefined\_circular}\), \(\texttt{below\_threshold}\), \(\texttt{collapsed}\), \(\texttt{absent}\), and \(\texttt{outside\_scope}\);
\item
  pair-observable statuses, classifying unordered pair observables, with values \(\texttt{real}\), \(\texttt{real\_provisional}\), \(\texttt{real\_duplicated}\), and \(\texttt{blocked}\);
\item
  promotion statuses, classifying promotion-bridge judgments, with values \(\texttt{candidate}\), \(\texttt{accepted}\), \(\texttt{strict}\), \(\texttt{non\_strict}\), \(\texttt{failed\_descent}\), \(\texttt{failed\_stability}\), \(\texttt{failed\_audit}\), \(\texttt{failed\_no\_smuggling}\), \(\texttt{local\_only}\), \(\texttt{globally\_obstructed}\), and \(\texttt{outside\_scope}\);
\item
  claim statuses, classifying instrument-indexed claim records, with values \(\texttt{accepted}\), \(\texttt{rejected}\), \(\texttt{provisional}\), \(\texttt{blocked}\), \(\texttt{outside\_scope}\), \(\texttt{absent\_with\_record}\), \(\texttt{below\_threshold}\), \(\texttt{failed\_audit}\), and \(\texttt{undefined\_circular}\);
\item
  gate statuses, classifying individual gate checks during promotion or claim acceptance, with values \(\texttt{pass}\), \(\texttt{fail}\), \(\texttt{not\_required}\), \(\texttt{not\_checked}\), and \(\texttt{outside\_scope}\), together with the strictness-gate refinements defined later;
\item
  square statuses, classifying decorated-square fillers in the secondary high-structure module of Section 14, with values \(\texttt{exact}\), \(\texttt{lax}\), \(\texttt{obstructed}\), \(\texttt{blocked}\), \(\texttt{outside\_scope}\), \(\texttt{failed\_audit}\), \(\texttt{undefined\_circular}\), and \(\texttt{provisional}\).
\end{itemize}

The discipline of the paper is that these families are kept apart. No theorem in the paper infers a value in one family from a value in another without an explicit bridge rule. In particular, role-channel activity does not by itself imply directed-cell action, directed-cell action does not imply pair-observable realness, claim acceptance under one instrument does not imply instrument-free truth or acceptance under another instrument, and gate passage does not imply theorem-grade strength of a surrounding statement. Section 6 includes the formal status-family separation theorem that closes off the most common collapses.

\subsubsection{No-Overreading Rule}

The visibility discipline of the paper is summarized informally as the no-overreading rule:

\begin{quote}
An accepted claim under instrument \(I\) may use only the content that \(I\) marks visible, together with content that an admissible bridge under \(I\) has connected to visible content. Suppressed content not connected by an admissible bridge cannot support an accepted claim, even when the suppressed content is in fact correct.
\end{quote}

The rule is not a metaphysical statement about truth. Suppressed content is not asserted to be false, and content outside the scope of \(I\) is not asserted to be unknowable. The rule is a record-level constraint: an instrument's acceptance verdict is admissible only when the use-set of the claim has been audited against the instrument's visibility map and bridge data, and any suppressed dependency without a bridge invalidates acceptance under \(I\). The formal statement, with the definition of the overreading defect and the precise acceptance rule, appears in Section 11; the rule is invoked in earlier sections only at this informal level.

\subsection{Paper Scope Statement}

\subsubsection{Formal Scope}

All judgments, theorems, propositions, schemas, model realizations, and countermodels in this paper are interpreted inside

\[
\mathbf{BirdInt}^{\mathrm{aud}}_{\mathrm{fin}}
\]

or inside one of the explicitly named finite extension hosts of Subsection 2.1.2. Every theorem statement carries explicit host, instrument, package, profile, visibility, threshold, audit, and nonclaim assumptions, and these assumptions are part of the theorem rather than an ambient context that may be relaxed without comment. The shorthand

\[
I\vdash\varphi(\mathcal T)
\]

abbreviates

\[
\Gamma;\mathcal T;I\vdash\varphi:\texttt{accepted},
\]

with \(\Gamma\) a finite typed context, \(\mathcal T\) an admissible theory package, \(I\) an instrument, and \(\varphi\) a claim whose use-set has been audited against the visibility, threshold, audit, and nonclaim records attached to the judgment. The shorthand never abbreviates instrument-free truth, and a statement of the form \(I\vdash\varphi(\mathcal T)\) does not transfer to a different instrument \(I'\) except through an explicit instrument bridge.

\subsubsection{Main Positive Scope}

Within the formal scope above, the paper proves and presents:

\begin{itemize}
\tightlist
\item
  a finite typed audited interaction calculus with role-channel, directed-cell, pair-observable, promotion, dependency, and instrument-indexed claim judgments, together with the typed witness, update, profile, level, visibility, threshold, audit, and defect data those judgments carry;
\item
  a finite theorem spine including a no-total-six-symbol-algebra theorem, typed non-collapse, cell well-formedness decidability, covered-cell status uniqueness, status-family separation, descent through a finite quotient, the Markov closure-deficit theorem under finite kernel and partition assumptions, noncommuting completions, idempotent saturation, strict-extension nonfactorization, and a fixed-interface definability bound;
\item
  a graph/cohomology drive module containing forest no-drive, exact-form null-drive, nonzero-affinity equivalence, and gating-affinity suppression, scoped to finite cycle-supported cochain data rather than used as a universal theory of drive;
\item
  a finite countermodel atlas blocking the most common overclaims, with the construction, supported no-go, blocked overclaim, and any required positive bridge recorded for each entry;
\item
  finite promotion-bridge theory with promotion-gate soundness, strict-promotion-implies-nonfactorization, and non-strict-promotion-implies-factorization;
\item
  instrument-indexed claim semantics with no-overreading suppression, same-level self-audit failure, and a finite rotating-audit theorem;
\item
  scoped model-realization theorems for finite stochastic PICA, the audited Cantor shell, and finite abstract interpretation, each tagged model-only and accompanied by the host, the realized fragment, the realization map, and the explicit nonclaims;
\item
  secondary high-structure semantics in a finite decorated double-category fragment, used only as an alternative organization of finite descent, completion, and strict-extension data, never as a replacement for the typed/statused calculus.
\end{itemize}

\subsubsection{Main Negative Scope}

The paper does not prove, and does not assume as background, any of the following:

\begin{itemize}
\tightlist
\item
  a universal exact-six theorem stating that every phenomenon, every theory, or all of mathematics reduces to the six primitive labels;
\item
  a total binary algebra \(*:\mathbb P^2\to\mathbb P\) on the primitive labels, or any equivalent reduction of typed interaction to primitive-label arithmetic;
\item
  a rule by which the primitive pair \((P_i,P_j)\) alone determines directed-cell status;
\item
  collapse of role-channel, directed-cell, pair-observable, promotion, claim, gate, or square status families;
\item
  an implication from role-channel activity to directed-cell action, or from directed-cell action to pair-observable realness;
\item
  a unique decomposition theorem stating that every theory, package, or interaction admits a canonical decomposition into the six primitives;
\item
  package-implies-closure, idempotence-implies-commutation, repeated fixed completion as strict theory growth, local-implies-global packaging, or refinement-implies-improvement;
\item
  strictness-implies-closure, strictness-implies-drive, or route-mismatch-implies-drive;
\item
  audit-is-truth, instrument-free acceptance, accepted claims over unbridged suppressed content, or rotating-audit finality;
\item
  a same-level complete self-certification of any instrument or theory, including the interpreting instrument of this paper;
\item
  empirical realization of the calculus by any external physical, biological, or computational system; empirical-bridge assertions about such systems require a finite admissible record on a named host and do not become realization theorems;
\item
  model-realization-as-universal-theorem, PICA universality, or a PICA \(6\times6\) table as a universal interaction algebra;
\item
  complete model coverage by the model families of Section 13, that is, a theorem stating that every fragment of the calculus is realized by PICA, Cantor, abstract interpretation, or graph/cohomology data;
\item
  proof-assistant certification beyond the finite audited declarations listed in \texttt{lean/manifest.toml}. Section 15 and Appendix F record the completed Lean 4 validation track for that manifest, not a claim about unlisted statements, empirical implementations, or unrestricted host extensions.
\end{itemize}

The negative scope is not a defensive afterthought. It is a working constraint on every later section: a theorem that drifts into any of these excluded claims, or that quietly relaxes its host, instrument, or audit assumptions to make a stronger statement appear to follow, is rejected on scope grounds before its proof is examined. The limitations and nonclaims section collects the dropped and deferred statements in their final form.
 \section{The Finite Audited Domain}

This section declares the ambient domain of the calculus, before the central judgment object is defined in Section 4. The domain is a finite tuple of finite typed components: a primitive set, a level set, a host inventory, a content universe, instruments, packages, profiles, witness and update families, defect records, judgment families, status families, gates, dependencies, audits, visibility data, source records, nonclaim records, and model-realization modules. Each component is introduced as a finite schema; the theorems and classifiers that depend on these components are stated in later sections and never quoted as already established here.

\subsection{Domain Declaration}

\subsubsection{Full Tuple}

The finite audited interaction calculus is the tuple

\[
\mathbf{BirdInt}^{\mathrm{aud}}_{\mathrm{fin}}
\;=\;
(\,\mathbb P,\;\mathsf{Lev},\;\mathsf{Host},\;\mathsf{Cont},\;\mathsf{Instr},\;\mathsf{Pkg},\;\mathsf{Prof},\;\mathsf{Wit},\;\mathsf{Upd},\;\mathsf{Def},\;\mathsf{Judg},\;\mathsf{Status},\;\mathsf{Gate},\;\mathsf{Dep},\;\mathsf{Audit},\;\mathsf{Vis},\;\mathsf{Source},\;\mathsf{Nonclaim},\;\mathsf{Real}\,).
\]

The tuple is a finite record: each component is either a finite set, a finite family of finite records, or a finite typed map between such families. The components are introduced in Subsections 3.2 through 3.7 and Section 4; full definitions of the status, defect, gate, audit, and model-realization components appear in Sections 4, 5, 10, 11, and 13, respectively. The tuple itself is the schema-grade declaration of what counts as the ambient domain of the calculus.

The notation \(\mathbf{BirdInt}^{\mathrm{aud}}_{\mathrm{fin}}\) refers to this declared domain throughout the paper. It is distinct from the inherited Foundations II domain \(\mathsf{FATCD}\) of finite audited typed closure descriptions: \(\mathsf{FATCD}\) is the description domain for the scoped exact-six program of that paper, and \(\mathbf{BirdInt}^{\mathrm{aud}}_{\mathrm{fin}}\) is the interaction calculus of this paper. Where Foundations II structures are reused inside the calculus, the paper records the embedding rather than renaming \(\mathsf{FATCD}\).

\subsubsection{Interpretation of Components}

The components of the tuple have the following roles. Each description is schema-grade: it fixes what kind of object the component is, not what theorems it satisfies.

\begin{itemize}
\tightlist
\item
  \(\mathbb P\) is the fixed finite set of primitive labels \(\{P_1,\ldots,P_6\}\) defined in Subsection 2.2.1. It is the underlying label alphabet of every role-channel, directed-cell, pair-observable, promotion, and dependency judgment.
\item
  \(\mathsf{Lev}=\{\mathsf{Beh},\mathsf{Ver},\mathsf{Str}\}\) is the finite level set defined in Subsection 2.2.3, used to tag the kind of content a judgment refers to and to gate cross-level claims through translation, lift, forgetting, or reflection bridges.
\item
  \(\mathsf{Host}\) is the finite inventory of permitted hosts declared in Subsection 2.1.2: \(\mathsf H_{\mathrm{fin}}\), \(\mathsf H_{\mathrm{prob}}\), \(\mathsf H_{\mathrm{graph}}\), \(\mathsf H_{\mathrm{AI}}\), \(\mathsf H_{\mathrm{PICA}}\), \(\mathsf H_{\mathrm{CantorShell}}\), \(\mathsf H_{\mathrm{DC}}\), \(\mathsf H_{\mathrm{instr}}\). Every theorem statement names its host.
\item
  \(\mathsf{Cont}\) is the finite content universe of the domain, defined in Subsection 3.2. It collects the records that claims may use: theory packages, instruments, witnesses, updates, profiles, defect records, threshold records, audit records, visibility records, level/lens/interface records, package and promotion interfaces, quotients, lenses, completion endomaps, kernels, graphs, finite families, refinements, and other finite typed records introduced as the calculus is built.
\item
  \(\mathsf{Instr}\) is the finite family of instrument records, defined in Subsection 3.4. An instrument fixes the visible/suppressed map, scope, claim types, threshold and audit policies, level assignment, bridge policy, source policy, and nonclaim list under which a claim is to be evaluated.
\item
  \(\mathsf{Pkg}\) is the finite family of theory packages, defined in Subsection 3.3, each carrying the inherited Foundations I package data \((Z,f,\Sigma_f,E,\mathcal A)\) together with the additional Foundations III typed records (visibility, source, defect, promotion, model-realization) attached around it.
\item
  \(\mathsf{Prof}\) is the finite family of profile records, defined in Subsection 3.6. A profile fixes per-judgment data such as actor and informant level tags, scale, regime, bridge type, instrument mode, and locality, and is required by every directed-cell, pair-observable, promotion, and claim judgment.
\item
  \(\mathsf{Wit}\) assigns to each primitive \(P_i\in\mathbb P\) its typed witness family \(\mathsf{Wit}(P_i)\), defined in Subsection 3.7. The informant-side data of a directed cell is required to belong to one of these families.
\item
  \(\mathsf{Upd}\) assigns to each primitive \(P_i\in\mathbb P\) its typed update family \(\mathsf{Upd}(P_i)\), defined in Subsection 3.7. The actor-side data of a directed cell is required to belong to one of these families.
\item
  \(\mathsf{Def}\) is the finite family of defect records, defined in Section 5. A defect record carries a defect type, a domain, a witness, a threshold, a measured value, a status consequence, and an audit/source record, and feeds the defect-to-status rules of Subsection 5.6.
\item
  \(\mathsf{Judg}\) is the finite family of judgment forms used by the calculus: role-channel, directed-cell, pair-observable, promotion, dependency, and instrument-indexed claim judgments, all defined in Section 4.
\item
  \(\mathsf{Status}\) is the finite family of status sets, one per judgment family: role-channel, directed-cell, pair-observable, promotion, claim, gate, and square statuses. The sets are listed in Subsection 2.3.2 and defined fully in Section 5.
\item
  \(\mathsf{Gate}\) is the finite family of gate specifications used during promotion and claim acceptance, defined in Section 10. A gate is a finite check whose \(\texttt{pass}\), \(\texttt{fail}\), \(\texttt{not\_required}\), \(\texttt{not\_checked}\), or \(\texttt{outside\_scope}\) result feeds promotion-status and claim-status decisions.
\item
  \(\mathsf{Dep}\) is the finite family of dependency relations between primitives, including positive dependencies, blocked dependencies, bridge-required dependencies, and countermodel-supported no-go dependencies, defined in Subsection 4.6.
\item
  \(\mathsf{Audit}\) is the finite family of audit records, defined in Subsection 4.8. An audit record carries a target, an instrument, replay or check data, and admissibility data; it is distinct from the inherited audit functional \(\mathcal A\) that lives inside a theory package \(\mathcal T\).
\item
  \(\mathsf{Vis}\) is the finite family of visibility maps and visibility bridges, defined in Subsection 3.5. A visibility map sends content to one of \(\texttt{visible}\), \(\texttt{suppressed}\), \(\texttt{outside}\), \(\texttt{unknown}\) under an instrument, and a visibility bridge is the admissible record by which a claim may use suppressed content.
\item
  \(\mathsf{Source}\) is the finite family of source records, defined in Subsection 4.8. A source record fixes source identity, source type, trace, and source validity, and is consulted by pair-observable and claim acceptance rules.
\item
  \(\mathsf{Nonclaim}\) is the finite family of nonclaim records carried by an instrument or attached to a judgment. A nonclaim records what the surrounding statement does not assert, in the same finite typed format as the rest of the calculus.
\item
  \(\mathsf{Real}\) is the finite family of model-realization modules, defined in Section 13: the PICA, audited Cantor shell, abstract-interpretation, and graph/cohomology realizations of named fragments of the calculus.
\end{itemize}

The tuple is finite at every component: each of \(\mathbb P\), \(\mathsf{Lev}\), \(\mathsf{Host}\), the per-instrument and per-package records, the witness and update families per primitive, and the per-judgment status sets is finite. No component is defined by an unbounded reflection scheme or a closure under arbitrary semantic operations. A claim that requires an extension of any component is not interpreted in \(\mathbf{BirdInt}^{\mathrm{aud}}_{\mathrm{fin}}\) as defined here and lies outside the formal scope of this paper.

\subsection{Content Universe}

\subsubsection{Finite Content Objects}

The content universe of the calculus is the finite family

\[
\mathsf{Cont}
\]

of records that claims, judgments, and audits may refer to. Every member of \(\mathsf{Cont}\) is a finite typed record built from the base host \(\mathsf H_{\mathrm{fin}}\), possibly with components drawn from a declared extension host. The schema is extended only by the finite record types introduced in this paper's declared components; each such record type is required to be a finite typed object on a named host, and no record type is defined by closure under an unbounded operation.

The standard members of \(\mathsf{Cont}\) used by this paper are the following.

\begin{longtable}[]{@{}ll@{}}
\toprule
\begin{minipage}[b]{0.47\columnwidth}\raggedright
Record kind\strut
\end{minipage} & \begin{minipage}[b]{0.47\columnwidth}\raggedright
Reference\strut
\end{minipage}\tabularnewline
\midrule
\endhead
\begin{minipage}[t]{0.47\columnwidth}\raggedright
theory package \(\mathcal T=(Z,f,\Sigma_f,E,\mathcal A)\) together with attached Foundations III records\strut
\end{minipage} & \begin{minipage}[t]{0.47\columnwidth}\raggedright
Subsection 3.3\strut
\end{minipage}\tabularnewline
\begin{minipage}[t]{0.47\columnwidth}\raggedright
instrument record \(I\)\strut
\end{minipage} & \begin{minipage}[t]{0.47\columnwidth}\raggedright
Subsection 3.4\strut
\end{minipage}\tabularnewline
\begin{minipage}[t]{0.47\columnwidth}\raggedright
primitive label \(P_i\in\mathbb P\)\strut
\end{minipage} & \begin{minipage}[t]{0.47\columnwidth}\raggedright
Subsection 2.2.1\strut
\end{minipage}\tabularnewline
\begin{minipage}[t]{0.47\columnwidth}\raggedright
witness \(W_i\in\mathsf{Wit}(P_i)\)\strut
\end{minipage} & \begin{minipage}[t]{0.47\columnwidth}\raggedright
Subsection 3.7\strut
\end{minipage}\tabularnewline
\begin{minipage}[t]{0.47\columnwidth}\raggedright
update \(U_i\in\mathsf{Upd}(P_i)\)\strut
\end{minipage} & \begin{minipage}[t]{0.47\columnwidth}\raggedright
Subsection 3.7\strut
\end{minipage}\tabularnewline
\begin{minipage}[t]{0.47\columnwidth}\raggedright
profile \(\lambda\)\strut
\end{minipage} & \begin{minipage}[t]{0.47\columnwidth}\raggedright
Subsection 3.6\strut
\end{minipage}\tabularnewline
\begin{minipage}[t]{0.47\columnwidth}\raggedright
level/lens/interface record \(L\)\strut
\end{minipage} & \begin{minipage}[t]{0.47\columnwidth}\raggedright
Subsection 4.3\strut
\end{minipage}\tabularnewline
\begin{minipage}[t]{0.47\columnwidth}\raggedright
visibility record \(V\), suppressed-content record, visibility bridge\strut
\end{minipage} & \begin{minipage}[t]{0.47\columnwidth}\raggedright
Subsection 3.5\strut
\end{minipage}\tabularnewline
\begin{minipage}[t]{0.47\columnwidth}\raggedright
threshold record \(\Theta\)\strut
\end{minipage} & \begin{minipage}[t]{0.47\columnwidth}\raggedright
Subsection 5.2\strut
\end{minipage}\tabularnewline
\begin{minipage}[t]{0.47\columnwidth}\raggedright
audit record \(A\), source record, replay trace\strut
\end{minipage} & \begin{minipage}[t]{0.47\columnwidth}\raggedright
Subsection 4.8\strut
\end{minipage}\tabularnewline
\begin{minipage}[t]{0.47\columnwidth}\raggedright
defect record \(\delta\) with type, domain, witness, threshold, measured value, status consequence\strut
\end{minipage} & \begin{minipage}[t]{0.47\columnwidth}\raggedright
Section 5\strut
\end{minipage}\tabularnewline
\begin{minipage}[t]{0.47\columnwidth}\raggedright
object map or interface \(\pi_j\), \(\pi_{j+1}\) in a promotion bridge\strut
\end{minipage} & \begin{minipage}[t]{0.47\columnwidth}\raggedright
Subsection 4.5\strut
\end{minipage}\tabularnewline
\begin{minipage}[t]{0.47\columnwidth}\raggedright
quotient \(q\), finite map \(F\), descended map \(F^\sharp\), completion endomap \(E\)\strut
\end{minipage} & \begin{minipage}[t]{0.47\columnwidth}\raggedright
Sections 6 and 7\strut
\end{minipage}\tabularnewline
\begin{minipage}[t]{0.47\columnwidth}\raggedright
finite stochastic kernel \(K\) on a finite state space, finite information quantities\strut
\end{minipage} & \begin{minipage}[t]{0.47\columnwidth}\raggedright
Subsection 7.3 (host \(\mathsf H_{\mathrm{prob}}\))\strut
\end{minipage}\tabularnewline
\begin{minipage}[t]{0.47\columnwidth}\raggedright
finite graph \(G\), cycle basis, cochain family \(\mathcal F\)\strut
\end{minipage} & \begin{minipage}[t]{0.47\columnwidth}\raggedright
Section 8 (host \(\mathsf H_{\mathrm{graph}}\))\strut
\end{minipage}\tabularnewline
\begin{minipage}[t]{0.47\columnwidth}\raggedright
refinement, staging, gluing record \(\Delta\)\strut
\end{minipage} & \begin{minipage}[t]{0.47\columnwidth}\raggedright
Subsection 5.3 and Section 7\strut
\end{minipage}\tabularnewline
\begin{minipage}[t]{0.47\columnwidth}\raggedright
nonclaim record \(\mathcal N\)\strut
\end{minipage} & \begin{minipage}[t]{0.47\columnwidth}\raggedright
Subsection 4.8\strut
\end{minipage}\tabularnewline
\bottomrule
\end{longtable}

Every claim \(\varphi\) in the calculus has a finite use-set

\[
\mathsf{Use}(\varphi)\;\subseteq\;\mathsf{Cont}
\]

consisting of the records that the statement of \(\varphi\) refers to. Visibility, suppression, source-of-truth, and overreading checks in later sections operate on \(\mathsf{Use}(\varphi)\), not on the entire content universe. A claim whose use-set contains a record that is not finite, not typed, or not on a permitted host is not admissible in \(\mathbf{BirdInt}^{\mathrm{aud}}_{\mathrm{fin}}\).

\subsubsection{Typed Content}

Every record in \(\mathsf{Cont}\) is typed in four respects.

\begin{itemize}
\tightlist
\item
  \textbf{Host.} Each record carries a host tag drawn from \(\mathsf{Host}\). The host tag fixes which finite mathematical universe the record lives in: \(\mathsf H_{\mathrm{fin}}\) for the base host, \(\mathsf H_{\mathrm{prob}}\) for finite stochastic content, \(\mathsf H_{\mathrm{graph}}\) for finite graph and cohomology content, \(\mathsf H_{\mathrm{AI}}\) for finite abstract-interpretation content, \(\mathsf H_{\mathrm{PICA}}\) and \(\mathsf H_{\mathrm{CantorShell}}\) for the model-realization hosts, \(\mathsf H_{\mathrm{DC}}\) for the high-structure fragment, and \(\mathsf H_{\mathrm{instr}}\) for instrument-side content.
\item
  \textbf{Level.} Each record carries a level tag drawn from \(\mathsf{Lev}=\{\mathsf{Beh},\mathsf{Ver},\mathsf{Str}\}\). The level tag fixes whether the record describes behavioral/trace content, verification/audit content, or structural/object content. A judgment that uses records at more than one level is admissible only with an explicit level bridge (translation, lift, forgetting, or reflection).
\item
  \textbf{Instrument.} Each record is interpreted under the instrument \(I\) that the surrounding judgment names. The instrument fixes whether the record is visible, suppressed, outside-scope, or unknown for the purpose of the claim, and what bridges may connect it to visible content. A record may be in \(\mathsf{Cont}\) globally and yet be suppressed under a particular instrument; suppression is a relation between a record and an instrument, not a property of the record alone, and it does not assert that the suppressed record is false.
\item
  \textbf{Profile.} Each judgment is tagged by a profile \(\lambda\in\mathsf{Prof}\) which fixes per-judgment data including actor and informant level tags, scale, regime, bridge type, instrument mode, and locality. The profile of a judgment together with its host, level, and instrument tags determines the admissibility of every record in its use-set.
\end{itemize}

Typed content is the working unit of the calculus: every later definition, theorem, classifier, gate, and audit operates on records whose host, level, instrument, and profile tags are explicit. A statement that refers to ``a finite map'' or ``a record'' without an admissible host, level, instrument, and profile context is a schema, not a judgment, and does not by itself support a claim.

\subsection{Theory Packages}

\subsubsection{Package Record}

A theory package is the inherited Foundations I object

\[
\mathcal T \;=\; (\,Z,\;f,\;\Sigma_f,\;E,\;\mathcal A\,)
\]

together with the additional Foundations III typed records that the calculus attaches to it. The inherited components have the inherited meaning: \(Z\) is a finite carrier or scoped comparison domain, \(f\) is a finite lens or coarse-description map of \(Z\), \(\Sigma_f\) is the expressive content (finite definability) induced by \(f\), \(E\) is a completion or packaging endomap, and \(\mathcal A\) is the audit functional inherited from Foundations I. The inherited form is preserved verbatim throughout the paper; no Foundations III usage redefines \(\mathcal T\), \(f\), \(\Sigma_f\), \(E\), or \(\mathcal A\) as different objects.

In Foundations III, a package record carries the inherited five-tuple together with finite typed records for the additional content the interaction calculus depends on:

\begin{itemize}
\tightlist
\item
  a host tag \(h_{\mathcal T}\in\mathsf{Host}\) fixing the host on which the package lives;
\item
  a level tag \(\ell_{\mathcal T}\in\mathsf{Lev}\) fixing whether the package describes behavioral, verification, or structural content;
\item
  a finite visibility profile \(V_{\mathcal T}\) recording which fields of the package are designated visible, suppressed, outside-scope, or unknown by the instruments that consult it;
\item
  a finite source record \(\mathrm{src}_{\mathcal T}\) recording source identity, source type, trace, and source validity;
\item
  a finite audit record family \(A_{\mathcal T}\) supplying provenance, replay, and threshold-check entries used by the package's audit functional \(\mathcal A\);
\item
  a finite defect record family \(\delta_{\mathcal T}\) tracking package-level closure deficits, gluing obstructions, idempotence residuals, and other defects;
\item
  finite promotion data when the package participates in a promotion bridge: an old interface, a target package, a target interface, and the bridge gate record;
\item
  finite model-realization data when the package realizes a fragment of \(\mathbf{BirdInt}^{\mathrm{aud}}_{\mathrm{fin}}\) inside a named model host;
\item
  a finite nonclaim record \(\mathcal N_{\mathcal T}\) recording what claims about the package are explicitly not asserted.
\end{itemize}

The visibility, suppression, source, audit, defect, promotion, model-realization, and nonclaim records are extra structure attached around the inherited five-tuple. They never silently replace any inherited component, and a Foundations III theorem about a package always restates which inherited fields it uses.

The canonical lens-induced expressive content is the inherited family

\[
\Sigma_f \;=\; \{\,f^{-1}(B)\;:\;B\subseteq f(Z)\,\},
\]

and the canonical lens completion is the inherited endomap

\[
E_f(S)\;=\;f^{-1}(f(S)).
\]

Where Foundations III uses a different completion, the choice is named explicitly and the relationship to \(E_f\) is recorded.

The word ``closure'' is used only with its declared sense. The completion endomap \(E\) is a packaging/completion operation on a declared content space; an order-closure operator is used only on a poset host satisfying the relevant closure axioms; and macro closure refers to closed induced dynamics such as a descended update law. None of these follows from the others without an explicit theorem or bridge.

\subsubsection{Package Morphisms}

The calculus uses four kinds of finite morphism between theory packages, all of which are records on the base host (or on an extension host when a package lives there).

\begin{itemize}
\tightlist
\item
  \textbf{Package map.} A package map \(\Phi:\mathcal T\to\mathcal T'\) is a finite typed record consisting of a finite map \(\phi_Z:Z\to Z'\) on carriers together with the induced data on lenses, expressive content, completion endomaps, audit functionals, and Foundations III attached records. A package map is admissible when \(\phi_Z\) commutes with the lens data up to the threshold and visibility recorded by both packages, and when the audit and source records of the codomain extend those of the domain.
\item
  \textbf{Refinement.} A package refinement is a package map \(\Phi:\mathcal T\to\mathcal T'\) for which the domain lens factors through the codomain lens, so the codomain expressive content extends the domain expressive content, and the refinement is recorded in the level/lens/interface field of any judgment that uses it. A refinement is not assumed to improve any closure deficit; refinement-implies-improvement is excluded by the negative scope of Subsection 2.4.3.
\item
  \textbf{Quotient.} A package quotient is a package map \(\Phi\) whose carrier component \(\phi_Z\) is a finite surjection \(q:Z\to Z/\!\sim\) for a finite equivalence relation \(\sim\) on \(Z\), together with the matching quotient data on lens, expressive content, and completion. Quotients are the finite maps used by the descent theorems of Section 7.
\item
  \textbf{Completion-induced map.} A completion-induced map is a finite package map or package record induced by applying a named completion endomap \(E:\mathcal V\to\mathcal V\) on its declared content space. When the host supplies inclusions or fixed-point records, those maps are recorded explicitly; they are not inferred from the symbol \(E\) alone. Repeated application of an idempotent completion is not treated as strict package growth unless a new interface, bridge, refinement, or nonfactorization witness is recorded.
\end{itemize}

Package maps compose as records, but composition is not assumed to preserve every Foundations III attached record without an explicit compatibility condition: in particular, audit and source records compose only when the audit functionals of the codomain extend those of the domain, and a composition that violates an inherited audit condition is a defect, not a new package map.

\subsubsection{Package Admissibility}

A theory package \(\mathcal T\) is admissible in \(\mathbf{BirdInt}^{\mathrm{aud}}_{\mathrm{fin}}\), written

\[
\mathsf{AdmPkg}(\mathcal T),
\]

when all of the following hold.

\begin{itemize}
\tightlist
\item
  The inherited five-tuple \((Z,f,\Sigma_f,E,\mathcal A)\) is present and finite. The carrier \(Z\) is a finite set, the lens \(f\) is a finite map from \(Z\) to a finite codomain, the expressive content \(\Sigma_f\) is a finite family of subsets of \(Z\), the completion endomap \(E\) is a finite endomap on the relevant content space, and the audit functional \(\mathcal A\) is a finite functional whose domain is determined by the package's host.
\item
  The host tag \(h_{\mathcal T}\) lies in \(\mathsf{Host}\), and every component of the package lives in or under that host.
\item
  The level tag \(\ell_{\mathcal T}\) lies in \(\mathsf{Lev}\). Cross-level use of the package requires a level bridge.
\item
  The visibility profile \(V_{\mathcal T}\) partitions the package's fields into visible, suppressed, outside-scope, and unknown classes for each instrument that may consult \(\mathcal T\), and the partition is finite.
\item
  The source record \(\mathrm{src}_{\mathcal T}\) is present, finite, and consistent with the source policy of any instrument that consults \(\mathcal T\).
\item
  The audit record family \(A_{\mathcal T}\) is finite and consistent with \(\mathcal A\). The audit record family is distinct from the audit functional: \(A_{\mathcal T}\) supplies the finite checked entries that \(\mathcal A\) certifies.
\item
  The defect record family \(\delta_{\mathcal T}\) is finite, with each entry conforming to the defect record schema of Section 5.
\item
  Any attached promotion data, model-realization data, or level/lens/interface data is finite, hosted, assigned an explicit visibility status under the relevant instrument, and accompanied by its own source, audit, and nonclaim records.
\item
  The nonclaim record \(\mathcal N_{\mathcal T}\) is present and finite.
\end{itemize}

A package that fails any of these conditions is not admissible. Subsequent sections only consider admissible packages, and a theorem statement that uses a package implicitly carries \(\mathsf{AdmPkg}(\mathcal T)\) as a hypothesis. The finite shape of these checks is part of the input to the later well-formedness results; this subsection only records the package schema and admissibility conditions.

\subsection{Instruments}

\subsubsection{Instrument Record}

An instrument is a finite typed record

\[
I\;=\;(\,\mathrm{id}_I,\;\mathsf{Scope}_I,\;\mathsf{ClaimTypes}_I,\;\mathsf{Visible}_I,\;\mathsf{Suppressed}_I,\;\mathsf{CheckRules}_I,\;\mathsf{EvidencePolicy}_I,\;\mathsf{ThresholdPolicy}_I,\;\mathsf{AuditPolicy}_I,\;\mathsf{Level}_I,\;\mathsf{BridgePolicy}_I,\;\mathsf{SourcePolicy}_I,\;\mathcal N_I\,).
\]

The components have the following roles, all finite:

\begin{itemize}
\tightlist
\item
  \(\mathrm{id}_I\): a finite identifier under which the instrument is registered. Two distinct instruments have distinct identifiers, and an instrument identifier is part of the data of every claim accepted under that instrument.
\item
  \(\mathsf{Scope}_I\): the finite scope record of the instrument, fixing the host, level, and content classes that fall inside its scope. A claim whose use-set is not contained in \(\mathsf{Scope}_I\) has status \(\texttt{outside\_scope}\) under \(I\).
\item
  \(\mathsf{ClaimTypes}_I\): the finite family of claim types that \(I\) is configured to evaluate, including role-channel, directed-cell, pair-observable, promotion, dependency, and instrument-indexed claim types as appropriate.
\item
  \(\mathsf{Visible}_I,\;\mathsf{Suppressed}_I\): the finite visibility partition of \(\mathsf{Scope}_I\) into visible and suppressed content under \(I\). The complement of \(\mathsf{Visible}_I\cup\mathsf{Suppressed}_I\) within \(\mathsf{Scope}_I\) is the unknown class; content outside \(\mathsf{Scope}_I\) is in the outside-scope class. The full four-way partition is given in Subsection 3.5.
\item
  \(\mathsf{CheckRules}_I\): the finite family of check rules used to classify a claim under \(I\). Check rules are typed finite records: each rule states which fields of a claim it consults, what comparison or test it performs, and what classifier verdict it contributes.
\item
  \(\mathsf{EvidencePolicy}_I\): the finite policy for what counts as admissible evidence under \(I\), distinguishing direct evidence, bridge-relayed evidence, and inadmissible evidence.
\item
  \(\mathsf{ThresholdPolicy}_I\): the finite family of thresholds used by \(I\), including exact, approximate, tolerance, rank/dimension, visibility, and source thresholds. The threshold record schema is fixed in Subsection 5.2.
\item
  \(\mathsf{AuditPolicy}_I\): the finite policy for which audit records \(I\) requires, in what order, and against which sources. The policy is checked by the gate semantics of Section 10 and the rotating-audit semantics of Section 11. Passing an audit policy contributes to an instrument-indexed status; it is not instrument-free truth.
\item
  \(\mathsf{Level}_I\in\mathsf{Lev}\): the level annotation of \(I\). An instrument operates at a single level; cross-level operation requires an explicit level bridge inside the instrument's claim record.
\item
  \(\mathsf{BridgePolicy}_I\): the finite policy for which bridges \(I\) admits, including visibility bridges, level bridges, instrument-transfer bridges, drive bridges, and model-realization bridges. The policy fixes which suppressed content can be brought into the use-set of an accepted claim. A generic audit bridge is not a drive certificate unless the bridge records the required \(\mathrm{P6}_{\mathrm{drive}}\) data.
\item
  \(\mathsf{SourcePolicy}_I\): the finite policy for which source-of-truth records \(I\) accepts, distinguishing real sources, fallback sources, dirty sources, proxy sources, and missing sources.
\item
  \(\mathcal N_I\): the finite nonclaim record of \(I\). The nonclaims of \(I\) are part of the instrument and propagate to every claim accepted under \(I\); an accepted claim under \(I\) is interpreted with the nonclaims of \(I\) attached.
\end{itemize}

Two instruments \(I_1\) and \(I_2\) may have the same identifier components in some fields and yet be distinct; equality of instruments is a finite record comparison across all fields. An instrument is admissible when every component is finite and consistent with \(\mathsf{Host}\), \(\mathsf{Lev}\), the visibility schema of Subsection 3.5, the threshold schema of Subsection 5.2, and the gate and audit schemas of Sections 10 and 11.

\subsubsection{Instrument-Relative Judgments}

Every nontrivial claim in the calculus is interpreted under an instrument. The general form of an instrument-indexed claim is

\[
\Gamma\,;\;\mathcal T\,;\;I\;\vdash\;\varphi:\chi,
\]

where \(\Gamma\) is a finite typed context of attached records (witnesses, updates, profiles, level/lens/interface records, visibility records, threshold records, audit records, defect records, and nonclaim records), \(\mathcal T\) is an admissible theory package, \(I\) is an admissible instrument, \(\varphi\) is a claim with a finite use-set \(\mathsf{Use}(\varphi)\subseteq\mathsf{Cont}\), and \(\chi\) is the claim status assigned to \(\varphi\) under \(\Gamma,\mathcal T,I\).

The shorthand

\[
I\;\vdash\;\varphi(\mathcal T)
\]

abbreviates

\[
\Gamma\,;\;\mathcal T\,;\;I\;\vdash\;\varphi:\texttt{accepted}.
\]

A claim that holds in this sense is accepted by \(I\) under the package \(\mathcal T\) and the surrounding context \(\Gamma\). The shorthand never abbreviates instrument-free truth, and a statement of the form \(I\vdash\varphi(\mathcal T)\) does not transfer to a different instrument \(I'\) except through an explicit instrument-transfer bridge admitted by both \(I\) and \(I'\).

The full classifier under which \(\chi\) is determined consults the visibility partition \(\mathsf{Visible}_I/\mathsf{Suppressed}_I\), the bridge policy \(\mathsf{BridgePolicy}_I\), the threshold policy \(\mathsf{ThresholdPolicy}_I\), the source policy \(\mathsf{SourcePolicy}_I\), the audit policy \(\mathsf{AuditPolicy}_I\), and the nonclaim record \(\mathcal N_I\). The classifier is specified at schema-grade in Subsection 11.2; the no-overreading and same-level self-audit theorems that depend on it are stated and proved at theorem-grade in Sections 11.4 and 11.5. This subsection only fixes the form of the judgment and the meaning of the shorthand. No theorem about acceptance or overreading is asserted here.

A claim that names a context and package but does not name an instrument is not a judgment in the calculus. A claim that names an instrument but whose context, package, or use-set is incomplete fails the well-formedness step of Section 6 before the classifier of Section 11 is reached.

\subsection{Visibility and Suppression}

Each instrument \(I\) partitions the content universe according to what \(I\) adjudicates and how it adjudicates it. The partition has four classes; this subsection fixes the partition map, the two named classes (visible and suppressed), and the schema of admissible visibility bridges. No theorem about acceptance or overreading is asserted here; the formal no-overreading theorem and the bridge-aware acceptance rule are stated in Section 11.

Each instrument \(I\) carries a finite visibility map

\[
\mathsf{vis}_I:\;\mathsf{Cont}\;\to\;\{\,\texttt{visible},\;\texttt{suppressed},\;\texttt{outside},\;\texttt{unknown}\,\},
\]

inducing a finite four-way partition of the content universe of Subsection 3.2:

\[
\mathsf{Cont}\;=\;\mathsf{Visible}_I\;\sqcup\;\mathsf{Suppressed}_I\;\sqcup\;\mathsf{Outside}_I\;\sqcup\;\mathsf{Unknown}_I.
\]

The four classes correspond, respectively, to content \(I\) may use directly, content within \(I\)'s scope that \(I\) may not use directly, content outside \(I\)'s declared scope, and content within \(I\)'s scope whose visibility verdict has not yet been resolved. The classes are part of \(I\)'s record (Subsection 3.4.1); changing the partition is changing the instrument.

\subsubsection{Visible Content}

The visible class under \(I\) is

\[
\mathsf{Visible}_I\;=\;\{\,x\in\mathsf{Cont}\;:\;\mathsf{vis}_I(x)=\texttt{visible}\,\}.
\]

A claim whose use-set \(\mathsf{Use}(\varphi)\) lies entirely in \(\mathsf{Visible}_I\), and whose witness, update, profile, threshold, audit, and source records are themselves in \(\mathsf{Visible}_I\), is a visible claim under \(I\). The classifier of Section 11 may consult visible content directly. Visibility under \(I\) is not a property of \(x\) alone: the same record can be visible under one instrument and suppressed, outside, or unknown under another.

\subsubsection{Suppressed Content}

The suppressed class under \(I\) is

\[
\mathsf{Suppressed}_I\;=\;\{\,x\in\mathsf{Cont}\;:\;\mathsf{vis}_I(x)=\texttt{suppressed}\,\}.
\]

Suppression is a record-level relation between \(x\) and \(I\), not a truth verdict on \(x\):

\[
x\in\mathsf{Suppressed}_I\;\;\not\Rightarrow\;\;x\;\text{is false}.
\]

Suppressed content may be exactly correct under another instrument or in an external description. The relation \(x\in\mathsf{Suppressed}_I\) asserts only that \(I\) does not have admissible visibility into \(x\); under \(I\), a claim that depends on \(x\) without an admissible bridge over-reads \(I\) (in the informal sense fixed in Subsection 2.3.3).

The remaining two classes are recorded for completeness:

\[
\mathsf{Outside}_I\;=\;\{\,x\in\mathsf{Cont}\;:\;\mathsf{vis}_I(x)=\texttt{outside}\,\},\qquad \mathsf{Unknown}_I\;=\;\{\,x\in\mathsf{Cont}\;:\;\mathsf{vis}_I(x)=\texttt{unknown}\,\}.
\]

A claim whose use-set meets \(\mathsf{Outside}_I\) has status \(\texttt{outside\_scope}\) under \(I\); a claim whose use-set meets \(\mathsf{Unknown}_I\) requires the later classifier to consult the corresponding unknown-content rule. No acceptance status for unknown content is asserted in this subsection.

\subsubsection{Visibility Bridges}

Suppressed content can become admissibly usable by a claim only via a recorded visibility bridge. A visibility bridge under \(I\) is a finite typed record

\[
L\;=\;(\,\mathrm{id}_L,\;c,\;c',\;\mathsf{BridgeType}_L,\;\mathsf{Guarantee}_L,\;\mathsf{Strength}_L,\;A_L,\;V_L,\;\mathcal N_L\,),
\]

with \(c\in\mathsf{Suppressed}_I\), \(c'\in\mathsf{Visible}_I\), \(\mathsf{BridgeType}_L\) drawn from the bridge inventory of \(\mathsf{BridgePolicy}_I\) (Subsection 3.4.1), and \(\mathsf{Guarantee}_L\), \(\mathsf{Strength}_L\), \(A_L\), \(V_L\), \(\mathcal N_L\) the finite guarantee, strength, audit, visibility, and nonclaim records that the bridge carries. A bridge is admissible under \(I\) when it has been registered in \(\mathsf{BridgePolicy}_I\) and when \(A_L\), \(V_L\), \(\mathcal N_L\) pass the gates of Section 10 with the strength required by the surrounding claim.

Let \(\mathcal L\) denote a finite set of admissible visibility bridges under \(I\). The schema-level use of \(\mathcal L\) is to record which suppressed content has been brought into admissible reach for which claims; the formal acceptance rule, including the constraint that every suppressed dependency of an accepted claim is bridged by some \(L\in\mathcal L\), is stated and proved in Section 11. This subsection only fixes the schema; no acceptance or no-overreading theorem is asserted here.

The witness and update records introduced in Subsection 3.7 are content records in this sense: when a later directed-cell or pair-observable judgment uses a witness \(W_i\) or update \(U_i\), that record has a visibility status under the named instrument and may require a visibility bridge if it is suppressed.

\subsection{Profiles}

A profile is the finite typed record attached to a directed-cell, pair-observable, promotion, dependency, claim, role-channel, or model-realization judgment that fixes the per-judgment scale, regime, bridge type, instrument mode, and level data. Profiles are part of every nontrivial judgment in Section 4 and beyond; their finite admissibility schema and their dependence on judgment status are recorded here. No theorem is asserted in this subsection; the no-total-six-symbol-algebra theorem of Section 6 uses the profile schema, but is not proved here.

\subsubsection{Profile Grammar}

A profile is a finite typed record

\[
\lambda\;=\;(\,\ell_i,\;\ell_j,\;\mathsf{Scale}_\lambda,\;\mathsf{Regime}_\lambda,\;\mathsf{BridgeType}_\lambda,\;\mathsf{InstrumentMode}_\lambda,\;\mathsf{Locality}_\lambda\,)
\]

with the following finite components.

\begin{itemize}
\tightlist
\item
  \(\ell_i,\ell_j\in\mathsf{Lev}\) are the actor-side and informant-side level tags. For a directed cell \(P_i\leftarrow P_j\), \(\ell_i\) is the level on which the actor update \(U_i\) lives and \(\ell_j\) is the level on which the informant witness \(W_j\) lives. For a judgment that is not directed (e.g., a pair observable or a claim), the two tags coincide or are determined by the judgment family's convention.
\item
  \(\mathsf{Scale}_\lambda\) is a finite scale label, distinguishing the scale of the carriers, partitions, kernels, or graphs the judgment depends on (for example, microscale or macroscale on a stochastic host, or fine or coarse on a finite carrier).
\item
  \(\mathsf{Regime}_\lambda\) is a finite regime label, distinguishing dynamical, structural, and audit regimes; specific regimes used by the paper are named where the relevant judgment is introduced.
\item
  \(\mathsf{BridgeType}_\lambda\) is drawn from the finite bridge inventory of \(\mathsf{BridgePolicy}_I\) (Subsection 3.4.1): visibility, level, instrument-transfer, drive, or model-realization. The profile records which bridge type the surrounding judgment uses.
\item
  \(\mathsf{InstrumentMode}_\lambda\) is a finite instrument-mode label, distinguishing the surrounding instrument's default, audit-only, transfer, and rotating-audit modes.
\item
  \(\mathsf{Locality}_\lambda\) is a finite locality label, distinguishing local, semi-global, and global use of the judgment data on a cover or staging.
\end{itemize}

Specific named profiles used elsewhere in the paper include \(\lambda_{\mathrm{cell}}\), \(\lambda_{\mathrm{prom}}\), \(\lambda_{\mathrm{AI}}\), \(\lambda_{\mathrm{PICA}}\), and \(\lambda_{\mathrm{Cantor}}\). Each is a finite record with the components above; their values are fixed where the corresponding judgment is introduced and are reused thereafter without redefinition.

\subsubsection{Profile Admissibility}

A profile \(\lambda\) is admissible under instrument \(I\) and theory package \(\mathcal T\), written

\[
\mathsf{AdmProfile}(\lambda,\,I,\,\mathcal T),
\]

when every component of \(\lambda\) is finite and consistent with \(I\), \(\mathcal T\), and the surrounding judgment family. The required checks are:

\begin{itemize}
\tightlist
\item
  \(\ell_i,\ell_j\in\mathsf{Lev}\), with cross-level \((\ell_i\neq\ell_j)\) profiles requiring a level bridge of one of the four types from Subsection 2.2.3 (translation, lift, forgetting, reflection);
\item
  \(\mathsf{BridgeType}_\lambda\) admitted by \(\mathsf{BridgePolicy}_I\);
\item
  \(\mathsf{InstrumentMode}_\lambda\) matching \(I\)'s declared mode for the surrounding judgment;
\item
  \(\mathsf{Scale}_\lambda\), \(\mathsf{Regime}_\lambda\), and \(\mathsf{Locality}_\lambda\) each consistent with the host of \(\mathcal T\) and with the specific data of the surrounding judgment.
\end{itemize}

A profile that fails any of these conditions is not admissible. When the failure is that the candidate judgment is not covered by the declared profile, host, or instrument, the later judgment classifier records the candidate as outside scope rather than accepted. Profile admissibility is a finite check on finite records; like package admissibility (Subsection 3.3.3) and instrument admissibility (Subsection 3.4.1), it feeds the well-formedness classifier of Section 6 rather than producing its own theorem here.

Profile mutation is not allowed silently. A judgment whose statement carries one profile and whose proof or downstream use depends on a different profile requires an explicit profile-update record, itself a typed entry in \(\mathsf{Cont}\) subject to the visibility, source, and audit discipline of Subsections 3.5, 4.8, and 10.

\subsubsection{Profile Dependence}

The status assigned to a directed-cell, pair-observable, promotion, dependency, or claim record depends on the profile attached to the judgment. The same primitive labels and the same witness, update, level, visibility, threshold, audit, and defect data --- when attached to a different profile --- can produce a different status. Concretely:

\begin{itemize}
\tightlist
\item
  a directed cell \(P_i\leftarrow P_j\) with profile \(\lambda\) may receive directed-cell status \(\texttt{action}\) from the classifier of Subsection 5.6.2, while the same primitive pair under a different profile \(\lambda'\) may receive \(\texttt{outside\_scope}\), \(\texttt{below\_threshold}\), or \(\texttt{collapsed}\);
\item
  a pair-observable judgment on the same primitive pair, with the same source data, may be classified \(\texttt{real}\) under one profile and \(\texttt{blocked}\) under another;
\item
  a promotion bridge with the same source and target packages may be \(\texttt{strict}\) under one profile and \(\texttt{non\_strict}\) under another.
\end{itemize}

This is a schema-grade observation about which fields the later classifiers must inspect: status is not a function of primitive labels alone, because the classifier's input includes the profile (along with witness, update, threshold, audit, visibility, defect, and source records). The no-total-six-symbol-algebra theorem of Section 6 is the theorem-grade form of this observation, and it is not proved here.

\subsection{Witness and Update Families}

\subsubsection{Witness Map}

The witness map of the calculus is the finite typed family

\[
\mathsf{Wit}\colon\mathbb P\;\to\;\mathrm{Type},\qquad P_i\;\mapsto\;\mathsf{Wit}(P_i),
\]

assigning to each primitive label a finite collection of admissible witness types. Each \(\mathsf{Wit}(P_i)\) is a finite family of finite typed records: for any host tag drawn from \(\mathsf{Host}\), the witness types specialize to finite records on that host. A witness \(W_i\in\mathsf{Wit}(P_i)\) is the informant-side data of a directed cell, or one side of a pair-observable record, in which \(P_i\) supplies the witness family; its host, level, instrument, and profile tags are determined by the surrounding judgment.

\subsubsection{Update Map}

The update map of the calculus is the finite typed family

\[
\mathsf{Upd}\colon\mathbb P\;\to\;\mathrm{Type},\qquad P_i\;\mapsto\;\mathsf{Upd}(P_i),
\]

assigning to each primitive label a finite collection of admissible update types. Each \(\mathsf{Upd}(P_i)\) is a finite family of finite typed records describing the effect that the primitive \(P_i\) produces when it occupies the actor role in a directed cell or promotion bridge. As with witnesses, every update record carries host, level, instrument, and profile tags consistent with the surrounding judgment.

\subsubsection{Witness/Update Constraint}

A directed cell with informant primitive \(P_j\) and actor primitive \(P_i\) is admissible only when its witness and update fields satisfy the typed actor/informant rule

\[
W_j\in\mathsf{Wit}(P_j)\quad\text{and}\quad U_i\in\mathsf{Upd}(P_i).
\]

That is, the informant-side data \(W_j\) must belong to the witness family of the informant's primitive, and the actor-side data \(U_i\) must belong to the update family of the actor's primitive. The constraint is checked at the well-formedness step of Section 6 together with the package, instrument, profile, level, visibility, threshold, audit, and defect conditions of the surrounding judgment.

The notation \(P_i\leftarrow P_j\) for a directed cell is a typed bridge judgment that records this rule. It is not an expression of a binary product on \(\mathbb P\), and a statement that uses the notation but violates the witness/update rule is a malformed cell, not a primitive product.

\subsubsection{Primitive Witness Table}

The primitive witness families used in this paper are listed below. Each entry is the finite witness family of the corresponding primitive: the entry names the kinds of records that count as witnesses, and each named record type is itself a finite typed record on the relevant host.

\begin{longtable}[]{@{}ll@{}}
\toprule
\begin{minipage}[b]{0.47\columnwidth}\raggedright
Primitive\strut
\end{minipage} & \begin{minipage}[b]{0.47\columnwidth}\raggedright
\(\mathsf{Wit}(P_i)\)\strut
\end{minipage}\tabularnewline
\midrule
\endhead
\begin{minipage}[t]{0.47\columnwidth}\raggedright
\(P_1\)\strut
\end{minipage} & \begin{minipage}[t]{0.47\columnwidth}\raggedright
descent success/failure records, rewrite obstruction records, closure-deficit records\strut
\end{minipage}\tabularnewline
\begin{minipage}[t]{0.47\columnwidth}\raggedright
\(P_2\)\strut
\end{minipage} & \begin{minipage}[t]{0.47\columnwidth}\raggedright
representability and nonrepresentability records, gate decisions, support and capacity records\strut
\end{minipage}\tabularnewline
\begin{minipage}[t]{0.47\columnwidth}\raggedright
\(P_3\)\strut
\end{minipage} & \begin{minipage}[t]{0.47\columnwidth}\raggedright
route-mismatch records, holonomy records, commutator records, critical-pair records, pasting-mismatch records\strut
\end{minipage}\tabularnewline
\begin{minipage}[t]{0.47\columnwidth}\raggedright
\(P_4\)\strut
\end{minipage} & \begin{minipage}[t]{0.47\columnwidth}\raggedright
stage records, refinement records, filtration records, cover records, local/global obstruction records\strut
\end{minipage}\tabularnewline
\begin{minipage}[t]{0.47\columnwidth}\raggedright
\(P_5\)\strut
\end{minipage} & \begin{minipage}[t]{0.47\columnwidth}\raggedright
package map records, quotient records, completion records, idempotence records, fixed-point records\strut
\end{minipage}\tabularnewline
\begin{minipage}[t]{0.47\columnwidth}\raggedright
\(P_6\)\strut
\end{minipage} & \begin{minipage}[t]{0.47\columnwidth}\raggedright
audit records, provenance records, source-of-truth records, threshold-check records, nonclaim records, drive certificates\strut
\end{minipage}\tabularnewline
\bottomrule
\end{longtable}

The drive-certificate entry under \(\mathsf{Wit}(P_6)\) is reserved for \(\mathrm{P6}_{\mathrm{drive}}\) data, which is the drive face of \(P_6\) rather than the generic audit face. A drive certificate is admissible only when the surrounding host supports it, and Section 8 lists the finite cycle-supported and gating-suppression cases in which the certificate may or may not be supported. The presence of generic audit records under \(\mathsf{Wit}(P_6)\) does not by itself supply a drive certificate.

\subsubsection{Primitive Update Table}

The primitive update families used in this paper are listed below.

\begin{longtable}[]{@{}ll@{}}
\toprule
\begin{minipage}[b]{0.47\columnwidth}\raggedright
Primitive\strut
\end{minipage} & \begin{minipage}[b]{0.47\columnwidth}\raggedright
\(\mathsf{Upd}(P_i)\)\strut
\end{minipage}\tabularnewline
\midrule
\endhead
\begin{minipage}[t]{0.47\columnwidth}\raggedright
\(P_1\)\strut
\end{minipage} & \begin{minipage}[t]{0.47\columnwidth}\raggedright
mark descent failure, install descended update, rewrite, repair\strut
\end{minipage}\tabularnewline
\begin{minipage}[t]{0.47\columnwidth}\raggedright
\(P_2\)\strut
\end{minipage} & \begin{minipage}[t]{0.47\columnwidth}\raggedright
mark representable or nonrepresentable, gate support, update admissibility record\strut
\end{minipage}\tabularnewline
\begin{minipage}[t]{0.47\columnwidth}\raggedright
\(P_3\)\strut
\end{minipage} & \begin{minipage}[t]{0.47\columnwidth}\raggedright
record route mismatch, holonomy, commutator, critical pair\strut
\end{minipage}\tabularnewline
\begin{minipage}[t]{0.47\columnwidth}\raggedright
\(P_4\)\strut
\end{minipage} & \begin{minipage}[t]{0.47\columnwidth}\raggedright
refine lens, advance stage, record local/global obstruction\strut
\end{minipage}\tabularnewline
\begin{minipage}[t]{0.47\columnwidth}\raggedright
\(P_5\)\strut
\end{minipage} & \begin{minipage}[t]{0.47\columnwidth}\raggedright
install package map, apply completion endomap, register fixed-point record, saturate\strut
\end{minipage}\tabularnewline
\begin{minipage}[t]{0.47\columnwidth}\raggedright
\(P_6\)\strut
\end{minipage} & \begin{minipage}[t]{0.47\columnwidth}\raggedright
add audit, check provenance, commit source, record nonclaim, certify drive or no-drive\strut
\end{minipage}\tabularnewline
\bottomrule
\end{longtable}

Each entry above is a finite typed update record on the host of the surrounding judgment. The certify-drive-or-no-drive entry under \(\mathsf{Upd}(P_6)\) is the update counterpart to the drive certificate witness type; it is admissible only when the surrounding host and instrument support drive judgments, and the drive or no-drive records of Section 8 are recorded through this update family.

The witness and update tables together fix the typed actor/informant data for directed cells and the related witness/update records used by pair observables, promotion bridges, and instrument-indexed claims. Section 4 introduces the judgment families that consume this data; Section 5 introduces the defect records that the data may produce; and Sections 6 through 11 use these records in the theorem-grade results.
 \section{\texorpdfstring{4. The Central Object \(\mathbf{BirdInt}\)}{4. The Central Object \textbackslash mathbf\{BirdInt\}}}

Section 3 declared the ambient domain \(\mathbf{BirdInt}^{\mathrm{aud}}_{\mathrm{fin}}\) as a finite tuple of finite typed components. This section names that tuple as the central object of the paper and defines the judgment families that the rest of the paper consumes: role-channel, directed-cell, pair-observable, promotion, dependency, and instrument-indexed claim judgments, together with the audit, source, and admissibility records that the calculus carries on every judgment line. The judgment families are introduced at schema-grade; their classifier rules and theorem-grade results appear in Sections 5 through 11.

\subsection{Central Object Definition}

\subsubsection{Object Form}

The central object of this paper is the finite typed structure

\[
\mathbf{BirdInt}^{\mathrm{aud}}_{\mathrm{fin}}
\;=\;
(\,\mathbb P,\;\mathsf{Lev},\;\mathsf{Host},\;\mathsf{Cont},\;\mathsf{Instr},\;\mathsf{Pkg},\;\mathsf{Prof},\;\mathsf{Wit},\;\mathsf{Upd},\;\mathsf{Def},\;\mathsf{Judg},\;\mathsf{Status},\;\mathsf{Gate},\;\mathsf{Dep},\;\mathsf{Audit},\;\mathsf{Vis},\;\mathsf{Source},\;\mathsf{Nonclaim},\;\mathsf{Real}\,),
\]

declared in Subsection 3.1.1, equipped with the judgment families introduced in Subsections 4.2 through 4.7 and the source, audit, and admissibility records of Subsections 4.8 and 4.9. The same notation \(\mathbf{BirdInt}^{\mathrm{aud}}_{\mathrm{fin}}\) refers to this object; no second name is introduced. It is distinct from the inherited Foundations II domain \(\mathsf{FATCD}\). The qualifier \(\mathrm{fin}\) records that every component is finite, and the qualifier \(\mathrm{aud}\) records that every accepted claim is interpreted under an explicit audit record.

The object is not a category, an algebra, or a closure operator. It is a finite typed record over a finite host inventory, carrying finite typed records in each component. The judgment families of this section operate on the components directly, and a theorem about \(\mathbf{BirdInt}^{\mathrm{aud}}_{\mathrm{fin}}\) is a theorem about a finite check on those records, under explicit instrument, package, profile, level, visibility, threshold, audit, and nonclaim assumptions.

\subsubsection{Internal Versus Model-Realization Components}

The components of \(\mathbf{BirdInt}^{\mathrm{aud}}_{\mathrm{fin}}\) split into two groups.

The \textbf{internal components} are \(\mathbb P\), \(\mathsf{Lev}\), \(\mathsf{Host}\), \(\mathsf{Cont}\), \(\mathsf{Instr}\), \(\mathsf{Pkg}\), \(\mathsf{Prof}\), \(\mathsf{Wit}\), \(\mathsf{Upd}\), \(\mathsf{Def}\), \(\mathsf{Judg}\), \(\mathsf{Status}\), \(\mathsf{Gate}\), \(\mathsf{Dep}\), \(\mathsf{Audit}\), \(\mathsf{Vis}\), \(\mathsf{Source}\), and \(\mathsf{Nonclaim}\). These are the components on which the calculus's syntax, classifier rules, defect-to-status rules, gate semantics, and instrument-indexed claim semantics operate. Theorems stated in Sections 5 through 11 are theorems about the internal components.

The \textbf{model-realization component} is \(\mathsf{Real}\). This is the family of finite model-realization modules through which scoped finite hosts (PICA, the audited Cantor shell, finite abstract interpretation, finite graph/cohomology fragments) realize named fragments of the internal calculus. Statements about the model-realization component are tagged model-only when they assert a realization; they do not transfer to internal universal claims, and the negative scope of Subsection 2.4.3 records this explicitly.

The split is enforced by the discipline of the paper: a theorem about an internal component does not silently use a model-realization fact, and a model-realization theorem is always tagged model-only and accompanied by the host, the realized fragment, the realization map, and the explicit nonclaims. Section 13 carries the model-realization material; the present section defines only the internal judgment families.

\subsection{Role-Channel Judgments}

\subsubsection{Judgment Form}

A role-channel judgment fixes the standing of a primitive role channel under an instrument and a theory package. Its form is

\[
\Gamma\,;\;\mathcal T\,;\;I\;\vdash\;\mathrm{Chan}_i:\rho_i,
\]

where \(\Gamma\) is a finite typed context, \(\mathcal T\) is an admissible theory package, \(I\) is an admissible instrument, \(\mathrm{Chan}_i\) is the role channel for primitive \(P_i\in\mathbb P\), and \(\rho_i\) is the role-channel status assigned to \(\mathrm{Chan}_i\) under \(\Gamma,\mathcal T,I\).

The notation \(\mathrm{Chan}_i\) records that the subject of the judgment is the role channel of \(P_i\), not the primitive label \(P_i\) itself. Channels are typed slots through which role-aligned witness, threshold, level, visibility, and audit data are recorded; the primitive label only names the channel.

\subsubsection{Channel Statuses}

The role-channel status set is inherited from Foundations II:

\[
\mathrm{RoleStatus}\;=\;\{\,\texttt{active\_projection},\;\texttt{below\_threshold},\;\texttt{collapsed\_boundary\_case},\;\texttt{absent\_with\_record},\;\texttt{inapplicable\_outside\_scope}\,\}.
\]

The five values have the inherited meanings:

\begin{itemize}
\tightlist
\item
  \(\texttt{active\_projection}\): an active derived role projection exists with the required role-aligned witness, threshold, level, audit, and relevant visibility/support data.
\item
  \(\texttt{below\_threshold}\): candidate role evidence is present, but the threshold data required for activation is absent or not met.
\item
  \(\texttt{collapsed\_boundary\_case}\): a candidate role feature cluster collapses under the relevant witness schema and cannot support an active projection.
\item
  \(\texttt{absent\_with\_record}\): no relevant raw features or witness data for the role are present, and the absence is explicitly recorded.
\item
  \(\texttt{inapplicable\_outside\_scope}\): the role-\(P_i\) content lies outside the scope covered by \(\mathcal T\) under \(I\).
\end{itemize}

Foundations III does not introduce additional role-channel statuses. The new status families introduced for directed cells, pair observables, promotions, claims, gates, and squares (Subsections 2.3.2 and 5.1) extend the calculus; they do not replace or relabel the inherited five role-channel statuses.

\subsubsection{Channel/Activation Discipline}

A role channel is a typed slot. An active witness is a recorded value passing the threshold, level, witness, and audit conditions of the surrounding instrument and package. The two are distinct:

\begin{itemize}
\tightlist
\item
  a primitive may have an open role channel without an active witness, with status \(\texttt{absent\_with\_record}\) or \(\texttt{below\_threshold}\);
\item
  a primitive may have an active witness on a particular profile and yet be inapplicable on a different profile of the same package and instrument, with status \(\texttt{inapplicable\_outside\_scope}\) on the second profile;
\item
  a primitive may have status \(\texttt{active\_projection}\) on its role channel and yet not produce a directed cell with status \(\texttt{action}\), because directed-cell action additionally requires admissible actor/informant data, profile, visibility, threshold, audit, and defect records on the cell.
\end{itemize}

No theorem in this paper infers an active witness from the existence of a channel, and no theorem infers a directed-cell status from a role-channel status. The status-family separation theorem of Section 6 makes the second of these formal.

\subsection{Directed-Cell Judgments}

Directed-cell judgments are the central typed bridge judgments of \(\mathbf{BirdInt}^{\mathrm{aud}}_{\mathrm{fin}}\). A directed cell records that, under an admissible context, package, and instrument, an actor primitive consumes typed witness data from an informant primitive through a typed update, with the level/lens/interface, visibility, threshold, audit, and defect data of the surrounding judgment. The cell is the working unit of typed interaction; the no-total-six-symbol-algebra theorem of Section 6 records that cell status does not factor through the primitive pair.

\subsubsection{Judgment Form}

The directed-cell judgment form is

\[
\Gamma\,;\;\mathcal T\,;\;I\;\vdash\;(\,P_i\leftarrow P_j,\;W_j,\;U_i,\;L,\;V,\;\Theta,\;A,\;\delta\,)^{\lambda}\;:\;\sigma,
\]

where \(\Gamma\) is a finite typed context, \(\mathcal T\) is an admissible theory package, \(I\) is an admissible instrument, \(\lambda\in\mathsf{Prof}\) is an admissible profile, and \(\sigma\in\mathrm{CellStatus}\) is the directed-cell status assigned by the classifier of Subsection 5.6.2. The cell record itself is the finite typed tuple

\[
\mathsf{Cell}_{ij}^{\lambda}\;=\;(\,P_i\leftarrow P_j,\;W_j,\;U_i,\;L,\;V,\;\Theta,\;A,\;\delta\,),
\]

so the judgment may equivalently be written \(\Gamma;\mathcal T;I\vdash\mathsf{Cell}_{ij}^{\lambda}:\sigma\). The notation \(P_i\leftarrow P_j\) is a typed bridge designation, not a binary product on \(\mathbb P\); it records the actor and informant roles for the surrounding witness, update, and status data.

\subsubsection{Actor and Informant Roles}

In the cell \(P_i\leftarrow P_j\), the actor primitive is \(P_i\) and the informant primitive is \(P_j\). The actor receives the typed update \(U_i\in\mathsf{Upd}(P_i)\), and the informant supplies the typed witness \(W_j\in\mathsf{Wit}(P_j)\), with \(\mathsf{Wit}\) and \(\mathsf{Upd}\) as in Subsection 3.7. The arrow direction is part of the data: the cells \(P_i\leftarrow P_j\) and \(P_j\leftarrow P_i\) are distinct judgments with different witness/update typing requirements and may receive different statuses under the same package, instrument, and profile.

The actor/informant typing rule is

\[
W_j\in\mathsf{Wit}(P_j)\quad\text{and}\quad U_i\in\mathsf{Upd}(P_i),
\]

inherited from Subsection 3.7.3. A cell whose witness or update fails this rule is malformed and fails the well-formedness step of Section 6 before any classifier verdict is assigned.

\subsubsection{Cell Components}

The cell judgment uses the following finite typed records, each hosted and instrument-indexed as required by \(\mathcal T\), \(I\), and the profile \(\lambda\):

\begin{itemize}
\tightlist
\item
  \(W_j\in\mathsf{Wit}(P_j)\): the informant-side witness data.
\item
  \(U_i\in\mathsf{Upd}(P_i)\): the actor-side update data.
\item
  \(L\): the level/lens/interface record. \(L\) specifies the level tags on which the actor and informant operate, the lens through which the cell content is read, and the interface against which the cell is judged. \(L\) is consistent with the level data of \(\mathcal T\) and \(I\); cross-level cells require a level bridge inside \(L\).
\item
  \(V\): the visibility record under \(I\) for the fields of the cell, drawn from the visibility partition of Subsection 3.5. \(V\) specifies which of \(W_j\), \(U_i\), \(L\), \(\Theta\), \(A\), \(\delta\) are visible, suppressed, outside-scope, or unknown under \(I\), and lists any visibility bridges admitted under \(\mathsf{BridgePolicy}_I\) that bring suppressed fields into admissible reach.
\item
  \(\Theta\): the threshold record. \(\Theta\) specifies the exact, approximate, tolerance, rank, visibility, and source thresholds the cell must meet, drawn from the eleven-element threshold family of Subsection 5.2.2.
\item
  \(A\): the audit record of the cell, in the schema of Subsection 4.8.2. \(A\) is distinct from the inherited audit functional \(\mathcal A\) of \(\mathcal T\): \(A\) is the finite audit record, while compatibility with \(\mathcal A\) and \(\mathsf{AuditPolicy}_I\) is an additional typed condition when the package audit functional is relevant.
\item
  \(\delta\): the defect record family attached to the cell. Each entry of \(\delta\) is a typed defect record in the schema of Subsection 5.2.1, drawn from the primitive defect families of Subsection 5.3 and the visibility defects of Subsection 5.5.
\item
  \(\lambda\in\mathsf{Prof}\): the profile of the cell, fixing the actor and informant level tags, scale, regime, bridge type, instrument mode, and locality, in the schema of Subsection 3.6.1.
\item
  \(\sigma\in\mathrm{CellStatus}\): the directed-cell status assigned to the cell by the classifier. This is part of the judgment line, not an extra field of the cell tuple \(\mathsf{Cell}_{ij}^{\lambda}\).
\end{itemize}

A cell with the actor/informant typing rule satisfied and with finite, admissible \(L,V,\Theta,A,\delta,\lambda\) records is well-formed. The well-formedness step is a finite decidable check (Theorem 3, Subsection 6.3); the resulting status \(\sigma\) is then determined by the classifier of Subsection 5.6.2 from the \(L,V,\Theta,A,\delta\) data and the profile.

\subsubsection{Directed-Cell Statuses}

The directed-cell status set is

\[
\mathrm{CellStatus}\;=\;\{\,\texttt{action},\;\texttt{implicit},\;\texttt{trivial},\;\texttt{blocked},\;\texttt{undefined\_circular},\;\texttt{below\_threshold},\;\texttt{collapsed},\;\texttt{absent},\;\texttt{outside\_scope}\,\}.
\]

The classifier assigns at most one of these to a covered well-formed cell (Theorem 4, Subsection 6.4). The classifier priority order, frozen for the rest of the paper, is

\[
\texttt{outside\_scope}\;\succ\;\texttt{absent}\;\succ\;\texttt{blocked}\;\succ\;\texttt{undefined\_circular}\;\succ\;\texttt{collapsed}\;\succ\;\texttt{below\_threshold}\;\succ\;\texttt{trivial}\;\succ\;\texttt{implicit}\;\succ\;\texttt{action}.
\]

Under this priority, a cell whose data simultaneously matches more than one status condition is assigned the highest-priority match; the per-status conditions are specified in Subsection 5.6.2 and the uniqueness theorem is Theorem 4. The priority ordering reflects the working discipline of the calculus: out-of-scope and absence verdicts are recorded before any further check on the cell's content; blocking, circularity, and collapse are recorded before threshold or activity verdicts; and the \(\texttt{action}\) verdict is reached only after every higher-priority condition has been ruled out.

\subsection{Pair-Observable Judgments}

Pair-observable judgments record the standing of an unordered pair of primitives as a joint observable, distinct from the directed-cell judgments of Subsection 4.3. The pair-observable family has its own status set, its own classifier inputs (notably a source-of-truth record), and its own admissibility conditions. Directed-cell action does not by itself imply pair-observable realness; the schema-level distinction is recorded here, and the formal status-family separation appears in Theorem 5 of Section 6.

\subsubsection{Judgment Form}

The pair-observable judgment form is

\[
\Gamma\,;\;\mathcal T\,;\;I\;\vdash\;\mathsf{PairObs}_{\{i,j\}}^{\mu}\;:\;\omega,
\]

where \(\{i,j\}\subseteq\{1,\ldots,6\}\) is an unordered pair of primitive indices, \(\mu\in\mathsf{Prof}\) is a profile drawn from the family of profiles for pair-observable judgments, and \(\omega\in\mathrm{PairStatus}\) is the pair-observable status assigned to the record by the classifier of Subsection 5.6.3. The unordered pair \(\{i,j\}\) equals \(\{j,i\}\); the judgment is the same for either ordering.

A pair-observable record is a finite typed tuple recording, in addition to \(\{i,j\}\) and \(\mu\), the joint witness data, the source-of-truth record, the branch-compatibility record, the visibility record, the threshold record, the audit record, the defect record, and the nonclaim record that the surrounding instrument and package require. The full schema is registered in \(\mathsf{Cont}\) and consumed at the well-formedness step of Section 6.

\subsubsection{Pair Versus Directed Cell}

A pair observable is not a directed cell. The directed cell \(P_i\leftarrow P_j\) records typed actor/informant bridge data; the pair observable \(\mathsf{PairObs}_{\{i,j\}}^{\mu}\) records joint source-of-truth content for the unordered pair. The two are distinct judgment families with distinct status sets, distinct classifier inputs, and distinct admissibility schemas.

In particular,

\[
\Gamma\,;\;\mathcal T\,;\;I\;\vdash\;\mathsf{Cell}_{ij}^{\lambda}\;:\;\texttt{action}\;\;\not\Rightarrow\;\;\Gamma\,;\;\mathcal T\,;\;I\;\vdash\;\mathsf{PairObs}_{\{i,j\}}^{\mu}\;:\;\texttt{real}.
\]

A directed cell with status \(\texttt{action}\) records that the typed actor/informant bridge passes the well-formedness, threshold, visibility, and audit conditions of the cell record. A pair observable with status \(\texttt{real}\) additionally records that the unordered pair has a real source of truth, an admissible branch-compatibility record, and an admissible source policy under \(I\). The countermodel sheet \(\mathsf{CM}_4\) of Section 9 supplies a finite construction in which a directed cell is \(\texttt{action}\) and the corresponding pair observable is \(\texttt{blocked}\). Section 6 contains the theorem-grade form of this separation as part of Theorem 5; this subsection records only the schema-level distinction.

\subsubsection{Source-of-Truth Conditions}

Pair-observable acceptance depends on a source-of-truth record drawn from a finite typed family. The source classes used by this paper, drawn from the source-policy values of Subsection 4.8.1, are:

\begin{itemize}
\tightlist
\item
  \textbf{real source.} A finite source record, certified under \(\mathsf{SourcePolicy}_I\) with admissible audit and bridge data, that supports the unordered pair under profile \(\mu\).
\item
  \textbf{fallback source.} A finite source record certified as a designated fallback under \(\mathsf{SourcePolicy}_I\). A fallback source supports a pair observable at \(\texttt{real\_provisional}\) at most by default; it does not support \(\texttt{real}\) unless the profile explicitly upgrades the fallback.
\item
  \textbf{dirty source.} A finite source record whose audit or threshold checks fail. A dirty source does not support an accepted pair observable; the resulting status is \(\texttt{blocked}\).
\item
  \textbf{proxy source.} A finite source record obtained via a typed bridge from another source, admissible only when the bridge is recorded in \(\mathsf{BridgePolicy}_I\) and audited; a proxy source supports a pair observable at most at \(\texttt{real\_provisional}\).
\item
  \textbf{missing source.} No source record is supplied. The pair observable does not pass acceptance: the status is \(\texttt{blocked}\) when a source was required by \(\mathsf{SourcePolicy}_I\).
\end{itemize}

The pair-observable status set used by this paper is

\[
\mathrm{PairStatus}\;=\;\{\,\texttt{real},\;\texttt{real\_provisional},\;\texttt{real\_duplicated},\;\texttt{blocked}\,\}.
\]

The source-condition names above--fallback source, dirty source, proxy source, and missing source--are not additional pair-observable statuses. They are source-policy conditions that the classifier routes into one of the four values in \(\mathrm{PairStatus}\).

The classifier of Subsection 5.6.3 consumes the pair record together with its source-of-truth record, branch-compatibility record, visibility record, threshold record, audit record, and defect record, and assigns a value in \(\mathrm{PairStatus}\) to a covered well-formed pair-observable record. As with directed cells, no theorem in this paper infers a pair-observable status from a directed-cell status, a role-channel status, or a primitive pair alone; the schema discipline of this subsection together with the status-family separation theorem of Section 6 closes that gap.

\subsection{Promotion Judgments}

Promotion judgments record the standing of a finite promotion bridge from a source theory package to a target theory package under an instrument. They are the typed/statused machinery on which the strict-extension theorems of Section 7 and the promotion-gate-soundness theorem of Section 10 are built. Like directed-cell and pair-observable judgments, promotion judgments have their own status set and their own classifier inputs; in particular, the strict and non-strict refinements of an accepted promotion are properties of the bridge's object-map data and do not by themselves imply closure or drive on the target package.

\subsubsection{Judgment Form}

The promotion judgment form is

\[
\Gamma\,;\;\mathcal T_j\,;\;I\;\vdash\;\mathrm{Promote}(B_{j\to j+1})\;:\;\sigma_{\mathrm{promote}},
\]

where \(B_{j\to j+1}\) is the promotion bridge from source theory package \(\mathcal T_j\) to target theory package \(\mathcal T_{j+1}\), \(I\) is an admissible instrument with \(\mathsf{ClaimTypes}_I\) that includes promotion claims, and \(\sigma_{\mathrm{promote}}\in\mathrm{PromotionStatus}\) is the promotion status assigned by the classifier of Subsection 5.6.4.

\subsubsection{Promotion Bridge Object}

A promotion bridge is the finite typed tuple

\[
B_{j\to j+1}\;=\;(\,\mathcal T_j,\;\mathcal T_{j+1},\;\pi_j,\;\pi_{j+1},\;L,\;V,\;\Theta,\;A,\;\delta,\;\mathcal N\,),
\]

with the following components, all finite:

\begin{itemize}
\tightlist
\item
  \(\mathcal T_j\): the source theory package, admissible in the sense of Subsection 3.3.3.
\item
  \(\mathcal T_{j+1}\): the target theory package, admissible in the same sense.
\item
  \(\pi_j\): the old object map or interface, recording how content is read out of \(\mathcal T_j\).
\item
  \(\pi_{j+1}\): the new object map or interface, recording how content is read out of \(\mathcal T_{j+1}\).
\item
  \(L\): the level/lens/interface record of the bridge, fixing the levels at which \(\pi_j\) and \(\pi_{j+1}\) operate.
\item
  \(V\): the visibility record of the bridge under \(I\), in the schema of Subsection 3.5.
\item
  \(\Theta\): the threshold record of the bridge, drawn from the threshold family of Subsection 5.2.2.
\item
  \(A\): the audit record of the bridge, distinct from the audit functionals \(\mathcal A_j\), \(\mathcal A_{j+1}\) of the source and target packages.
\item
  \(\delta\): the defect record family attached to the bridge, including the strictness, no-smuggling, stability, descent, and audit defects defined in Subsections 5.3 and 5.4.
\item
  \(\mathcal N\): the nonclaim record attached to the bridge.
\end{itemize}

The expanded form of \(B_{j\to j+1}\) recorded in \(\mathsf{Cont}\) also carries a bridge identifier, the carrier \(S\) on which the object maps act, the codomains \(O_j\) and \(O_{j+1}\), the gate-results record, the host tag, and the bridge profile:

\[
\begin{aligned}
B_{j\to j+1}=\bigl(&\mathrm{id},\mathcal T_j,\mathcal T_{j+1},S,O_j,O_{j+1},\pi_j,\pi_{j+1},\\
&L,V,\Theta,A,\delta,\mathsf{GateResults},\mathcal N,\mathsf{HostTag},\lambda_{\mathrm{prom}}\bigr).
\end{aligned}
\]

A bridge is admissible when every component is finite, hosted, and consistent with the package and instrument data of the surrounding judgment, and when the gate record passes the gate semantics of Section 10.

A promotion bridge is \textbf{strict} when

\[
\pi_{j+1}\;\not\factor\;\pi_j,
\]

equivalently when there exist \(s,s'\in S\) with \(\pi_j(s)=\pi_j(s')\) and \(\pi_{j+1}(s)\neq\pi_{j+1}(s')\). A promotion bridge is \textbf{non-strict} when \(\pi_{j+1}\) factors through \(\pi_j\), that is, when there exists a finite map \(\phi:\mathrm{im}(\pi_j)\to O_{j+1}\) with \(\pi_{j+1}=\phi\,\pi_j\). Strictness is a property of the typed object-map data; it is not by itself a closure or drive claim, and the negative scope of Subsection 2.4.3 records the corresponding nonclaims (NC-10 and NC-11).

\subsubsection{Promotion Statuses}

The promotion status set is

\[
\mathrm{PromotionStatus}\;=\;\{\,\texttt{candidate},\;\texttt{accepted},\;\texttt{strict},\;\texttt{non\_strict},\;\texttt{failed\_descent},\;\texttt{failed\_stability},\;\texttt{failed\_audit},\;\texttt{failed\_no\_smuggling},\;\texttt{local\_only},\;\texttt{globally\_obstructed},\;\texttt{outside\_scope}\,\}.
\]

The values \(\texttt{strict}\) and \(\texttt{non\_strict}\) are accepted-status refinements: a strict promotion is an accepted promotion in which the strictness condition \(\pi_{j+1}\not\factor\pi_j\) holds, and a non-strict promotion is an accepted promotion in which \(\pi_{j+1}\) factors through \(\pi_j\). The classifier of Subsection 5.6.4 produces an accepted, strict, or non-strict verdict only when every required core gate in \(\mathsf{GateResults}(B)\) passes; the failure-tagged values record specific gate failures (descent, stability, audit, no-smuggling), and \(\texttt{local\_only}\) and \(\texttt{globally\_obstructed}\) record the optional local-to-global gate's verdict. The verdict \(\texttt{candidate}\) records a bridge whose required core gates have not yet been resolved, typically because at least one core gate has status \(\texttt{not\_checked}\).

A promotion bridge with status \(\texttt{strict}\) does not by that fact alone imply macro closure, descended dynamics, or drive: NC-10 and NC-11 of Subsection 2.4.3 record the corresponding strictness nonclaims, and the later no-go results are stated against these would-be implications.

\subsection{Dependency and No-Go Judgments}

Dependency judgments record typed relationships between content classes under an instrument. They are the working format for the no-go and bridge-required claims that appear throughout the paper, including the route-mismatch-requires-bridge dependency that keeps the drive face \(\mathrm{P6}_{\mathrm{drive}}\) separate from generic \(P_3\) holonomy. Like the other judgment families, dependencies have their own admissibility schema; theorem-grade no-go statements supported by countermodels are recorded in Sections 7 through 10 and Appendix C.

\subsubsection{Judgment Form}

The dependency judgment form is

\[
\Gamma\,;\;\mathcal C\,;\;I\;\vdash\;S\;\mathrel{\mathsf{dep}_r}\;X\;[\,H,\;A,\;V,\;\Theta\,],
\]

where \(\Gamma\) is a finite typed context, \(\mathcal C\) is a finite content context (typically a theory package, a model-realization module, or a content fragment of \(\mathsf{Cont}\)), \(I\) is an admissible instrument, \(S\) and \(X\) are content classes (primitive labels, primitive role channels, named structural records, drive certificates, gating data, or any record drawn from \(\mathsf{Cont}\)), \(r\) is the relation type, and \(H,A,V,\Theta\) are the host, audit, visibility, and threshold records under which the dependency is asserted.

The relation types admitted by this paper are

\[
\{\,\texttt{suffices},\;\texttt{insufficient},\;\texttt{destroys},\;\texttt{blocks},\;\texttt{enables},\;\texttt{equivalent},\;\texttt{requires\_bridge},\;\texttt{outside\_scope}\,\}.
\]

The relation type fixes the polarity and form of the dependency; the host, audit, visibility, and threshold records fix the scope under which the dependency is asserted.

\subsubsection{Dependency Claims}

The dependency judgment family supports four categories of statement, each appearing throughout the paper.

\textbf{Positive dependency.} A statement of the form \(S\mathrel{\mathsf{dep}_{\texttt{suffices}}}X\) or \(S\mathrel{\mathsf{dep}_{\texttt{enables}}}X\) records that, under the named scope, \(S\) supplies or enables \(X\). Positive dependencies in this paper are typically theorem-grade results stated under explicit host, instrument, and audit assumptions: the descent theorem of Section 7, for example, gives a positive dependency from a vanishing descent defect to the existence of a descended map.

\textbf{Blocked dependency.} A statement of the form \(S\mathrel{\mathsf{dep}_{\texttt{insufficient}}}X\), \(S\mathrel{\mathsf{dep}_{\texttt{blocks}}}X\), or \(S\mathrel{\mathsf{dep}_{\texttt{destroys}}}X\) records that \(S\) does not support \(X\) under the named scope, or actively obstructs it. Blocked dependencies are the natural form of the paper's no-go results: the gating-affinity-suppression theorem of Section 8, for instance, records a \(\texttt{destroys}\) dependency from \(P_2\)-style support deletion to cohomological drive on the resulting forest.

\textbf{Bridge-required dependency.} A statement of the form \(S\mathrel{\mathsf{dep}_{\texttt{requires\_bridge}}}X\) records that \(S\) does not by itself yield \(X\), but does so via a typed bridge admitted by \(\mathsf{BridgePolicy}_I\). The canonical instance in this paper is the route-mismatch dependency

\[
P_3\;\mathrel{\mathsf{dep}_{\texttt{requires\_bridge}}}\;\mathrm{P6}_{\mathrm{drive}},
\]

recording that route mismatch or holonomy alone does not establish a drive certificate; an explicit drive bridge is required for a positive drive claim.

\textbf{Countermodel-supported no-go.} A statement of the form \(S\mathrel{\mathsf{dep}_{\texttt{insufficient}}}X\) or \(S\mathrel{\mathsf{dep}_{\texttt{blocks}}}X\) supported by a finite explicit countermodel construction. Sections 7, 8, 9, and 10 use this category for their no-go and separation results; the countermodel atlas of Section 9 and Appendix C records the constructions in full.

A dependency judgment is admissible when every component is finite and consistent with the surrounding instrument and content context, and when its relation label \(r\) is one of the eight allowed values above. The label \(\texttt{outside\_scope}\) records the case in which the surrounding instrument does not adjudicate the dependency.

\subsection{Structural Downward Influence Versus Top-Down Causal Channels}

The calculus records several forms of downward influence across packages, profiles, levels, and hosts. A macro package may select a representative substrate record by a completion or lift; a macro admissibility condition may exclude substrate states or histories; a phase, profile, or protocol record may index which substrate update rule is active. These are structural downward influences: they are paths, restrictions, or context-indexing records from a macro record to a substrate record.

They are not, by themselves, top-down causal channels. In this paper the word ``channel'' is reserved for intervention-respecting difference-making records. A completion, lift, representative map, feasibility restriction, or context index is a structural path unless the top-down channel gates pass. Such a record must exhibit an admissible finite family of macro interventions and an audited comparison showing that changing the macro intervention changes a substrate-future response under matched controls.

A macro record is a finite level-tagged or lens-related record \(M\) in the current package: for example a package object, profile parameter, phase, protocol, admissibility gate, policy, abstract value, or promoted object-map value. A substrate record \(X\) is a lower-level or finer-lens future observable, state predicate, update trace, or output record in the same host or a connected host. The macro/substrate relation is therefore relative to a package, lens, level profile, host, and bridge record, rather than an external micro/macro ontology.

A top-down channel record is a finite tuple

\[
\mathsf{TopDownChannelRecord}_{M\to X}^{\lambda}
=
(M,X,\mathsf{Interv}_M,\mathsf{Ctrl},\mathsf{Resp},D,\Theta,L,V,A,
\delta_{\mathrm{td}},\mathcal N).
\]

Here \(M\) is a macro record relative to \((\mathcal T,I,\lambda)\), \(X\) is a substrate-future record, \(\mathsf{Interv}_M\) is a finite family of admissible macro interventions, \(\mathsf{Ctrl}\) records matched controls, \(\mathsf{Resp}\) records the substrate response under each intervention, \(D\) is the host-supplied effect comparator, \(\Theta\) is the effect threshold, \(L,V,A\) are bridge, visibility, and audit records, \(\delta_{\mathrm{td}}\) is the measured difference-making defect or certificate, and \(\mathcal N\) records the required nonclaims.

The calculus does not fix a universal probability metric for top-down channels. The host supplies the comparator \(D_{\mathsf H}\) and threshold \(\Theta_{\mathrm{td}}\). A probabilistic host may instantiate the comparator by total variation, KL divergence, or a channel-capacity quantity; a deterministic host may instantiate it by finite output inequality. The present calculus requires only that the comparator, threshold, source records, and audit be finite and visible or bridged.

The claim

\[
\Gamma;\mathcal T;I
\vdash
\mathsf{TopDownChannel}(R):\texttt{accepted}
\]

is accepted only if the following gates pass: macro-record well-formedness, admissible intervention, matched controls, feasibility, source-of-truth, visibility, audit, no-smuggling, and effect threshold. In particular, the record must contain at least two admissible macro interventions and a host-supplied comparison satisfying

\[
\Theta_{\mathrm{td}}(\delta_{\mathrm{td}})=\texttt{pass}.
\]

Thus accepted top-down channel claims carry a positive structural dependency,

\[
\mathsf{TopDownChannel}(R)
\;\mathrel{\mathsf{dep}_{\texttt{suffices}}}\;
\mathsf{StructDown}(M,X),
\]

but structural downward influence does not imply top-down causal-channel acceptance:

\[
\mathsf{StructDown}(M,X)
\;\mathrel{\mathsf{dep}_{\texttt{insufficient}}}\;
\mathsf{TopDownChannel}(R):\texttt{accepted}.
\]

A completion, lift, feasibility gate, or context index is a structural path. It becomes a top-down causal channel only when the intervention and matched-response gates pass. The required top-down gates are

\[
\mathsf{TDGate}(R)
=
(G_{\mathrm{macro}},G_{\mathrm{interv}},G_{\mathrm{matched}},G_{\mathrm{feas}},
G_{\mathrm{source}},G_{\mathrm{vis}},G_{\mathrm{audit}},G_{\mathrm{nosmuggle}},
G_{\mathrm{effect}}).
\]

The macro gate requires that \(M\) be a well-defined package/profile record rather than a post hoc label. The intervention gate requires at least two admissible macro choices. The matched-control gate fixes background, initial condition class, substrate support, package/lens, level profile, and source-policy conditions. The feasibility gate requires the interventions to be feasible in the declared layer. The source, visibility, audit, and no-smuggling gates apply the existing source-of-truth, no-overreading, audit, and anti-smuggling discipline to every record used by the top-down claim. The effect gate is the host-specific threshold check \(\Theta_{\mathrm{td}}\).

At schema grade, top-down channel gate soundness states that

\[
\Gamma;\mathcal T;I
\vdash
\mathsf{TopDownChannel}(R):\texttt{accepted}
\quad\Longrightarrow\quad
\bigwedge_{G\in\mathsf{ReqTDGates}(R)}\mathsf{GatePass}(R,G).
\]

The record is not assigned to a single primitive. Its macro package or representative path may use \(P_5\), its feasible intervention family may use \(P_2\), its layer/profile bridge may use \(P_4\), its response or transition-law comparison may use \(P_1\) when an induced operator is at issue, and its provenance, source, threshold, visibility, and no-smuggling checks are \(P_6\)-audited. \(P_3\) route data may enter the record when protocol or order matters, but it is not by itself a channel certificate.

\subsection{Instrument-Indexed Claims}

Instrument-indexed claims are the most general judgment family in the calculus. Every nontrivial assertion in the paper --- directed-cell, pair-observable, promotion, dependency, theorem, model-realization, or countermodel --- is interpreted as an instrument-indexed claim under an admissible instrument. This subsection fixes the claim form at the level of generality required throughout the paper, and records the per-claim status set used by the classifier of Subsection 5.6.5; the classifier rules and the no-overreading and same-level self-audit theorems are stated in Section 11.

\subsubsection{Claim Form}

The general instrument-indexed claim form is

\[
\Gamma\,;\;\mathcal T\,;\;I\;\vdash\;\varphi\;:\;\chi,
\]

introduced at schema-grade in Subsection 3.4.2 and reused throughout the paper. Here \(\Gamma\) is a finite typed context, \(\mathcal T\) is an admissible theory package, \(I\) is an admissible instrument, \(\varphi\) is a claim with a finite use-set \(\mathsf{Use}(\varphi)\subseteq\mathsf{Cont}\), and \(\chi\in\mathrm{ClaimStatus}\) is the claim status assigned by the classifier of Subsection 5.6.5.

A claim record under \(I\) is the finite typed tuple consisting of the claim content \(\varphi\), its declared claim type drawn from \(\mathsf{ClaimTypes}_I\), the target package \(\mathcal T\), the surrounding context \(\Gamma\), the source records consulted, the threshold and audit records, the visibility bridges \(\mathcal L\) used to bring suppressed content into admissible reach, and the nonclaim record attached to the claim. The shorthand

\[
I\;\vdash\;\varphi(\mathcal T)
\]

abbreviates \(\Gamma;\mathcal T;I\vdash\varphi:\texttt{accepted}\) and is used in cross-references where the surrounding context is fixed; the shorthand never abbreviates instrument-free truth, and a statement of the form \(I\vdash\varphi(\mathcal T)\) does not transfer to a different instrument \(I'\) except through an explicit instrument-transfer bridge.

\subsubsection{Claim Strength}

The claim status set used for instrument-indexed claims is

\[
\mathrm{ClaimStatus}\;=\;\{\,\texttt{accepted},\;\texttt{rejected},\;\texttt{provisional},\;\texttt{blocked},\;\texttt{outside\_scope},\;\texttt{absent\_with\_record},\;\texttt{below\_threshold},\;\texttt{failed\_audit},\;\texttt{undefined\_circular}\,\}.
\]

The classifier of Subsection 5.6.5 assigns one of these to a covered well-formed claim record, in the priority order specified in Subsection 11.2.1. The values record specific verdicts:

\begin{itemize}
\tightlist
\item
  \(\texttt{accepted}\): the claim is accepted by \(I\) under \(\Gamma,\mathcal T\) with admissible visibility, threshold, source, audit, and nonclaim data.
\item
  \(\texttt{rejected}\): the instrument has an admissible refuting witness for the claim.
\item
  \(\texttt{provisional}\): partial evidence exists and the profile permits provisional status.
\item
  \(\texttt{blocked}\): the claim is in scope but unadjudicable because required evidence, a bridge, a source, or visibility support is missing.
\item
  \(\texttt{outside\_scope}\): the claim's use-set is not contained in \(\mathsf{Scope}_I\), or a host or level mismatch removes the claim from \(I\)'s adjudication.
\item
  \(\texttt{absent\_with\_record}\): required evidence is absent and the absence is audited.
\item
  \(\texttt{below\_threshold}\): a threshold check required by \(\Theta\) or by \(\mathsf{ThresholdPolicy}_I\) has failed.
\item
  \(\texttt{failed\_audit}\): an audit policy check has failed.
\item
  \(\texttt{undefined\_circular}\): the claim depends on an unstratified circular support or audit loop, including the same-level self-audit case of Subsection 11.5.
\end{itemize}

The paper-facing claim strengths of Subsection 2.3.1 (theorem-grade, proposition-grade, schema-grade, model-only, countermodel-grade, future-work, dropped) and the per-claim instrument statuses above are different layers. Claim strength records how a numbered statement is presented in the manuscript; claim status records how an instrument classifies a finite claim record. A theorem-grade or proposition-grade statement, under the hypotheses recorded in its statement, is represented by an accepted instrument-indexed claim; a schema-grade record is a definition or rule that is not adjudicated as a claim unless separately wrapped as one; a model-only statement is adjudicated only inside its declared model host; a countermodel-grade statement is adjudicated with its finite countermodel record attached; future-work and dropped statements are not asserted as accepted claims in this paper.

\subsection{Source, Provenance, and Audit Records}

Source and audit records are the working format for the calculus's reproducibility and scope discipline. They are not themselves proofs, and admissibility under an instrument is not the same as instrument-free truth. This subsection records the source-record and audit-record schemas used throughout the paper, with explicit separation between the audit record \(A\) and the inherited audit functional \(\mathcal A\) of Foundations I.

\subsubsection{Source Records}

A source record is a finite typed tuple

\[
\mathrm{src}\;=\;(\,\mathrm{id}_{\mathrm{src}},\;\mathrm{type}_{\mathrm{src}},\;\mathrm{trace}_{\mathrm{src}},\;\mathrm{valid}_{\mathrm{src}}\,),
\]

with the following finite components.

\begin{itemize}
\item
  \(\mathrm{id}_{\mathrm{src}}\): a finite identifier under which the source is registered. Two records with the same identifier refer to the same source.
\item
  \(\mathrm{type}_{\mathrm{src}}\): the source-type tag, drawn from the finite source-of-truth family

  \[
  \mathsf{SourceOfTruth}\;=\;\{\,\texttt{committed\_state},\;\texttt{audited\_cell\_records},\;\texttt{independent\_pair\_witness},\;\texttt{simulation\_trace},\;\texttt{ablation\_record},\;\texttt{fallback},\;\texttt{unknown},\;\texttt{contradictory}\,\},
  \]

  inherited from packet 149 §20. The first five tags (\(\texttt{committed\_state}\), \(\texttt{audited\_cell\_records}\), \(\texttt{independent\_pair\_witness}\), \(\texttt{simulation\_trace}\), \(\texttt{ablation\_record}\)) are candidate source-of-truth tags whose force is determined by \(\mathsf{SourcePolicy}_I\) and the surrounding profile; \(\texttt{fallback}\) supports at most provisional pair-observable acceptance unless explicitly upgraded by an audited bridge, and \(\texttt{unknown}\) or \(\texttt{contradictory}\) does not supply a full source-of-truth record.
\item
  \(\mathrm{trace}_{\mathrm{src}}\): a finite trace recording how the source was constructed and how it has been audited; entries are typed records consistent with \(\mathsf{AuditPolicy}_I\) for the surrounding instrument.
\item
  \(\mathrm{valid}_{\mathrm{src}}\): a finite validity record, recording whether the source is currently within its declared validity window or condition.
\end{itemize}

A claim that depends on a source records the identifier of that source in its use-set. A claim whose source is \(\texttt{unknown}\) or \(\texttt{contradictory}\) does not supply a full source-of-truth record under any instrument that requires one; the resulting judgment is routed by the relevant source policy, typically to \(\texttt{blocked}\) or \(\texttt{outside\_scope}\). Provisional support is the fallback-source case above, not a consequence of an unknown or contradictory source by itself.

\subsubsection{Audit Records}

An audit record is a finite typed tuple

\[
A\;=\;(\,\mathsf{ClaimRef},\;\mathsf{EvidenceRefs},\;\mathsf{CheckRules},\;\mathsf{CheckResults},\;\mathsf{SourceRefs},\;\mathsf{VisibilityCheck},\;\mathsf{ThresholdCheck},\;\mathsf{CircularityCheck},\;\mathsf{Nonclaims},\;\mathsf{Provenance}\,),
\]

with the following components.

\begin{itemize}
\tightlist
\item
  \(\mathsf{ClaimRef}\): a finite reference to the object or claim being audited. Targets include theory packages, directed cells, pair observables, promotion bridges, instrument-indexed claim records, dependency records, and source records.
\item
  \(\mathsf{EvidenceRefs}\): the finite evidence references consulted by the audit.
\item
  \(\mathsf{CheckRules}\): the finite check rules, including the instrument and audit policy under which the audit is performed. An audit record without an instrumented check policy is malformed.
\item
  \(\mathsf{CheckResults}\): the finite replay or check results produced by applying the check rules.
\item
  \(\mathsf{SourceRefs}\): the source records used by the audited claim or object.
\item
  \(\mathsf{VisibilityCheck}\): the finite visibility check against the named instrument's visibility map and bridge policy.
\item
  \(\mathsf{ThresholdCheck}\): the finite threshold check against the relevant threshold record.
\item
  \(\mathsf{CircularityCheck}\): the check that no unstratified circular audit/support loop occurs.
\item
  \(\mathsf{Nonclaims}\): the nonclaim records required by the audited claim or object.
\item
  \(\mathsf{Provenance}\): the finite provenance trace of the audit record.
\end{itemize}

The audit passes only when its target resolves, its evidence resolves, the source-of-truth policy passes, visibility passes, thresholds pass, no unstratified circularity occurs, and the required nonclaims are present. The audit record family \(\mathsf{Audit}\) of \(\mathbf{BirdInt}^{\mathrm{aud}}_{\mathrm{fin}}\) is the finite collection of all such records used by judgments in the paper.

The audit record \(A\) is distinct from the inherited audit functional \(\mathcal A\) that lives inside a theory package \(\mathcal T=(Z,f,\Sigma_f,E,\mathcal A)\): \(\mathcal A\) is the package-internal audit functional, while \(A\) is a finite audit record, certificate, trace, or evidence object. Compatibility with \(\mathcal A\) is a separate typed condition when the inherited package audit functional is used. The two are not interchangeable, and a theorem that uses one does not by that fact use the other.

A drive certificate is a specialization of an audit record whose check data records the drive face \(\mathrm{P6}_{\mathrm{drive}}\) --- a nonzero cycle affinity, a nontrivial cohomology class, or another declared drive bridge --- rather than a generic \(P_6\) audit. A generic admissible audit record does not by itself supply a drive certificate; Subsection 5.3.6 records the drive-certificate schema and Section 8 states the no-drive theorems that constrain it.

\subsubsection{Evidence Hygiene}

Source and audit records support reproducibility and scope discipline. They are not themselves proofs, and admissibility under an instrument is not the same as the truth of a claim that uses them. The discipline of the paper is:

\begin{itemize}
\tightlist
\item
  a source record fixes which source a claim depends on, so that the claim's reliance on a real, fallback, dirty, proxy, unknown, or contradictory source is identifiable at the classifier step;
\item
  an audit record fixes which checks a claim has passed under a specific instrument, so that the claim is reproducible by anyone with access to the same records and the same audit policy;
\item
  neither record asserts instrument-free truth. An audited record may be exact under one instrument and suppressed or outside-scope under another, and the no-overreading theorem of Section 11 prevents an accepted claim under one instrument from being silently transferred to a different instrument.
\end{itemize}

NC-15 of Subsection 2.4.3 records the corresponding nonclaim: \(A\) passing under \(I\), \(\mathcal T\), and the surrounding scope yields acceptance under that scope, not instrument-free truth. Audits and source records record the hygiene of the calculus, not its conclusions; the conclusions are theorems whose proofs use the records together with the typed witness, update, profile, level, visibility, and threshold data of the surrounding judgments.

\subsection{Domain Instance Admissibility}

This subsection records the admissibility conditions on a full domain instance \(\mathfrak D\in\mathbf{BirdInt}^{\mathrm{aud}}_{\mathrm{fin}}\). Sections 5 through 14 work with admissible instances throughout, and the hypothesis ``let \(\mathfrak D\) be admissible'' is implicit in every theorem statement.

\subsubsection{Admissible Instance}

A full domain instance \(\mathfrak D\in\mathbf{BirdInt}^{\mathrm{aud}}_{\mathrm{fin}}\) is admissible, written

\[
\mathsf{AdmDomain}(\mathfrak D),
\]

when every component of the instance is finite and every record in every component meets its admissibility schema:

\begin{itemize}
\tightlist
\item
  \(\mathbb P=\{P_1,\ldots,P_6\}\) is the fixed finite primitive set of Subsection 2.2.1.
\item
  \(\mathsf{Lev}=\{\mathsf{Beh},\mathsf{Ver},\mathsf{Str}\}\) is the fixed finite level set of Subsection 2.2.3.
\item
  \(\mathsf{Host}\) is a finite subfamily of the host inventory of Subsection 2.1.2; only hosts named in \(\mathsf{Host}\) may appear on the host tag of any record in the instance.
\item
  \(\mathsf{Cont}\) is the finite content universe of the instance (Subsection 3.2): finite typed records, each with a host tag from \(\mathsf{Host}\), a level tag from \(\mathsf{Lev}\), an instrument-relative visibility class drawn from Subsection 3.5, and a profile attachment when required.
\item
  All references among content records resolve, or else the missing reference is represented by an explicit absence record in the appropriate status family.
\item
  \(\mathsf{Instr}\) is a finite family of admissible instruments, each meeting the conditions of Subsection 3.4.1.
\item
  \(\mathsf{Pkg}\) is a finite family of admissible theory packages, each meeting the conditions of Subsection 3.3.3.
\item
  \(\mathsf{Prof}\) is a finite family of admissible profiles, each meeting the conditions of Subsection 3.6.2.
\item
  \(\mathsf{Wit}\) and \(\mathsf{Upd}\) are the witness and update maps of Subsection 3.7.1-3.7.2; their values on each \(P_i\) are finite typed families, with the actor/informant typing rule \(W_j\in\mathsf{Wit}(P_j)\), \(U_i\in\mathsf{Upd}(P_i)\) enforced on every directed-cell record.
\item
  \(\mathsf{Def}\) is a finite family of admissible defect records in the schema of Subsection 5.2.1.
\item
  \(\mathsf{Judg}\) is the finite family of judgment forms introduced in Subsections 4.2 through 4.7 (role-channel, directed-cell, pair-observable, promotion, dependency, instrument-indexed claim).
\item
  \(\mathsf{Status}\) is the finite family of status sets associated to those judgment forms (role, cell, pair, promotion, claim, gate, square), listed in Subsection 5.1.
\item
  Threshold records are finite and typed in the threshold schema of Subsection 5.2.
\item
  \(\mathsf{Gate}\) is the finite family of gate specifications used in Sections 10 and 11, with the gate status sets \(\mathrm{GateStatus}\) and \(\mathrm{StrictGateStatus}\) from Subsection 5.1.6; promotion gates are explicit whenever a promotion judgment is claimed.
\item
  \(\mathsf{Dep}\) is a finite family of admissible dependency records in the schema of Subsection 4.6.
\item
  \(\mathsf{Audit}\) and \(\mathsf{Source}\) are finite families of admissible audit and source records in the schemas of Subsections 4.8.2 and 4.8.1.
\item
  \(\mathsf{Vis}\) is the finite family of visibility maps and visibility bridges of Subsection 3.5.
\item
  \(\mathsf{Nonclaim}\) is a finite family of admissible nonclaim records, attached either to instruments (as \(\mathcal N_I\)) or to individual judgments.
\item
  \(\mathsf{Real}\) is a finite family of admissible model-realization modules; admissibility of a model-realization module is the schema fixed in Section 13.
\end{itemize}

A full instance that satisfies all of the above is the working object of every theorem in Sections 5 through 11 and of every model-only realization statement in Section 13. Admissibility is a finite check on finite records: it composes the per-component admissibility checks of Sections 3 and 4 into a single instance-level condition and feeds the well-formedness, classifier, gate, and audit machinery of Sections 5 through 11.

\subsubsection{Exclusion Conditions}

A candidate instance is rejected as outside the scope of \(\mathbf{BirdInt}^{\mathrm{aud}}_{\mathrm{fin}}\) when any of the following holds:

\begin{itemize}
\tightlist
\item
  a component of the tuple is missing or infinite;
\item
  a host tag on any record points to a host outside \(\mathsf{Host}\) (in particular, outside the host inventory of Subsection 2.1.2);
\item
  a level tag on any record points outside \(\mathsf{Lev}\);
\item
  a theory package fails \(\mathsf{AdmPkg}\) from Subsection 3.3.3, including missing inherited Foundations I data \((Z,f,\Sigma_f,E,\mathcal A)\) or missing Foundations III attached records;
\item
  an instrument fails the instrument admissibility conditions of Subsection 3.4.1;
\item
  a profile fails \(\mathsf{AdmProfile}\) from Subsection 3.6.2;
\item
  a witness or update record violates the actor/informant typing rule of Subsection 3.7.3;
\item
  a directed-cell, pair-observable, promotion, dependency, or claim record has unresolved references, missing required fields, or status-family collapse where the record uses a status from one family in place of another;
\item
  an audit record is malformed, has unresolved evidence/source references, fails visibility or threshold checks, contains an unstratified circularity, or lacks required nonclaims;
\item
  a source record has type \(\texttt{unknown}\) or \(\texttt{contradictory}\) where the surrounding instrument's source policy requires a real or fallback source;
\item
  a nonclaim record is missing where one is required by the surrounding instrument or judgment;
\item
  a model-realization module is asserted as a universal theorem rather than as a scoped model-only realization (NC-19 of Subsection 2.4.3).
\end{itemize}

A claim, judgment, or theorem stated against a rejected candidate instance is not interpreted in \(\mathbf{BirdInt}^{\mathrm{aud}}_{\mathrm{fin}}\) and is excluded from the formal scope of this paper. Sections 5 through 14 work with admissible instances throughout, and a hypothesis ``let \(\mathfrak D\) be admissible'' is implicit in every theorem statement.
 \section{Status and Defect Semantics}

This section gathers the full status families used by the calculus and the defect record schemas that feed the defect-to-status rules consumed by Sections 6 through 11. The role-channel, directed-cell, pair-observable, promotion, and claim status families were introduced per judgment in Section 4; the gate and square families were referenced there. This section consolidates all seven families, fixes the defect record schema, lists the primitive and promotion-specific defect families used by the paper, and states the defect-to-status rules that drive the classifiers.

\subsection{Status-Family Separation}

The seven status families used in this paper are

\[
\mathrm{RoleStatus},\quad\mathrm{CellStatus},\quad\mathrm{PairStatus},\quad\mathrm{PromotionStatus},\quad\mathrm{ClaimStatus},\quad\mathrm{GateStatus},\quad\mathrm{SqStatus}.
\]

They are kept apart at the schema level: a value in one family is not by definition a value in another, and a status assigned to a judgment in one family does not by itself fix a status in another. The status-family separation theorem of Subsection 6.5 records this discipline at theorem-grade.

\subsubsection{Role-Channel Statuses}

The role-channel status set is the inherited Foundations II five:

\[
\mathrm{RoleStatus}\;=\;\{\,\texttt{active\_projection},\;\texttt{below\_threshold},\;\texttt{collapsed\_boundary\_case},\;\texttt{absent\_with\_record},\;\texttt{inapplicable\_outside\_scope}\,\}.
\]

The meanings are recorded in Subsection 4.2.2.

\subsubsection{Directed-Cell Statuses}

The directed-cell status set is

\[
\mathrm{CellStatus}\;=\;\{\,\texttt{action},\;\texttt{implicit},\;\texttt{trivial},\;\texttt{blocked},\;\texttt{undefined\_circular},\;\texttt{below\_threshold},\;\texttt{collapsed},\;\texttt{absent},\;\texttt{outside\_scope}\,\}.
\]

The classifier priority order is fixed: \(\texttt{outside\_scope}\succ\texttt{absent}\succ\texttt{blocked}\succ\texttt{undefined\_circular}\succ\texttt{collapsed}\succ\texttt{below\_threshold}\succ\texttt{trivial}\succ\texttt{implicit}\succ\texttt{action}\), as recorded in Subsection 4.3.4. Subsection 5.6.2 specifies the per-status conditions that the classifier checks.

\subsubsection{Pair-Observable Statuses}

The pair-observable status set is

\[
\mathrm{PairStatus}\;=\;\{\,\texttt{real},\;\texttt{real\_provisional},\;\texttt{real\_duplicated},\;\texttt{blocked}\,\}.
\]

The meanings and source-of-truth conditions are recorded in Subsection 4.4.

\subsubsection{Promotion Statuses}

The promotion status set is

\[
\mathrm{PromotionStatus}\;=\;\{\,\texttt{candidate},\;\texttt{accepted},\;\texttt{strict},\;\texttt{non\_strict},\;\texttt{failed\_descent},\;\texttt{failed\_stability},\;\texttt{failed\_audit},\;\texttt{failed\_no\_smuggling},\;\texttt{local\_only},\;\texttt{globally\_obstructed},\;\texttt{outside\_scope}\,\}.
\]

The values \(\texttt{strict}\) and \(\texttt{non\_strict}\) are accepted-status refinements; the failure-tagged values record specific gate failures, as in Subsection 4.5.3.

\subsubsection{Claim Statuses}

The instrument-indexed claim status set is

\[
\mathrm{ClaimStatus}\;=\;\{\,\texttt{accepted},\;\texttt{rejected},\;\texttt{provisional},\;\texttt{blocked},\;\texttt{outside\_scope},\;\texttt{absent\_with\_record},\;\texttt{below\_threshold},\;\texttt{failed\_audit},\;\texttt{undefined\_circular}\,\}.
\]

Subsection 4.7.2 records the meanings; Subsection 5.6.5 specifies the classifier, and Section 11 proves the no-overreading and same-level self-audit theorems against it.

\subsubsection{Gate Statuses}

A promotion or claim gate has status drawn from the finite set

\[
\mathrm{GateStatus}\;=\;\{\,\texttt{pass},\;\texttt{fail},\;\texttt{not\_required},\;\texttt{not\_checked},\;\texttt{outside\_scope}\,\}.
\]

A strictness gate has the refined set

\[
\mathrm{StrictGateStatus}\;=\;\{\,\texttt{strict\_pass},\;\texttt{non\_strict\_pass},\;\texttt{not\_required},\;\texttt{not\_checked},\;\texttt{bad\_audit},\;\texttt{outside\_scope}\,\}.
\]

The core promotion gates used by the paper are

\[
\mathsf{CoreGates}\;=\;\{\,G_{\mathrm{suff}},\;G_{\mathrm{desc}},\;G_{\mathrm{stab}},\;G_{\mathrm{ctrl}},\;G_{\mathrm{nosmuggle}},\;G_{\mathrm{vis}},\;G_{\mathrm{audit}}\,\},
\]

with the additional strictness gate \(G_{\mathrm{strict}}\) and the optional local-to-global gate \(G_{\mathrm{locglob}}\). Their roles are fixed in Section 10.

\subsubsection{Square Statuses}

The high-structure module of Section 14 uses the square status set

\[
\mathrm{SqStatus}\;=\;\{\,\texttt{exact},\;\texttt{lax},\;\texttt{obstructed},\;\texttt{blocked},\;\texttt{outside\_scope},\;\texttt{failed\_audit},\;\texttt{undefined\_circular},\;\texttt{provisional}\,\}.
\]

The values have the meanings of an exact square (\(\texttt{exact}\)), an ordered or approximate square comparison (\(\texttt{lax}\)), a square in scope but with nonzero defect (\(\texttt{obstructed}\)), and the standard outside, failed-audit, circular, and provisional cases. Section 14 records the per-value conditions.

\subsection{General Defect Record}

\subsubsection{Record Form}

A typed defect record is the finite tuple

\[
\mathsf{Def}_i^{\lambda}(\mathcal T,I)\;=\;(\,\delta_i,\;\Omega_i,\;\Theta_i,\;\mathrm{polarity}_i,\;W_i,\;A_i,\;V_i,\;\sigma_i\,)
\]

attached to a primitive index \(i\in\{1,\ldots,6\}\) (or to a non-primitive defect family in Subsections 5.4 and 5.5), under profile \(\lambda\), theory package \(\mathcal T\), and instrument \(I\). The fields are finite and have the following roles:

\begin{itemize}
\tightlist
\item
  \(\delta_i\): the defect value, a finite typed value drawn from the value type of \(\Omega_i\).
\item
  \(\Omega_i\): the value object or defect type, naming the kind of value \(\delta_i\) takes (numeric, Boolean, categorical, set-valued, cohomological, audit-valued).
\item
  \(\Theta_i\): the threshold or pass condition against which \(\delta_i\) is evaluated.
\item
  \(\mathrm{polarity}_i\): the polarity tag, naming whether the defect is an obstruction, certificate, gap, residual, capacity loss, or other typed polarity drawn from a finite family fixed by the surrounding family.
\item
  \(W_i\): the witness, certificate, or counterexample supporting \(\delta_i\).
\item
  \(A_i\): the audit record of the defect, in the schema of Subsection 4.8.2.
\item
  \(V_i\): the visibility record of the defect under \(I\), in the schema of Subsection 3.5.
\item
  \(\sigma_i\): the induced status or status annotation, drawn from the status family appropriate to the surrounding judgment.
\end{itemize}

The general threshold rule for evaluating a defect is

\[
\mathsf{Eval}_{\Theta_i}(\delta_i)\;\in\;\{\,\texttt{pass},\;\texttt{fail},\;\texttt{not\_checked},\;\texttt{outside\_scope}\,\},
\]

producing one of four verdicts. The verdict feeds the defect-to-status rules of Subsection 5.6.

\subsubsection{Threshold Types}

The threshold field \(\Theta_i\) is drawn from the finite family

\[
\Theta_i\;\in\;\{\,\Theta_{=0},\;\Theta_{\neq0},\;\Theta_{\le\varepsilon},\;\Theta_{\ge\varepsilon},\;\Theta_{\in S},\;\Theta_{\notin S},\;\Theta_{\mathrm{pass}},\;\Theta_{\mathrm{source}},\;\Theta_{\mathrm{bridge}},\;\Theta_{\mathrm{audit}},\;\Theta_{\mathrm{vis}}\,\}.
\]

The eleven types have the following meanings, used by the per-defect rules of Subsections 5.3 through 5.5:

\begin{longtable}[]{@{}ll@{}}
\toprule
Threshold & Meaning\tabularnewline
\midrule
\endhead
\(\Theta_{=0}\) & exact defect vanishes\tabularnewline
\(\Theta_{\neq0}\) & nontrivial witness exists\tabularnewline
\(\Theta_{\le\varepsilon}\) & approximate defect tolerated below \(\varepsilon\)\tabularnewline
\(\Theta_{\ge\varepsilon}\) & activation passes a lower bound \(\varepsilon\)\tabularnewline
\(\Theta_{\in S}\) & representability or membership condition holds\tabularnewline
\(\Theta_{\notin S}\) & exclusion or nonrepresentability holds\tabularnewline
\(\Theta_{\mathrm{pass}}\) & a Boolean audit or gate pass\tabularnewline
\(\Theta_{\mathrm{source}}\) & the source-of-truth requirement is met\tabularnewline
\(\Theta_{\mathrm{bridge}}\) & a required bridge exists and is strong enough\tabularnewline
\(\Theta_{\mathrm{audit}}\) & the audit policy passes\tabularnewline
\(\Theta_{\mathrm{vis}}\) & the visibility check passes\tabularnewline
\bottomrule
\end{longtable}

Each threshold type has a fixed evaluation rule on values of the appropriate kind: \(\Theta_{=0}\) checks \(\delta=0\), \(\Theta_{\le\varepsilon}\) checks \(|\delta|\le\varepsilon\), and so on. The threshold types are the finite vocabulary used throughout Sections 5 through 11; no theorem in the paper introduces a threshold type outside this list without explicitly extending the family.

\subsection{Primitive Defect Families}

For each primitive \(P_i\), the calculus declares a finite family of defect record types. Each entry below uses the named defects fixed in the final status/defect packet. Full constructions and proofs that depend on these defects appear in Sections 6, 7, and 8; this subsection records only the schema and the defect-to-status link.

\subsubsection{\texorpdfstring{\(P_1\): Descent, Closure, Rewrite Defects}{P\_1: Descent, Closure, Rewrite Defects}}

The \(P_1\) defect family contains three named entries.

\begin{itemize}
\item
  \textbf{Exact quotient descent defect.} Given a finite quotient \(q:X\to Y\) and a finite map \(F:X\to X\), the split-pair set is

  \[
  \delta_1^{\mathrm{split}}(q,F)\;=\;\{\,(x,x')\;:\;q(x)=q(x')\;\wedge\;qF(x)\neq qF(x')\,\}.
  \]

  Threshold: \(\Theta_1^{\mathrm{desc}}:\delta_1^{\mathrm{split}}=\varnothing\). Polarity: obstruction. The status rule is \(\delta_1^{\mathrm{split}}=\varnothing\Rightarrow\) descent passes, while \(\delta_1^{\mathrm{split}}\neq\varnothing\) records an active \(P_1\) obstruction.
\item
  \textbf{Candidate descended-map defect.} For a candidate \(F^\sharp\), the defect is

  \[
  \delta_1^{\mathrm{cand}}(q,F,F^\sharp)\;=\;\{\,x:F^\sharp(q_{\mathrm{im}}(x))\neq q_{\mathrm{im}}F(x)\,\}.
  \]

  Threshold: \(\Theta_{=0}\). Polarity: obstruction.
\item
  \textbf{Markov closure deficit.} For finite Markov abstraction, the closure deficit is

  \[
  \mathrm{CD}_{\tau}^{\mu}(\Pi)\;=\;I(X_t;\Pi(X_{t+\tau})\mid\Pi(X_t)).
  \]

  Threshold: \(\Theta_1^{\mathrm{closure}}:\mathrm{CD}_{\tau}^{\mu}(\Pi)=0\) for exact closure, or \(\mathrm{CD}_{\tau}^{\mu}(\Pi)\le\varepsilon\) for approximate closure. Polarity: gap.
\end{itemize}

\subsubsection{\texorpdfstring{\(P_2\): Representability and Gating Defects}{P\_2: Representability and Gating Defects}}

The \(P_2\) defect family contains two named entries.

\begin{itemize}
\item
  \textbf{Representability defect.} For a lens \(f:Z\to O\) and predicate \(s\subseteq Z\),

  \[
  \delta_2^{\mathrm{rep}}(s;f)\;=\;
  \begin{cases}
  0, & s\in\Sigma_f,\\
  1, & s\notin\Sigma_f.
  \end{cases}
  \]

  In full finite-lens mode, \(s\in\Sigma_f\) iff \(s=f^{-1}(B)\) for some \(B\subseteq\mathrm{im}(f)\). A nonrepresentability witness is a pair \(z,z'\) with \(f(z)=f(z')\) and \(\mathbf 1_s(z)\neq\mathbf 1_s(z')\). Threshold: \(\Theta_{\in S}\) for representable, \(\Theta_{\notin S}\) for nonrepresentable. Polarity: obstruction.
\item
  \textbf{Gating / support-loss defect.} For graph/support gating \(G\leadsto G'\),

  \[
  \delta_2^{\mathrm{cycle}}(G,G')\;=\;\beta_1(G)-\beta_1(G').
  \]

  The drive-suppression threshold is \(\Theta_2^{\mathrm{forest}}:\beta_1(G')=0\), and \(\beta_1(G')=0\Rightarrow H^1(G')=0\). Polarity: capacity loss.
\end{itemize}

\subsubsection{\texorpdfstring{\(P_3\): Route-Mismatch, Holonomy, Pasting Defects}{P\_3: Route-Mismatch, Holonomy, Pasting Defects}}

The \(P_3\) defect family contains the formula-level entries fixed in the status/defect packet.

\begin{itemize}
\item
  \textbf{Completion commutator defect.} For completions \(E_1,E_2:C\to C\),

  \[
  \delta_3^{\mathrm{comp}}(E_1,E_2)\;=\;\{\,c:E_1E_2(c)\neq E_2E_1(c)\,\}.
  \]

  Route mismatch is active when \(\delta_3^{\mathrm{comp}}\neq\varnothing\), and route coherence holds when \(\delta_3^{\mathrm{comp}}=\varnothing\). Polarity: obstruction.
\item
  \textbf{Holonomy defect.} The holonomy defect is

  \[
  \delta_3^{\mathrm{hol}}(\gamma)\;=\;\operatorname{Hol}(\gamma)-\mathrm{id}.
  \]

  It is active when \(\operatorname{Hol}(\gamma)\neq\mathrm{id}\). The associated nonclaim is

  \[
  P_3^{\mathrm{holonomy}}\not\Rightarrow \mathrm{P6}_{\mathrm{drive}}
  \]

  without an explicit bridge.
\end{itemize}

Pasting and critical-pair material belong to the \(P_3\) role vocabulary and to the high-structure square module, but no additional primitive \(P_3\) pasting-defect formula is introduced here beyond the named packet entries.

\subsubsection{\texorpdfstring{\(P_4\): Refinement, Staging, Gluing Defects}{P\_4: Refinement, Staging, Gluing Defects}}

The \(P_4\) defect family contains two named entries.

\begin{itemize}
\item
  \textbf{Refinement factor defect.} For \(f_0:Z\to O_0\), \(f_1:Z\to O_1\), and forgetting map \(\rho:O_1\to O_0\),

  \[
  \delta_4^{\mathrm{ref}}(f_0,f_1,\rho)\;=\;\{\,z:f_0(z)\neq\rho f_1(z)\,\}.
  \]

  Refinement passes iff \(\delta_4^{\mathrm{ref}}=\varnothing\). Polarity: obstruction.
\item
  \textbf{Gluing defect.} For local data over an overlap,

  \[
  \delta_4^{\mathrm{glue}}\;=\;\{\,x\in U\cap V:\mathrm{res}_{U,U\cap V}(x_U)\neq\mathrm{res}_{V,U\cap V}(x_V)\,\}.
  \]

  Gluing passes iff \(\delta_4^{\mathrm{glue}}=\varnothing\). If local data exist but gluing fails, the induced promotion status is \(\texttt{globally\_obstructed}\).
\end{itemize}

\subsubsection{\texorpdfstring{\(P_5\): Packaging, Completion, Objecthood Defects}{P\_5: Packaging, Completion, Objecthood Defects}}

The \(P_5\) defect family contains three named entries.

\begin{itemize}
\item
  \textbf{Idempotence defect.} Given an endomap \(E:X\to X\), the idempotence defect is

  \[
  \delta_5^{\mathrm{idem}}(E)\;=\;\{\,x:E^2(x)\neq E(x)\,\}.
  \]

  Threshold: \(\Theta_5^{\mathrm{idem}}:\delta_5^{\mathrm{idem}}=\varnothing\). If \(E^2=E\), then \(E^n=E\) for \(n\ge1\), and the fixed exact completion profile may assign \(P_5\leftarrow P_5:\texttt{trivial}\).
\item
  \textbf{Fixed-point residual.} The fixed-point residual is

  \[
  \delta_5^{\mathrm{fix}}(o,E)\;=\;E(o)-o,
  \]

  or, in set-valued form, \(\{\,o:E(o)\neq o\,\}\). Threshold: \(\Theta_{=0}\) or \(\Theta_{\le\varepsilon}\).
\item
  \textbf{Package collapse defect.} Package collapse examples include \(\texttt{identity\_package}\), \(\texttt{total\_collapse}\), and \(\texttt{singleton\_collapse}\). These may induce \(\texttt{collapsed}\) or \(\texttt{trivial}\) depending on profile.
\end{itemize}

\subsubsection{\texorpdfstring{\(P_6\): Audit, Provenance, Drive Defects}{P\_6: Audit, Provenance, Drive Defects}}

The \(P_6\) defect family contains three named entries.

\begin{itemize}
\item
  \textbf{Audit defect.} The audit defect has values

  \[
  \delta_6^{\mathrm{audit}}\in\{\,\texttt{audit\_passes},\;\texttt{missing\_audit},\;\texttt{failed\_audit},\;\texttt{source\_failure},\;\texttt{visibility\_failure},\;\texttt{threshold\_failure},\;\texttt{nonclaim\_failure},\;\texttt{circular\_audit}\,\}.
  \]

  The audit threshold is \(\Theta_6^{\mathrm{audit}}:\delta_6^{\mathrm{audit}}=\texttt{audit\_passes}\). Circular audit induces \(\texttt{undefined\_circular}\).
\item
  \textbf{Provenance / source defect.} The source defect has values

  \[
  \delta_6^{\mathrm{source}}\in\{\,\texttt{committed},\;\texttt{dirty},\;\texttt{fallback\_forbidden},\;\texttt{unknown},\;\texttt{contradictory},\;\texttt{missing}\,\}.
  \]

  Real pair claims require \(\delta_6^{\mathrm{source}}=\texttt{committed}\).
\item
  \textbf{Drive certificate.} The drive certificate is the separate typed record

  \[
  \kappa_6^{\mathrm{drive}},
  \]

  kept distinct from \(\delta_6^{\mathrm{audit}}\). In the finite graph/cohomology host, the cohomological drive-certificate schema is

  \[
  [a]\neq0\in H^1(G)\;\Longleftrightarrow_{\mathrm{schema}}\;\exists\gamma:\oint_\gamma a\neq0.
  \]

  This schema is available only in the declared finite graph/cohomology host with the relevant cohomology-drive bridge. It is not supplied by a generic audit record.
\end{itemize}

The separation between the audit face \(P_6\) and the drive face \(\mathrm{P6}_{\mathrm{drive}}\) is preserved at the defect level: a generic audit defect is in the \(P_6\) audit subfamily and does not by itself imply a drive certificate failure, and a missing drive certificate does not by itself imply an audit defect. The route-mismatch dependency \(P_3\mathrel{\mathsf{dep}_{\texttt{requires\_bridge}}}\mathrm{P6}_{\mathrm{drive}}\) of Subsection 4.6.2 records this separation at the dependency level.

\subsection{Promotion-Specific Defects}

Promotion bridges carry a finite family of defect records that are not attached to any single primitive but to the bridge as a whole. Each entry below names the defect, the threshold it is checked against, and the promotion status it induces on failure, in the schema of Subsections 5.2 and 4.5.

\subsubsection{Strictness / Factorization Defect}

The strictness defect of a promotion bridge \(B_{j\to j+1}\) with object maps \(\pi_j,\pi_{j+1}\) on a finite carrier \(S\) is

\[
\Delta_{\mathrm{fact}}(\pi_j,\pi_{j+1})\;=\;\{\,(s,s')\in S\times S\;:\;\pi_j(s)=\pi_j(s')\;\wedge\;\pi_{j+1}(s)\neq\pi_{j+1}(s')\,\}.
\]

The strict threshold \(\Theta_{\mathrm{strict}}\) asks \(\Delta_{\mathrm{fact}}\neq\varnothing\); the non-strict threshold \(\Theta_{\mathrm{nonstrict}}\) asks \(\Delta_{\mathrm{fact}}=\varnothing\). The strictness gate \(G_{\mathrm{strict}}\) of Subsection 4.5 is then assigned

\[
\Delta_{\mathrm{fact}}\neq\varnothing\;\;\Longrightarrow\;\;G_{\mathrm{strict}}=\texttt{strict\_pass},\qquad \Delta_{\mathrm{fact}}=\varnothing\;\;\Longrightarrow\;\;G_{\mathrm{strict}}=\texttt{non\_strict\_pass}.
\]

Strictness is a property of the typed object-map data \(\pi_j,\pi_{j+1}\) alone; it is not by itself a closure or drive claim, and the negative scope of Subsection 2.4.3 records NC-10 and NC-11 against any such inference. Theorem 13 of Section 7 (factorization defect equivalence) states the equivalence between \(\Delta_{\mathrm{fact}}=\varnothing\) and the existence of a factorization \(\pi_{j+1}=\phi\,\pi_j\), and is the theorem-grade form of this defect.

\subsubsection{No-Smuggling Defect}

The no-smuggling defect \(\delta_{\mathrm{nosmuggle}}(B)\) records violations of the bridge's declared scope. The packet enumerates the following finite list of violation kinds:

\begin{enumerate}
\def\labelenumi{\arabic{enumi}.}
\tightlist
\item
  target-layer data used to prove target novelty,
\item
  suppressed content used without an admissible bridge,
\item
  simulation used as theorem,
\item
  local evidence used as global without a gluing record,
\item
  a \(P_3\)-witness used as a \(\mathrm{P6}_{\mathrm{drive}}\) certificate without a drive bridge,
\item
  a lossy summary used as exact data.
\end{enumerate}

The threshold is \(\Theta_{\mathrm{nosmuggle}}:\delta_{\mathrm{nosmuggle}}=\varnothing\). Failure induces

\[
\mathrm{Promote}(B)\;:\;\texttt{failed\_no\_smuggling}
\]

in the promotion classifier of Subsection 5.6.4, and the no-smuggling gate \(G_{\mathrm{nosmuggle}}\) of Subsection 4.5 records the gate verdict.

\subsubsection{Stability Defect}

For a target object \(o\in O_{j+1}\) formed by completion endomap \(E\) on the target package's content space, the stability defect is

\[
\delta_{\mathrm{stab}}\;=\;\{\,o\;:\;E(o)\neq o\,\},
\]

or its approximate-residual analogue under threshold \(\Theta_{\le\varepsilon}\) when an approximate completion is in scope. The stability gate \(G_{\mathrm{stab}}\) is required when the bridge claims that the target objects are fixed points of \(E\). Failure induces

\[
\mathrm{Promote}(B)\;:\;\texttt{failed\_stability}.
\]

When the bridge does not claim fixed-point structure on the target, the stability gate is recorded as \(\texttt{not\_required}\) rather than enforced.

\subsubsection{Descent Defect Inside Promotion}

When a promotion bridge claims macro closure on the target package --- that is, when the bridge claims a descended dynamics \(F^\sharp\) on the target's quotient or interface --- the descent gate \(G_{\mathrm{desc}}\) is required. The gate uses the descent defect of Subsection 5.3.1 (\(\delta_{\mathrm{split}}(q,F)\)) on the relevant quotient and update data of the target package. Failure induces

\[
\mathrm{Promote}(B)\;:\;\texttt{failed\_descent}.
\]

If no closure or descended-dynamics claim is made on the bridge, the descent gate is recorded as \(G_{\mathrm{desc}}=\texttt{not\_required}\) rather than enforced; in particular, the strictness gate alone does not require descent, since \(\pi_{j+1}\not\factor\pi_j\) does not by itself entail any claim about a descended dynamics on the target. NC-10 of Subsection 2.4.3 is the corresponding nonclaim.

\subsection{Claim-Level Visibility Defects}

Instrument-indexed claim records carry a finite family of visibility defects, all aggregated into a single visibility defect record. The threshold and induced statuses are recorded here at schema-grade; the formal no-overreading theorem is Theorem 23 in Section 11.

The aggregated visibility defect of a claim \(\varphi\) under instrument \(I\) relative to a finite admissible bridge family \(\mathcal L\) is

\[
\delta_{\mathrm{vis}}(I,\varphi,\mathcal L)\;=\;\delta_{\mathrm{overread}}\;\cup\;\delta_{\mathrm{weakbridge}}\;\cup\;\delta_{\mathrm{outside}}\;\cup\;\delta_{\mathrm{unknown}}.
\]

The threshold is \(\Theta_{\mathrm{vis}}:\delta_{\mathrm{vis}}=\varnothing\). Each of the four components is a finite set of records, defined below.

\subsubsection{Hidden-Content Overread}

The hidden-content overread defect is

\[
\delta_{\mathrm{overread}}(I,\varphi,\mathcal L)\;=\;\{\,x\in\mathsf{Use}(\varphi)\;:\;x\in\mathsf{Suppressed}_I\;\text{ and no admissible bridge in }\mathcal L\text{ permits }x\,\}.
\]

A nonempty \(\delta_{\mathrm{overread}}\) records that the claim depends on suppressed content for which no bridge has been registered or audited under \(\mathsf{BridgePolicy}_I\). The no-overreading theorem (Theorem 23 of Section 11) later uses the condition \(\delta_{\mathrm{overread}}=\varnothing\) in its acceptance argument; this subsection only records the defect schema.

\subsubsection{Missing Bridge}

The weak-bridge defect \(\delta_{\mathrm{weakbridge}}(I,\varphi,\mathcal L)\) records suppressed content that has a registered bridge in \(\mathcal L\), but where the bridge fails to provide the strength required by the claim:

\[
\delta_{\mathrm{weakbridge}}(I,\varphi,\mathcal L)\;=\;\{\,(x,L)\;:\;x\in\mathsf{Use}(\varphi),\;L\in\mathcal L\text{ bridges }x,\;\mathsf{Strength}(L)<_{\varphi}\mathsf{RequiredStrength}(\varphi)\,\}.
\]

Each entry pairs a use-set element \(x\) with the bridge \(L\) that admissibly relates \(x\) to visible content but whose declared strength, threshold, audit, or nonclaim profile does not meet the required strength of the surrounding claim. The promotion-gate semantics of Section 10 give the later promotion-specific use of this comparison; this subsection records only the claim-level visibility schema. A nonempty \(\delta_{\mathrm{weakbridge}}\) is reported separately from \(\delta_{\mathrm{overread}}\) so that ``no bridge'' and ``bridge failed'' failures are distinguishable in the classifier output and the audit record.

\subsubsection{Claim-Strength Mismatch}

The outside and unknown components of \(\delta_{\mathrm{vis}}\) record use-set entries that lie outside \(I\)'s scope or in the \(\texttt{unknown}\) visibility class:

\[
\delta_{\mathrm{outside}}(I,\varphi)\;=\;\mathsf{Use}(\varphi)\cap\mathsf{Outside}_I,\qquad \delta_{\mathrm{unknown}}(I,\varphi)\;=\;\mathsf{Use}(\varphi)\cap\mathsf{Unknown}_I.
\]

A nonempty \(\delta_{\mathrm{outside}}\) records that the claim refers to content the surrounding instrument does not adjudicate at all; a nonempty \(\delta_{\mathrm{unknown}}\) records that the claim refers to content within \(I\)'s scope whose visibility verdict has not yet been resolved. Both are visibility-side mismatches between the claim's declared strength and what \(I\) can adjudicate; the classifier of Subsection 5.6.5 routes such records away from \(\texttt{accepted}\).

The induced status on a claim with a nonempty visibility defect is recorded by judgment family in the table below:

\begin{longtable}[]{@{}ll@{}}
\toprule
\begin{minipage}[b]{0.47\columnwidth}\raggedright
Claim family\strut
\end{minipage} & \begin{minipage}[b]{0.47\columnwidth}\raggedright
Induced status\strut
\end{minipage}\tabularnewline
\midrule
\endhead
\begin{minipage}[t]{0.47\columnwidth}\raggedright
ordinary instrument-indexed claim\strut
\end{minipage} & \begin{minipage}[t]{0.47\columnwidth}\raggedright
\(\texttt{blocked}\)\strut
\end{minipage}\tabularnewline
\begin{minipage}[t]{0.47\columnwidth}\raggedright
directed cell\strut
\end{minipage} & \begin{minipage}[t]{0.47\columnwidth}\raggedright
\(\texttt{blocked}\)\strut
\end{minipage}\tabularnewline
\begin{minipage}[t]{0.47\columnwidth}\raggedright
pair observable\strut
\end{minipage} & \begin{minipage}[t]{0.47\columnwidth}\raggedright
\(\texttt{blocked}\)\strut
\end{minipage}\tabularnewline
\begin{minipage}[t]{0.47\columnwidth}\raggedright
promotion bridge\strut
\end{minipage} & \begin{minipage}[t]{0.47\columnwidth}\raggedright
\(\texttt{failed\_no\_smuggling}\) or \(\texttt{outside\_scope}\)\strut
\end{minipage}\tabularnewline
\begin{minipage}[t]{0.47\columnwidth}\raggedright
same-level audit\strut
\end{minipage} & \begin{minipage}[t]{0.47\columnwidth}\raggedright
\(\texttt{undefined\_circular}\) when the failure is circular\strut
\end{minipage}\tabularnewline
\bottomrule
\end{longtable}

The promotion case routes through the no-smuggling defect of Subsection 5.4.2 because a visibility failure on a promotion bridge is recorded as a no-smuggling violation in the gate semantics of Section 10; the same-level audit case is the special-case overreading of self-soundness recorded as circularity in Subsection 11.5.

\subsection{Defect-to-Status Rules}

The classifier rules below convert finite defect-record verdicts into status assignments per judgment family. Each rule is a finite check on the surrounding judgment, and each rule is consumed by the corresponding theorem in Sections 6, 10, or 11. The rules are stated as priority orders: when more than one condition matches, the highest-priority match is recorded.

\subsubsection{Role-Status Rules}

The role-channel classifier \(\mathsf{RoleClassify}(P_i,\mathcal T,I)\) assigns a value in \(\mathrm{RoleStatus}\) in the priority order

\begin{enumerate}
\def\labelenumi{\arabic{enumi}.}
\tightlist
\item
  \(\texttt{inapplicable\_outside\_scope}\), when \(P_i\) is outside the host or instrument scope under \(\mathcal T\);
\item
  \(\texttt{absent\_with\_record}\), when a required witness for the role is absent and the absence is audited;
\item
  \(\texttt{collapsed\_boundary\_case}\), when the candidate role configuration collapses under the relevant witness schema and cannot support an active projection;
\item
  \(\texttt{below\_threshold}\), when role evidence is present but the activation threshold is not met;
\item
  \(\texttt{active\_projection}\), otherwise, when role evidence is visible and audited.
\end{enumerate}

\subsubsection{Cell-Status Rules}

The directed-cell classifier consumes the cell record together with the per-cell defect records of Subsection 5.3 and the visibility defect of Subsection 5.5. The priority order is

\[
\texttt{outside\_scope}\;\succ\;\texttt{absent}\;\succ\;\texttt{blocked}\;\succ\;\texttt{undefined\_circular}\;\succ\;\texttt{collapsed}\;\succ\;\texttt{below\_threshold}\;\succ\;\texttt{trivial}\;\succ\;\texttt{implicit}\;\succ\;\texttt{action}.
\]

The per-status conditions are:

\begin{itemize}
\tightlist
\item
  \(\texttt{outside\_scope}\): the cell's host, level, or content lies outside \(\mathsf{Scope}_I\).
\item
  \(\texttt{absent}\): a required witness or update is missing and the absence is audited.
\item
  \(\texttt{blocked}\): the cell uses content with the wrong type, a missing or weak bridge, hidden suppressed content, an unaudited source, or a failed visibility check.
\item
  \(\texttt{undefined\_circular}\): the cell's audit is circular at the same level.
\item
  \(\texttt{collapsed}\): the cell's content collapses under a quotient or completion to a degenerate record.
\item
  \(\texttt{below\_threshold}\): the cell's evidence is present but a threshold check fails.
\item
  \(\texttt{trivial}\): the cell is well-formed but its update is the identity, or its content is no-op under the fixed profile; in particular, under a fixed completion profile, \(E^2=E\) gives \(P_5\leftarrow P_5:\texttt{trivial}\).
\item
  \(\texttt{implicit}\): the cell records typed bridge data that is implicit in the surrounding package and is not asserted as an active witness.
\item
  \(\texttt{action}\): the cell is well-formed, all required defects pass their thresholds, the visibility defect is empty (with admissible bridges in \(\mathcal L\) for any suppressed content used), and the audit and source policies pass.
\end{itemize}

The covered-cell uniqueness theorem (Theorem 4 of Section 6) later proves exact uniqueness for covered well-formed cells using this priority order.

\subsubsection{Pair-Status Rules}

The pair-observable classifier consumes the pair record together with its source-of-truth record and visibility defect, and assigns a value in \(\mathrm{PairStatus}=\{\texttt{real},\texttt{real\_provisional},\texttt{real\_duplicated},\texttt{blocked}\}\). The priority order, drawn from packet 151 §16, is

\begin{enumerate}
\def\labelenumi{\arabic{enumi}.}
\tightlist
\item
  \(\texttt{blocked}\), when any of the source, visibility, threshold, or audit defects fails its threshold; or when the source-of-truth record has type \(\texttt{unknown}\) or \(\texttt{contradictory}\) where a real or fallback source is required; or when the pair contract forbids the available source class entirely;
\item
  \(\texttt{real\_duplicated}\), when the pair record carries an explicit duplicate-pair record together with an admissible source under the bridge policy of \(I\);
\item
  \(\texttt{real}\), when the source-of-truth record is real (one of the five admissible real-source classes from Subsection 4.8.1: \(\texttt{committed\_state}\), \(\texttt{audited\_cell\_records}\), \(\texttt{independent\_pair\_witness}\), \(\texttt{simulation\_trace}\), or \(\texttt{ablation\_record}\) admitted by \(\mathsf{SourcePolicy}_I\)), the visibility defect is empty, and the audit and threshold policies pass;
\item
  \(\texttt{real\_provisional}\), when the source-of-truth record is fallback or proxy and admissible under \(\mathsf{SourcePolicy}_I\), the visibility defect is empty, and the bridge to a real source required by the policy is recorded.
\end{enumerate}

\subsubsection{Promotion-Status Rules}

The promotion classifier consumes the bridge \(B_{j\to j+1}\), its gate record \(\mathsf{GateResults}(B)\), its strictness defect, no-smuggling defect, stability defect, and descent defect, and assigns a value in \(\mathrm{PromotionStatus}\). The priority order is

\begin{enumerate}
\def\labelenumi{\arabic{enumi}.}
\tightlist
\item
  \(\texttt{outside\_scope}\), when the bridge or its content lies outside \(\mathsf{Scope}_I\);
\item
  \(\texttt{failed\_audit}\), when the audit gate fails;
\item
  \(\texttt{failed\_no\_smuggling}\), when the no-smuggling gate or visibility gate fails;
\item
  \(\texttt{failed\_descent}\), when the descent gate fails on a bridge that requires descent;
\item
  \(\texttt{failed\_stability}\), when the stability gate fails on a bridge that requires stability;
\item
  \(\texttt{globally\_obstructed}\), when the optional local-to-global gate \(G_{\mathrm{locglob}}\) fails;
\item
  \(\texttt{local\_only}\), when \(G_{\mathrm{locglob}}\) records local validity without admissible global packaging;
\item
  \(\texttt{candidate}\), when at least one required core gate has status \(\texttt{not\_checked}\);
\item
  \(\texttt{strict}\), when accepted core gates pass and \(\Delta_{\mathrm{fact}}\neq\varnothing\);
\item
  \(\texttt{non\_strict}\), when accepted core gates pass, \(\Delta_{\mathrm{fact}}=\varnothing\), and the factorization is witnessed;
\item
  \(\texttt{accepted}\), when accepted core gates pass and strictness is not required or not classified.
\end{enumerate}

The promotion-gate-soundness theorem (Theorem 19 of Section 10) later proves the accepted-status consequence for verdicts in \(\{\texttt{accepted},\texttt{strict},\texttt{non\_strict}\}\); this rule supplies only the classifier schema.

\subsubsection{Claim-Status Rules}

The claim classifier consumes the claim record, its use-set, the visibility defect of Subsection 5.5, the threshold and source records, the audit record, and the bridge policy of \(I\). The priority order is

\[
\texttt{outside\_scope}\;\succ\;\texttt{blocked}\;\succ\;\texttt{absent\_with\_record}\;\succ\;\texttt{failed\_audit}\;\succ\;\texttt{undefined\_circular}\;\succ\;\texttt{below\_threshold}\;\succ\;\texttt{rejected}\;\succ\;\texttt{provisional}\;\succ\;\texttt{accepted}.
\]

The per-status conditions are:

\begin{itemize}
\tightlist
\item
  \(\texttt{outside\_scope}\): the claim's use-set is not contained in \(\mathsf{Scope}_I\), or its claim type is not in \(\mathsf{ClaimTypes}_I\).
\item
  \(\texttt{blocked}\): a visibility check, missing bridge, wrong evidence type, or source check fails; threshold failures are routed separately to \(\texttt{below\_threshold}\).
\item
  \(\texttt{absent\_with\_record}\): the claim refers to records that are absent in \(\Gamma\) or \(\mathcal T\), with the absence audited.
\item
  \(\texttt{failed\_audit}\): the audit policy check fails.
\item
  \(\texttt{undefined\_circular}\): the claim refers to its own acceptance verdict at the same level.
\item
  \(\texttt{below\_threshold}\): a threshold check fails.
\item
  \(\texttt{rejected}\): the claim is well-formed but the check rules of \(I\) assign a rejection verdict.
\item
  \(\texttt{provisional}\): partial evidence is present and the profile permits provisional status.
\item
  \(\texttt{accepted}\): the claim is well-formed, the visibility defect is empty (with all required bridges admissible), the threshold and source policies pass, the audit policy passes, and the claim is not circular at the same level.
\end{itemize}

The no-overreading theorem (Theorem 23 of Section 11) later uses the constraint \(\delta_{\mathrm{overread}}=\varnothing\) in the acceptance argument, and the same-level self-audit failure theorem (Theorem 24 of Section 11) later applies the circularity rule to complete same-level self-audit attempts without a level shift.
 \section{Core Finite Interaction Theorems}

Sections 2 through 5 declared the finite audited interaction calculus, its judgment families, and the status and defect schemas. This section proves the core theorems of \(\mathbf{BirdInt}^{\mathrm{aud}}_{\mathrm{fin}}\): directed-cell status does not factor through the primitive pair, active typed interaction does not identify primitive roles, well-formedness of directed cells is decidable, every covered well-formed directed cell receives a unique status, and the seven status families are not collapsed into one. All five theorems are theorem-grade, stated under explicit host, instrument, and admissibility assumptions, and proved finitely.

\subsection{Theorem 1: No Total Six-Symbol Algebra}

\subsubsection{Statement}

\textbf{Theorem 1 (NoTotalAlgebra; theorem-grade).} Let \(\mathsf{CellData}\) be the class of typed directed-cell records of \(\mathbf{BirdInt}^{\mathrm{aud}}_{\mathrm{fin}}\), with the projection

\[
\mathsf{pair}\colon\mathsf{CellData}\to\mathbb P\times\mathbb P,\qquad C\mapsto(P_i,P_j)
\]

extracting the actor and informant primitive labels of \(C\), and the status assignment

\[
\mathsf{status}\colon\mathsf{CellData}\to\mathrm{CellStatus}
\]

returning the directed-cell status of \(C\) under the classifier of Subsection 5.6.2. If there exist \(C_1,C_2\in\mathsf{CellData}\) with

\[
\mathsf{pair}(C_1)\;=\;\mathsf{pair}(C_2)\qquad\text{and}\qquad\mathsf{status}(C_1)\;\neq\;\mathsf{status}(C_2),
\]

then \(\mathsf{status}\) does not factor through \(\mathsf{pair}\):

\[
\mathsf{status}\;\not\factor\;\mathsf{pair}.
\]

In particular, no total operation \(*\colon\mathbb P\times\mathbb P\to\mathbb P\) together with any decoder \(d\colon\mathbb P\to\mathrm{CellStatus}\) can faithfully encode the typed, statused, profile-relative interaction recorded by directed cells of \(\mathbf{BirdInt}^{\mathrm{aud}}_{\mathrm{fin}}\).

\subsubsection{Hypotheses}

The theorem is stated under the following finite admissibility assumptions:

\begin{itemize}
\tightlist
\item
  \(\mathbb P=\{P_1,\ldots,P_6\}\) is the fixed finite primitive set of Subsection 2.2.1.
\item
  \(\mathsf{CellData}\) is the class of admissible directed-cell records as defined in Subsection 4.3, that is, finite typed tuples \((P_i\leftarrow P_j,W_j,U_i,L,V,\Theta,A,\delta)^{\lambda}\) together with \(\Gamma,\mathcal T,I\), satisfying the actor/informant typing rule \(W_j\in\mathsf{Wit}(P_j)\), \(U_i\in\mathsf{Upd}(P_i)\) of Subsection 3.7.3 and the well-formedness conditions of Subsection 4.3.3.
\item
  \(\mathrm{CellStatus}\) is the directed-cell status set fixed in Subsection 4.3.4 and Subsection 5.1.2.
\item
  \(\mathsf{status}\) is the classifier of Subsection 5.6.2, applied to covered records.
\item
  The hypothesis posits two directed-cell records sharing the same primitive pair but receiving distinct statuses. The proof packet supplies such a finite profile-dependent witness: \(P_6\leftarrow P_3:\texttt{action}\) under an audit-update profile and \(P_6\leftarrow P_3:\texttt{blocked}\) under a drive-certification profile with the required bridge missing.
\end{itemize}

\subsubsection{Proof Strategy}

The argument is the finite factorization obstruction stated in the theorem.

Suppose for contradiction that \(\mathsf{status}\) factored through \(\mathsf{pair}\): that is, suppose there were a map

\[
g\colon\mathbb P\times\mathbb P\to\mathrm{CellStatus}
\]

such that \(\mathsf{status}(C)=g(\mathsf{pair}(C))\) for every \(C\in\mathsf{CellData}\). Apply this to the two records \(C_1,C_2\) supplied by the hypothesis. Since \(\mathsf{pair}(C_1)=\mathsf{pair}(C_2)\), evaluating \(g\) at this common pair gives

\[
\mathsf{status}(C_1)\;=\;g\bigl(\mathsf{pair}(C_1)\bigr)\;=\;g\bigl(\mathsf{pair}(C_2)\bigr)\;=\;\mathsf{status}(C_2),
\]

contradicting the hypothesis \(\mathsf{status}(C_1)\neq\mathsf{status}(C_2)\). Hence no such \(g\) exists, and \(\mathsf{status}\) does not factor through \(\mathsf{pair}\).

A total operation \(*\colon\mathbb P\times\mathbb P\to\mathbb P\) followed by a decoder \(d\colon\mathbb P\to\mathrm{CellStatus}\) is the special case \(g=d\circ *\), which is itself a map \(\mathbb P\times\mathbb P\to\mathrm{CellStatus}\). The same argument applies, and no such \((*,d)\) pair can encode \(\mathsf{status}\) faithfully.

The instrument-indexed form of the theorem is

\[
\Gamma\,;\;\mathcal T\,;\;I_{\mathrm{fin}}\;\vdash\;\textsc{NoTotalAlgebra}\;:\;\texttt{accepted},
\]

where \(\mathcal T\) is the finite admissible package carried by the instance and all cell records are interpreted in \(\mathbf{BirdInt}^{\mathrm{aud}}_{\mathrm{fin}}\).

\subsubsection{Consequence}

Directed cells \(P_i\leftarrow P_j\) are typed bridge judgments, not primitive products. The notation is closed against algebraic interpretation: there is no underlying multiplication of primitives whose values can determine cell status, and the typed witness, update, profile, level, visibility, threshold, audit, and defect data carried by every cell record are essential for the classifier's verdict. NC-1 of Subsection 2.4.3 is the corresponding nonclaim. Subsequent sections of this paper use directed-cell records, not primitive-pair labels, as the working unit of typed interaction; in particular, Section 13's PICA model-realization sheet does not interpret the \(6\times 6\) table as a universal interaction algebra.

\subsection{Theorem 2: Typed Non-Collapse}

\subsubsection{Statement}

\textbf{Theorem 2 (TypedNonCollapse; theorem-grade).} For every admissible directed cell

\[
\mathsf{Cell}_{ij}^{\lambda}\;=\;(\,P_i\leftarrow P_j,\;W_j,\;U_i,\;L,\;V,\;\Theta,\;A,\;\delta\,)
\]

in \(\mathbf{BirdInt}^{\mathrm{aud}}_{\mathrm{fin}}\), active typed interaction does not identify the actor and informant primitives:

\[
\mathsf{Cell}_{ij}^{\lambda}\;:\;\texttt{action}\;\;\not\Rightarrow\;\;P_i\;=\;P_j.
\]

More generally, the directed-cell judgment \(P_i\leftarrow P_j\) does not entail \(P_i=P_j\), regardless of the assigned cell status.

\subsubsection{Proof Strategy}

The proof uses the type separation of the witness and update families of Subsection 3.7.

A directed cell is a record

\[
(P_i\leftarrow P_j,\;W_j,\;U_i,\;L,\;V,\;\Theta,\;A,\;\delta)
\]

with the typing rules of Subsection 3.7.3:

\[
W_j\;\in\;\mathsf{Wit}(P_j),\qquad U_i\;\in\;\mathsf{Upd}(P_i).
\]

The rules carry both primitive labels as part of the data of the cell; they do not assert any equality relation between \(P_i\) and \(P_j\). The record constructor permits distinct primitive labels, for example \(P_1\leftarrow P_5\) for descent through a package or \(P_3\leftarrow P_5\) for noncommuting completions, with witness data drawn from the informant family and update data drawn from the actor family. Hence active interaction does not entail \(P_i=P_j\).

The same type-separation argument applies independently of the assigned cell status. The directed-cell judgment is a typed bridge with two slots; identifying the slots would conflate \(\mathsf{Wit}(P_j)\) and \(\mathsf{Upd}(P_i)\), two distinct typed families on distinct primitive indices, and is not sanctioned by the typing rule.

The diagonal case \(i=j\) is consistent with the typing rule but does not yield a collapse argument either way: a cell \(P_i\leftarrow P_i\) records typed witness data on the informant slot and typed update data on the actor slot, both indexed by \(P_i\) but drawn from distinct typed families (\(\mathsf{Wit}(P_i)\) and \(\mathsf{Upd}(P_i)\)).

The instrument-indexed form of the theorem is

\[
\Gamma\,;\;\mathcal T\,;\;I_{\mathrm{fin}}\;\vdash\;\textsc{TypedNonCollapse}\;:\;\texttt{accepted}.
\]

\subsubsection{Corollary}

Pair labels alone do not determine cell status. The same primitive pair \((P_i,P_j)\), with different witness, update, profile, threshold, visibility, audit, or defect data, can produce different cell statuses; this was the working content of Theorem 1, and Theorem 2 records the underlying typed-record discipline that licenses it. No theorem in this paper infers the equality of two primitive labels from the existence of a directed-cell judgment between them; in particular, no theorem in Sections 7, 10, or 13 does so even when the cell receives status \(\texttt{action}\).

\subsection{Theorem 3: Cell Well-Formedness Decidability}

\subsubsection{Statement}

\textbf{Theorem 3 (CellWellFormednessDecidability; theorem-grade).} In \(\mathbf{BirdInt}^{\mathrm{aud}}_{\mathrm{fin}}\), well-formedness of a candidate directed cell is decidable. For each candidate cell record

\[
C\;=\;(\,P_i\leftarrow P_j,\;W,\;U,\;L,\;V,\;\Theta,\;A,\;\delta\,)^{\lambda},
\]

together with \(\Gamma,\mathcal T,I\), there is a finite decision procedure \(\mathsf{WFCell}(C)\in\{\texttt{wf},\texttt{not\_wf}\}\) that returns \(\texttt{wf}\) exactly when \(C\) is admissible in the sense of Subsection 4.3.

\subsubsection{Algorithm}

The procedure runs the following finite checks in order; the cell is well-formed iff every step returns \(\texttt{wf}\).

\begin{enumerate}
\def\labelenumi{\arabic{enumi}.}
\tightlist
\item
  \textbf{Primitive labels.} Verify \(P_i,P_j\in\mathbb P\). This is a finite membership test on the six-element set.
\item
  \textbf{Witness type.} Verify \(W\in\mathsf{Wit}(P_j)\). The family \(\mathsf{Wit}(P_j)\) is a finite typed family (Subsection 3.7.1); the test is a finite membership check.
\item
  \textbf{Update type.} Verify \(U\in\mathsf{Upd}(P_i)\). The family \(\mathsf{Upd}(P_i)\) is finite (Subsection 3.7.2); the test is a finite membership check.
\item
  \textbf{Profile admissibility.} Verify \(\mathsf{AdmProfile}(\lambda,I,\mathcal T)\). The check is the finite per-field comparison of Subsection 3.6.2.
\item
  \textbf{Resolution of \(L,V,\Theta,A,\delta,\lambda\).} Verify that each component is a finite typed record on the host of \(\mathcal T\), with all references in its fields resolved against \(\mathsf{Cont}\) or recorded as audited absences.
\item
  \textbf{Reference resolution.} Verify that every reference in \(W,U,\delta,A\) --- to a quotient, a kernel, a defect record, an audit entry, a source identifier, or any other element of \(\mathsf{Cont}\) --- resolves either to an admissible record or to an admissibly recorded absence.
\item
  \textbf{Host and instrument adjudication.} Verify that the surrounding host \(h_{\mathcal T}\in\mathsf{Host}\) admits the cell type and that \(\mathsf{Scope}_I\) contains the cell record's content classes.
\end{enumerate}

\subsubsection{Proof}

Each step of the procedure is a finite decidable check: a membership test in a finite set, a finite per-field record comparison, or a finite policy comparison against a finite policy record. The composition of finitely many decidable checks is decidable, hence \(\mathsf{WFCell}\) is decidable. The instrument-indexed form of the theorem is

\[
\Gamma\,;\;\mathcal T\,;\;I_{\mathrm{fin}}\;\vdash\;\forall\,C,\;\mathsf{WFCell}(C)\text{ decidable}\;:\;\texttt{accepted}.
\]

The procedure is the well-formedness step that the classifier of Subsection 5.6.2 invokes before assigning a cell status. Theorem 4 will use this decidability step when proving uniqueness of the status assigned to a covered well-formed cell.

\subsection{Theorem 4: Covered Cell Status Uniqueness}

\subsubsection{Statement}

\textbf{Theorem 4 (CoveredCellUniqueStatus; theorem-grade).} For every covered candidate directed cell \(C\in\mathsf{CoveredCell}\), the finite priority classifier \(\mathsf{CellClassify}\) of Subsection 5.6.2 assigns exactly one cell status:

\[
\forall\,C\in\mathsf{CoveredCell},\quad\exists!\,\sigma\in\mathrm{CellStatus}\;\;:\;\;\mathsf{CellClassify}(C)\;=\;\sigma.
\]

A covered cell is a well-formed directed-cell record (Theorem 3) whose finite per-status conditions of Subsection 5.6.2 are jointly exhaustive on the record.

\subsubsection{Classifier Priority}

The status set is fixed in Subsection 5.1.2, and the priority order is frozen in Subsections 4.3.4 and 5.6.2:

\[
\texttt{outside\_scope}\;\succ\;\texttt{absent}\;\succ\;\texttt{blocked}\;\succ\;\texttt{undefined\_circular}\;\succ\;\texttt{collapsed}\;\succ\;\texttt{below\_threshold}\;\succ\;\texttt{trivial}\;\succ\;\texttt{implicit}\;\succ\;\texttt{action}.
\]

A cell whose data simultaneously matches more than one per-status condition is assigned the highest-priority match; the per-status conditions are specified in full in Subsection 5.6.2.

\subsubsection{Proof}

The classifier is a finite first-match procedure over the fixed priority order. The covered-cell hypothesis ensures that the per-status conditions of Subsection 5.6.2 --- outside-scope, absent, blocked, undefined-circular, collapsed, below-threshold, trivial, implicit, action --- are jointly exhaustive on \(C\): at least one of the nine conditions matches.

Existence of a status follows: the classifier walks the priority list, evaluates each per-status condition in order, and stops at the first match. Since at least one condition matches (covered-cell hypothesis), the classifier returns some \(\sigma\in\mathrm{CellStatus}\).

Uniqueness follows from the priority discipline: if more than one condition matches a cell, the classifier records only the highest-priority one. The output is therefore a single value in \(\mathrm{CellStatus}\). Each per-status check in Subsection 5.6.2 is a finite test on finite records; the entire classifier is a finite first-match procedure. Hence \(\mathsf{CellClassify}\) is a total function on \(\mathsf{CoveredCell}\) with codomain \(\mathrm{CellStatus}\), and assigns exactly one status to each covered well-formed cell.

The instrument-indexed form of the theorem is

\[
\Gamma\,;\;\mathcal T\,;\;I_{\mathrm{fin}}\;\vdash\;\textsc{CoveredCellUniqueStatus}\;:\;\texttt{accepted}.
\]

The result depends on the priority order having been frozen before the proof: the freeze is recorded in Subsection 4.3.4 and reused in Subsection 5.6.2, and the proof above does not introduce a separate ordering. Subsection 5.6.2 supplies the per-status conditions; this subsection only records that the priority assignment is well-defined and produces a unique value on each covered cell.

\subsection{Theorem 5: Status-Family Separation}

\subsubsection{Statement}

\textbf{Theorem 5 (StatusFamilySeparation; theorem-grade).} The five judgment-attached status families

\[
\mathrm{RoleStatus},\quad\mathrm{CellStatus},\quad\mathrm{PairStatus},\quad\mathrm{PromotionStatus},\quad\mathrm{ClaimStatus}
\]

are distinct judgment-indexed families and are not collapsed into a single status algebra. In particular, none of the following implications holds in \(\mathbf{BirdInt}^{\mathrm{aud}}_{\mathrm{fin}}\):

\[
\mathrm{Chan}_i\;:\;\texttt{active\_projection}\;\;\not\Rightarrow\;\;\mathsf{Cell}_{ij}^{\lambda}\;:\;\texttt{action},
\]

\[
\mathsf{Cell}_{ij}^{\lambda}\;:\;\texttt{action}\;\;\not\Rightarrow\;\;\mathsf{PairObs}_{\{i,j\}}^{\mu}\;:\;\texttt{real},
\]

\[
\mathrm{Promote}(B)\;:\;\texttt{strict}\;\;\not\Rightarrow\;\;\mathsf{Cell}_{ij}^{\lambda}\;:\;\texttt{action}.
\]

\subsubsection{Proof}

Each non-implication is established by exhibiting finite admissible records, or a finite admissible record schema, satisfying the antecedent without satisfying the consequent. The channel and promotion separations use the classifier-input distinctions of Sections 4 and 5; the cell-to-pair separation is the countermodel-backed case, with the explicit construction recorded as \(\mathsf{CM}_4\) in Section 9 and Appendix C.

\textbf{Channel does not imply cell.} Take a primitive \(P_i\) whose role channel is classified \(\texttt{active\_projection}\) under \(\mathcal T,I\) --- that is, an admissible role-channel record with a witness that passes the threshold, level, witness, and audit conditions of Subsection 5.6.1. Now consider a directed cell \(P_i\leftarrow P_j\) for some \(j\neq i\) whose informant data \(W_j\in\mathsf{Wit}(P_j)\) is missing and whose absence is audited. The cell is well-formed, but the classifier of Subsection 5.6.2 assigns it status \(\texttt{absent}\), not \(\texttt{action}\). Thus an active role channel does not by itself produce a directed cell with status \(\texttt{action}\); the channel and the cell are adjudicated by separate classifiers with separate inputs.

\textbf{Cell does not imply pair.} Take a directed cell \(\mathsf{Cell}_{ij}^{\lambda}\) that is well-formed and classified \(\texttt{action}\) under \(\Gamma,\mathcal T,I\). Now consider the unordered pair observable \(\mathsf{PairObs}_{\{i,j\}}^{\mu}\) on the same primitives, with a missing source-of-truth record where \(\mathsf{SourcePolicy}_I\) requires a real source. The cell is \(\texttt{action}\), but the pair classifier of Subsection 5.6.3 assigns the pair record status \(\texttt{blocked}\). Countermodel \(\mathsf{CM}_4\) of Section 9 supplies the explicit construction; the cell-to-pair non-implication is therefore witnessed by a finite admissible pair of records.

\textbf{Strict promotion does not imply cell action.} Take a promotion bridge \(B_{j\to j+1}\) classified \(\texttt{strict}\) under \(\Gamma,\mathcal T_j,I\) --- that is, an admissible bridge whose strictness gate \(G_{\mathrm{strict}}\) is \(\texttt{strict\_pass}\), equivalently \(\Delta_{\mathrm{fact}}(\pi_j,\pi_{j+1})\neq\varnothing\). Strictness is a property of the typed object-map data alone (Subsection 5.4.1) and does not by itself produce a directed cell of any specific primitive pair. A directed-cell record on the target package, with its own witness, update, profile, threshold, audit, visibility, and defect data, is required to obtain a \(\texttt{action}\) verdict on a cell. Hence strict promotion does not imply directed-cell action.

The three separation examples taken together establish that the five judgment-attached status families have different classifier inputs and non-equivalent admissibility conditions; no single status algebra is part of the domain. The gate and square families of Subsections 5.1.6 and 5.1.7 are likewise distinct from these five: their classifier inputs are gate records (per the gate semantics of Section 10) and square fillers (per the high-structure semantics of Section 14), not the cell, pair, promotion, or claim records.

The instrument-indexed form of the theorem is

\[
\Gamma\,;\;\mathcal T\,;\;I_{\mathrm{fin}}\;\vdash\;\textsc{StatusFamilySeparation}\;:\;\texttt{accepted}.
\]

\subsubsection{Consequence}

The paper does not infer pair-observable realness from directed-cell action, promotion strictness from package or channel activity, or claim acceptance from raw visibility alone. The status discipline of the calculus is the working content of the separation theorem: each status family is consulted only for its own judgments, and a status-level inference between families is recorded only through an explicit bridge or theorem. NC-3, NC-4, NC-5, and NC-6 of Subsection 2.4.3 are the corresponding nonclaims, and the no-overreading theorem of Section 11 gives the claim-level analogue at theorem-grade.
 \section{Descent, Packaging, Completion, and Strictness}

This section establishes the main finite theorem spine beyond the well-formedness, classifier, and status-family results of Section 6. Theorem 6 records the descent condition for a finite update through a quotient or finite map. Theorem 7 records the defect equivalence that makes the descent obstruction a finite typed record. Theorem 8 is the Markov closure-deficit theorem, the only theorem in this paper that lives on the stochastic host \(\mathsf H_{\mathrm{prob}}\). Theorems 9 and 10 record noncommuting completions and the resulting pasting defect. Theorem 11 is idempotent saturation. Theorems 12 and 13 record strict-extension nonfactorization and its defect equivalence. Theorem 14 is the fixed-interface definability bound.

\subsection{Theorem 6: Descent Through Quotient / Finite Map}

\subsubsection{Statement}

\textbf{Theorem 6 (DescentThroughQuotient; theorem-grade).} Let \(X\) and \(Y\) be finite sets on the host \(\mathsf H_{\mathrm{fin}}\), let \(q:X\to Y\) be a finite package, quotient, or abstraction map, and let \(F:X\to X\) be a finite update. Let \(q_{\mathrm{im}}:X\to\mathrm{im}(q)\) denote the effective-image map. There exists a finite descended update \(F^\sharp:\mathrm{im}(q)\to\mathrm{im}(q)\) with

\[
F^\sharp\,q_{\mathrm{im}}\;=\;q_{\mathrm{im}}\,F
\]

if and only if

\[
\forall\,x,x'\in X,\quad q(x)=q(x')\;\Longrightarrow\;qF(x)=qF(x').
\]

Equivalently, for a quotient \(X/\!\sim\) with quotient map \(q\),

\[
\exists\,F^\sharp\colon\;qF=F^\sharp q\;\;\Longleftrightarrow\;\;\bigl(x\sim x'\;\Longrightarrow\;F(x)\sim F(x')\bigr).
\]

\subsubsection{Data}

The hypotheses are: a finite set \(X\) on the host \(\mathsf H_{\mathrm{fin}}\), a finite map \(q:X\to Y\) with effective image \(\mathrm{im}(q)\subseteq Y\), a finite update \(F:X\to X\), and a candidate descended update \(F^\sharp:\mathrm{im}(q)\to\mathrm{im}(q)\). The judgment is registered under an admissible finite profile whose lens data include the effective-image map \(q_{\mathrm{im}}\).

\subsubsection{Proof}

For the forward direction, suppose \(F^\sharp\) exists with \(F^\sharp q_{\mathrm{im}}=q_{\mathrm{im}}F\). Take \(x,x'\in X\) with \(q(x)=q(x')\). Then \(q_{\mathrm{im}}(x)=q_{\mathrm{im}}(x')\), so

\[
q_{\mathrm{im}}F(x)\;=\;F^\sharp q_{\mathrm{im}}(x)\;=\;F^\sharp q_{\mathrm{im}}(x')\;=\;q_{\mathrm{im}}F(x').
\]

Taking the underlying values in the subset \(\mathrm{im}(q)\subseteq Y\) gives \(qF(x)=qF(x')\), establishing the fiber-constancy condition.

For the reverse direction, suppose \(q(x)=q(x')\Rightarrow qF(x)=qF(x')\). Define \(F^\sharp\) on \(\mathrm{im}(q)\) by

\[
F^\sharp(q_{\mathrm{im}}(x))\;=\;q_{\mathrm{im}}(F(x)),\qquad x\in X.
\]

The fiber-constancy hypothesis ensures this is well-defined: if \(q(x)=q(x')\), then \(qF(x)=qF(x')\), so the value does not depend on the chosen representative in the fiber. Restricting the codomain to \(\mathrm{im}(q)\) gives a finite map \(F^\sharp:\mathrm{im}(q)\to\mathrm{im}(q)\), and the relation \(F^\sharp q_{\mathrm{im}}=q_{\mathrm{im}}F\) holds by construction.

Both directions are finite checks on finite sets, so the proof is finitary throughout. The instrument-indexed form is

\[
\Gamma\,;\;\mathsf H_{\mathrm{fin}}\,;\;I_{\mathrm{fin}}\;\vdash\;\textsc{DescentThroughQuotient}\;:\;\texttt{accepted}.
\]

\subsubsection{Role Interpretation}

The theorem records the \(P_1/P_5/P_6\) descent condition: \(P_1\) supplies the update \(F\), \(P_5\) supplies the package \(q\), and \(P_6\) records the audit trail under which the descended map \(F^\sharp\) is accepted. Descent through \(q\) is admissible exactly when the descent defect \(\delta_1^{\mathrm{split}}(q,F)\) of Subsection 5.3.1 vanishes; Theorem 7 records this equivalence at the defect level.

\subsection{Theorem 7: Descent Defect Equivalence}

\subsubsection{Statement}

\textbf{Theorem 7 (DescentDefectEquivalence; theorem-grade).} With \(X,Y,q,F\) as in Theorem 6, define the descent defect

\[
\delta_1^{\mathrm{split}}(q,F)\;=\;\{\,(x,x')\in X\times X\;:\;q(x)=q(x')\;\wedge\;qF(x)\neq qF(x')\,\}.
\]

Then

\[
\delta_1^{\mathrm{split}}(q,F)\;=\;\varnothing\;\;\Longleftrightarrow\;\;F\;\text{descends through}\;q,
\]

and equivalently

\[
\delta_1^{\mathrm{split}}(q,F)\;\neq\;\varnothing\;\;\Longleftrightarrow\;\;\nexists\,F^\sharp\colon F^\sharp\,q_{\mathrm{im}}=q_{\mathrm{im}}\,F.
\]

\subsubsection{Proof}

By definition, \(\delta_1^{\mathrm{split}}(q,F)\) collects exactly the pairs \((x,x')\) that violate the fiber-constancy condition of Theorem 6. The set is empty iff the fiber-constancy condition holds for every pair, which by Theorem 6 is equivalent to the existence of a descended update \(F^\sharp\). The set is nonempty iff there exists a violating pair, which by Theorem 6 is equivalent to the nonexistence of any \(F^\sharp\).

Both directions are finite checks on the finite set \(X\times X\). The instrument-indexed form is

\[
\Gamma\,;\;\mathsf H_{\mathrm{fin}}\,;\;I_{\mathrm{fin}}\;\vdash\;\textsc{DescentDefectEquivalence}\;:\;\texttt{accepted}.
\]

The theorem makes the descent obstruction a finite typed record: a non-empty \(\delta_1^{\mathrm{split}}\) is the witness used by the descent gate of Subsection 5.4.4 in promotion records that claim macro closure on the target package, and a vanishing \(\delta_1^{\mathrm{split}}\) is the corresponding passing condition for that gate.

\subsection{Theorem 8: Markov Closure-Deficit Theorem}

\subsubsection{Statement}

\textbf{Theorem 8 (FiniteMarkovClosureDeficit; theorem-grade with assumptions).} On the finite stochastic host \(\mathsf H_{\mathrm{prob}}\), fix a finite state space \(X\), a finite stochastic kernel on \(X\), a fixed distribution, a finite abstraction \(\Pi:X\to Y\), a time lag \(\tau\), the declared KL predictive closure loss \(\mathcal L_\tau\), and a declared admissible class \(\mathcal K_{\mathrm{adm}}\) of macro kernels \(K\) on \(Y\) containing the conditional macro kernel \(K^\ast(\cdot\mid y)=P(\Pi(X_{t+\tau})\mid\Pi(X_t)=y)\) on the support of \(\Pi(X_t)\). Then the minimum predictive closure loss over admissible macro kernels is

\[
\min_{K\in\mathcal K_{\mathrm{adm}}}\;\mathcal L_\tau(K)\;=\;I\bigl(X_t\,;\,\Pi(X_{t+\tau})\,\big|\,\Pi(X_t)\bigr),
\]

where \(I(\,\cdot\,;\,\cdot\,|\,\cdot\,)\) is the conditional mutual information.

\subsubsection{Assumptions}

The theorem requires the following finite assumptions, all carried explicitly on the judgment line:

\begin{itemize}
\tightlist
\item
  \(\mathsf H_{\mathrm{prob}}\) is the host: a finite state space \(X\) with a finite stochastic kernel on \(X\) and a fixed distribution.
\item
  \(\Pi:X\to Y\) is a finite abstraction to a finite macro space \(Y\).
\item
  \(\tau\) is a finite time lag.
\item
  \(\mathcal L_\tau\) is the declared closure loss; in this theorem statement it is the KL predictive loss
  \[\mathcal L_\tau(K)\;=\;\mathbb E\Bigl[\,D_{\mathrm{KL}}\bigl(P(\Pi(X_{t+\tau})\mid X_t)\,\big\|\,K(\,\cdot\,\mid\Pi(X_t))\bigr)\,\Bigr].\]
\item
  The admissible macro-kernel class \(\mathcal K_{\mathrm{adm}}\) is a declared finite-dimensional class of stochastic kernels \(K:Y\to\Delta(Y)\) over which the minimization is taken, and it contains the conditional macro kernel \(K^\ast(\cdot\mid y)=P(Y_{t+\tau}\mid Y_t=y)\) on the support of \(Y_t\).
\end{itemize}

The theorem is theorem-grade only under these declared assumptions; it is not a universal stochastic-emergence statement.

\subsubsection{Proof}

Write \(Y_t=\Pi(X_t)\) and \(Y_{t+\tau}=\Pi(X_{t+\tau})\). For each macro state \(y\in Y\) in the support of \(Y_t\), decompose the conditional KL loss into a sum of two non-negative terms:

\[
\mathbb E\Bigl[\,D_{\mathrm{KL}}\bigl(P(Y_{t+\tau}\mid X_t)\,\big\|\,K(\,\cdot\,\mid y)\bigr)\;\Big|\;Y_t=y\,\Bigr]\;=\;\mathbb E\Bigl[\,D_{\mathrm{KL}}\bigl(P(Y_{t+\tau}\mid X_t)\,\big\|\,P(Y_{t+\tau}\mid Y_t=y)\bigr)\;\Big|\;Y_t=y\,\Bigr]\;+\;D_{\mathrm{KL}}\bigl(P(Y_{t+\tau}\mid Y_t=y)\,\big\|\,K(\,\cdot\,\mid y)\bigr).
\]

The first term does not depend on the choice of \(K\). The second term is non-negative and is minimized by the admissible choice \(K^\ast(\,\cdot\,\mid y)\;=\;P(Y_{t+\tau}\mid Y_t=y)\), at which it vanishes. Taking expectation over \(Y_t=y\) and minimizing over \(K\in\mathcal K_{\mathrm{adm}}\),

\[
\min_{K\in\mathcal K_{\mathrm{adm}}}\;\mathcal L_\tau(K)\;=\;\mathbb E\Bigl[\,D_{\mathrm{KL}}\bigl(P(Y_{t+\tau}\mid X_t)\,\big\|\,P(Y_{t+\tau}\mid Y_t)\bigr)\,\Bigr]\;=\;I\bigl(X_t\,;\,Y_{t+\tau}\,\big|\,Y_t\bigr).
\]

Substituting back \(Y_t=\Pi(X_t)\) and \(Y_{t+\tau}=\Pi(X_{t+\tau})\) gives the stated equality. Support-zero cases are handled by restricting to the support of the finite distribution; the finite-state assumption keeps the optimization a finite-dimensional convex problem with closed-form minimizer \(K^\ast\).

The instrument-indexed form is

\[
\Gamma\,;\;\mathsf H_{\mathrm{prob}}\,;\;I_{\mathrm{prob}}\;\vdash\;\textsc{FiniteMarkovClosureDeficit}\;:\;\texttt{accepted}.
\]

\subsubsection{Caution}

The theorem is assumption-scoped: the equality holds only when the host is \(\mathsf H_{\mathrm{prob}}\) and the assumptions of Subsection 7.3.2 are met. It is not a universal stochastic-emergence theorem and does not by itself imply package-implies-closure or strictness-implies-closure; NC-7 of Subsection 2.4.3 records the corresponding nonclaim. The closure deficit is a finite information-theoretic quantity on a declared finite Markov host; its interpretation outside that host requires a separate bridge.

\subsection{Theorem 9: Noncommuting Completions}

\subsubsection{Statement}

\textbf{Theorem 9 (NoncommutingCompletions; theorem-grade).} There exist finite idempotent completion operators \(E_1,E_2:C\to C\) on a finite carrier \(C\) with

\[
E_1^2\;=\;E_1,\qquad E_2^2\;=\;E_2,
\]

but

\[
E_1\,E_2\;\neq\;E_2\,E_1.
\]

Equivalently, there exists \(c\in C\) with \(E_1E_2(c)\neq E_2E_1(c)\).

\subsubsection{Construction}

Take \(U=\{a,b,c\}\) and \(C=\mathcal P(U)\), the finite power set of \(U\). Define

\[
E_1(S)\;=\;\begin{cases}\,S\cup\{b\}, & a\in S,\\[2pt] S, & a\notin S,\end{cases}\qquad E_2(S)\;=\;\begin{cases}\,S\cup\{c\}, & b\in S,\\[2pt] S, & b\notin S.\end{cases}
\]

Each map is a finite endomap of \(C\).

\textbf{Idempotence.} For \(E_1\): if \(a\in S\), then \(E_1(S)=S\cup\{b\}\) still contains \(a\), so \(E_1(E_1(S))=(S\cup\{b\})\cup\{b\}=S\cup\{b\}=E_1(S)\). If \(a\notin S\), then \(E_1(S)=S\), hence \(E_1(E_1(S))=E_1(S)\). So \(E_1^2=E_1\). The argument for \(E_2\) is symmetric (with \(b\) in place of \(a\) and \(c\) in place of \(b\)): \(E_2^2=E_2\).

\textbf{Noncommutation.} Take \(S_0=\{a\}\). Then

\[
E_1\,E_2(\{a\})\;=\;E_1(\{a\})\;=\;\{a,b\},
\]

since \(E_2\) leaves \(\{a\}\) fixed (because \(b\notin\{a\}\)). On the other hand,

\[
E_2\,E_1(\{a\})\;=\;E_2(\{a,b\})\;=\;\{a,b,c\},
\]

since \(E_1\) adds \(b\) (because \(a\in\{a\}\)) and then \(E_2\) adds \(c\) (because \(b\in\{a,b\}\)). Hence \(E_1E_2(\{a\})\neq E_2E_1(\{a\})\).

\subsubsection{Proof}

The construction above is a finite witness: \(E_1\) and \(E_2\) are finite idempotent endomaps of a finite carrier \(C\) with \(E_1E_2(\{a\})\neq E_2E_1(\{a\})\). The theorem is established.

The instrument-indexed form is

\[
\Gamma\,;\;\mathsf H_{\mathrm{fin}}\,;\;I_{\mathrm{fin}}\;\vdash\;\textsc{NoncommutingCompletions}\;:\;\texttt{accepted}.
\]

\subsubsection{Interpretation}

Route mismatch is a real defect even when each individual completion is exact and idempotent. The result rules out the overclaim targeted by NC-8 (no idempotence-implies-commutation) of Subsection 2.4.3 at theorem-grade: \(E_1^2=E_1\wedge E_2^2=E_2\) does not imply \(E_1E_2=E_2E_1\). Theorem 10 will record the equivalence between non-vanishing route-mismatch defect \(\Delta_3^{\mathrm{comp}}(E_1,E_2)\) and noncommutation; the construction here is the finite witness behind that equivalence.

\subsection{Theorem 10: Completion-Pasting Defect}

\subsubsection{Statement}

\textbf{Theorem 10 (CompletionPastingDefect; theorem-grade).} For finite completion endomaps \(E_1,E_2:C\to C\) on a finite carrier \(C\), define the route-mismatch defect

\[
\Delta_3^{\mathrm{comp}}(E_1,E_2)\;=\;\{\,c\in C\;:\;E_1E_2(c)\;\neq\;E_2E_1(c)\,\}.
\]

Then

\[
\Delta_3^{\mathrm{comp}}(E_1,E_2)\;\neq\;\varnothing\;\;\Longleftrightarrow\;\;E_1E_2\;\neq\;E_2E_1.
\]

The defect \(\Delta_3^{\mathrm{comp}}\) is a \(P_3\)-route/pasting defect (Subsection 5.3.3) informed by \(P_5\)-completion witnesses (Subsection 5.3.5).

\subsubsection{Proof}

By definition, \(\Delta_3^{\mathrm{comp}}(E_1,E_2)\) collects exactly the points \(c\in C\) at which the two composites disagree.

If \(\Delta_3^{\mathrm{comp}}=\varnothing\), then for every \(c\in C\), \(E_1E_2(c)=E_2E_1(c)\), so \(E_1E_2=E_2E_1\) as maps on \(C\).

If \(\Delta_3^{\mathrm{comp}}\neq\varnothing\), then there exists some \(c\in C\) with \(E_1E_2(c)\neq E_2E_1(c)\), so \(E_1E_2\neq E_2E_1\) as maps on \(C\).

Both directions are finite checks on the finite carrier \(C\). The instrument-indexed form is

\[
\Gamma\,;\;\mathsf H_{\mathrm{fin}}\,;\;I_{\mathrm{fin}}\;\vdash\;\textsc{CompletionPastingDefect}\;:\;\texttt{accepted}.
\]

Together with Theorem 9, the result records that the route-mismatch defect is a finite typed record: any noncommutation of finite completions is witnessed by a non-empty \(\Delta_3^{\mathrm{comp}}\), and the construction of Theorem 9 is a finite explicit instance with \(\{a\}\in\Delta_3^{\mathrm{comp}}(E_1,E_2)\). The high-structure pasting-square fragment of Section 14 later uses this route-mismatch record as the finite input for decorated-square pasting defects.

\subsection{Theorem 11: Idempotent Saturation}

\subsubsection{Statement}

\textbf{Theorem 11 (IdempotentSaturation; theorem-grade).} Let \(X\) be a finite set on \(\mathsf H_{\mathrm{fin}}\), and let \(E:X\to X\) be a finite self-map. If \(E^2=E\), then for every \(n\ge 1\),

\[
E^n\;=\;E.
\]

Pointwise: \(\forall n\ge 1\;\forall x\in X,\;E^n(x)=E(x)\).

\subsubsection{Proof}

Induct on \(n\).

\textbf{Base case} (\(n=1\)). \(E^1=E\) by definition.

\textbf{Inductive step.} Assume \(E^n=E\) for some \(n\ge 1\). Then

\[
E^{n+1}\;=\;E\circ E^n\;=\;E\circ E\;=\;E^2\;=\;E.
\]

At each induction step, the equality \(E^n=E\) is an equality of finite functions on the finite carrier \(X\); the argument is the ordinary induction on positive integers using the idempotence equation \(E^2=E\). The instrument-indexed form is

\[
\Gamma\,;\;\mathsf H_{\mathrm{fin}}\,;\;I_{\mathrm{fin}}\;\vdash\;\textsc{IdempotentSaturation}\;:\;\texttt{accepted}.
\]

\subsubsection{Consequence}

Repeated application of a fixed exact completion does not produce theory growth: the saturation \(E,E^2,E^3,\ldots\) collapses to \(E\) once \(E^2=E\). Strict theory growth requires a new object map, bridge, refinement, or nonfactorization witness, not iteration of a fixed completion; NC-9 of Subsection 2.4.3 is the corresponding nonclaim. Under a fixed exact completion profile, the directed-cell consequence is

\[
E^2\;=\;E\;\;\Longrightarrow\;\;P_5\leftarrow P_5\;:\;\texttt{trivial}
\]

in the cell classifier of Subsection 5.6.2, recording that an actor cell with this profile and update is admissible but trivial in the saturated case.

\subsection{Theorem 12: Strict-Extension Nonfactorization}

\subsubsection{Statement}

\textbf{Theorem 12 (StrictExtensionNonfactorization; theorem-grade).} Let \(S\) be a finite set on \(\mathsf H_{\mathrm{fin}}\), and let

\[
\pi_0\colon S\to O_0,\qquad \pi_1\colon S\to O_1
\]

be finite object maps with effective codomain \(O_0^{\mathrm{eff}}=\mathrm{im}(\pi_0)\subseteq O_0\). Then \(\pi_1\) does not factor through \(\pi_0\), written \(\pi_1\not\factor\pi_0\), meaning

\[
\nexists\,\phi\colon O_0^{\mathrm{eff}}\to O_1\;\;\text{with}\;\;\pi_1\;=\;\phi\,\pi_0,
\]

if and only if

\[
\exists\,s,s'\in S\;:\;\pi_0(s)=\pi_0(s')\;\wedge\;\pi_1(s)\neq\pi_1(s').
\]

\subsubsection{Data}

The hypotheses are: a finite carrier \(S\), finite object maps \(\pi_0,\pi_1\), the effective codomain \(O_0^{\mathrm{eff}}=\mathrm{im}(\pi_0)\), and a candidate factorization \(\phi:O_0^{\mathrm{eff}}\to O_1\). The theorem treats \(\pi_0,\pi_1\) as the old and new interface of a promotion bridge or extension; the strict-extension condition is the absence of \(\phi\).

\subsubsection{Proof}

For the forward direction, suppose there exist \(s,s'\in S\) with \(\pi_0(s)=\pi_0(s')\) and \(\pi_1(s)\neq\pi_1(s')\). Suppose for contradiction that some \(\phi:O_0^{\mathrm{eff}}\to O_1\) satisfies \(\pi_1=\phi\,\pi_0\). Then

\[
\phi(\pi_0(s))\;=\;\pi_1(s)\;\neq\;\pi_1(s')\;=\;\phi(\pi_0(s')),
\]

contradicting that \(\phi\) is well-defined as a function on the common value \(\pi_0(s)=\pi_0(s')\). Hence no \(\phi\) exists, and \(\pi_1\not\factor\pi_0\).

For the reverse direction, suppose no such pair \((s,s')\) exists, that is, \(\pi_0(s)=\pi_0(s')\Rightarrow\pi_1(s)=\pi_1(s')\) for every \(s,s'\in S\). Define \(\phi\) on \(O_0^{\mathrm{eff}}\) by \(\phi(\pi_0(s))=\pi_1(s)\) for \(s\in S\). The fiber-constancy condition just stated ensures that \(\phi\) does not depend on the choice of \(s\) in the fiber, so \(\phi\) is well-defined. Thus there is a finite map \(\phi:O_0^{\mathrm{eff}}\to O_1\) with \(\pi_1=\phi\,\pi_0\). Therefore, if no witness pair exists, \(\pi_1\) factors through \(\pi_0\); equivalently, nonfactorization forces such a witness pair.

Both directions are finite checks on the finite set \(S\times S\). The instrument-indexed form is

\[
\Gamma\,;\;\mathsf H_{\mathrm{fin}}\,;\;I_{\mathrm{fin}}\;\vdash\;\textsc{StrictExtensionNonfactorization}\;:\;\texttt{accepted}.
\]

\subsubsection{Consequence}

The strict-extension condition is exactly nonfactorization of the new interface through the old interface; this is the strictness-gate comparison of Subsection 5.4.1 at theorem grade. Strictness does not imply closure (NC-10) or drive (NC-11); the theorem records only the nonfactorization equivalence on the typed object-map data, and the no-go countermodels of Section 9 supply finite explicit instances in which strictness coexists with absent macro closure or absent drive certificates.

\subsection{Theorem 13: Factorization Defect Equivalence}

\subsubsection{Statement}

\textbf{Theorem 13 (FactorizationDefectEquivalence; theorem-grade).} With \(S,\pi_0,\pi_1\) as in Theorem 12, define the factorization defect

\[
\Delta_{\mathrm{fact}}(\pi_0,\pi_1)\;=\;\{\,(s,s')\in S\times S\;:\;\pi_0(s)=\pi_0(s')\;\wedge\;\pi_1(s)\neq\pi_1(s')\,\}.
\]

Then

\[
\Delta_{\mathrm{fact}}(\pi_0,\pi_1)\;\neq\;\varnothing\;\;\Longleftrightarrow\;\;\pi_1\;\not\factor\;\pi_0,
\]

and equivalently

\[
\Delta_{\mathrm{fact}}(\pi_0,\pi_1)\;=\;\varnothing\;\;\Longleftrightarrow\;\;\pi_1\;\factor\;\pi_0.
\]

\subsubsection{Proof}

By definition, \(\Delta_{\mathrm{fact}}(\pi_0,\pi_1)\) collects exactly the witness pairs of Theorem 12. The equivalence then follows immediately: nonfactorization holds iff a witness pair exists iff \(\Delta_{\mathrm{fact}}\) is nonempty; factorization holds iff no witness pair exists iff \(\Delta_{\mathrm{fact}}\) is empty.

Both directions are finite checks on the finite set \(S\times S\). The instrument-indexed form is

\[
\Gamma\,;\;\mathsf H_{\mathrm{fin}}\,;\;I_{\mathrm{fin}}\;\vdash\;\textsc{FactorizationDefectEquivalence}\;:\;\texttt{accepted}.
\]

The theorem makes the strictness condition a finite typed record: \(\Delta_{\mathrm{fact}}\) is the strictness defect of Subsection 5.4.1 read in the present setting, and the strictness gate \(G_{\mathrm{strict}}\) is \(\texttt{strict\_pass}\) exactly when \(\Delta_{\mathrm{fact}}\neq\varnothing\).

\subsection{Theorem 14: Fixed-Interface Definability Bound}

\subsubsection{Statement}

\textbf{Theorem 14 (FixedInterfaceDefinabilityBound; theorem-grade).} Let \(f:X\to Y\) be a finite map on \(\mathsf H_{\mathrm{fin}}\), and define the full finite-lens definability of \(f\) as

\[
\mathrm{Def}(f)\;=\;\{\,f^{-1}(B)\;:\;B\subseteq\mathrm{im}(f)\,\}.
\]

Then

\[
|\,\mathrm{Def}(f)\,|\;=\;2^{|\mathrm{im}(f)|}.
\]

\subsubsection{Proof}

Consider the map \(\Phi:\mathcal P(\mathrm{im}(f))\to\mathrm{Def}(f)\) defined by \(\Phi(B)=f^{-1}(B)\). The map is well-defined and surjective onto \(\mathrm{Def}(f)\) by construction.

\textbf{Injectivity.} Suppose \(B,B'\subseteq\mathrm{im}(f)\) with \(B\neq B'\). Without loss of generality, take \(y\in B\setminus B'\). Since \(y\in\mathrm{im}(f)\), there is some \(x\in X\) with \(f(x)=y\). Then \(x\in f^{-1}(B)\) but \(x\notin f^{-1}(B')\), so \(f^{-1}(B)\neq f^{-1}(B')\), that is, \(\Phi(B)\neq\Phi(B')\).

\textbf{Surjectivity.} By definition, every element of \(\mathrm{Def}(f)\) is of the form \(f^{-1}(B)\) for some \(B\subseteq\mathrm{im}(f)\), hence in the image of \(\Phi\).

Therefore \(\Phi\) is a bijection \(\mathcal P(\mathrm{im}(f))\to\mathrm{Def}(f)\), and

\[
|\,\mathrm{Def}(f)\,|\;=\;|\mathcal P(\mathrm{im}(f))|\;=\;2^{|\mathrm{im}(f)|}.
\]

The instrument-indexed form is

\[
\Gamma\,;\;\mathsf H_{\mathrm{fin}}\,;\;I_{\mathrm{fin}}\;\vdash\;\textsc{FixedInterfaceDefinabilityBound}\;:\;\texttt{accepted}.
\]

\subsubsection{Consequence}

For a fixed finite interface \(f\), the number of definable package distinctions is bounded by \(2^{|\mathrm{im}(f)|}\). Open-ended definability requires changing the interface, the package, the instrument, or the host; it does not arise from repeated application of a fixed completion (Theorem 11) or from the same lens at a deeper level. Combined with the strict-extension theorems (Theorems 12 and 13), the bound makes the discipline of Section 10 explicit: strict promotion is recorded as an interface change rather than as theory growth at fixed interface.
 \section{Graph/Cohomology Drive Module}

This section proves the scoped no-drive and drive-separation results that the calculus uses to keep the drive face \(\mathrm{P6}_{\mathrm{drive}}\) distinct from generic \(P_6\) audit. All theorems in this section live on the host \(\mathsf H_{\mathrm{graph}}\) of finite graphs and finite cycle/cochain data, and all theorems are scoped to cycle-supported cohomological drive on a declared finite graph. The module is not a universal theory of drive: it does not classify all thermodynamic force, all entropy production, or all directionality. The negative scope of Subsection 2.4.3 records the corresponding nonclaim (NC-12 and NC-17 in particular).

\subsection{Finite Graph/Cohomology Support Model}

This subsection fixes the finite graph/cohomology support model used by Theorems 15 through 18. It is a schema-grade declaration of the host data, the drive certificate format, and the gating operations admitted on the host. No theorem is asserted in this subsection; the no-drive and gating results are stated and proved in Subsections 8.2 through 8.5.

\subsubsection{Graph Data}

The host \(\mathsf H_{\mathrm{graph}}\) provides finite graph data:

\begin{itemize}
\tightlist
\item
  a finite vertex set \(V(G)\);
\item
  a finite edge set \(E(G)\) (oriented or unoriented as the surrounding judgment requires);
\item
  the finite first Betti number \(\beta_1(G)\;=\;|E(G)|\;-\;|V(G)|\;+\;c(G)\), where \(c(G)\) is the number of connected components of \(G\);
\item
  the finite cycle complex \(Z_*(G)\) and the finite cochain complex \(C^*(G)\) with coboundary \(d:C^0\to C^1\), taken over a fixed coefficient field or another explicitly declared free coefficient host;
\item
  the finite first cohomology \(H^1(G)\;=\;\ker(d_1)/\mathrm{im}(d_0)\).
\end{itemize}

The cycle integral of a finite \(1\)-cochain \(a\in C^1(G)\) around a finite cycle \(\gamma\in Z_1(G)\) is the standard sum

\[
\oint_\gamma a\;=\;\sum_{e\in\gamma}\,a(e),
\]

with the orientation conventions of \(\gamma\) applied to each edge \(e\in\gamma\). Every record in this section carries a host tag \(\mathsf H_{\mathrm{graph}}\), a level tag drawn from \(\mathsf{Lev}\), and an instrument annotation. Cycle-supported drive is the only drive notion adjudicated here.

\subsubsection{Drive Certificate}

The drive face \(\mathrm{P6}_{\mathrm{drive}}^{H^1}\) on \(\mathsf H_{\mathrm{graph}}\) is a finite typed record stating that some cohomology class \([a]\in H^1(G)\) is nonzero. A positive cycle-supported cohomological drive certificate consists of

\begin{itemize}
\tightlist
\item
  a finite \(1\)-cochain \(a\in C^1(G)\);
\item
  a finite cycle \(\gamma\in Z_1(G)\) with \(\oint_\gamma a\neq 0\);
\item
  the audit and source records required by \(\mathsf{AuditPolicy}_I\) and \(\mathsf{SourcePolicy}_I\) for the surrounding instrument.
\end{itemize}

A no-drive certificate is the corresponding \(\oint_\gamma a=0\) record on every cycle of \(G\), together with the same audit and source data.

A generic \(P_6\) audit record is not a drive certificate: it does not record cycle integrals or cohomology data. The separation between the audit face \(P_6\) and the drive face \(\mathrm{P6}_{\mathrm{drive}}^{H^1}\) is preserved at the certificate level throughout this section.

\subsubsection{Gating}

A \(P_2\)-style gate on \(\mathsf H_{\mathrm{graph}}\) is a finite admissible map \(G\mapsto G'\) obtained by deleting a finite set of edges or by restricting to a finite induced subgraph. A forest-producing gate is a gate whose output \(G'\) has \(\beta_1(G')=0\). Gates are recorded as finite typed records on the surrounding judgment; an admissible gate is one that conforms to \(\mathsf{ScopeAndPolicy}\) data of the surrounding instrument.

\subsection{Theorem 15: Forest No-Drive}

\subsubsection{Statement}

\textbf{Theorem 15 (ForestNoDrive; theorem-grade).} For a finite graph \(G\) on \(\mathsf H_{\mathrm{graph}}\), if \(\beta_1(G)=0\), then

\[
H^1(G)\;=\;0.
\]

In particular, no cycle-supported cohomological drive certificate \(\mathrm{P6}_{\mathrm{drive}}^{H^1}\) is admissible on \(G\).

\subsubsection{Proof}

The first Betti number \(\beta_1(G)\) is the dimension over a coefficient field, or the corresponding free rank over a declared free coefficient host, of the cycle space \(Z_1(G)\) from Subsection 8.1.1. If \(\beta_1(G)=0\), then \(Z_1(G)=0\): the graph has no nontrivial cycles in the cycle space. The first cohomology

\[
H^1(G)\;=\;\ker(d_1)/\mathrm{im}(d_0)
\]

has dimension, or free rank in the declared free-host case, equal to \(\beta_1(G)\) by the standard finite cycle/cochain duality on a finite graph. Hence \(H^1(G)=0\).

With \(H^1(G)=0\) and \(Z_1(G)=0\), every cycle integral \(\oint_\gamma a\) is over an empty cycle (or a sum of trivial cycles) and vanishes, so no nonzero cycle affinity is recorded. Hence no cycle-supported drive certificate is admissible: a positive certificate would require some \(\gamma\in Z_1(G)\) with \(\oint_\gamma a\neq 0\), but no such cycle exists.

The instrument-indexed form is

\[
\Gamma\,;\;\mathsf H_{\mathrm{graph}}\,;\;I_{\mathrm{graph}}\;\vdash\;\textsc{ForestNoDrive}\;:\;\texttt{accepted}.
\]

\subsection{Theorem 16: Exact-Form Null-Drive}

\subsubsection{Statement}

\textbf{Theorem 16 (ExactFormNullDrive; theorem-grade).} Let \(G\) be a finite graph on \(\mathsf H_{\mathrm{graph}}\). If a finite \(1\)-cochain \(a\in C^1(G)\) is exact, that is,

\[
a\;=\;d\phi\quad\text{for some }0\text{-cochain }\phi\in C^0(G),
\]

then for every cycle \(\gamma\in Z_1(G)\),

\[
\oint_\gamma a\;=\;0.
\]

\subsubsection{Proof}

Let \(\gamma\in Z_1(G)\) be a cycle with oriented edges \(e_1,\ldots,e_n\), where \(e_k\) goes from vertex \(v_{k-1}\) to vertex \(v_k\), and the cycle closes up at \(v_n=v_0\). By definition of the coboundary \(d\phi(e)=\phi(\mathrm{head}(e))-\phi(\mathrm{tail}(e))\), the cycle integral telescopes:

\[
\oint_\gamma a\;=\;\sum_{k=1}^{n}\,a(e_k)\;=\;\sum_{k=1}^{n}\bigl(\phi(v_k)\;-\;\phi(v_{k-1})\bigr)\;=\;\phi(v_n)\;-\;\phi(v_0)\;=\;0,
\]

since \(v_n=v_0\) on the cycle. The instrument-indexed form is

\[
\Gamma\,;\;\mathsf H_{\mathrm{graph}}\,;\;I_{\mathrm{graph}}\;\vdash\;\textsc{ExactFormNullDrive}\;:\;\texttt{accepted}.
\]

The result records that exact \(1\)-cochains carry no cycle drive: any cohomology class supported on an exact cochain is zero, and the corresponding drive certificate is empty on every cycle.

\subsection{Theorem 17: Nonzero-Affinity Equivalence}

\subsubsection{Statement}

\textbf{Theorem 17 (NonzeroAffinityEquivalence; theorem-grade).} Let \(G\) be a finite graph on \(\mathsf H_{\mathrm{graph}}\) with the fixed coefficient field or declared free coefficient host of Subsection 8.1.1, and let \(a\in C^1(G)\) with cohomology class \([a]\in H^1(G)\). Then

\[
[a]\;\neq\;0\;\;\Longleftrightarrow\;\;\exists\,\gamma\in Z_1(G)\;:\;\oint_\gamma a\;\neq\;0.
\]

\subsubsection{Proof}

For the reverse direction, suppose some \(\gamma\in Z_1(G)\) satisfies \(\oint_\gamma a\neq 0\). By the contrapositive of Theorem 16, \(a\) is not exact: if \(a=d\phi\) for some \(\phi\), then \(\oint_\gamma a=0\), contradicting \(\oint_\gamma a\neq 0\). Hence \([a]\neq 0\in H^1(G)\).

For the forward direction, suppose \([a]\neq 0\). The cycle/cochain pairing on a finite graph with the fixed coefficient field or declared free coefficient host is non-degenerate: the cycle space \(Z_1(G)\) and the cohomology \(H^1(G)\) have matching finite dimension or free rank \(\beta_1(G)\), and the bilinear pairing

\[
\langle\,\cdot\,,\,\cdot\,\rangle\colon Z_1(G)\times H^1(G)\to\text{coefficient host},\qquad \langle\gamma,[b]\rangle\;=\;\oint_\gamma b,
\]

is well-defined (independent of representative \(b\) of \([b]\), by Theorem 16) and non-degenerate. If every cycle integral were zero --- that is, \(\oint_\gamma a=0\) for all \(\gamma\in Z_1(G)\) --- then \([a]\) would pair trivially with every cycle class, hence \([a]=0\) by non-degeneracy of the pairing, contradicting the hypothesis. Therefore some cycle \(\gamma\) has \(\oint_\gamma a\neq 0\).

The instrument-indexed form is

\[
\Gamma\,;\;\mathsf H_{\mathrm{graph}}\,;\;I_{\mathrm{graph}}\;\vdash\;\textsc{NonzeroAffinityEquivalence}\;:\;\texttt{accepted}.
\]

The theorem records the equivalence between the cohomological drive certificate \([a]\neq 0\) and the cycle-affinity drive certificate \(\oint_\gamma a\neq 0\) on \(\mathsf H_{\mathrm{graph}}\) with the declared coefficients; both carry the same cycle-supported cohomological drive face \(\mathrm{P6}_{\mathrm{drive}}^{H^1}\).

\subsection{Theorem 18: Gating-Affinity Suppression}

\subsubsection{Statement}

\textbf{Theorem 18 (GatingAffinitySuppression; theorem-grade).} Let \(G\) be a finite graph on \(\mathsf H_{\mathrm{graph}}\), and let \(G\mapsto G'\) be a \(P_2\)-style gate that produces a forest, that is,

\[
\beta_1(G')\;=\;0.
\]

Then

\[
H^1(G')\;=\;0,
\]

and the cycle-supported cohomological drive certificate \(\mathrm{P6}_{\mathrm{drive}}^{H^1}\) is absent-with-record on \(G'\).

\subsubsection{Proof}

By Theorem 15 applied to \(G'\), the cohomology \(H^1(G')\) vanishes when \(\beta_1(G')=0\). Since the drive certificate \(\mathrm{P6}_{\mathrm{drive}}^{H^1}\) requires a nonzero cohomology class \([a]\in H^1(G')\), no such certificate is admissible on \(G'\). The drive certificate record on \(G'\) has status \(\texttt{absent\_with\_record}\) in the relevant role-channel sense: the gate produced a graph on which no cycle supports a drive.

The instrument-indexed form is

\[
\Gamma\,;\;\mathsf H_{\mathrm{graph}}\,;\;I_{\mathrm{graph}}\;\vdash\;\textsc{GatingAffinitySuppression}\;:\;\texttt{accepted}.
\]

\subsubsection{Nonclaim}

The theorem concerns cycle-supported cohomological drive on a declared finite graph and gate. It does not classify all drive, all thermodynamic force, all entropy production, or all directionality. A drive notion that does not reduce to a cycle-cohomology certificate on a finite graph is outside the scope of this theorem; its adjudication requires a separate drive bridge.

The dependency

\[
P_3\;\mathrel{\mathsf{dep}_{\texttt{requires\_bridge}}}\;\mathrm{P6}_{\mathrm{drive}}
\]

of Subsection 4.6.2 remains active: route mismatch alone does not establish a drive certificate, even when route mismatch is recorded as a \(P_3\) defect. NC-12 of Subsection 2.4.3 records the corresponding nonclaim.
 \section{Countermodel Atlas}

This section records the six countermodels of the first-paper atlas. Each countermodel is a finite explicit construction together with the no-go statement it supports, the overclaim it blocks, and any positive bridge that would have to be supplied for a positive claim. The full atlas of twelve countermodels is recorded in Appendix C; the six countermodels in this section are the load-bearing constructions used by Sections 6 through 11.

\subsection{Atlas Convention}

This subsection fixes the sheet form for countermodels and selects which countermodels appear in the main text. It is a schema-grade declaration, with the actual constructions in Subsections 9.2 through 9.7.

\subsubsection{Countermodel Sheet Form}

Each countermodel sheet contains the following five fields, in this order:

\begin{itemize}
\tightlist
\item
  \textbf{Construction.} A finite explicit construction on a named host, with the data of the calculus (carriers, maps, packages, profiles, witnesses, updates, defects, audit and source records) made explicit. The construction is finite at every step.
\item
  \textbf{Supported no-go.} The dependency judgment \(S\mathrel{\mathsf{dep}_r}X\) that the construction supports, in the form fixed by Subsection 4.6.1, with relation type \(r\) drawn from the eight admissible types \(\texttt{suffices}\), \(\texttt{insufficient}\), \(\texttt{destroys}\), \(\texttt{blocks}\), \(\texttt{enables}\), \(\texttt{equivalent}\), \(\texttt{requires\_bridge}\), and \(\texttt{outside\_scope}\).
\item
  \textbf{What it blocks.} The overclaim that the no-go rules out, named together with the corresponding nonclaim NC-\(n\) from Subsection 2.4.3.
\item
  \textbf{What it does not block.} The positive content that the construction is consistent with, so that the no-go is not over-read into a stronger negative statement.
\item
  \textbf{Required positive bridge.} When a positive claim of the form \(S\Rightarrow X\) would need an additional bridge, the bridge type is named (for example, a drive bridge \(B_{3\to 6}^{\mathrm{drive}}\) or a gluing bridge \(B_{\mathrm{glue}}\)).
\end{itemize}

Internally, the packet records a countermodel sheet as a finite tuple

\[
(\mathrm{id},H,\mathsf{Data},\mathsf{ActiveClaim},\mathsf{FailedClaim},W,A,V,\mathcal N),
\]

with the five prose fields above unpacking that tuple for the reader.

\subsubsection{Main-Text Selection}

The six countermodels of this section are the load-bearing first-paper atlas:

\begin{itemize}
\tightlist
\item
  \(\mathsf{CM}_1\) --- Active package but no macro closure (Subsection 9.2);
\item
  \(\mathsf{CM}_2\) --- Active route/holonomy but no drive (Subsection 9.3);
\item
  \(\mathsf{CM}_3\) --- Local packages fail globally (Subsection 9.4);
\item
  \(\mathsf{CM}_4\) --- Directed cell active but pair observable blocked (Subsection 9.5);
\item
  \(\mathsf{CM}_5\) --- Refinement worsens closure (Subsection 9.6);
\item
  \(\mathsf{CM}_6\) --- Same-level self-audit circularity (Subsection 9.7).
\end{itemize}

The full twelve-countermodel atlas, including same-primitive-pair-different-statuses (\(\mathsf{CM}_7\)), idempotent-completions-do-not-commute (\(\mathsf{CM}_8\)), strict-extension-without-closure (\(\mathsf{CM}_9\)), strict-extension-without-drive (\(\mathsf{CM}_{10}\)), fallback-source (\(\mathsf{CM}_{11}\)), and gating-destroys-cycle-supported-drive (\(\mathsf{CM}_{12}\)), is recorded with full constructions in Appendix C. Subsection 9.8 names these appendix-only countermodels and points to Appendix C for the constructions.

\subsection{Countermodel 1: Active Package but No Macro Closure}

\subsubsection{Construction}

Take \(X=\{a,b,c\}\) and \(Y=\{A,B\}\), both finite sets on \(\mathsf H_{\mathrm{fin}}\). Define a finite package map

\[
q:X\to Y,\qquad q(a)=A,\quad q(b)=A,\quad q(c)=B.
\]

The map \(q\) is admissible as a \(P_5\)-package on \(X\) in the schema of Subsection 3.3: it is a finite map on a finite carrier with a finite codomain.

Define a finite update

\[
F:X\to X,\qquad F(a)=a,\quad F(b)=c,\quad F(c)=c.
\]

Then \(q(a)=q(b)=A\), while \(qF(a)=q(a)=A\) and \(qF(b)=q(c)=B\), so

\[
qF(a)\;\neq\;qF(b).
\]

By Theorem 7, \((a,b)\in\delta_1^{\mathrm{split}}(q,F)\), so

\[
\delta_1^{\mathrm{split}}(q,F)\;\neq\;\varnothing,
\]

and no finite descended update \(F^\sharp:Y\to Y\) with \(qF=F^\sharp q\) exists.

\subsubsection{Supported No-Go}

The construction supports the dependency judgment

\[
\Gamma\,;\;\mathsf H_{\mathrm{fin}}\,;\;I_{\mathrm{fin}}\;\vdash\;P_5(q)\;\mathrel{\mathsf{dep}_{\texttt{insufficient}}}\;P_1(qF=F^\sharp q),
\]

that is, \(P_5\not\Rightarrow P_1\text{ macro closure}\) on the host \(\mathsf H_{\mathrm{fin}}\) under the declared admissible instrument.

\subsubsection{What It Blocks}

The overclaim ``a package automatically closes the dynamics'' is rejected. NC-7 of Subsection 2.4.3 records the corresponding nonclaim. The construction is the working content of NC-7: package admissibility on a finite carrier does not by itself produce a descended update.

\subsubsection{What It Does Not Block}

The construction does not say packages are useless or that descent is generically impossible. It says only that package existence alone does not guarantee descent of a given update, and a positive descent claim requires an explicit witness: in this construction, replacing \(F\) with a different finite update consistent with the fibers of \(q\) would produce an admissible descended map.

\subsubsection{Required Positive Bridge}

A positive descent claim \(P_5\Rightarrow P_1\) on a specific package \(q\) and update \(F\) requires an explicit descent record: the empty defect \(\delta_1^{\mathrm{split}}(q,F)=\varnothing\) together with the audit and visibility data required by the surrounding judgment. Theorem 7 makes the vanishing of \(\delta_1^{\mathrm{split}}(q,F)\) necessary and sufficient for the descended update, while the audit and visibility records remain separate admissibility requirements.

\subsection{Countermodel 2: Active Route / Holonomy but No Drive}

\subsubsection{Construction}

Take a finite cycle graph \(G\) on \(\mathsf H_{\mathrm{graph}}\), for example the triangle \(0\to 1\to 2\to 0\) with \(\beta_1(G)=1\) and a single nontrivial cycle \(\gamma\). Equip the surrounding judgment with a \(P_3\)-route record asserting nontrivial holonomy around \(\gamma\):

\[
\operatorname{Hol}(\gamma)\;\neq\;\mathrm{id},
\]

so that the role channel of \(P_3\) is admissible with active route data.

For the drive face, take the zero \(1\)-cochain \(a=0\in C^1(G)\). By Theorem 16 (or directly from \(a=0\)), every cycle integral vanishes:

\[
\oint_\gamma a\;=\;0\quad\text{for every cycle }\gamma\in Z_1(G).
\]

Hence \([a]=0\in H^1(G)\), and no positive cycle-supported cohomological drive certificate \(\mathrm{P6}_{\mathrm{drive}}^{H^1}\) is admissible on \(G\) for this \(a\). The route/holonomy data is active, but the drive certificate is empty.

\subsubsection{Supported No-Go}

The construction supports the dependency judgment

\[
\Gamma\,;\;\mathsf H_{\mathrm{graph}}\,;\;I_{\mathrm{graph}}\;\vdash\;P_3^{\mathrm{holonomy}}\;\mathrel{\mathsf{dep}_{\texttt{requires\_bridge}}}\;\mathrm{P6}_{\mathrm{drive}}^{H^1},
\]

recording that route mismatch or holonomy alone does not establish a cohomological drive certificate.

\subsubsection{What It Blocks}

The overclaim ``route mismatch or holonomy automatically certifies drive'' is rejected. NC-12 of Subsection 2.4.3 records the corresponding nonclaim. The construction is consistent with active \(P_3\) data and absent \(\mathrm{P6}_{\mathrm{drive}}^{H^1}\) data on the same graph.

\subsubsection{What It Does Not Block}

The construction does not say drive certificates are unattainable on \(G\) or that route mismatch is incompatible with drive. A different choice of \(1\)-cochain \(a\) on \(G\) --- one with \(\oint_\gamma a\neq 0\) --- would supply a positive drive certificate; the construction shows only that the choice of \(a\) is independent data from the route/holonomy record.

\subsubsection{Required Positive Bridge}

A positive inference of the form \(P_3\Rightarrow\mathrm{P6}_{\mathrm{drive}}^{H^1}\) requires an explicit drive bridge

\[
B_{3\to 6}^{\mathrm{drive}},
\]

a finite typed record in \(\mathsf{BridgePolicy}_I\) that connects the route/holonomy data to a specific cochain \(a\) with \([a]\neq 0\). The bridge supplies the missing data: the route witness alone does not produce \(a\), and the choice of \(a\) is the substantive content of the bridge.

\subsection{Countermodel 3: Local Packages Fail Globally}

\subsubsection{Construction}

Take a finite carrier \(Z=\{a,b,c\}\) on \(\mathsf H_{\mathrm{fin}}\) with a finite cover \(\{U,V\}\) given by

\[
U\;=\;\{a,b\},\qquad V\;=\;\{b,c\},
\]

so that the overlap is \(U\cap V=\{b\}\). Let local packages be finite maps

\[
q_U:U\to O_U,\qquad q_V:V\to O_V,
\]

each admissible on its own patch. Suppose the overlap target distinguishes two incompatible values \(0\) and \(1\), and that the restrictions to the overlap satisfy

\[
\mathrm{res}_{U,U\cap V}(q_U(b))\;=\;0,\qquad \mathrm{res}_{V,U\cap V}(q_V(b))\;=\;1,
\]

with no admissible identification of \(0\) and \(1\) in the overlap target. Each local package is admissible on its own patch, but no global package \(q:Z\to O\) can extend both: any admissible global \(q\) would assign a single value to \(b\) on the overlap, but the local data demand two incompatible values there.

The gluing defect \(\delta_4^{\mathrm{glue}}\) on the cover \(\{U,V\}\) is therefore nonempty, with \(b\) as a witness.

\subsubsection{Supported No-Go}

The construction supports the dependency judgment

\[
\Gamma\,;\;\mathsf H_{\mathrm{fin}}\,;\;I_{\mathrm{fin}}\;\vdash\;\bigl\{P_5(q_U),\,P_5(q_V)\bigr\}\;\mathrel{\mathsf{dep}_{\texttt{insufficient}}}\;P_5(q_{\mathrm{global}}),
\]

recording that local-\(P_5\)-validity does not by itself imply global package existence on the cover.

\subsubsection{What It Blocks}

The overclaim ``local validity implies global validity'' is rejected. NC-14 of Subsection 2.4.3 records the corresponding nonclaim. The construction is consistent with two admissible local packages and an absent global package on the same finite carrier.

\subsubsection{What It Does Not Block}

The construction does not say global packaging is generically impossible or that covers are useless. Different overlap data --- with admissibly identifiable values on \(U\cap V\) --- would produce a global package; the construction shows only that the overlap data is part of the global packaging input.

\subsubsection{Required Positive Repair}

A positive global-packaging claim from local data requires a \(P_4\)-gluing bridge

\[
B_{\mathrm{glue}}\quad\text{with}\quad \delta_4^{\mathrm{glue}}\;=\;\varnothing,
\]

together with a \(P_6\)-audit compatibility record on the overlaps. Absent these, the gluing defect blocks the global claim, and Subsection 5.3.4 records the gluing-defect schema used by \(B_{\mathrm{glue}}\).

\subsection{Countermodel 4: Directed Cell Active but Pair Observable Blocked}

\subsubsection{Construction}

Take the primitive pair \(P_6\leftarrow P_3\) in \(\mathbf{BirdInt}^{\mathrm{aud}}_{\mathrm{fin}}\), with the directed-cell schema of Subsection 4.3. Equip the cell with witness and update data

\[
W_3\;=\;\mathsf{RouteMismatchWit},\qquad U_6\;=\;\mathsf{AddAuditRecord},
\]

drawn from \(\mathsf{Wit}(P_3)\) and \(\mathsf{Upd}(P_6)\) respectively. Suppose the witness is visible under the surrounding instrument \(I\) and audited under \(\mathsf{AuditPolicy}_I\), and that the cell's profile \(\lambda_{\mathrm{audit}}\) records audit-update mode. The classifier of Subsection 5.6.2 then assigns

\[
\mathsf{Cell}_{63}^{\lambda_{\mathrm{audit}}}\;:\;\texttt{action}.
\]

Now consider the unordered pair-observable judgment

\[
\mathsf{PairObs}_{\{3,6\}}^{\mu_{\mathrm{drive}}}
\]

with profile \(\mu_{\mathrm{drive}}\) that asserts the \(P_3\)-witness certifies a drive observable. Suppose no drive bridge is present: \(B_{3\to 6}^{\mathrm{drive}}\) is absent from \(\mathcal L\), and no admissible drive source-of-truth certifying \([a]\neq0\) is recorded under \(\mu_{\mathrm{drive}}\). The cell is well-formed and active, but the pair record fails the source-and-bridge check; the pair classifier of Subsection 5.6.3 assigns

\[
\mathsf{PairObs}_{\{3,6\}}^{\mu_{\mathrm{drive}}}\;:\;\texttt{blocked}.
\]

The same primitive pair \(\{3,6\}\) thus carries an active directed cell and a blocked pair observable simultaneously.

\subsubsection{Supported No-Go}

The construction supports the dependency judgment

\[
\Gamma\,;\;\mathbf{BirdInt}^{\mathrm{aud}}_{\mathrm{fin}}\,;\;I_{\mathrm{fin}}\;\vdash\;\mathsf{Cell}_{63}^{\lambda_{\mathrm{audit}}}:\texttt{action}\;\mathrel{\mathsf{dep}_{\texttt{insufficient}}}\;\mathsf{PairObs}_{\{3,6\}}^{\mu_{\mathrm{drive}}}:\texttt{real},
\]

recording that directed-cell action does not by itself imply pair-observable realness on the same primitive pair.

\subsubsection{What It Blocks}

The overclaim ``directed-cell action implies pair-observable realness'' is rejected. NC-5 of Subsection 2.4.3 records the corresponding nonclaim, and Theorem 5 of Section 6 makes the cell-to-pair non-implication formal as part of the status-family separation theorem.

\subsubsection{What It Does Not Block}

The construction does not say pair observables are unattainable on \(\{3,6\}\) or that audit cells are incompatible with pair realness. A different pair record on the same primitives, with an admissible source and the required drive bridge, may classify \(\texttt{real}\) or \(\texttt{real\_provisional}\) in the schema of Subsection 4.4.3.

\subsubsection{Required Positive Bridge}

A positive pair-observable claim requires its own typed records: a real or fallback source under \(\mathsf{SourcePolicy}_I\), an admissible branch-compatibility record, the visibility, threshold, and audit data of the pair record, and (when the pair claims drive content) the drive bridge \(B_{3\to 6}^{\mathrm{drive}}\) under \(\mathsf{BridgePolicy}_I\). The directed-cell record contributes none of these; the pair record carries them as separate data.

\subsection{Countermodel 5: Refinement Worsens Closure}

\subsubsection{Construction}

Take a finite carrier \(X=\{a,b,c,d\}\) on \(\mathsf H_{\mathrm{fin}}\). Define a coarse abstraction

\[
q_0\colon X\to\{*\},\qquad q_0(x)\;=\;*\;\text{ for every }x\in X,
\]

so that \(q_0\) has a single fiber and every finite update on \(X\) descends through \(q_0\) trivially.

Now define a refinement

\[
q_1\colon X\to\{A,B,C\},\qquad q_1(a)=A,\quad q_1(b)=A,\quad q_1(c)=B,\quad q_1(d)=C.
\]

The coarse abstraction \(q_0\) factors through \(q_1\) via the unique forgetting map \(r:\{A,B,C\}\to\{*\}\), so \(q_0=rq_1\). Thus the package data of \(q_1\) is more refined than that of \(q_0\) in the sense of Subsection 3.3.

Now choose a finite update \(F:X\to X\) that creates a \(q_1\)-split. For instance, take

\[
F(a)\;=\;c,\quad F(b)\;=\;d,\quad F(c)\;=\;c,\quad F(d)\;=\;d.
\]

Then \(q_1(a)=q_1(b)=A\), but \(q_1F(a)=q_1(c)=B\) and \(q_1F(b)=q_1(d)=C\), so

\[
q_1F(a)\;\neq\;q_1F(b),
\]

and \((a,b)\in\delta_1^{\mathrm{split}}(q_1,F)\neq\varnothing\). By Theorem 7, no descended update through \(q_1\) exists.

By contrast, every update on \(X\) --- including this one --- descends through \(q_0\), since \(q_0\) has a single fiber and the descended map on \(\{*\}\) is the identity. Thus the closure deficit on \(q_0\) is empty, while the closure deficit on \(q_1\) is non-empty: the refinement from \(q_0\) to \(q_1\) has worsened the closure status of the same update.

\subsubsection{Supported No-Go}

The construction supports the dependency judgment

\[
\Gamma\,;\;\mathsf H_{\mathrm{fin}}\,;\;I_{\mathrm{fin}}\;\vdash\;P_4(q_0\leadsto q_1)\;\mathrel{\mathsf{dep}_{\texttt{insufficient}}}\;\bigl(\delta_1^{\mathrm{split}}(q_1,F)\;\le\;\delta_1^{\mathrm{split}}(q_0,F)\bigr),
\]

recording that \(P_4\)-refinement does not by itself improve closure: a refinement can introduce a non-empty \(\delta_1^{\mathrm{split}}\) where there was none before.

\subsubsection{What It Blocks}

The overclaim ``finer is always better'' is rejected. NC-13 of Subsection 2.4.3 records the corresponding nonclaim. The construction is the working content of NC-13: refinement can introduce closure defects rather than improve them.

\subsubsection{What It Does Not Block}

The construction does not say refinement is always harmful or that closure improvement under refinement is impossible. A different finite update \(F'\) --- one whose action is consistent with the fibers of \(q_1\) --- would descend through \(q_1\) as well as \(q_0\). The construction shows only that compatibility of \(F\) with \(q_1\) is independent of compatibility with \(q_0\).

\subsubsection{Required Positive Bridge}

A positive refinement-improvement theorem requires monotonicity hypotheses on the refinement, the update, and the closure functional, recorded as a \(P_4\)-monotonicity bridge. Such hypotheses are recorded in the future-work boundary of Subsection 16.3.1 and are not asserted here.

\subsection{Countermodel 6: Same-Level Self-Audit Circularity}

\subsubsection{Construction}

Let \(I_0\) be a finite admissible instrument under an admissible package \(\mathcal T_0\) on the host \(\mathsf H_{\mathrm{instr}}\), with the instrument record schema of Subsection 3.4.1. Consider the candidate same-level self-soundness claim

\[
\Gamma\,;\;\mathcal T_0\,;\;I_0\;\vdash\;\mathrm{Sound}(I_0)\;:\;?,
\]

asserting that \(I_0\) certifies its own soundness without a level shift. Two cases arise from the classifier of Subsection 5.6.5.

\subsubsection{Cases}

\textbf{Outside-scope case.} Suppose the soundness claim is not in \(\mathsf{ClaimTypes}(I_0)\), that is, \(I_0\) is not configured to evaluate claims about its own soundness. The classifier assigns

\[
\Gamma\,;\;\mathcal T_0\,;\;I_0\;\vdash\;\mathrm{Sound}(I_0)\;:\;\texttt{outside\_scope}.
\]

This is the case for instruments configured to evaluate object-level claims only.

\textbf{Circular case.} Suppose the soundness claim is syntactically in \(\mathsf{ClaimTypes}(I_0)\) but its acceptance under \(I_0\) depends on \(I_0\)'s own soundness verdict --- the audit policy required to certify \(\mathrm{Sound}(I_0)\) consults the very verdict being computed. The classifier assigns

\[
\Gamma\,;\;\mathcal T_0\,;\;I_0\;\vdash\;\mathrm{Sound}(I_0)\;:\;\texttt{undefined\_circular}.
\]

This is the case when an instrument is asked to certify a property of itself whose verification cannot be carried out without already knowing the verdict.

In neither case does \(I_0\) certify its own soundness at \(\texttt{accepted}\) without a level shift.

\subsubsection{Supported No-Go}

The construction supports the dependency judgment

\[
\Gamma\,;\;\mathsf H_{\mathrm{instr}}\,;\;I_{\mathrm{instr}}\;\vdash\;I_0\text{ audits }\mathcal T_0\;\mathrel{\mathsf{dep}_{\texttt{insufficient}}}\;\bigl(I_0\vdash\mathrm{Sound}(I_0):\texttt{accepted}\bigr),
\]

recording that an instrument's audit capacity for object-level content is insufficient to certify its own soundness at the same level.

\subsubsection{What It Blocks}

The overclaim ``an instrument can certify its own complete soundness without a level shift'' is rejected. NC-16 of Subsection 2.4.3 records the corresponding nonclaim, and Theorem 24 of Section 11 makes the same-level self-audit failure formal.

\subsubsection{What It Does Not Block}

The construction does not say self-audit data is unattainable or that instruments cannot record any self-referential information. Object-level audit updates within \(P_6\leftarrow P_6\) (auditing a specific object-level record) and fixed audit refreshes (re-running an existing audit policy on existing records) remain admissible: Subsection 11.6 records the three-case separation between object-level audit, fixed audit refresh, and complete same-level self-audit. Only the third case is blocked or circular.

\subsubsection{Required Positive Repair}

A positive same-level soundness claim requires a level shift: a higher instrument \(I_1\) that audits \(I_0\) via an admissible audit bridge \(B^{\mathrm{audit}}_{0\to 1}\) in the rotating-audit chain of Subsection 11.7. The chain is finite and admissible; Theorem 25 of Section 11 records that an extension instrument can audit the lower stack but does not certify its own soundness, so the rotating audit is non-final (NC-17 of Subsection 2.4.3).

\subsection{Additional Countermodels for Appendix}

The remaining six countermodels of the full atlas are recorded with their full constructions and dependency judgments in Appendix C; they are named here for cross-reference. Each uses the sheet form of Subsection 9.1.1.

\subsubsection{Same Primitive Pair, Different Statuses}

\(\mathsf{CM}_7\) exhibits two directed-cell records on the same primitive pair \(P_6\leftarrow P_3\): one with audit-update profile classified \(\texttt{action}\), the other with drive-certification profile and missing drive bridge classified \(\texttt{blocked}\). Supports NC-3 (no primitive-pair status determination); the construction is the working content of Theorem 1 of Section 6. Full sheet in Appendix C.

\subsubsection{Idempotent Completions Do Not Commute}

\(\mathsf{CM}_8\) exhibits two finite idempotent completion operators \(E_1,E_2\) on the power set \(\mathcal P(\{a,b,c\})\) that fail to commute. Supports NC-8 (no idempotence-implies-commutation); the construction is the working content of Theorem 9 of Section 7. Full sheet in Appendix C.

\subsubsection{Strict Extension Without Macro Closure}

\(\mathsf{CM}_9\) exhibits a finite promotion bridge with strict object-map nonfactorization \(\pi_{j+1}\not\factor\pi_j\) and no descended dynamics on the target package. Supports NC-10 (no strictness-implies-closure). Full sheet in Appendix C.

\subsubsection{Strict Extension Without Drive}

\(\mathsf{CM}_{10}\) exhibits a finite strict promotion bridge whose target host is a forest, so \(H^1=0\) and no admissible cohomological drive certificate exists. Supports NC-11 (no strictness-implies-drive). Full sheet in Appendix C.

\subsubsection{Fallback Source Cannot Support Real Pair}

\(\mathsf{CM}_{11}\) exhibits a pair-observable record whose only available source has type \(\texttt{fallback}\) under a profile that does not upgrade fallback. The classifier of Subsection 5.6.3 assigns at most \(\texttt{real\_provisional}\), never \(\texttt{real}\). Supports the source-of-truth discipline of Subsection 4.4.3. Full sheet in Appendix C.

\subsubsection{Gating Destroys Cycle-Supported Drive}

\(\mathsf{CM}_{12}\) exhibits a triangle graph with a nonzero drive cochain and a \(P_2\)-style gate that deletes one edge, leaving a path with \(\beta_1=0\) and hence \(H^1=0\). Supports the gating-suppression discipline of Section 8 and is the working content of Theorem 18. Full sheet in Appendix C.
 \section{Promotion and Stacking}

This section gives the finite promotion-bridge theory in scoped form. Subsection 10.1 records the promotion-bridge object as introduced at schema-grade in Subsection 4.5 and adds the gate record. Subsection 10.2 specifies the gate semantics. Subsections 10.3, 10.4, and 10.5 prove the three promotion theorems: gate soundness, strict-implies-nonfactorization, and non-strict-implies-factorization. Subsection 10.6 records the toy strict-bridge example used as a compact illustration. Subsection 10.7 lists what this paper does not prove about stacking.

\subsection{Promotion Bridge Object}

\subsubsection{Bridge Data}

A promotion bridge from a source theory package \(\mathcal T_j\) to a target theory package \(\mathcal T_{j+1}\) is the finite typed tuple

\[
B_{j\to j+1}\;=\;(\,\mathcal T_j,\;\mathcal T_{j+1},\;\pi_j,\;\pi_{j+1},\;L,\;V,\;\Theta,\;A,\;\delta,\;\mathcal N\,),
\]

introduced at schema-grade in Subsection 4.5.2. The expanded form recorded in \(\mathsf{Cont}\) also carries the bridge identifier, the carrier \(S\) on which the object maps act, the codomains \(O_j\) and \(O_{j+1}\), the gate-results record \(\mathsf{GateResults}(B)\), the host tag \(h_B\in\mathsf{Host}\), and the bridge profile \(\lambda_{\mathrm{prom}}\). In the packet's expanded order this is

\[
\begin{aligned}
B_{j\to j+1}\;=\;(\,&\mathrm{id},\;\mathcal T_j,\;\mathcal T_{j+1},\;S,\;O_j,\;O_{j+1},\;\pi_j,\;\pi_{j+1},\\[2pt]
&L,\;V,\;\Theta,\;A,\;\delta,\;\mathsf{GateResults},\;\mathcal N,\;\mathsf{HostTag},\;\lambda_{\mathrm{prom}}\,).
\end{aligned}
\]

Every component is finite. The source and target packages are admissible in the sense of Subsection 3.3.3, the object maps \(\pi_j,\pi_{j+1}\) are finite maps with \(\pi_j:S\to O_j\) and \(\pi_{j+1}:S\to O_{j+1}\), the level/lens/interface record \(L\) fixes the levels at which the bridge operates, the visibility record \(V\) partitions the bridge's content under \(\mathsf{Visible}_I/\mathsf{Suppressed}_I/\mathsf{Outside}_I/\mathsf{Unknown}_I\), the threshold record \(\Theta\) is drawn from the threshold family of Subsection 5.2.2, the audit record \(A\) is in the schema of Subsection 4.8.2, the defect record \(\delta\) collects the strictness, no-smuggling, stability, and descent defects of Subsection 5.4, and the nonclaim record \(\mathcal N\) is finite.

\subsubsection{Strict Versus Non-Strict Promotion}

A promotion bridge is \textbf{strict} when

\[
\pi_{j+1}\;\not\factor\;\pi_j,
\]

equivalently when there exist \(s,s'\in S\) with \(\pi_j(s)=\pi_j(s')\) and \(\pi_{j+1}(s)\neq\pi_{j+1}(s')\). A promotion bridge is \textbf{non-strict} when \(\pi_{j+1}\) factors through \(\pi_j\), that is, when there exists a finite map \(\phi:\mathrm{im}(\pi_j)\to O_{j+1}\) with \(\pi_{j+1}=\phi\,\pi_j\). For the object-map data of an accepted promotion, the factorization dichotomy is exhaustive: by Theorem 13 of Section 7, the factorization defect \(\Delta_{\mathrm{fact}}(\pi_j,\pi_{j+1})\) is empty exactly when \(\pi_{j+1}\) factors through \(\pi_j\), and non-empty exactly when it does not.

Strictness is a property of the typed object-map data \(\pi_j,\pi_{j+1}\) alone; it is not by itself a closure or drive claim. The negative scope of Subsection 2.4.3 records NC-10 (no strictness-implies-closure) and NC-11 (no strictness-implies-drive); the no-go countermodels \(\mathsf{CM}_9\) and \(\mathsf{CM}_{10}\) of Section 9 supply finite explicit instances in which strictness coexists with absent closure or absent drive.

\subsection{Promotion Gates}

\subsubsection{Gate List}

The required gate family of a promotion bridge is

\[
\mathsf{PromGate}(B)\;=\;(\,G_{\mathrm{suff}},\;G_{\mathrm{desc}},\;G_{\mathrm{stab}},\;G_{\mathrm{ctrl}},\;G_{\mathrm{nosmuggle}},\;G_{\mathrm{vis}},\;G_{\mathrm{audit}},\;G_{\mathrm{strict}}\,),
\]

with the optional local-to-global gate

\[
G_{\mathrm{locglob}}.
\]

Each gate has a fixed role:

\begin{itemize}
\tightlist
\item
  \(G_{\mathrm{suff}}\): well-formedness gate, checking that \(B\) as a record satisfies the typed schema of Subsection 10.1.1.
\item
  \(G_{\mathrm{desc}}\): descent gate, required when the bridge claims macro closure or descended dynamics on the target package; it is fed by the descent defect of Subsection 5.3.1 (\(\delta_1^{\mathrm{split}}\)) and Theorem 7.
\item
  \(G_{\mathrm{stab}}\): stability gate, required when the bridge claims target objects are fixed points of a completion endomap; it is fed by the stability defect \(\delta_{\mathrm{stab}}\) of Subsection 5.4.3.
\item
  \(G_{\mathrm{ctrl}}\): control-of-content gate, checking that no target-layer data has been used to certify target novelty.
\item
  \(G_{\mathrm{nosmuggle}}\): no-smuggling gate, fed by the no-smuggling defect of Subsection 5.4.2.
\item
  \(G_{\mathrm{vis}}\): visibility gate, checking that the bridge's content is admissible under \(\mathsf{Visible}_I\) and the bridge policy of \(I\).
\item
  \(G_{\mathrm{audit}}\): audit gate, checking that the audit record \(A\) passes \(\mathsf{AuditPolicy}_I\).
\item
  \(G_{\mathrm{strict}}\): strictness gate, with refined status set \(\mathrm{StrictGateStatus}\) from Subsection 5.1.6, fed by the factorization defect \(\Delta_{\mathrm{fact}}\) of Subsection 5.4.1.
\item
  \(G_{\mathrm{locglob}}\): optional local-to-global gate, present when the bridge claims a global package built from local packages on a cover.
\end{itemize}

The required core gates of a bridge are the subset

\[
\mathsf{ReqCoreGates}(B)\;\subseteq\;\mathsf{CoreGates}\;=\;\{\,G_{\mathrm{suff}},\,G_{\mathrm{desc}},\,G_{\mathrm{stab}},\,G_{\mathrm{ctrl}},\,G_{\mathrm{nosmuggle}},\,G_{\mathrm{vis}},\,G_{\mathrm{audit}}\,\}
\]

whose roles the bridge actually claims. The strictness gate \(G_{\mathrm{strict}}\) is part of the promotion-gate record but sits outside \(\mathsf{CoreGates}\); when strictness is required and classified, it records the strict/non-strict refinement, while a bridge whose strictness is not required or not classified may remain at the unrefined status \(\texttt{accepted}\).

\subsubsection{Gate Statuses}

Each ordinary gate carries a status from the packet status set

\[
\mathrm{GateStatus}\;=\;\{\,\texttt{pass},\;\texttt{fail},\;\texttt{not\_required},\;\texttt{not\_checked},\;\texttt{outside\_scope}\,\}.
\]

The strictness gate carries a refined status from

\[
\mathrm{StrictGateStatus}\;=\;\{\,\texttt{strict\_pass},\;\texttt{non\_strict\_pass},\;\texttt{not\_required},\;\texttt{not\_checked},\;\texttt{bad\_audit},\;\texttt{outside\_scope}\,\}.
\]

The gate-results record \(\mathsf{GateResults}(B)\) is the finite assignment of these statuses to the gates in the bridge record. A gate may be marked \(\texttt{not\_required}\) when the bridge does not claim the corresponding content; the descent gate, for example, is \(\texttt{not\_required}\) for a bridge that records only an interface change with no closure claim. A required core gate marked \(\texttt{not\_checked}\) does not yield an accepted promotion; the promotion classifier of Subsection 5.6.4 routes such a bridge to \(\texttt{candidate}\).

\subsection{Theorem 19: Promotion Gate Soundness}

\subsubsection{Statement}

\textbf{Theorem 19 (PromotionGateSoundness; theorem-grade).} Let \(B_{j\to j+1}\) be a promotion bridge. If the promotion classifier of Subsection 5.6.4 assigns

\[
\mathrm{Promote}(B):\sigma\quad\text{with}\quad\sigma\in\{\,\texttt{accepted},\,\texttt{strict},\,\texttt{non\_strict}\,\},
\]

then every required core gate exists in the gate-results record and passes:

\[
\forall\,G\in\mathsf{ReqCoreGates}(B),\quad\mathsf{GateExists}(B,G)\;\wedge\;\mathsf{GatePass}(B,G).
\]

\subsubsection{Proof}

By the definition of the promotion classifier in Subsection 5.6.4, the accepted core of a promotion bridge is exactly the condition that the bridge is admissible and every required core gate exists and passes. The verdicts \(\texttt{strict}\) and \(\texttt{non\_strict}\) are refinements of the accepted core by the strictness gate: \(\texttt{strict}\) adds \(G_{\mathrm{strict}}=\texttt{strict\_pass}\), and \(\texttt{non\_strict}\) adds \(G_{\mathrm{strict}}=\texttt{non\_strict\_pass}\). The remaining promotion-status values record specific failure or unresolved conditions and are routed by the classifier to non-accepted statuses such as \(\texttt{failed\_audit}\), \(\texttt{failed\_no\_smuggling}\), \(\texttt{failed\_descent}\), \(\texttt{failed\_stability}\), \(\texttt{globally\_obstructed}\), \(\texttt{local\_only}\), \(\texttt{candidate}\), or \(\texttt{outside\_scope}\). Hence whenever the verdict lies in the accepted family, every required core gate exists and passes.

The instrument-indexed form is

\[
\Gamma\,;\;\mathbf{BirdInt}^{\mathrm{aud}}_{\mathrm{fin}}\,;\;I_{\mathrm{prom}}\;\vdash\;\textsc{PromotionGateSoundness}\;:\;\texttt{accepted}.
\]

The result records the basic discipline of promotion: an acceptance verdict in the calculus is not informal; it is precisely the condition that every required core gate has passed its finite admissibility check.

\subsection{Theorem 20: Strict Promotion Implies Nonfactorization}

\subsubsection{Statement}

\textbf{Theorem 20 (StrictPromotionImpliesNonfactorization; theorem-grade).} Let \(B_{j\to j+1}\) be a promotion bridge. If the promotion classifier assigns

\[
\mathrm{Promote}(B_{j\to j+1})\;:\;\texttt{strict},
\]

then \(\pi_{j+1}\not\factor\pi_j\). Equivalently, there exist \(s,s'\in S\) with \(\pi_j(s)=\pi_j(s')\) and \(\pi_{j+1}(s)\neq\pi_{j+1}(s')\).

\subsubsection{Proof}

By the strictness-gate semantics of Subsection 5.4.1, the strict-pass refinement \(G_{\mathrm{strict}}=\texttt{strict\_pass}\) is recorded exactly when \(\Delta_{\mathrm{fact}}(\pi_j,\pi_{j+1})\neq\varnothing\). The promotion classifier of Subsection 5.6.4 produces the verdict \(\texttt{strict}\) only when \(G_{\mathrm{strict}}=\texttt{strict\_pass}\), so \(\Delta_{\mathrm{fact}}\neq\varnothing\). By Theorem 13 of Section 7, \(\Delta_{\mathrm{fact}}\neq\varnothing\) is equivalent to \(\pi_{j+1}\not\factor\pi_j\), and a witness pair \((s,s')\) exists by definition of \(\Delta_{\mathrm{fact}}\).

The instrument-indexed form is

\[
\Gamma\,;\;\mathbf{BirdInt}^{\mathrm{aud}}_{\mathrm{fin}}\,;\;I_{\mathrm{prom}}\;\vdash\;\textsc{StrictPromotionImpliesNonfactorization}\;:\;\texttt{accepted}.
\]

\subsection{Theorem 21: Non-Strict Promotion Implies Factorization}

\subsubsection{Statement}

\textbf{Theorem 21 (NonStrictPromotionImpliesFactorization; theorem-grade).} Let \(B_{j\to j+1}\) be a promotion bridge. If the promotion classifier assigns

\[
\mathrm{Promote}(B_{j\to j+1})\;:\;\texttt{non\_strict},
\]

then \(\pi_{j+1}\factor\pi_j\). There exists a finite map \(\phi:\mathrm{im}(\pi_j)\to O_{j+1}\) with \(\pi_{j+1}=\phi\,\pi_j\).

\subsubsection{Proof}

By the strictness-gate semantics of Subsection 5.4.1, the non-strict-pass refinement \(G_{\mathrm{strict}}=\texttt{non\_strict\_pass}\) is recorded exactly when \(\Delta_{\mathrm{fact}}(\pi_j,\pi_{j+1})=\varnothing\). The promotion classifier of Subsection 5.6.4 produces the verdict \(\texttt{non\_strict}\) only when \(G_{\mathrm{strict}}=\texttt{non\_strict\_pass}\), so \(\Delta_{\mathrm{fact}}=\varnothing\). By Theorem 13 of Section 7, \(\Delta_{\mathrm{fact}}=\varnothing\) is equivalent to \(\pi_{j+1}\factor\pi_j\), and Theorem 12 supplies the explicit factorization \(\phi\) on \(\mathrm{im}(\pi_j)\).

The instrument-indexed form is

\[
\Gamma\,;\;\mathbf{BirdInt}^{\mathrm{aud}}_{\mathrm{fin}}\,;\;I_{\mathrm{prom}}\;\vdash\;\textsc{NonStrictPromotionImpliesFactorization}\;:\;\texttt{accepted}.
\]

\subsection{Example Theorem 22: Toy 22 Strict Bridge}

\subsubsection{Data}

\textbf{Theorem 22 (Toy22StrictBridge; theorem-grade example).} There exists a finite two-layer promotion bridge \(B_{0\to 1}^{\mathrm{Toy22}}\) on the host \(\mathsf H_{\mathrm{fin}}\) with carrier \(S=\{a,b\}\) and object maps

\[
\pi_0\colon S\to O_0=\{*\},\qquad \pi_1\colon S\to O_1=\{u,v\}
\]

defined by \(\pi_0(a)=\pi_0(b)=*\) and \(\pi_1(a)=u\), \(\pi_1(b)=v\). The bridge is admissible under a finite instrument \(I_{\mathrm{prom}}\), with audit and visibility records that pass the required core gates and with no closure or drive claim attached to the bridge.

\subsubsection{Gate Results}

The gate-results record \(\mathsf{GateResults}(B_{0\to 1}^{\mathrm{Toy22}})\) assigns:

\begin{itemize}
\tightlist
\item
  \(G_{\mathrm{suff}}=\texttt{pass}\), since the bridge data conforms to the schema of Subsection 10.1.1;
\item
  \(G_{\mathrm{vis}}=\texttt{pass}\), since the bridge content lies in \(\mathsf{Visible}_{I_{\mathrm{prom}}}\);
\item
  \(G_{\mathrm{audit}}=\texttt{pass}\), since the audit record passes \(\mathsf{AuditPolicy}_{I_{\mathrm{prom}}}\);
\item
  \(G_{\mathrm{ctrl}}=\texttt{pass}\), since no target-layer data has been used to certify target novelty;
\item
  \(G_{\mathrm{nosmuggle}}=\texttt{pass}\), since the no-smuggling defect of Subsection 5.4.2 is empty;
\item
  \(G_{\mathrm{desc}}=\texttt{not\_required}\), since no closure or descended-dynamics claim is made on the bridge;
\item
  \(G_{\mathrm{stab}}=\texttt{not\_required}\), since no fixed-point structure is claimed;
\item
  \(G_{\mathrm{locglob}}=\texttt{not\_required}\), since no local-to-global packaging claim is made.
\end{itemize}

The strictness gate is fed by the factorization defect \(\Delta_{\mathrm{fact}}(\pi_0,\pi_1)\). We have \(\pi_0(a)=\pi_0(b)=*\) but \(\pi_1(a)=u\neq v=\pi_1(b)\), so \((a,b)\in\Delta_{\mathrm{fact}}(\pi_0,\pi_1)\neq\varnothing\). By Theorems 12 and 13 of Section 7 and the strictness-gate semantics, \(G_{\mathrm{strict}}=\texttt{strict\_pass}\).

The promotion classifier of Subsection 5.6.4 then assigns

\[
\mathrm{Promote}(B_{0\to 1}^{\mathrm{Toy22}})\;:\;\texttt{strict}.
\]

The instrument-indexed form is

\[
\Gamma\,;\;\mathsf H_{\mathrm{fin}}\,;\;I_{\mathrm{prom}}\;\vdash\;\textsc{Toy22StrictBridge}\;:\;\texttt{accepted}.
\]

\subsubsection{Interpretation}

The toy bridge is a compact illustrative example of strict promotion, not a universal stacking theorem. It records that strictness is the working content of object-map nonfactorization on a finite carrier; it does not by itself imply closure, drive, or any deeper structure on the target package. Section 13's Cantor model-realization sheet supplies a more substantive scoped strict-bridge realization on the audited Cantor shell, and Subsection 10.7 records that no general stacking composition theorem follows from the toy example.

\subsection{Stacking Claims Not Proved}

This subsection records two claims that this paper does not prove. They are deferred to future work; the limitations and future-work boundary of Subsection 16.3.2 lists them in the formal future-work register.

\subsubsection{No General Stacking Composition Theorem}

The paper does not prove a general stacking composition theorem of the form

\[
\mathrm{Promote}(B_{j\to j+1}):\texttt{strict}\;\wedge\;\mathrm{Promote}(B_{j+1\to j+2}):\texttt{strict}\;\Longrightarrow\;\mathrm{Promote}(B_{j\to j+2}):\texttt{strict},
\]

with \(B_{j\to j+2}\) a composed bridge derived from \(B_{j\to j+1}\) and \(B_{j+1\to j+2}\). A general composition theorem would require results on the composition of strictness defects across bridges, the joint gate semantics on the composed bridge, and the joint nonclaim discipline that propagates through composition. None of these is proved here. Subsection 16.3.2 records the corresponding future-work item.

\subsubsection{No Local-to-Global Promotion Theorem}

The paper does not prove a local-to-global promotion theorem of the form

\[
\bigl(\,\forall\,U_\alpha\in\text{cover},\;\mathrm{Promote}(B_\alpha\text{ on }U_\alpha):\texttt{strict}\,\bigr)\;\Longrightarrow\;\mathrm{Promote}(B_{\mathrm{global}}):\texttt{strict}.
\]

The local-to-global gate \(G_{\mathrm{locglob}}\) of Subsection 10.2 is optional precisely because the conditions under which local strict promotions glue into a global strict promotion are not classified at theorem-grade in this paper. The countermodel \(\mathsf{CM}_3\) of Section 9 records the failure case for ordinary packages without a gluing bridge on a finite cover; a positive gluing theorem for compatible local promotion bridges is recorded as future work in Subsection 16.3.2.
 \section{Instruments, Visibility, and Rotating Audit}

This section specifies the instrument-indexed claim semantics, the visibility discipline that the calculus uses to prevent overreading, and the rotating-audit machinery that records self-reference honestly. Subsections 11.1 and 11.2 fix the claim semantics and the classifier priority. Subsection 11.3 records the visibility and suppression discipline at the level required for the proofs. Subsections 11.4 and 11.5 prove Theorems 23 and 24. Subsection 11.6 isolates the \(P_6\leftarrow P_6\) status separation. Subsection 11.7 records the rotating-audit chain. Subsection 11.8 proves Theorem 25. Subsection 11.9 discusses instrument transfer and marks its full classifier as future work.

\subsection{Instrument-Indexed Claim Semantics}

This subsection records the instrument record schema, the claim record schema, and the form of the acceptance judgment. The classifier rules and the no-overreading and same-level self-audit theorems that depend on them are stated in Subsections 11.4 and 11.5.

\subsubsection{Instrument Record}

The instrument record \(I\) has the form fixed in Subsection 3.4.1, with finite components

\[
I\;=\;(\,\mathrm{id}_I,\;\mathsf{Scope}_I,\;\mathsf{ClaimTypes}_I,\;\mathsf{Visible}_I,\;\mathsf{Suppressed}_I,\;\mathsf{CheckRules}_I,\;\mathsf{EvidencePolicy}_I,\;\mathsf{ThresholdPolicy}_I,\;\mathsf{AuditPolicy}_I,\;\mathsf{Level}_I,\;\mathsf{BridgePolicy}_I,\;\mathsf{SourcePolicy}_I,\;\mathcal N_I\,).
\]

Every component is finite, and admissibility of \(I\) is decidable on finite records by the conditions of Subsection 3.4.1; write \(\mathsf{AdmInstrument}(I)\) for this finite admissibility predicate.

\subsubsection{Claim Record}

A claim record under instrument \(I\) is a finite typed tuple

\[
\varphi(\mathcal T)\;=\;(\,\mathrm{id}_{\varphi},\;\mathcal T,\;\mathsf{ClaimType},\;\mathsf{Content},\;\mathsf{LevelProfile},\;\mathsf{HostTag},\;\mathsf{EvidenceRefs},\;\Theta,\;A,\;V,\;\mathcal L,\;\mathcal N_{\varphi}\,).
\]

The tuple fields record:

\begin{itemize}
\tightlist
\item
  the claim content \(\varphi\), with finite use-set \(\mathsf{Use}(\varphi)\subseteq\mathsf{Cont}\);
\item
  the declared claim type drawn from \(\mathsf{ClaimTypes}_I\);
\item
  the target package \(\mathcal T\) and surrounding context \(\Gamma\);
\item
  the source records consulted, drawn from the source-of-truth schema of Subsection 4.8.1;
\item
  the threshold and audit records, drawn from the threshold and audit schemas of Subsections 5.2 and 4.8.2;
\item
  the visibility record \(V\) and the visibility bridges \(\mathcal L\) used to bring suppressed content into admissible reach, drawn from \(\mathsf{BridgePolicy}_I\);
\item
  the nonclaim record \(\mathcal N_\varphi\) attached to the claim.
\end{itemize}

A claim record is well-formed when every field is present, finite, and consistent with the instrument's policies; well-formedness is a finite check on finite records.

\subsubsection{Acceptance Rule}

The general claim form is

\[
\Gamma\,;\;\mathcal T\,;\;I\;\vdash\;\varphi\;:\;\chi,
\]

with \(\chi\in\mathrm{ClaimStatus}\), where

\[
\mathrm{ClaimStatus}\;=\;\{\,\texttt{accepted},\;\texttt{rejected},\;\texttt{provisional},\;\texttt{blocked},\;\texttt{outside\_scope},\;\texttt{absent\_with\_record},\;\texttt{below\_threshold},\;\texttt{failed\_audit},\;\texttt{undefined\_circular}\,\}.
\]

The classifier of Subsection 5.6.5 consumes the claim record, the visibility partition \(\mathsf{Visible}_I/\mathsf{Suppressed}_I/\mathsf{Outside}_I/\mathsf{Unknown}_I\), the threshold and source records, the audit record, the bridge family \(\mathcal L\), and the nonclaim records, and assigns a value \(\chi\in\mathrm{ClaimStatus}\) to the claim. The verdict \(\chi=\texttt{accepted}\) records that the required in-scope, visibility, evidence, threshold, audit, circularity, and nonclaim checks have passed under \(I\); the other statuses record the specific failure or scope condition that prevented acceptance.

The shorthand \(I\vdash\varphi(\mathcal T)\) abbreviates \(\Gamma;\mathcal T;I\vdash\varphi:\texttt{accepted}\) and never abbreviates instrument-free truth, in the sense fixed in Subsection 3.4.2.

\subsection{Claim Classifier Priority}

\subsubsection{Priority Order}

The claim classifier of Subsection 5.6.5 is a finite first-match procedure with the priority order

\[
\texttt{outside\_scope}\;\succ\;\texttt{blocked}\;\succ\;\texttt{absent\_with\_record}\;\succ\;\texttt{failed\_audit}\;\succ\;\texttt{undefined\_circular}\;\succ\;\texttt{below\_threshold}\;\succ\;\texttt{rejected}\;\succ\;\texttt{provisional}\;\succ\;\texttt{accepted}.
\]

A claim record matching multiple per-status conditions is assigned the highest-priority match. The order encodes the intended discipline: an outside-scope claim is recorded outside scope rather than rejected; a claim whose use-set contains suppressed unbridged content is blocked rather than rejected; audited absence and failed audits are recorded before threshold or refutation outcomes; circularity is recorded once the earlier scope, visibility, source, and audit checks do not already determine the status; and the accepted verdict is the lowest-priority match, reached only when every higher condition has been ruled out.

\subsubsection{Soundness Role}

The classifier prevents suppressed content from supporting overstrong claims. Its priority order first routes claims with outside-scope use-set entries to \(\texttt{outside\_scope}\); for in-scope visibility failures, a nonempty visibility defect \(\delta_{\mathrm{vis}}\) of Subsection 5.5 routes the claim to \(\texttt{blocked}\) before \(\texttt{accepted}\) can be reached. Likewise, a same-level circular self-reference routes to \(\texttt{undefined\_circular}\) once the earlier scope, visibility, source, and audit checks have not already determined the status. The no-overreading and same-level self-audit failure theorems of Subsections 11.4 and 11.5 are the theorem-grade form of these routing rules.

\subsection{Visibility and Suppression}

This subsection records the visibility data required for Theorems 23 and 24. The full visibility schema is fixed in Subsections 3.5 and 5.5; this subsection collects the definitions used immediately by the proofs.

\subsubsection{Visible Content}

The visible class under \(I\) is

\[
\mathsf{Visible}_I\;=\;\{\,x\in\mathsf{Cont}\;:\;\mathsf{vis}_I(x)=\texttt{visible}\,\},
\]

with \(\mathsf{vis}_I:\mathsf{Cont}\to\{\texttt{visible},\texttt{suppressed},\texttt{outside},\texttt{unknown}\}\) as in Subsection 3.5. The classifier of Subsection 5.6.5 may use visible content directly when classifying a claim.

\subsubsection{Suppressed Content}

The suppressed class under \(I\) is

\[
\mathsf{Suppressed}_I\;=\;\{\,x\in\mathsf{Cont}\;:\;\mathsf{vis}_I(x)=\texttt{suppressed}\,\}.
\]

Suppression is a relation between content and instrument, not a truth verdict on the content: \(x\in\mathsf{Suppressed}_I\) does not imply that \(x\) is false, only that \(I\) does not have admissible visibility into \(x\) without a bridge.

\subsubsection{Visibility Bridges}

A visibility bridge under \(I\) is the finite typed record of Subsection 3.5.3:

\[
L\;=\;(\,\mathrm{id}_L,\;c,\;c',\;\mathsf{BridgeType}_L,\;\mathsf{Guarantee}_L,\;\mathsf{Strength}_L,\;A_L,\;V_L,\;\mathcal N_L\,),
\]

with \(c\in\mathsf{Suppressed}_I\), \(c'\in\mathsf{Visible}_I\), and \(\mathsf{BridgeType}_L\in\mathsf{BridgePolicy}_I\). A bridge is admissible for a claim \(\varphi\) under \(I\), written \(\mathsf{AdmVisBridge}(L,I,\varphi)\), when it targets suppressed content used by \(\varphi\), has a visible proxy, is allowed by \(I\), is strong enough for the declared claim strength, is audited, and is noncircular. Let \(\mathcal L\) denote a finite set of admissible visibility bridges under \(I\).

\subsubsection{Visibility Defects}

The visibility defect family of a claim \(\varphi\) under instrument \(I\) relative to a bridge family \(\mathcal L\) is the union of four finite typed sets

\[
\delta_{\mathrm{vis}}(I,\varphi,\mathcal L)\;=\;\delta_{\mathrm{overread}}\;\cup\;\delta_{\mathrm{weakbridge}}\;\cup\;\delta_{\mathrm{outside}}\;\cup\;\delta_{\mathrm{unknown}},
\]

defined in Subsection 5.5. The overreading defect is

\[
\delta_{\mathrm{overread}}(I,\varphi,\mathcal L)\;=\;\{\,x\in\mathsf{Use}(\varphi)\;:\;x\in\mathsf{Suppressed}_I\;\text{ and no admissible bridge in }\mathcal L\text{ permits }x\,\},
\]

and the threshold for the visibility check is \(\Theta_{\mathrm{vis}}:\delta_{\mathrm{vis}}=\varnothing\).

The weak-bridge defect is

\[
\delta_{\mathrm{weakbridge}}(I,\varphi,\mathcal L)\;=\;\{\,(x,L)\;:\;x\in\mathsf{Use}(\varphi),\;L\in\mathcal L\text{ bridges }x,\;\mathsf{Strength}(L)<_{\varphi}\mathsf{RequiredStrength}(\varphi)\,\}.
\]

The outside and unknown defects are

\[
\delta_{\mathrm{outside}}(I,\varphi)\;=\;\mathsf{Use}(\varphi)\cap\mathsf{Outside}_I,\qquad
\delta_{\mathrm{unknown}}(I,\varphi)\;=\;\mathsf{Use}(\varphi)\cap\mathsf{Unknown}_I.
\]

\subsection{Theorem 23: No-Overreading Suppression}

\subsubsection{Statement}

\textbf{Theorem 23 (NoOverreadingSuppression; theorem-grade).} Let \(I\) be an admissible instrument, \(\varphi\) a finite claim, and \(\mathcal L\) a finite set of admissible visibility bridges under \(I\). If

\[
\Gamma\,;\;\mathcal T\,;\;I\;\vdash\;\varphi\;:\;\texttt{accepted},
\]

then

\[
\delta_{\mathrm{overread}}(I,\varphi,\mathcal L)\;=\;\varnothing.
\]

Equivalently, every \(x\in\mathsf{Use}(\varphi)\) with \(x\in\mathsf{Suppressed}_I\) is bridged by some admissible \(L\in\mathcal L\).

\subsubsection{Proof}

By the classifier priority order of Subsection 11.2.1, the verdict \(\texttt{accepted}\) is reached only after every higher-priority condition has been ruled out. The \(\texttt{blocked}\) status sits above \(\texttt{accepted}\) in the priority order, and one of its triggering conditions in Subsection 5.6.5 is a non-empty visibility defect \(\delta_{\mathrm{vis}}\neq\varnothing\). Since \(\delta_{\mathrm{overread}}\subseteq\delta_{\mathrm{vis}}\), a non-empty overread defect forces \(\delta_{\mathrm{vis}}\neq\varnothing\) and the classifier returns \(\texttt{blocked}\) before reaching \(\texttt{accepted}\). Contrapositively, \(\texttt{accepted}\) implies \(\delta_{\mathrm{overread}}=\varnothing\).

The instrument-indexed form is

\[
\Gamma\,;\;\mathsf H_{\mathrm{instr}}\,;\;I_{\mathrm{instr}}\;\vdash\;\textsc{NoOverreadingSuppression}\;:\;\texttt{accepted}.
\]

\subsubsection{Corollaries}

The no-overreading constraint specializes to three corollaries on the judgment families of Section 4.

\textbf{Directed-cell corollary.} A directed-cell judgment with status \(\texttt{action}\) has empty overread defect: every suppressed field of the cell record (witness, update, level/lens/interface, threshold, audit, or defect record) is bridged by an admissible \(L\in\mathcal L\) on the cell's instrument.

\textbf{Pair-observable corollary.} A pair-observable judgment with status \(\texttt{real}\) has empty overread defect on the source-of-truth record and the branch-compatibility record. By the classifier of Subsection 5.6.3, the source must be a real-source class admissible under \(\mathsf{SourcePolicy}_I\), and any suppressed dependency must be bridged.

\textbf{Promotion corollary.} A promotion bridge with status \(\texttt{strict}\) has empty overread defect; the no-smuggling gate \(G_{\mathrm{nosmuggle}}\) of Section 10 enforces the bridge-level analogue at the gate level.

\subsection{Theorem 24: Same-Level Self-Audit Failure}

\subsubsection{Target}

Let \(I_0\) be an admissible instrument under an admissible package \(\mathcal T_0\) on the host \(\mathsf H_{\mathrm{instr}}\). Consider the candidate same-level self-soundness claim \(\mathrm{Sound}(I_0)\), presented as the attempted judgment

\[
\Gamma\,;\;\mathcal T_0\,;\;I_0\;\vdash\;\mathrm{Sound}(I_0):\chi.
\]

If \(\chi=\texttt{accepted}\), this would be written in shorthand as \(I_0\vdash\mathrm{Sound}(I_0)\). The target claim asserts that \(I_0\) certifies its own soundness without a level shift to a higher instrument.

\subsubsection{Outside-Scope Case}

Suppose \(\mathrm{Sound}(I_0)\notin\mathsf{ClaimTypes}(I_0)\), that is, \(I_0\) is not configured to evaluate claims about its own soundness. By the classifier of Subsection 5.6.5, the claim's use-set contains a claim type outside \(I_0\)'s scope, and the priority routing of Subsection 11.2 returns

\[
\Gamma\,;\;\mathcal T_0\,;\;I_0\;\vdash\;\mathrm{Sound}(I_0)\;:\;\texttt{outside\_scope}.
\]

This is the case for instruments configured to evaluate object-level claims only.

\subsubsection{Circular Case}

Suppose \(\mathrm{Sound}(I_0)\in\mathsf{ClaimTypes}(I_0)\) but the audit policy required to certify \(\mathrm{Sound}(I_0)\) consults the very verdict being computed --- the acceptance of \(\mathrm{Sound}(I_0)\) depends on \(I_0\)'s own soundness verdict. The classifier records this circular self-reference at the same level and returns

\[
\Gamma\,;\;\mathcal T_0\,;\;I_0\;\vdash\;\mathrm{Sound}(I_0)\;:\;\texttt{undefined\_circular}.
\]

This is the case when an instrument is asked to certify a property of itself whose verification cannot be carried out without already knowing the verdict.

\subsubsection{Statement and Proof}

\textbf{Theorem 24 (SameLevelSelfAuditFailure; theorem-grade).} Without an admissible level-shift or reflection bridge to a higher instrument, the same-level self-soundness claim does not receive verdict \(\texttt{accepted}\) under \(I_0\). The classifier returns either \(\texttt{outside\_scope}\) or \(\texttt{undefined\_circular}\) depending on whether the candidate claim is in \(I_0\)'s claim types and whether its acceptance is self-dependent.

\textbf{Proof.} For the complete same-level self-audit target, the classifier has the two relevant cases of Subsections 11.5.2 and 11.5.3. If the target claim is not in scope, the priority rule returns \(\texttt{outside\_scope}\). If it is syntactically in scope but its audit support depends on the same verdict or on an unstratified same-level audit cycle, the priority rule returns \(\texttt{undefined\_circular}\). Both statuses sit above \(\texttt{accepted}\), so \(\texttt{accepted}\) is not reached. Hence the same-level self-soundness claim is not accepted under \(I_0\) in the absence of a level shift.

The instrument-indexed form is

\[
\Gamma\,;\;\mathsf H_{\mathrm{instr}}\,;\;I_{\mathrm{instr}}\;\vdash\;\textsc{SameLevelSelfAuditFailure}\;:\;\texttt{accepted}.
\]

NC-16 of Subsection 2.4.3 is the corresponding nonclaim. A positive soundness claim about a lower instrument requires a level shift via a higher instrument \(I_1\) that audits \(I_0\) through an admissible audit bridge; Subsection 11.7 records the rotating-audit chain that supplies such bridges, and Subsection 11.8 records that a rotating audit can certify the lower stack but does not certify itself, so the rotating audit is non-final (NC-17).

\subsection{\texorpdfstring{\(P_6\leftarrow P_6\) Status Separation}{P\_6\textbackslash leftarrow P\_6 Status Separation}}

This subsection records a schema-level distinction between three judgments that all carry the same primitive pair \(P_6\leftarrow P_6\): the object-level audit update, the fixed audit refresh, and the complete same-level self-audit. They differ in profile, update, and directed-cell status; Theorem 1 of Section 6 covers the formal non-factorization of cell status through the primitive pair.

\subsubsection{Object-Level Audit Update}

A directed cell \(P_6\leftarrow P_6\) in which the actor update \(U_6\) records an audit on a specific object-level record --- a directed cell, a pair observable, a promotion bridge, or a claim record other than \(I\)'s own soundness --- is admissible. Its profile \(\lambda_{\mathrm{obj}}\) records audit-update mode at object level. The cell may receive directed-cell status \(\texttt{action}\) in the classifier of Subsection 5.6.2 when its witness, threshold, audit, and visibility records pass the standard cell admissibility checks.

\subsubsection{Fixed Audit Refresh}

A directed cell \(P_6\leftarrow P_6\) in which the actor update is a fixed-instrument audit refresh --- re-running an existing audit policy on an already present, unchanged audit record under a profile \(\lambda_{\mathrm{refresh}}\) that records refresh mode --- is admissible and classified \(\texttt{trivial}\) in the priority order of Subsection 5.6.2. If the audit step produces new admissible audit content, it is no longer this fixed-refresh case and must be classified under its own object-level audit-update profile.

\subsubsection{Complete Same-Level Self-Audit}

A directed cell \(P_6\leftarrow P_6\) in which the actor update asserts a complete same-level self-soundness verdict on the surrounding instrument \(I_0\) is the case covered by Theorem 24 of Subsection 11.5. By that theorem, the cell-level reflection of the soundness claim is either \(\texttt{outside\_scope}\) or \(\texttt{undefined\_circular}\) in the corresponding claim record; the directed-cell status of the cell is \(\texttt{outside\_scope}\) or \(\texttt{undefined\_circular}\) accordingly.

The three cases together separate object-level audit, fixed audit refresh, and complete same-level self-certification within the same primitive pair \((P_6,P_6)\). They are an instance of the typed/profile-relative discipline of the calculus: the same primitive pair receives different statuses on different profiles, as Theorem 1 records.

\subsection{Rotating Audit}

This subsection records the rotating-audit discipline that replaces forbidden complete same-level self-audit with a finite, level-shifted, bounded audit of one instrument by a strictly higher instrument. It records the rotating-audit chain, the rotating-audit bridge, the admissibility predicate for that bridge, and the rotating-audit defect schema. Theorem 25 of Subsection 11.8 then states the finite rotating-audit theorem and its non-finality.

\subsubsection{Rotating-Audit Chain}

A rotating-audit chain is a finite ordered family of admissible instruments

\[
\mathbb I_{\le N}\;=\;(I_0,I_1,\ldots,I_N),
\]

with strictly increasing instrument levels \(\mathsf{Level}(I_0)<\mathsf{Level}(I_1)<\cdots<\mathsf{Level}(I_N)\), together with a finite ordered family of rotating-audit bridges

\[
B^{\mathrm{audit}}_{n\to n+1}\quad\text{for each } 0\le n<N,
\]

each pointing one level upward. Each \(I_n\) is an instrument in the sense of Subsection 11.1.1, with finite logs, declared visibility policy, declared source policy, declared audit policy, declared nonclaim register, and declared scope. The chain is finite by construction: \(N\) is a fixed natural number and there is no implicit limit instrument at the top. The chain has no final self-certifying instrument; Subsection 11.8 records this non-finality formally.

\subsubsection{Rotating-Audit Bridge}

For each adjacent pair \((I_n,I_{n+1})\) with \(\mathsf{Level}(I_{n+1})>\mathsf{Level}(I_n)\), the rotating-audit bridge

\[
B^{\mathrm{audit}}_{n\to n+1}\;=\;(I_n,\;I_{n+1},\;\Phi_n,\;L,\;V,\;\Theta,\;A,\;\delta,\;\mathcal N)
\]

with expanded finite record

\[
\begin{aligned}
B^{\mathrm{audit}}_{n\to n+1}\;=\;(&\mathrm{id},\;I_n,\;I_{n+1},\;\Phi_n,\;\mathsf{TargetRecords}_n,\;L_{n\to n+1},\;V_{n\to n+1},\\
&\Theta_{n\to n+1},\;A_{n\to n+1},\;\delta_{n\to n+1},\;\mathcal N_{n\to n+1},\;\lambda_{\mathrm{audit}}).
\end{aligned}
\]

It is a directed bridge from \(I_n\) into \(I_{n+1}\) whose bridge type is \(\texttt{reflection/audit}\). Its components are:

\begin{itemize}
\tightlist
\item
  \(\Phi_n\), the rotating-audit target scope: a finite, declared subset of \(I_n\)'s records (code, check rules, claim log, source policy, visibility policy, nonclaim register, status classifier, trace records, bounded-soundness records) on which \(I_{n+1}\) undertakes to compute a finite audit defect, in the sense of Subsection 11.7.4;
\item
  \(L_{n\to n+1}\), the language map that translates \(I_n\)-records into \(I_{n+1}\)'s claim language;
\item
  \(V_{n\to n+1}\), the visibility map that exposes the audited records to \(I_{n+1}\) either directly or through visibility bridges of Subsection 11.3.3;
\item
  \(\Theta_{n\to n+1}\), the audit thresholds (witness, audit, visibility) that \(I_{n+1}\) applies when classifying the bridge;
\item
  \(A_{n\to n+1}\), the audit record of the bridge itself, recording that the bridge is auditable in \(I_{n+1}\);
\item
  \(\delta_{n\to n+1}\), the rotating-audit defect of Subsection 11.7.4;
\item
  \(\mathcal N_{n\to n+1}\), the bridge-level nonclaim register that records, at minimum, \(\mathrm{NC}\text{-}16\), \(\mathrm{NC}\text{-}17\), and \(\mathrm{NC}\text{-}18\);
\item
  \(\lambda_{\mathrm{audit}}\), the bridge profile recording rotating-audit mode.
\end{itemize}

The rotating-audit bridge is the instrument-level vehicle by which \(I_{n+1}\) reflects \(I_n\)'s declared records into its own claim language under a declared, finite scope \(\Phi_n\). It is a \(P_6\leftarrow P_6\) object at the bridge layer, distinct from the directed-cell \(P_6\leftarrow P_6\) judgments separated in Subsection 11.6: a rotating-audit bridge audits an instrument across a level shift; the cell-level judgments of Subsection 11.6 audit object-level records under a single instrument.

\subsubsection{Rotating-Audit Admissibility}

The admissibility predicate

\[
\mathsf{AdmRotAuditBridge}(B^{\mathrm{audit}}_{n\to n+1})
\]

holds iff every clause below holds:

\begin{enumerate}
\def\labelenumi{\arabic{enumi}.}
\tightlist
\item
  \(I_n\) and \(I_{n+1}\) are admissible instruments in the sense of Subsection 11.1.1;
\item
  \(\mathsf{Level}(I_{n+1})>\mathsf{Level}(I_n)\);
\item
  \(\Phi_n\) is finite and declared as part of \(B^{\mathrm{audit}}_{n\to n+1}\);
\item
  every record in \(\Phi_n\) is in scope of \(I_{n+1}\), either directly or through a visibility bridge of Subsection 11.3.3;
\item
  \(L_{n\to n+1}\) translates the records of \(\Phi_n\) into well-typed \(I_{n+1}\)-claims;
\item
  \(V_{n\to n+1}\) exposes the records required by the audit thresholds, with no overreading defect of Subsection 11.3.4 against the visible content;
\item
  \(\Theta_{n\to n+1}\) defines audit thresholds that \(I_{n+1}\) is committed to apply;
\item
  \(A_{n\to n+1}\) is a finite audit record of the bridge itself, recording that the bridge is auditable in \(I_{n+1}\) and that its records are consistent with \(I_{n+1}\)'s audit policy;
\item
  \(\delta_{n\to n+1}\) is typed by the rotating-audit defect schema of Subsection 11.7.4;
\item
  the audit dependency graph of the chain has no unstratified same-level audit cycle: no record in \(\Phi_n\) depends on the very acceptance verdict that \(I_{n+1}\) is computing for \(\mathrm{Audit}_{\Phi_n}(I_n)\);
\item
  \(\mathcal N_{n\to n+1}\) is present and contains the rotating-audit nonclaims of the calculus, in particular \(\mathrm{NC}\text{-}16\), \(\mathrm{NC}\text{-}17\), and \(\mathrm{NC}\text{-}18\).
\end{enumerate}

The clauses are finite and effectively decidable on the finite records of \(B^{\mathrm{audit}}_{n\to n+1}\); admissibility is a syntactic predicate on the bridge object, not an external soundness condition.

\subsubsection{Rotating-Audit Defect}

The rotating-audit defect schema for a rotating-audit bridge \(B^{\mathrm{audit}}_{n\to n+1}\) is

\[
\mathsf{Def}^{\mathrm{rot}}_{n\to n+1}\;=\;(\delta_{\mathrm{rot}},\;\Omega_{\mathrm{rot}},\;\Theta_{\mathrm{rot}},\;\mathrm{polarity}_{\mathrm{rot}},\;W_{\mathrm{rot}},\;A_{\mathrm{rot}},\;V_{\mathrm{rot}},\;\sigma_{\mathrm{rot}}),
\]

following the defect-record schema of Subsection 5.2. The defect value \(\delta_{\mathrm{rot}}\) ranges over the rotating-audit defect codes

\[
\Omega_{\mathrm{rot}}\;=\;\{\,\texttt{audit\_passes},\;\texttt{missing\_target\_scope},\;\texttt{missing\_instrument\_record},\;\texttt{visibility\_failure},\;\texttt{translation\_failure},\;\texttt{source\_failure},\;\texttt{threshold\_failure},\;\texttt{nonclaim\_failure},\;\texttt{same\_level\_cycle},\;\texttt{higher\_audit\_failure}\,\},
\]

with intended readings:

\begin{itemize}
\tightlist
\item
  \(\texttt{audit\_passes}\): the bridge audit succeeded under its declared thresholds;
\item
  \(\texttt{missing\_target\_scope}\): \(\Phi_n\) is not finite or not declared;
\item
  \(\texttt{missing\_instrument\_record}\): a record required by \(\Phi_n\) is not present in \(I_n\) or not visible in \(I_{n+1}\);
\item
  \(\texttt{visibility\_failure}\): a visibility-policy clause of \(V_{n\to n+1}\) fails;
\item
  \(\texttt{translation\_failure}\): \(L_{n\to n+1}\) fails to type a record in \(\Phi_n\);
\item
  \(\texttt{source\_failure}\): a source-of-truth clause of \(I_{n+1}\) fails on a record in \(\Phi_n\);
\item
  \(\texttt{threshold\_failure}\): an audit threshold in \(\Theta_{n\to n+1}\) is violated;
\item
  \(\texttt{nonclaim\_failure}\): a clause in \(\mathcal N_{n\to n+1}\) is violated;
\item
  \(\texttt{same\_level\_cycle}\): an unstratified same-level audit cycle is detected in the bridge dependency graph;
\item
  \(\texttt{higher\_audit\_failure}\): the bridge-level audit record \(A_{n\to n+1}\) fails.
\end{itemize}

The acceptance threshold is

\[
\Theta_{\mathrm{rot}}:\;\delta_{\mathrm{rot}}=\texttt{audit\_passes}.
\]

The value \(\texttt{audit\_passes}\) is the threshold value later consumed by the rotating-audit acceptance rule of Subsection 11.8. This subsection records only the bridge-level defect schema and the finite field on which the later rule operates.

\subsection{Theorem 25: Finite Rotating-Audit Theorem}

This subsection records the finite rotating-audit theorem of Foundations III. The theorem is theorem-grade: it states that, given any finite stack of admissible instruments, there is a one-level extension instrument and an audit bridge that can finitely audit the lower stack and accept its compliance iff the stack-audit defect is empty, while the higher instrument does not thereby certify itself. The statement, proof, and non-finality are recorded in Subsubsections 11.8.1--11.8.3. The theorem supports \(\mathrm{NC}\text{-}17\) of Subsection 2.4.3.

\subsubsection{Statement}

Let

\[
\mathbb I_{\le N}\;=\;(I_0,\,I_1,\,\ldots,\,I_N)
\]

be a finite admissible rotating-audit chain in the sense of Subsection 11.7.1, with strictly increasing levels and a finite family of admissible rotating-audit bridges \(\{B^{\mathrm{audit}}_{n\to n+1}\}_{0\le n<N}\). Let the lower-stack audit target be the finite union

\[
\mathsf{Target}_{\le N}\;=\;\bigcup_{n=0}^{N}\mathsf{InstrRecord}(I_n)\;\cup\;\bigcup_{n=0}^{N-1}\mathsf{BridgeRecord}\bigl(B^{\mathrm{audit}}_{n\to n+1}\bigr)\;\cup\;\bigcup_{n=0}^{N}\mathsf{ClaimLog}(I_n).
\]

Let the finite stack-audit defect be the finite typed union

\[
\delta_{\le N}^{\mathrm{stack}}\;=\;\delta_{\mathrm{instr}}\,\cup\,\delta_{\mathrm{bridge}}\,\cup\,\delta_{\mathrm{visibility}}\,\cup\,\delta_{\mathrm{source}}\,\cup\,\delta_{\mathrm{audit}}\,\cup\,\delta_{\mathrm{circular}}\,\cup\,\delta_{\mathrm{nonclaim}}\,\cup\,\delta_{\mathrm{status}},
\]

where each summand is the finite defect of the corresponding component, using the general defect-record schema of Subsection 5.2, the visibility-defect schema of Subsection 5.5, the source and audit record schemas of Subsection 4.8, the claim-status classifier of Subsection 11.2, and the rotating-audit defect schema of Subsection 11.7.4.

\textbf{Theorem 25 (FiniteRotatingAudit; theorem-grade).} There exists a one-level extension instrument

\[
I_{N+1}^{\mathrm{stack}}
\]

and an admissible audit bridge

\[
B^{\mathrm{audit}}_{\le N\to N+1}
\]

with \(\mathsf{Level}\bigl(I_{N+1}^{\mathrm{stack}}\bigr)>\max_{0\le n\le N}\mathsf{Level}(I_n)\) and target scope \(\Phi_{\le N}\supseteq\mathsf{Target}_{\le N}\) such that the following hold:

\begin{enumerate}
\def\labelenumi{\arabic{enumi}.}
\tightlist
\item
  (Audit acceptance.)
  \[
  \Gamma\,;\;\mathsf H_{\mathrm{instr}}\,;\;I_{N+1}^{\mathrm{stack}}\;\vdash\;\mathrm{Report}_{\Phi_{\le N}}(\mathbb I_{\le N})\;:\;\texttt{accepted}.
  \]
\item
  (Compliance characterization.)
  \[
  \Gamma\,;\;\mathsf H_{\mathrm{instr}}\,;\;I_{N+1}^{\mathrm{stack}}\;\vdash\;\mathrm{Compliant}_{\Phi_{\le N}}(\mathbb I_{\le N})\;:\;\texttt{accepted}\quad\Longleftrightarrow\quad\delta_{\le N}^{\mathrm{stack}}=\varnothing.
  \]
\item
  (Non-finality.)
  \[
  \Gamma\,;\;\mathsf H_{\mathrm{instr}}\,;\;I_{N+1}^{\mathrm{stack}}\;\vdash\;\mathrm{Audit}(\mathbb I_{\le N})\;:\;\texttt{accepted}\quad\not\Longrightarrow\quad\Gamma\,;\;\mathsf H_{\mathrm{instr}}\,;\;I_{N+1}^{\mathrm{stack}}\;\vdash\;\mathrm{Sound}\bigl(I_{N+1}^{\mathrm{stack}}\bigr)\;:\;\texttt{accepted}.
  \]
\end{enumerate}

The instrument-indexed form of the headline statement is

\[
\Gamma\,;\;\mathsf H_{\mathrm{instr}}\,;\;I_{\mathrm{instr}}\;\vdash\;\textsc{FiniteRotatingAudit}\;:\;\texttt{accepted}.
\]

\subsubsection{Proof}

The proof constructs the extension instrument, verifies the audit bridge is admissible, and reads off the three clauses from the rotating-audit acceptance condition and the rotating-audit non-finality statement.

\emph{Construction of \(I_{N+1}^{\mathrm{stack}}\).} Define \(I_{N+1}^{\mathrm{stack}}\) as a finite instrument record in the form of Subsection 11.1.1, with components fixed as follows:

\begin{itemize}
\tightlist
\item
  \(\mathsf{Scope}\bigl(I_{N+1}^{\mathrm{stack}}\bigr)\supseteq\mathsf{Target}_{\le N}\);
\item
  \(\mathsf{ClaimTypes}\bigl(I_{N+1}^{\mathrm{stack}}\bigr)\supseteq\{\,\mathrm{Report}_{\Phi_{\le N}}(\mathbb I_{\le N}),\;\mathrm{Compliant}_{\Phi_{\le N}}(\mathbb I_{\le N}),\;\mathrm{Audit}(\mathbb I_{\le N})\,\}\);
\item
  \(\mathsf{Visible}\bigl(I_{N+1}^{\mathrm{stack}}\bigr)\supseteq\mathsf{Target}_{\le N}\), with visibility bridges of Subsection 11.3.3 added where required to bring suppressed records into admissible reach;
\item
  \(\mathsf{CheckRules}\bigl(I_{N+1}^{\mathrm{stack}}\bigr)\) contains the finite rules that compute each of \(\delta_{\mathrm{instr}}\), \(\delta_{\mathrm{bridge}}\), \(\delta_{\mathrm{visibility}}\), \(\delta_{\mathrm{source}}\), \(\delta_{\mathrm{audit}}\), \(\delta_{\mathrm{circular}}\), \(\delta_{\mathrm{nonclaim}}\), \(\delta_{\mathrm{status}}\) on the finite records of \(\mathsf{Target}_{\le N}\);
\item
  \(\mathsf{ThresholdPolicy}\bigl(I_{N+1}^{\mathrm{stack}}\bigr)\) inherits the per-component thresholds from the source instruments and adds the rotating-audit threshold \(\Theta_{\mathrm{rot}}:\;\delta_{\mathrm{rot}}=\texttt{audit\_passes}\) at the bridge layer;
\item
  \(\mathsf{AuditPolicy}\bigl(I_{N+1}^{\mathrm{stack}}\bigr)\) contains the audit rules for the bridge \(B^{\mathrm{audit}}_{\le N\to N+1}\);
\item
  \(\mathsf{Level}\bigl(I_{N+1}^{\mathrm{stack}}\bigr)>\max_{0\le n\le N}\mathsf{Level}(I_n)\);
\item
  \(\mathcal N\bigl(I_{N+1}^{\mathrm{stack}}\bigr)\) contains \(\mathrm{NC}\text{-}16\), \(\mathrm{NC}\text{-}17\), and \(\mathrm{NC}\text{-}18\).
\end{itemize}

Each component is finite by construction: \(\mathsf{Target}_{\le N}\) is a finite union of finite records and the check rules, thresholds, audit rules, and nonclaim register are finite. Hence \(\mathsf{AdmInstrument}\bigl(I_{N+1}^{\mathrm{stack}}\bigr)\) holds.

\emph{Construction of the audit bridge.} Define \(B^{\mathrm{audit}}_{\le N\to N+1}\) to be the rotating-audit bridge of Subsection 11.7.2 with target scope \(\Phi_{\le N}\supseteq\mathsf{Target}_{\le N}\), language map \(L\) that translates each \(I_n\)-record into \(I_{N+1}^{\mathrm{stack}}\)'s claim language, visibility map \(V\) exposing each record in \(\Phi_{\le N}\) either directly or through visibility bridges, audit thresholds \(\Theta\) inherited from the per-component thresholds, audit record \(A\) that records the finite checks above are auditable in \(I_{N+1}^{\mathrm{stack}}\), defect \(\delta=\delta_{\le N}^{\mathrm{stack}}\), and nonclaim register containing \(\mathrm{NC}\text{-}16\), \(\mathrm{NC}\text{-}17\), and \(\mathrm{NC}\text{-}18\).

\emph{Admissibility of the bridge.} Each clause of \(\mathsf{AdmRotAuditBridge}\bigl(B^{\mathrm{audit}}_{\le N\to N+1}\bigr)\) of Subsection 11.7.3 is verified on the finite records: clause (1) is the admissibility of the source instruments and of \(I_{N+1}^{\mathrm{stack}}\); clause (2) is the strict level inequality fixed by construction; clause (3) is the finite declaration of \(\Phi_{\le N}\); clause (4) is satisfied because \(I_{N+1}^{\mathrm{stack}}\)'s scope contains \(\mathsf{Target}_{\le N}\); clause (5) is the language map; clauses (6)--(9) are the visibility, threshold, audit, and defect-typing requirements built into the construction; clause (10) is the no-unstratified-same-level-cycle requirement, satisfied because the audit dependency graph points strictly upward, from level \(\le N\) to level \(N+1\); clause (11) is the nonclaim presence. Hence the bridge is admissible.

\emph{Clause 1 (Audit acceptance).} The rotating-audit bridge admissibility clauses of Subsection 11.7.3 and the defect threshold of Subsection 11.7.4, applied to \(B^{\mathrm{audit}}_{\le N\to N+1}\) with \(\delta_{\mathrm{rot}}=\texttt{audit\_passes}\) for the report claim, yield

\[
\Gamma\,;\;\mathsf H_{\mathrm{instr}}\,;\;I_{N+1}^{\mathrm{stack}}\;\vdash\;\mathrm{Report}_{\Phi_{\le N}}(\mathbb I_{\le N})\;:\;\texttt{accepted}.
\]

The report claim records the value of \(\delta_{\le N}^{\mathrm{stack}}\) together with the per-component defect values; it is accepted iff the bridge audit itself passes, regardless of whether \(\delta_{\le N}^{\mathrm{stack}}\) is empty.

\emph{Clause 2 (Compliance characterization).} The compliance claim \(\mathrm{Compliant}_{\Phi_{\le N}}(\mathbb I_{\le N})\) asserts \(\delta_{\le N}^{\mathrm{stack}}=\varnothing\). Its acceptance under \(I_{N+1}^{\mathrm{stack}}\) requires the threshold policy to compute the finite stack-audit defect to \(\varnothing\). By definition of the per-component defects, \(\delta_{\le N}^{\mathrm{stack}}=\varnothing\) holds iff every component defect is empty, iff every clause of \(\mathsf{AdmRotAuditBridge}\) holds with no defect across the lower stack. The classifier of Subsection 11.2.1 then assigns \(\texttt{accepted}\) to \(\mathrm{Compliant}_{\Phi_{\le N}}(\mathbb I_{\le N})\) iff \(\delta_{\le N}^{\mathrm{stack}}=\varnothing\). The reverse implication is immediate from the priority order: if any component defect is nonempty, the classifier returns a higher-priority status (e.g.~\(\texttt{failed\_audit}\), \(\texttt{blocked}\), or \(\texttt{undefined\_circular}\)) on the corresponding underlying claim, which propagates to a non-\(\texttt{accepted}\) status for the compliance claim. Hence the equivalence

\[
\Gamma\,;\;\mathsf H_{\mathrm{instr}}\,;\;I_{N+1}^{\mathrm{stack}}\;\vdash\;\mathrm{Compliant}_{\Phi_{\le N}}(\mathbb I_{\le N})\;:\;\texttt{accepted}\quad\Longleftrightarrow\quad\delta_{\le N}^{\mathrm{stack}}=\varnothing.
\]

\emph{Clause 3 (Non-finality).} By Theorem 24 of Subsection 11.5, the same-level self-soundness claim \(\mathrm{Sound}\bigl(I_{N+1}^{\mathrm{stack}}\bigr)\) does not receive verdict \(\texttt{accepted}\) under \(I_{N+1}^{\mathrm{stack}}\) in the absence of an admissible level shift to a strictly higher instrument. The rotating-audit acceptance of \(\mathrm{Audit}(\mathbb I_{\le N})\) concerns the lower stack; it is not an audit of \(I_{N+1}^{\mathrm{stack}}\) by a still higher instrument, and so does not supply a level shift over \(I_{N+1}^{\mathrm{stack}}\). Hence

\[
\Gamma\,;\;\mathsf H_{\mathrm{instr}}\,;\;I_{N+1}^{\mathrm{stack}}\;\vdash\;\mathrm{Audit}(\mathbb I_{\le N})\;:\;\texttt{accepted}\quad\not\Longrightarrow\quad\Gamma\,;\;\mathsf H_{\mathrm{instr}}\,;\;I_{N+1}^{\mathrm{stack}}\;\vdash\;\mathrm{Sound}\bigl(I_{N+1}^{\mathrm{stack}}\bigr)\;:\;\texttt{accepted}.
\]

This completes the proof. \(\square\)

\subsubsection{Nonfinality}

Theorem 25 leaves the higher instrument \(I_{N+1}^{\mathrm{stack}}\) uncertified by itself. To audit \(I_{N+1}^{\mathrm{stack}}\), one treats

\[
\mathbb I_{\le N+1}\;=\;(I_0,\,I_1,\,\ldots,\,I_N,\,I_{N+1}^{\mathrm{stack}})
\]

as the next finite lower stack and adds a further admissible rotating-audit bridge \(B^{\mathrm{audit}}_{N+1\to N+2}\) to a strictly higher instrument \(I_{N+2}^{\mathrm{stack}}\), producing the extended chain

\[
\mathbb I_{\le N+2}\;=\;(I_0,\,I_1,\,\ldots,\,I_N,\,I_{N+1}^{\mathrm{stack}},\,I_{N+2}^{\mathrm{stack}}),
\]

and reapplies Theorem 25 with \(N\) replaced by \(N+1\). The construction is finite at every step but unbounded as a procedure: there is no finite tower with a final self-certifying top, because each one-level extension inherits Theorem 24's same-level self-audit failure for its own soundness claim. This is the rotating-audit non-finality clause \(\mathrm{NC}\text{-}17\) of Subsection 2.4.3.

The instrument-family completeness claim, in the sense of an instrument that audits every higher level of itself, is therefore not within the scope of \(\mathbf{BirdInt}^{\mathrm{aud}}_{\mathrm{fin}}\). Any positive soundness statement about an instrument \(I_n\) requires an admissible level-shifted audit by a strictly higher instrument \(I_{n+1}\) via a rotating-audit bridge of Subsection 11.7; the framework records the chain, the bridge, the admissibility predicate, and the defect, and accepts \(\mathrm{Compliant}_{\Phi_{\le N}}(\mathbb I_{\le N})\) iff \(\delta_{\le N}^{\mathrm{stack}}=\varnothing\), but never accepts \(\mathrm{Sound}(I_n)\) from the same instrument \(I_n\) itself.

\subsection{Instrument Transfer}

This subsection records the instrument-transfer rule that governs whether a claim accepted under one instrument may be transported to another, and marks the corresponding full classifier as future work. The rule supports the instrument-transfer nonclaim recorded in Subsection 2.4.3 and in the self-reference packet's NC-19 form: \(I\vdash\varphi\) does not imply \(J\vdash\varphi\).

\subsubsection{Transfer Rule}

Let \(I\) and \(J\) be admissible instruments in the sense of Subsection 11.1.1. Suppose

\[
\Gamma\,;\;\mathcal T\,;\;I\;\vdash\;\varphi\;:\;\texttt{accepted}.
\]

Acceptance under \(I\) does not, by itself, license any claim status under \(J\): the default is

\[
\Gamma\,;\;\mathcal T\,;\;I\;\vdash\;\varphi\;:\;\texttt{accepted}\quad\not\Longrightarrow\quad\Gamma\,;\;\mathcal T\,;\;J\;\vdash\;\varphi\;:\;\texttt{accepted}.
\]

A non-default transfer requires an instrument-transfer bridge

\[
B_{I\to J}\;=\;(\,\mathrm{id}_B,\;I,\;J,\;\varphi,\;L_{I\to J},\;V_{I\to J},\;\Theta_{I\to J},\;A_{I\to J},\;\delta_{I\to J},\;\mathcal N_{I\to J},\;\lambda_{\mathrm{transfer}}\,),
\]

a directed bridge from \(I\) to \(J\) admitted by the surrounding bridge policies, with bridge type \(\texttt{instrument-transfer}\) and components recording, respectively: the source claim, the language map that translates \(\varphi\) into \(J\)'s claim language, the visibility map exposing the records that supported \(\varphi\) under \(I\) to \(J\), the audit thresholds \(J\) applies to the transfer, the audit record of the bridge, the transfer defect, the bridge nonclaim register, and the bridge profile. Admissibility of \(B_{I\to J}\) for the source claim \(\varphi\), written

\[
\mathsf{AdmInstrBridge}(B_{I\to J},\varphi),
\]

requires \(I\) and \(J\) admissible, \(\varphi\) in the claim types of both instruments via \(L_{I\to J}\), the supporting records visible to \(J\) either directly or through visibility bridges of Subsection 11.3.3, audit thresholds declared, the bridge audited, no unstratified circular dependency between \(\varphi\) in \(I\) and the records that support its transfer to \(J\), and the nonclaim register present.

The instrument-transfer rule is then

\[
\frac{\Gamma\,;\;\mathcal T\,;\;I\;\vdash\;\varphi\;:\;\texttt{accepted}\quad\quad\mathsf{AdmInstrBridge}(B_{I\to J},\varphi)}{\Gamma\,;\;\mathcal T\,;\;J\;\vdash\;B_{I\to J}(\varphi)\;:\;\chi},
\]

where the conclusion's status \(\chi\) ranges over

\[
\chi\;\in\;\{\,\texttt{accepted},\;\texttt{provisional},\;\texttt{blocked}\,\}
\]

and is determined by the classifier of Subsection 11.2.1 applied to the transferred claim record under \(J\). Acceptance under \(I\) thus enters the transfer as a premise; the verdict under \(J\) is decided by \(J\)'s own classifier on the bridged claim, with \(\texttt{accepted}\), \(\texttt{provisional}\), or \(\texttt{blocked}\) depending on whether \(J\)'s visibility, threshold, audit, and source checks pass on the bridged content. The rule preserves \(\mathrm{NC}\text{-}19\): the source acceptance and the bridge admissibility are necessary but not sufficient conditions for \(\texttt{accepted}\) under \(J\).

\subsubsection{Future-Work Boundary}

A complete instrument-transfer classifier --- one that fixes, for every admissible \(B_{I\to J}\) and every claim type, the exact mapping from the source-side acceptance and the bridge defect record to the value \(\chi\in\{\texttt{accepted},\texttt{provisional},\texttt{blocked}\}\) on the target side --- is not stated as a theorem of Foundations III. Subsection 11.9.1 fixes the rule schema and the admissibility predicate; the per-claim-type classifier and its associated soundness theorem (the analogue of Theorem 23 for transferred claims) are deferred. This is consistent with the scope nonclaims of Subsection 2.4 and with the deferred-claims register of Appendix H: the calculus accepts only what is supported by an admissible bridge and a passing classifier verdict, and does not commit to a closed-form transfer-completeness statement.

The framework's default position is therefore the negative one: \(I\vdash\varphi\) does not imply \(J\vdash\varphi\). Positive transfer is recorded only by an admissible instrument-transfer bridge whose verdict under \(J\) is independently classified. Future-work strengthenings of Subsection 11.9.1 would supply the per-claim-type completeness theorem and the corresponding countermodel atlas; both are outside the scope of the present paper.
 \section{Abstract Interpretation Model Family}

This section records abstract interpretation as a finite model family for \(\mathbf{BirdInt}^{\mathrm{aud}}_{\mathrm{fin}}\) and proves the four model-family theorems that establish its scope. Subsection 12.1 fixes the finite host data and the realized role fragment. Subsections 12.2--12.4 prove Theorems 26, 27, and 28: closure as packaging, descent as sound transformer, and representability as fixedness. Subsection 12.5 proves Model Theorem 29, that a finite family of admissible abstract-interpretation hosts realizes the \(P_1/P_2/P_4/P_5/P_6\) fragment of the calculus. Subsection 12.6 records the abstract-interpretation nonclaims that fix the boundaries of this realization.

The discipline of the section is the model-realization discipline of Subsection 13.1: abstract interpretation realizes a declared fragment of the calculus and not the whole calculus. In particular, the section does not claim that abstract interpretation supplies the route mismatch primitive \(P_3\) outside the optional extension declared in Subsection 12.1.2, and does not claim that sound abstraction is complete representation, that closure implies transformer soundness, or that the AI fragment is a foundation for the calculus.

\subsection{Abstract-Interpretation Host}

This subsection fixes the finite host data --- concrete and abstract domains, abstraction and concretization maps, the closure operator, sound transformers, and audit records --- and identifies the role fragment that the host realizes inside \(\mathbf{BirdInt}^{\mathrm{aud}}_{\mathrm{fin}}\).

\subsubsection{Host Data}

A finite abstract-interpretation host is a tuple

\[
H\;=\;(\,C,\;A,\;\alpha,\;\gamma,\;\rho,\;F,\;F^{\sharp},\;A_{\mathrm{AI}}\,)
\]

with the following components:

\begin{itemize}
\tightlist
\item
  \((C,\le_C)\) is a finite concrete domain: a finite set \(C\) together with a partial order \(\le_C\). Concreteness and finiteness are imposed throughout: every supremum, infimum, and fixed point used by the host is taken in a finite poset and is decidable on finite records.
\item
  \((A,\le_A)\) is a finite abstract domain: a finite set \(A\) together with a partial order \(\le_A\).
\item
  \(\alpha:C\to A\) is a monotone abstraction map, and \(\gamma:A\to C\) is a monotone concretization map. The pair forms a Galois connection
  \[
  \alpha(c)\;\le_A\;a\quad\Longleftrightarrow\quad c\;\le_C\;\gamma(a),
  \]
  written \(\alpha\dashv\gamma\). Equivalently, \(\alpha\) and \(\gamma\) satisfy unit \(c\le_C\gamma\alpha(c)\) and counit \(\alpha\gamma(a)\le_A a\).
\item
  \(\rho:C\to C\) is the host closure operator, defined by \(\rho=\gamma\circ\alpha\). The fixedness set is \(\mathrm{Fix}(\rho)=\{c\in C:\rho(c)=c\}\).
\item
  \(F:C\to C\) is a monotone concrete transformer, and \(F^{\sharp}:A\to A\) is a monotone abstract transformer. The host carries the soundness condition that \(\alpha F\le_A F^{\sharp}\alpha\), equivalently \(F\gamma\le_C\gamma F^{\sharp}\) in the sense fixed in Subsection 12.3.
\item
  \(A_{\mathrm{AI}}\) is a finite audit record on the host data: it records the verifications that \(C\) and \(A\) are finite ordered domains, that \(\alpha\) and \(\gamma\) are monotone, that \(\alpha\dashv\gamma\), that \(\rho=\gamma\alpha\), that \(F\) and \(F^{\sharp}\) are monotone, and that the soundness inequality holds. The audit record is itself finite and indexed by a host-level instrument \(I_{\mathrm{AI}}\) in the sense of Subsection 11.1.
\end{itemize}

The host is admissible, written \(\mathsf{AdmAIHost}(H)\), when each clause above is recorded in \(A_{\mathrm{AI}}\) and verified by the finite checks of \(I_{\mathrm{AI}}\). Admissibility is a syntactic, finite predicate on \(H\).

The instrument-indexed claim form for any AI-host theorem in this section is

\[
\Gamma\,;\;\mathsf H_{\mathrm{AI}}\,;\;I_{\mathrm{AI}}\;\vdash\;\varphi\;:\;\chi,
\]

with \(\mathsf H_{\mathrm{AI}}\) the AI-host context recording \(H\) and \(I_{\mathrm{AI}}\) the instrument that audits the host.

\subsubsection{BirdInt Realized Fragment}

A finite admissible abstract-interpretation host \(H\) realizes the role fragment

\[
P_1\;/\;P_2\;/\;P_4\;/\;P_5\;/\;P_6
\]

of \(\mathbf{BirdInt}^{\mathrm{aud}}_{\mathrm{fin}}\), with the realization map fixed as follows:

\begin{itemize}
\tightlist
\item
  \(P_5\) (packaging) is realized by the closure operator \(\rho=\gamma\alpha\). Theorem 26 of Subsection 12.2 records that \(\rho\) is extensive, monotone, and idempotent, hence supplies the finite packaging/objecthood record for the AI host.
\item
  \(P_1\) (descent) is realized by the soundness inequality \(\alpha F\le_A F^{\sharp}\alpha\), equivalently \(F\gamma\le_C\gamma F^{\sharp}\). Theorem 27 of Subsection 12.3 records that the equivalence is the descent-square soundness condition for the AI host.
\item
  \(P_2\) (representability) is realized by exact representability \(p\,\text{representable}\iff\exists a\in A:\gamma(a)=p\). Theorem 28 of Subsection 12.4 records the equivalence with fixedness \(\rho(p)=p\) and with \(\mathrm{im}(\gamma)\).
\item
  \(P_4\) (refinement) is realized by audited refinement records between abstract domains within a finite family of hosts: a refinement record \(A_k\leadsto A_\ell\) is a finite audited bridge comparing the precision of two abstract domains within \(\mathsf H_{\mathrm{AI}}\).
\item
  \(P_6\) (audit) is realized by the host audit record \(A_{\mathrm{AI}}\) of Subsubsection 12.1.1, indexed by \(I_{\mathrm{AI}}\).
\end{itemize}

The role \(P_3\) (route mismatch) is not realized by the bare AI host. The route mismatch primitive of \(\mathbf{BirdInt}^{\mathrm{aud}}_{\mathrm{fin}}\) requires a route record and a route-mismatch defect, neither of which is supplied by a single Galois pair. An optional AI extension may add finite route records and route-mismatch checks on top of \(H\); in that extension and only in that extension does \(P_3\) enter the realized fragment. The present section restricts attention to the bare host and fixes the realized fragment as \(P_1/P_2/P_4/P_5/P_6\).

The fragment discipline of Subsection 13.1 then applies: realization of the fragment \(P_1/P_2/P_4/P_5/P_6\) does not imply realization of the whole calculus, and the abstract-interpretation host is not a foundation for \(\mathbf{BirdInt}^{\mathrm{aud}}_{\mathrm{fin}}\) but a model family within it.

\subsection{Theorem 26: Closure as Packaging}

This subsection proves Theorem 26: in any finite admissible abstract-interpretation host, the composite \(\rho=\gamma\alpha\) is a closure operator and so realizes the \(P_5\) packaging role of \(\mathbf{BirdInt}^{\mathrm{aud}}_{\mathrm{fin}}\).

\subsubsection{Statement}

Let \(H=(C,A,\alpha,\gamma,\rho,F,F^{\sharp},A_{\mathrm{AI}})\) be a finite admissible abstract-interpretation host of Subsection 12.1.1, with Galois connection \(\alpha\dashv\gamma\) and \(\rho=\gamma\alpha\).

\textbf{Theorem 26 (AIClosureAsPackaging; theorem-grade).} The map \(\rho:C\to C\) is a closure operator on \((C,\le_C)\): it is extensive, monotone, and idempotent, that is,

\[
c\;\le_C\;\rho(c)\quad\text{(extensive)},
\]

\[
c\;\le_C\;c'\;\Longrightarrow\;\rho(c)\;\le_C\;\rho(c')\quad\text{(monotone)},
\]

\[
\rho^{2}\;=\;\rho\quad\text{(idempotent)}.
\]

The instrument-indexed form is

\[
\Gamma\,;\;\mathsf H_{\mathrm{AI}}\,;\;I_{\mathrm{AI}}\;\vdash\;\textsc{AIClosureAsPackaging}\;:\;\texttt{accepted}.
\]

In particular, \(\rho\) is a finite packaging operator on \(C\) in the sense required by the \(P_5\) role of \(\mathbf{BirdInt}^{\mathrm{aud}}_{\mathrm{fin}}\).

\subsubsection{Proof}

The three closure-operator clauses follow directly from the Galois unit, the monotonicity of \(\alpha\) and \(\gamma\), and the Galois counit, applied to the finite ordered domains \((C,\le_C)\) and \((A,\le_A)\).

\emph{Extensivity.} The Galois unit asserts \(c\le_C\gamma\alpha(c)\) for every \(c\in C\). By definition \(\rho(c)=\gamma\alpha(c)\), so \(c\le_C\rho(c)\).

\emph{Monotonicity.} Suppose \(c\le_C c'\). Monotonicity of \(\alpha\) gives \(\alpha(c)\le_A\alpha(c')\). Monotonicity of \(\gamma\) then gives \(\gamma\alpha(c)\le_C\gamma\alpha(c')\), that is, \(\rho(c)\le_C\rho(c')\).

\emph{Idempotence.} Apply the Galois counit at \(a=\alpha(c)\): \(\alpha\gamma\alpha(c)\le_A\alpha(c)\). Monotonicity of \(\gamma\) gives \(\gamma\alpha\gamma\alpha(c)\le_C\gamma\alpha(c)\), that is, \(\rho^{2}(c)\le_C\rho(c)\). Conversely, extensivity applied at \(\rho(c)\) gives \(\rho(c)\le_C\rho(\rho(c))=\rho^{2}(c)\). Antisymmetry of \(\le_C\) on the finite poset \((C,\le_C)\) yields \(\rho^{2}(c)=\rho(c)\) for every \(c\in C\), hence \(\rho^{2}=\rho\).

Each clause is verified by a finite check on the finite host data \(H\) and recorded in the host audit record \(A_{\mathrm{AI}}\). Under the admissibility assumptions for \(H\) and the instrument \(I_{\mathrm{AI}}\), the in-scope, visibility, evidence, threshold, audit, and nonclaim conditions for the headline claim pass, so the instrument-indexed verdict is \(\texttt{accepted}\). \(\square\)

The closure operator \(\rho\) is the host's \(P_5\) packaging witness: it supplies the finite extensive, monotone, idempotent map on \(C\) that the realization map of Subsection 12.1.2 requires. Theorems 27 and 28 of Subsections 12.3 and 12.4 then state how the host's \(P_1\) and \(P_2\) roles are read off from \(\alpha\), \(\gamma\), \(F\), \(F^{\sharp}\), and \(\rho\).

\subsection{Theorem 27: Descent as Sound Transformer}

This subsection proves Theorem 27: in any finite admissible abstract-interpretation host, the abstract-side soundness inequality \(\alpha F\le_A F^{\sharp}\alpha\) is equivalent to the concrete-side soundness inequality \(F\gamma\le_C\gamma F^{\sharp}\). The equivalence is the descent-square soundness condition for the AI host and supplies the host's \(P_1\) descent witness.

\subsubsection{Statement}

Let \(H=(C,A,\alpha,\gamma,\rho,F,F^{\sharp},A_{\mathrm{AI}})\) be a finite admissible abstract-interpretation host of Subsection 12.1.1, with Galois connection \(\alpha\dashv\gamma\) and monotone transformers \(F:C\to C\) and \(F^{\sharp}:A\to A\).

\textbf{Theorem 27 (AIDescentAsSoundTransformer; theorem-grade).} The two soundness inequalities

\[
\alpha F\;\le_A\;F^{\sharp}\alpha\qquad\text{and}\qquad F\gamma\;\le_C\;\gamma F^{\sharp}
\]

are equivalent: each holds iff the other does. Equivalently,

\[
\alpha F(c)\;\le_A\;F^{\sharp}\alpha(c)\;\;\text{for every}\;c\in C\quad\Longleftrightarrow\quad F\gamma(a)\;\le_C\;\gamma F^{\sharp}(a)\;\;\text{for every}\;a\in A.
\]

The instrument-indexed form is

\[
\Gamma\,;\;\mathsf H_{\mathrm{AI}}\,;\;I_{\mathrm{AI}}\;\vdash\;\textsc{AIDescentAsSoundTransformer}\;:\;\texttt{accepted}.
\]

The common content of the two inequalities is the host's \(P_1\) descent record: the finite, audited soundness condition for the abstract transformer \(F^{\sharp}\) over the concrete transformer \(F\) through the Galois package \((\alpha,\gamma,\rho)\).

\subsubsection{Proof}

Each direction follows from the Galois adjunction together with monotonicity of \(\alpha\), \(\gamma\), \(F\), and \(F^{\sharp}\).

\emph{Forward direction.} Assume \(\alpha F\le_A F^{\sharp}\alpha\). Fix \(a\in A\); we show \(F\gamma(a)\le_C\gamma F^{\sharp}(a)\). By the Galois adjunction \(\alpha\dashv\gamma\), this is equivalent to \(\alpha F\gamma(a)\le_A F^{\sharp}(a)\). Apply the assumption at \(c=\gamma(a)\):

\[
\alpha F\gamma(a)\;\le_A\;F^{\sharp}\alpha\gamma(a).
\]

The Galois counit gives \(\alpha\gamma(a)\le_A a\), and monotonicity of \(F^{\sharp}\) gives \(F^{\sharp}\alpha\gamma(a)\le_A F^{\sharp}(a)\). Combining the two inequalities yields \(\alpha F\gamma(a)\le_A F^{\sharp}(a)\). Re-applying the adjunction returns \(F\gamma(a)\le_C\gamma F^{\sharp}(a)\), as required.

\emph{Reverse direction.} Assume \(F\gamma\le_C\gamma F^{\sharp}\). Fix \(c\in C\); we show \(\alpha F(c)\le_A F^{\sharp}\alpha(c)\). By the Galois adjunction, this is equivalent to \(F(c)\le_C\gamma F^{\sharp}\alpha(c)\). The Galois unit gives \(c\le_C\gamma\alpha(c)\), and monotonicity of \(F\) gives \(F(c)\le_C F\gamma\alpha(c)\). Apply the assumption at \(a=\alpha(c)\):

\[
F\gamma\alpha(c)\;\le_C\;\gamma F^{\sharp}\alpha(c).
\]

Combining the two inequalities yields \(F(c)\le_C\gamma F^{\sharp}\alpha(c)\). Re-applying the adjunction returns \(\alpha F(c)\le_A F^{\sharp}\alpha(c)\), as required.

Each step uses only the Galois unit, the Galois counit, monotonicity of \(\alpha\), \(\gamma\), \(F\), \(F^{\sharp}\), and antisymmetry/transitivity of \(\le_C\) and \(\le_A\) on the finite posets. Every check is finite and is recorded in the host audit record \(A_{\mathrm{AI}}\). Under the admissibility assumptions for \(H\) and the instrument \(I_{\mathrm{AI}}\), the in-scope, visibility, evidence, threshold, audit, and nonclaim conditions for the headline claim pass, so the instrument-indexed verdict is \(\texttt{accepted}\). \(\square\)

The equivalence supplies the host's \(P_1\) descent witness in either form: the abstract-side \(\alpha F\le_A F^{\sharp}\alpha\) is the upper-soundness inequality used when reasoning about \(F^{\sharp}\) as an over-approximation of \(F\) on the abstract side, and the concrete-side \(F\gamma\le_C\gamma F^{\sharp}\) is the descent inequality used when reasoning about the concretized abstract step as a finite over-approximation of the concrete step. The two inequalities are interchangeable under the Galois adjunction, and the host carries either form indifferently as its \(P_1\) record.

\subsection{Theorem 28: Representability as Fixedness}

This subsection proves Theorem 28: in any finite admissible abstract-interpretation host, exact representability of a concrete element coincides with fixedness under the closure operator \(\rho\) and with membership in the image of the concretization map \(\gamma\). The equivalence supplies the host's \(P_2\) representability witness.

\subsubsection{Statement}

Let \(H=(C,A,\alpha,\gamma,\rho,F,F^{\sharp},A_{\mathrm{AI}})\) be a finite admissible abstract-interpretation host of Subsection 12.1.1, with Galois connection \(\alpha\dashv\gamma\) and \(\rho=\gamma\alpha\). For \(p\in C\), define exact representability by

\[
p\;\text{representable}\quad\Longleftrightarrow\quad\exists\,a\in A:\;\gamma(a)=p.
\]

\textbf{Theorem 28 (AIRepresentabilityAsFixedness; theorem-grade).} For every \(p\in C\),

\[
p\;\text{representable}\quad\Longleftrightarrow\quad\rho(p)=p.
\]

Equivalently,

\[
\mathrm{Rep}_{\alpha}(C)\;=\;\mathrm{Fix}(\rho)\;=\;\mathrm{im}(\gamma),
\]

where \(\mathrm{Rep}_{\alpha}(C)=\{p\in C:p\text{ representable}\}\), \(\mathrm{Fix}(\rho)=\{p\in C:\rho(p)=p\}\), and \(\mathrm{im}(\gamma)=\{p\in C:\exists\,a\in A,\;\gamma(a)=p\}\).

The instrument-indexed form is

\[
\Gamma\,;\;\mathsf H_{\mathrm{AI}}\,;\;I_{\mathrm{AI}}\;\vdash\;\textsc{AIRepresentabilityAsFixedness}\;:\;\texttt{accepted}.
\]

The theorem identifies three finite, equally decidable subsets of \(C\) --- the representable elements, the \(\rho\)-fixed points, and the image of \(\gamma\) --- and so supplies the host's \(P_2\) representability witness as a finite, audit-checkable predicate on \(C\).

\subsubsection{Proof}

The two implications follow from the Galois adjunction \(\alpha\dashv\gamma\) and the definition \(\rho=\gamma\alpha\).

\emph{Forward direction.} Suppose \(p\) is representable, so there exists \(a\in A\) with \(\gamma(a)=p\). Then

\[
\rho(p)\;=\;\gamma\alpha\gamma(a).
\]

The standard Galois-adjunction identity \(\gamma\alpha\gamma=\gamma\) follows from the unit and counit. The counit gives \(\alpha\gamma(a)\le_A a\); applying the monotone map \(\gamma\) gives

\[
\gamma\alpha\gamma(a)\;\le_C\;\gamma(a).
\]

The unit at the concrete element \(\gamma(a)\) gives

\[
\gamma(a)\;\le_C\;\gamma\alpha\gamma(a).
\]

By antisymmetry, \(\gamma\alpha\gamma(a)=\gamma(a)\). Hence \(\rho(p)=\gamma\alpha\gamma(a)=\gamma(a)=p\).

\emph{Reverse direction.} Suppose \(\rho(p)=p\), so \(\gamma\alpha(p)=p\). Set \(a=\alpha(p)\in A\). Then \(\gamma(a)=\gamma\alpha(p)=p\), so \(p\) is representable.

The two directions together yield \(\mathrm{Rep}_{\alpha}(C)=\mathrm{Fix}(\rho)\). The remaining equality \(\mathrm{Fix}(\rho)=\mathrm{im}(\gamma)\) is now immediate: if \(p=\rho(p)\), then \(p=\gamma\alpha(p)\in\mathrm{im}(\gamma)\); conversely, if \(p=\gamma(a)\) for some \(a\in A\), then \(\rho(p)=\gamma\alpha\gamma(a)=\gamma(a)=p\), so \(p\in\mathrm{Fix}(\rho)\).

Each step uses only the Galois unit and counit and the identity \(\gamma\alpha\gamma=\gamma\) that is read off from them on the finite posets \((C,\le_C)\) and \((A,\le_A)\). Every check is finite and is recorded in the host audit record \(A_{\mathrm{AI}}\). Under the admissibility assumptions for \(H\) and the instrument \(I_{\mathrm{AI}}\), the in-scope, visibility, evidence, threshold, audit, and nonclaim conditions for the headline claim pass, so the instrument-indexed verdict is \(\texttt{accepted}\). \(\square\)

The fixed set \(\mathrm{Fix}(\rho)\) is the host's \(P_2\) representability witness. Together with Theorem 26 (which makes \(\rho\) a packaging operator) and Theorem 27 (which makes the soundness inequalities equivalent), Theorem 28 closes the local circle: representable elements are exactly the \(\rho\)-fixed elements, which are exactly the \(\gamma\)-images, and packaging, descent, and representability are read off from \(\rho\), \(F^{\sharp}\), and \(\gamma\) via the Galois adjunction.

\subsection{Model Theorem 29: AI Model-Family Theorem}

This subsection proves Model Theorem 29: a finite family of admissible abstract-interpretation hosts realizes the \(P_1/P_2/P_4/P_5/P_6\) fragment of \(\mathbf{BirdInt}^{\mathrm{aud}}_{\mathrm{fin}}\). The theorem assembles the per-host realizations of \(P_5\), \(P_1\), and \(P_2\) (Theorems 26, 27, 28) with the audited refinement records between hosts and the per-host audit records. The theorem is model-grade rather than theorem-grade: its acceptance verdict is \(\texttt{model\_realization}\), not \(\texttt{accepted}\), in the sense fixed in Subsection 13.1.

\subsubsection{Statement}

Let

\[
\mathsf{AIFam}_{\mathrm{fin}}\;=\;\{\,H_k\,\}_{k\in K}
\]

be a finite family, indexed by a finite set \(K\), of admissible abstract-interpretation hosts in the sense of Subsection 12.1.1:

\[
H_k\;=\;(\,C,\;A_k,\;\alpha_k,\;\gamma_k,\;\rho_k,\;F,\;F_k^{\sharp},\;A_{\mathrm{AI},k}\,),
\]

with \(\alpha_k\dashv\gamma_k\) and \(\rho_k=\gamma_k\alpha_k\) for each \(k\in K\), and with \(\mathsf{AdmAIHost}(H_k)\) holding for each \(k\). Let \(\mathsf{Refine}=\{(k,\ell,R_{k\to\ell})\}\) be a finite family of audited refinement records between selected hosts \(H_k\) and \(H_\ell\), each \(R_{k\to\ell}\) recording a finite, audited comparison of the precision of \(A_k\) and \(A_\ell\) with the same concrete domain \(C\).

\textbf{Model Theorem 29 (AIModelFamily; model-grade).} The finite family \((\mathsf{AIFam}_{\mathrm{fin}},\mathsf{Refine})\) realizes the role fragment \(P_1/P_2/P_4/P_5/P_6\) of \(\mathbf{BirdInt}^{\mathrm{aud}}_{\mathrm{fin}}\) via the realization map

\[
P_5\;\mapsto\;\rho_k,\qquad P_1\;\mapsto\;\alpha_k F\;\le_{A_k}\;F_k^{\sharp}\alpha_k,\qquad P_2\;\mapsto\;\bigl(p\;\text{representable}\;\Longleftrightarrow\;\rho_k(p)=p\bigr),
\]

\[
P_4\;\mapsto\;A_k\leadsto A_\ell\;\text{via}\;R_{k\to\ell},\qquad P_6\;\mapsto\;A_{\mathrm{AI},k}.
\]

In schematic form,

\[
\mathsf{AIFam}_{\mathrm{fin}}\;\models\;\mathbf{BirdInt}^{\mathrm{aud}}_{\mathrm{fin}}{}^{P_1/P_2/P_4/P_5/P_6}.
\]

The instrument-indexed form is

\[
\Gamma\,;\;\mathsf H_{\mathrm{AI}}\,;\;I_{\mathrm{AI}}\;\vdash\;\mathsf{AIFam}_{\mathrm{fin}}\;\models\;\mathbf{BirdInt}^{\mathrm{aud}}_{\mathrm{fin}}{}^{P_1/P_2/P_4/P_5/P_6}\;:\;\texttt{model\_realization}.
\]

The verdict \(\texttt{model\_realization}\) records that the family supplies an admissible model-realization map for the declared fragment, in the sense of the model-realization convention of Subsection 13.1, and is not the \(\texttt{accepted}\) verdict of a theorem-grade claim.

\subsubsection{Proof}

The proof verifies, host by host and bridge by bridge, that the components required by the fragment are present, finite, and audited.

\emph{\(P_5\) packaging.} For each \(k\in K\), Theorem 26 of Subsection 12.2 applied to \(H_k\) asserts that \(\rho_k=\gamma_k\alpha_k\) is a finite closure operator on \((C,\le_C)\). Hence the assignment \(P_5\mapsto\rho_k\) supplies a per-host packaging witness recorded in \(A_{\mathrm{AI},k}\).

\emph{\(P_1\) descent.} For each \(k\in K\), Theorem 27 of Subsection 12.3 applied to \(H_k\) asserts that the soundness inequality \(\alpha_k F\le_{A_k}F_k^{\sharp}\alpha_k\) is equivalent to its concrete-side form \(F\gamma_k\le_C\gamma_k F_k^{\sharp}\), and so supplies the host's \(P_1\) descent witness in either form. The assignment \(P_1\mapsto(\alpha_k F\le_{A_k}F_k^{\sharp}\alpha_k)\) records the abstract-side form on \(H_k\), recorded in \(A_{\mathrm{AI},k}\).

\emph{\(P_2\) representability.} For each \(k\in K\), Theorem 28 of Subsection 12.4 applied to \(H_k\) asserts \(p\,\text{representable}\iff\rho_k(p)=p\) for every \(p\in C\), with \(\mathrm{Rep}_{\alpha_k}(C)=\mathrm{Fix}(\rho_k)=\mathrm{im}(\gamma_k)\). The assignment \(P_2\mapsto(p\,\text{representable}\iff\rho_k(p)=p)\) records the host's \(P_2\) representability witness, recorded in \(A_{\mathrm{AI},k}\).

\emph{\(P_4\) refinement.} Each refinement record \(R_{k\to\ell}\in\mathsf{Refine}\) is a finite audited bridge of the form fixed in Subsection 12.1.2: it records a finite comparison of the abstract domains \(A_k\) and \(A_\ell\) with the same concrete domain \(C\), with audit conditions internal to \(\mathsf H_{\mathrm{AI}}\) and audit record contributing to the family-level audit aggregation. The assignment \(P_4\mapsto(A_k\leadsto A_\ell)\) thus supplies the \(P_4\) refinement witness on every recorded pair \((k,\ell)\).

\emph{\(P_6\) audit.} Each per-host audit record \(A_{\mathrm{AI},k}\) supplies the \(P_6\) audit witness for \(H_k\), indexed by the family-level instrument \(I_{\mathrm{AI}}\). The aggregation \(\bigcup_{k\in K}A_{\mathrm{AI},k}\) together with the bridge audit records of \(\mathsf{Refine}\) is the family-level audit witness.

\emph{Closure of the realization map.} The five assignments above are typed: each maps a role of the declared fragment to a finite, audited datum on \(\mathsf{AIFam}_{\mathrm{fin}}\) and \(\mathsf{Refine}\). No assignment is supplied for \(P_3\), in keeping with the bare-host scope fixed in Subsection 12.1.2. The instrument \(I_{\mathrm{AI}}\) classifies the model-realization claim under the model-realization classifier of Subsection 13.1 once the per-host admissibility, the per-host theorems, the bridge admissibility, and the family-level audit records have all been recorded.

The verdict is therefore \(\texttt{model\_realization}\), in the sense fixed in Subsection 13.1: the family realizes the declared fragment, and not the whole calculus. \(\square\)

The model-realization verdict is the model-grade analogue of the \(\texttt{accepted}\) verdict of Theorems 26, 27, and 28: it records that the host family supplies a finite, audited realization of the declared role fragment, while no realization datum is supplied for the unrealized role \(P_3\) on the bare host. The nonclaims of Subsection 12.6 record the boundaries of this realization: the family does not realize \(P_3\), sound abstraction is not complete representation, and closure does not imply transformer soundness.

\subsection{AI Nonclaims}

This subsection records the abstract-interpretation nonclaims that fix the boundaries of the realization established in Subsections 12.2--12.5. Each nonclaim is a scoped negative statement attached to the AI model family: it is a record of what the family does not establish, in the sense of the nonclaim discipline of Subsection 2.4. The nonclaims here support the global nonclaim discipline of Section 2 and record the three AI-specific exclusions: no AI totality, no soundness-as-completeness, and no closure-implies-soundness.

\subsubsection{No AI Totality}

Abstract interpretation does not realize the entire six-role calculus \(\mathbf{BirdInt}^{\mathrm{aud}}_{\mathrm{fin}}\). Theorem 29 of Subsection 12.5 asserts only the realization of the fragment \(P_1/P_2/P_4/P_5/P_6\). The bare AI host of Subsection 12.1.1 does not supply a route record or a route-mismatch defect, and so does not realize the route-mismatch primitive \(P_3\) described in Subsection 2.2.2. A full six-role AI model would require additional finite \(P_3\) data; such extensions are outside the scope of the present section and are not asserted here.

In particular, the family-level claim

\[
\Gamma\,;\;\mathsf H_{\mathrm{AI}}\,;\;I_{\mathrm{AI}}\;\vdash\;\mathsf{AIFam}_{\mathrm{fin}}\;\models\;\mathbf{BirdInt}^{\mathrm{aud}}_{\mathrm{fin}}\;:\;\chi
\]

does not receive verdict \(\texttt{model\_realization}\) for the whole calculus on the bare host. The verdict \(\texttt{model\_realization}\) is assigned only to the declared fragment \(P_1/P_2/P_4/P_5/P_6\).

\subsubsection{No Soundness-as-Completeness}

Sound abstraction does not imply complete representation. The \(P_1\) descent witness of Theorem 27 records the inequality \(\alpha F\le_A F^{\sharp}\alpha\) (equivalently \(F\gamma\le_C\gamma F^{\sharp}\)): the abstract transformer \(F^{\sharp}\) over-approximates the concrete transformer \(F\). The corresponding equality \(\alpha F=F^{\sharp}\alpha\), which would assert that \(F^{\sharp}\) is a complete (exact) representation of \(F\) through the Galois package \((\alpha,\gamma,\rho)\), is not implied by Theorem 27 and is not part of the host's \(P_1\) realization datum.

The formal distinction is the inequality/equality distinction itself: \(\alpha F\le_A F^{\sharp}\alpha\) is the soundness condition, while \(\alpha F=F^{\sharp}\alpha\) is the exactness condition. The host therefore records the descent inequality, not the descent equality, and the AI fragment does not commit to completeness.

\subsubsection{No Closure-Implies-Soundness}

The closure-operator status of \(\rho=\gamma\alpha\) from Theorem 26 does not imply the soundness of the abstract transformer \(F^{\sharp}\) from Theorem 27. The Galois adjunction \(\alpha\dashv\gamma\) is sufficient to establish that \(\rho\) is extensive, monotone, and idempotent (Theorem 26), and in particular this establishes the \(P_5\) packaging role of the host. It does not, by itself, establish \(\alpha F\le_A F^{\sharp}\alpha\), which is an additional, separate condition on the pair \((F,F^{\sharp})\).

The host records the closure operator \(\rho\) and the soundness inequality \(\alpha F\le_A F^{\sharp}\alpha\) as independent data. Theorem 26 commits only to the closure-operator properties of \(\rho\), and Theorem 27 commits only to the equivalence of the two soundness inequalities; neither theorem implies the other. The audit record \(A_{\mathrm{AI}}\) of Subsection 12.1.1 records both checks separately, and the realization map of Subsection 12.1.2 maps \(P_5\) to \(\rho\) and \(P_1\) to the soundness inequality, not from one to the other.

The three nonclaims above together fix the local boundary of the AI model family within \(\mathbf{BirdInt}^{\mathrm{aud}}_{\mathrm{fin}}\): the family does not realize the whole calculus, sound abstraction is not exact representation, and the packaging role does not imply the descent role. The model-realization verdict of Theorem 29 is therefore scoped: it records realization of the declared fragment \(P_1/P_2/P_4/P_5/P_6\), and is silent on the rest.
 \section{Model Realizations}

This section records the four model realizations of \(\mathbf{BirdInt}^{\mathrm{aud}}_{\mathrm{fin}}\) retained for the first paper: PICA, Cantor strict-bridge, the abstract-interpretation host of Section 12, and the graph/cohomology drive support of Section 8. Subsection 13.1 fixes the model-realization convention and the fragment discipline that governs every model in the family. Subsections 13.2, 13.3, and 13.4 record the PICA, Cantor, and graph/cohomology realizations and their model theorems (Theorems 33 and 34 plus the model-grade reading of Section 8). Subsection 13.5 organizes the four realizations into a combined model-realization matrix. Subsection 13.6 records the cross-model role-coverage observation and the corresponding nonclaim.

The discipline of the section is uniform: each realization is scoped, each declares its host, instrument, source-of-truth policy, audit policy, visibility policy, and nonclaim register, and each realizes a declared fragment of the calculus. None of the four realizations is claimed to be the whole calculus, and the section does not state a global model-coverage theorem.

\subsection{Model-Realization Convention}

This subsection fixes the form of a model-realization statement and the fragment discipline imposed on every model realization in the section.

\subsubsection{Realization Map}

A model-realization statement has the form

\[
\Gamma\,;\;\mathcal T\,;\;I\;\vdash\;\mathcal M\;\models\;\mathcal S\;:\;\texttt{model\_realization},
\]

where:

\begin{itemize}
\tightlist
\item
  \(\mathcal M\) is the realization object: a finite typed record carrying the host data, the instrument, the source-of-truth policy, the audit policy, the visibility policy, and the nonclaim register required by the host;
\item
  \(\mathcal S\) is the declared structure: a finite, named fragment of \(\mathbf{BirdInt}^{\mathrm{aud}}_{\mathrm{fin}}\), written either as a role fragment (for example, \(P_1/P_2/P_4/P_5/P_6\)) or as a named substructure (for example, \(\mathrm{cell/pair/prov}\) for the finite stochastic cell/pair/provenance fragment, or \(0\to1\) for an audited strict-bridge passage);
\item
  \(I\) is the instrument under which the realization claim is recorded, and \(\mathcal T\) is the surrounding target package, in the sense of Subsection 11.1.
\end{itemize}

The verdict \(\texttt{model\_realization}\) records that the realization map of \(\mathcal M\) has supplied a finite, audited assignment of host data to each component of \(\mathcal S\), and that the per-component admissibility, audit, and nonclaim checks of the host instrument have passed. The verdict is a model-grade claim-strength verdict and is distinct from the per-claim status \(\texttt{accepted}\) in \(\mathrm{ClaimStatus}\).

The realization map of \(\mathcal M\) to \(\mathcal S\) is a finite typed family of assignments

\[
\mathrm{real}_{\mathcal M\to\mathcal S}\;=\;\bigl\{\,\mathrm{real}_X\;:\;X\in\mathrm{Components}(\mathcal S)\,\bigr\},
\]

with each \(\mathrm{real}_X\) mapping a host datum from \(\mathcal M\) to the corresponding component \(X\) of \(\mathcal S\). The realization map is admissible when:

\begin{enumerate}
\def\labelenumi{\arabic{enumi}.}
\tightlist
\item
  every component of \(\mathcal S\) is assigned a finite host datum from \(\mathcal M\);
\item
  each assignment is audited by the host instrument \(I\);
\item
  the source-of-truth policy of \(\mathcal M\) is consistent with the source policy required by \(\mathcal S\);
\item
  the visibility policy of \(\mathcal M\) exposes the records required by \(\mathcal S\), with visibility bridges of Subsection 11.3.3 supplied where required;
\item
  the nonclaim register of \(\mathcal M\) contains the model-realization nonclaims of Subsection 13.6 and any host-specific nonclaims of the realization.
\end{enumerate}

Admissibility of the realization map is a finite predicate on \((\mathcal M,\mathcal S,I)\).

\subsubsection{Fragment Discipline}

Every model realization in the section is scoped to a declared fragment \(\mathcal S\) of \(\mathbf{BirdInt}^{\mathrm{aud}}_{\mathrm{fin}}\). The fragment discipline imposes the following rules on every realization claim:

\begin{itemize}
\tightlist
\item
  \emph{Fragment-only realization.} A model-realization statement \(\mathcal M\models\mathcal S:\texttt{model\_realization}\) asserts realization of \(\mathcal S\) and not of the whole calculus. The verdict \(\texttt{model\_realization}\) for \(\mathcal S\) does not imply a \(\texttt{model\_realization}\) verdict for any larger fragment \(\mathcal S'\supsetneq\mathcal S\).
\item
  \emph{No theorem grade by realization.} A model-realization verdict is not a theorem-grade verdict. In the absence of a separately admissible theorem-grade bridge, \(\mathcal M\models\mathcal S:\texttt{model\_realization}\) does not yield a \(\texttt{accepted}\) verdict on a theorem-grade claim about \(\mathcal S\). This is the model-realization nonclaim of Subsection 2.4.
\item
  \emph{Host-bound realization.} A model-realization verdict for \(\mathcal S\) under \(\mathcal M\) is bound to the host of \(\mathcal M\): it does not transfer to other hosts without an admissible instrument-transfer bridge in the sense of Subsection 11.9.
\item
  \emph{No model-coverage theorem.} The four realizations recorded in Subsections 13.2--13.4 and Section 12 do not, taken together, support a global model-coverage theorem for \(\mathbf{BirdInt}^{\mathrm{aud}}_{\mathrm{fin}}\). This is the model-coverage nonclaim of Subsection 13.6.
\end{itemize}

The fragment discipline is the central structural constraint of the section: it fixes the meaning of \(\models\), prevents over-reading of the model-realization verdict, and defines the boundary between model-grade and theorem-grade claims throughout Sections 13 and 12. The four realizations of the section are each scoped to a declared fragment, and the four fragments are recorded in the combined matrix of Subsection 13.5.

\subsection{PICA Realization}

This subsection records the PICA model-realization sheet and Model Theorem 33: an admissible PICA realization object induces a model of the finite stochastic cell/pair/provenance fragment of \(\mathbf{BirdInt}^{\mathrm{aud}}_{\mathrm{fin}}\). The PICA realization is host-bound and fragment-scoped: it does not realize the whole calculus, the recovered \(6\times 6\) table is a PICA-profile fact rather than a universal law, and raw PICA statuses recover into formal statuses only via the host classifier.

\subsubsection{Name and Host}

The PICA realization name is \(\mathsf{PICAReal}_{\mathrm{fin}}\). Its host is the finite stochastic implementation context \(\mathsf H_{\mathrm{PICA}}\), with components

\[
\mathsf H_{\mathrm{PICA}}\;=\;(\,\Omega,\;K,\;\mathbb P,\;\mathsf{Bus},\;\mathsf{CellRules},\;\mathsf{WitnessStore},\;\mathsf{UpdateStore},\;\mathsf{PairStore},\;\mathsf{DepGraph},\;A,\;V\,),
\]

where \(\Omega\) is the finite stochastic state space, \(K\) the finite Markov kernel on \(\Omega\), \(\mathbb P=\{P_1,\ldots,P_6\}\) the six primitive labels, and the remaining fields finite stores for the bus, cell rules, witness records, actor-update records, pair observables, the producer/consumer dependency graph, the host audit record \(A\), and the host visibility record \(V\). The host instrument is \(I_{\mathrm{PICA}}\), an admissible instrument in the sense of Subsection 11.1.1 whose visible class contains the host stores and whose audit policy and source policy are those of Subsubsections 13.2.6 and 13.2.8.

\subsubsection{Realization Object}

A PICA realization object is the finite typed tuple

\[
\mathcal P\;=\;(\,\Omega,\;K,\;\mathbb P,\;\mathsf{Bus},\;\mathsf{CellRules},\;\mathsf{WitnessStore},\;\mathsf{UpdateStore},\;\mathsf{StatusStore},\;\mathsf{PairStore},\;\mathsf{DepGraph},\;\mathsf{Abl},\;A,\;V,\;\mathcal N\,),
\]

extending the host data with \(\mathsf{StatusStore}\) (raw and recovered statuses), \(\mathsf{Abl}\) (finite ablation/model evidence), and \(\mathcal N\) (the PICA nonclaim register of Subsubsection 13.2.10). The object is admissible, written \(\mathsf{AdmPICAReal}(\mathcal P)\), when each store is finite, every entry is audited under \(I_{\mathrm{PICA}}\), and the directed-cell, pair-observable, source/audit, and status-recovery records satisfy the schemas used in Subsections 4.3, 4.4, 4.8, and 5.6. Admissibility is a finite predicate on the stores.

The associated realization map is

\[
\mathcal R_{\mathrm{PICA}}\;:\;\mathsf{PICAReal}_{\mathrm{fin}}\longrightarrow
\mathbf{BirdInt}^{\mathrm{aud}}_{\mathrm{fin\text{-}stoch}}{}^{\,\mathrm{cell}/\mathrm{pair}/\mathrm{prov}},
\]

sending \(\mathcal P\) to the realized finite stochastic model \(\mathcal M_{\mathcal P}\). The induced package is

\[
\mathcal T_{\mathrm{PICA}}\;=\;(\,\Omega,\;f_{\mathrm{PICA}},\;\Sigma_{f_{\mathrm{PICA}}},\;E_{\mathrm{PICA}},\;\mathcal A_{\mathrm{PICA}}\,),
\]

where \(f_{\mathrm{PICA}}\) is the bus/lens projection, \(\Sigma_{f_{\mathrm{PICA}}}\) is the finite visible-claim set, \(E_{\mathrm{PICA}}\) is the commit/update closure field, and \(\mathcal A_{\mathrm{PICA}}\) is the inherited package-level audit functional assembled from the finite PICA audit record \(A\).

\subsubsection{Actor-Update Map}

The PICA actor-update map sends each primitive \(P_i\) to a finite update sort drawn from \(\mathsf{UpdateStore}\):

\[
P_i\;\longmapsto\;\mathsf{Upd}_{\mathrm{PICA}}(P_i),
\]

with the assignments

\begin{longtable}[]{@{}ll@{}}
\toprule
primitive & PICA actor update sort\tabularnewline
\midrule
\endhead
\(P_1\) & descent repair / closure update\tabularnewline
\(P_2\) & support gate / admissibility update\tabularnewline
\(P_3\) & route-memory / mismatch update\tabularnewline
\(P_4\) & refinement / staging update\tabularnewline
\(P_5\) & package / completion / closure-summary update\tabularnewline
\(P_6\) & audit / provenance / source update\tabularnewline
\bottomrule
\end{longtable}

Each row records a finite update sort consumed by the directed-cell record of Subsubsection 13.2.5; the assignment is a typed map from \(\mathbb P\) to the PICA update sorts, audited by \(I_{\mathrm{PICA}}\).

\subsubsection{Informant-Witness Map}

The PICA informant-witness map sends each primitive \(P_j\) to a finite witness sort drawn from \(\mathsf{WitnessStore}\):

\[
P_j\;\longmapsto\;\mathsf{Wit}_{\mathrm{PICA}}(P_j),
\]

with the assignments

\begin{longtable}[]{@{}ll@{}}
\toprule
primitive & PICA informant witness sort\tabularnewline
\midrule
\endhead
\(P_1\) & descent/closure defect\tabularnewline
\(P_2\) & gate/support/admissibility witness\tabularnewline
\(P_3\) & route mismatch / route-memory witness\tabularnewline
\(P_4\) & refinement/stage witness\tabularnewline
\(P_5\) & package/closure/completion witness\tabularnewline
\(P_6\) & audit/provenance/source witness\tabularnewline
\bottomrule
\end{longtable}

The actor-update and informant-witness maps together specify the per-cell witness/update content for the directed-cell realization of Subsubsection 13.2.5.

\subsubsection{Directed-Cell Realization}

Each PICA directed cell becomes a typed claim in the calculus:

\[
\Gamma\,;\;\mathcal T_{\mathrm{PICA}}\,;\;I_{\mathrm{PICA}}\;\vdash\;\mathsf{Cell}_{ij}^{\lambda_{\mathrm{PICA}}}\;:\;\sigma_{ij},
\]

with the cell record

\[
(P_i\leftarrow P_j,\;W_j,\;U_i,\;L,\;V,\;\Theta,\;A,\;\delta),
\]

where \(W_j=\mathsf{Wit}_{\mathrm{PICA}}(P_j)\), \(U_i=\mathsf{Upd}_{\mathrm{PICA}}(P_i)\), \(L\) is the language data of Subsection 4.3, \(V\) is the visibility record, \(\Theta\) the threshold record, \(A\) the audit record, and \(\delta\) the cell defect record. The cell profile \(\lambda_{\mathrm{PICA}}\) is the PICA cell profile, recording stochastic-host mode and the per-cell source/visibility/audit thresholds inherited from \(\mathsf H_{\mathrm{PICA}}\).

The classifier of Subsection 5.6.2, applied across the \(6\times 6\) grid of \((i,j)\) pairs, returns a PICA-profile multiplicity record

\[
25\;\texttt{action},\quad 8\;\texttt{implicit},\quad 2\;\texttt{trivial},\quad 1\;\texttt{undefined\_circular}.
\]

This count is a PICA-profile realization fact: it records the multiplicities of the four directed-cell statuses within the PICA host under the recovery map of Subsubsection 13.2.6, and is not a universal table theorem about \(6\times 6\) tables in \(\mathbf{BirdInt}^{\mathrm{aud}}_{\mathrm{fin}}\).

\subsubsection{Status Recovery}

Raw PICA statuses recover as formal statuses through the recovery map

\[
\texttt{Action}\;\longmapsto\;\texttt{action},\qquad\texttt{Implicit}\;\longmapsto\;\texttt{implicit},\qquad\texttt{Trivial}\;\longmapsto\;\texttt{trivial},
\]

\[
\texttt{Undefined}\;\longmapsto\;\text{reason-dependent formal status},
\]

where the reason-dependent recovery of \(\texttt{Undefined}\) is one of \(\{\,\texttt{blocked},\;\texttt{undefined\_circular},\;\texttt{outside\_scope},\;\texttt{absent},\;\texttt{below\_threshold},\;\texttt{collapsed}\,\}\), decided by the formal classifier of Subsection 5.6 acting on the cell record, the threshold record, the audit record, and the source-of-truth record.

Raw PICA status never overrides the formal classifier verdict: if the recovered formal status assigned by the classifier differs from the raw status, the formal status is the authoritative cell status, and the raw/recovered discrepancy is recorded in the source/provenance correction record of \(\mathsf{StatusStore}\).

\subsubsection{Pair-Observable Realization}

PICA pair statuses recover as formal pair statuses through

\[
\texttt{real}\;\longmapsto\;\texttt{real},\qquad\texttt{provisional}\;\longmapsto\;\texttt{real\_provisional},\qquad\texttt{duplicated}\;\longmapsto\;\texttt{real\_duplicated},
\]

\[
\texttt{blocked}\;\longmapsto\;\texttt{blocked},
\]

with a PICA pair observable record

\[
\mathsf{PairObs}_{\{i,j\}}^{\mu}\;=\;(\,\{P_i,P_j\},\;\mu,\;\mathsf{ObservableType},\;\mathsf{BranchData},\;\mathsf{SourceOfTruth},\;L,\;V,\;\Theta,\;A,\;\delta,\;\mathcal N\,).
\]

Realness of the pair observable requires the conjunction

\[
\mathsf{SourceOfTruthOK}\;\land\;\mathsf{VisibilityPasses}\;\land\;\mathsf{AuditPasses}\;\land\;\mathsf{ThresholdPasses},
\]

with \(\mathsf{SourceOfTruthOK}\) the source policy of Subsubsection 13.2.8 and the remaining three predicates the host-level visibility, audit, and threshold checks of Subsections 11.3, 4.8, and 5.2. A pair observable for which any of the four conjuncts fails recovers as one of the non-real PICA statuses: \(\texttt{real\_provisional}\), \(\texttt{real\_duplicated}\), or \(\texttt{blocked}\), depending on which conjunct fails and on the threshold and audit data.

\subsubsection{Provenance Discipline}

The PICA source-of-truth predicate is

\[
\mathsf{SourceOfTruthOK}(R,Q)\;\;\Longleftrightarrow\;\;\mathsf{Committed}(R)\;\land\;\mathsf{CorrectProducer}(R,Q)\;\land\;\mathsf{CorrectProvenance}(R,Q)
\]

\[
\land\;\mathsf{NotFallbackUnlessAllowed}(R,Q)\;\land\;\mathsf{Visible}(R)\;\land\;\mathsf{Audited}(R),
\]

for a record \(R\) and a query/observable \(Q\). Each conjunct is a finite check on the host stores: \(\mathsf{Committed}\) on \(\mathsf{Bus}\) and \(\mathsf{StatusStore}\); \(\mathsf{CorrectProducer}\) on \(\mathsf{DepGraph}\); \(\mathsf{CorrectProvenance}\) on \(A\); \(\mathsf{NotFallbackUnlessAllowed}\) on the host's fallback policy; \(\mathsf{Visible}\) on \(V\); \(\mathsf{Audited}\) on \(A\).

The provenance discipline imposes two non-negotiable rules:

\begin{itemize}
\tightlist
\item
  \emph{Dirty exclusion.} \(R.\mathsf{dirty\_since\_step}\neq\varnothing\;\Longrightarrow\;\neg\,\mathsf{SourceOfTruthOK}(R,Q)\). A dirty record cannot be a source of truth for any query, and pair observables that depend on a dirty record cannot be \(\texttt{real}\).
\item
  \emph{Fallback-default cap.} \(R\;\text{is fallback}\;\Longrightarrow\;Q\;\text{recovers as}\;\texttt{real\_provisional}\;\text{at most}\), unless an explicit upgrade is recorded in \(A\) via an admissible source-upgrade bridge. A fallback record supports at most a provisional pair observation.
\end{itemize}

The discipline ensures that pair-realness is grounded in committed, correct, audited, visible, non-dirty, and not-merely-fallback records; under any failure of these conditions, the pair recovers as one of the non-real statuses of Subsubsection 13.2.7.

\subsubsection{Model Theorem 33: PICA Model-Realization}

\textbf{Model Theorem 33 (PICA Model-Realization; model-grade).} Let \(\mathcal P\in\mathsf{PICAReal}_{\mathrm{fin}}\) be admissible: \(\mathsf{AdmPICAReal}(\mathcal P)\) holds. Then the realization map \(\mathcal R_{\mathrm{PICA}}\) supplies a finite, audited assignment of host data to every component of the finite stochastic cell/pair/provenance fragment, so

\[
\mathcal R_{\mathrm{PICA}}(\mathcal P)\;\models\;\mathbf{BirdInt}^{\mathrm{aud}}_{\mathrm{fin\text{-}stoch}}{}^{\,\mathrm{cell}/\mathrm{pair}/\mathrm{prov}}.
\]

The instrument-indexed form is

\[
\Gamma\,;\;\mathsf H_{\mathrm{PICA}}\,;\;I_{\mathrm{PICA}}\;\vdash\;\mathsf{PICAReal}_{\mathrm{fin}}\;\models\;\mathbf{BirdInt}^{\mathrm{aud}}_{\mathrm{fin\text{-}stoch}}{}^{\,\mathrm{cell}/\mathrm{pair}/\mathrm{prov}}\;:\;\texttt{model\_realization}.
\]

\emph{Proof.} Each component of the realized fragment is assigned a finite host datum and audit record from \(\mathcal P\):

\begin{itemize}
\tightlist
\item
  \emph{Typed directed cells.} Subsubsections 13.2.3 and 13.2.4 supply the actor-update and informant-witness maps; Subsubsection 13.2.5 assembles them into the directed-cell record under the cell profile \(\lambda_{\mathrm{PICA}}\).
\item
  \emph{Cell-status recovery.} Subsubsection 13.2.6 supplies the recovery map from raw to formal statuses, with the formal classifier of Subsection 5.6 as the authoritative verdict.
\item
  \emph{Pair-observable contracts.} Subsubsection 13.2.7 supplies the pair-observable record and the four-conjunct realness predicate.
\item
  \emph{Source-of-truth and provenance discipline.} Subsubsection 13.2.8 supplies the \(\mathsf{SourceOfTruthOK}\) predicate and the dirty-exclusion and fallback-default rules.
\item
  \emph{Finite dependency and ablation records.} The host stores \(\mathsf{DepGraph}\) and \(\mathsf{Abl}\) supply finite dependency and ablation records, recorded in \(A\) under \(I_{\mathrm{PICA}}\).
\end{itemize}

Each assignment is a finite host datum, audited by \(I_{\mathrm{PICA}}\), and consistent with the source-of-truth and visibility policies of \(\mathsf H_{\mathrm{PICA}}\). The realization map is admissible in the sense of Subsection 13.1.1, so the scoped model-realization statement receives verdict \(\texttt{model\_realization}\). \(\square\)

\subsubsection{PICA Nonclaims}

The PICA nonclaim register \(\mathcal N_{\mathrm{PICA}}\) contains, in particular:

\begin{enumerate}
\def\labelenumi{\arabic{enumi}.}
\tightlist
\item
  PICA is not the universal six-bird algebra: \(\mathsf{PICAReal}_{\mathrm{fin}}\) realizes a declared finite stochastic fragment, not the whole calculus.
\item
  The recovered \(6\times 6\) multiplicity count is a PICA-profile model fact under the cell profile \(\lambda_{\mathrm{PICA}}\), not a universal \(6\times 6\) law.
\item
  Raw PICA status is not formal status without recovery and classifier check; the formal classifier verdict is the authoritative cell status.
\item
  Directed-cell status \(\texttt{action}\) does not imply pair-realness; pair-realness requires the source/visibility/audit/threshold conjunction of Subsubsection 13.2.7.
\item
  A fallback record does not by itself support a real pair source-of-truth; it supports at most \(\texttt{real\_provisional}\).
\item
  Ablation evidence in \(\mathsf{Abl}\) is model-realization evidence, not theorem-grade evidence, in the absence of an admissible bridge to a theorem-grade claim.
\item
  A \(P_3\) route mismatch is not a \(\mathrm{P6}_{\mathrm{drive}}\) assertion in the sense of Section 8: the two are connected only through an admissible drive bridge.
\item
  A dirty record cannot support a real claim under any host configuration of \(\mathsf H_{\mathrm{PICA}}\).
\item
  PICA realizes a finite stochastic fragment, and not the entirety of \(\mathbf{BirdInt}^{\mathrm{aud}}_{\mathrm{fin}}\); this restates the fragment discipline of Subsection 13.1.2.
\item
  Implementation correctness of any concrete PICA implementation is future work and is not certified by this realization theorem; the Lean validation record of Section 15 is a separate, manifest-scoped artifact.
\end{enumerate}

The register \(\mathcal N_{\mathrm{PICA}}\) is finite, recorded in \(\mathcal P\), and is part of the admissibility predicate \(\mathsf{AdmPICAReal}(\mathcal P)\) in the sense of Subsubsection 13.2.2.

\subsection{Cantor Strict-Bridge Realization}

This subsection records the Cantor strict-bridge realization sheet and Model Theorem 34: under the audited Cantor shell, the strict-bridge passage from a base cocycle-pressure package \(T_0\) to a completion-packaged target package \(T_1\) is admissibly classified \(\texttt{strict}\). The realization is host-bound to the audited shell and fragment-scoped to the strict-bridge passage; it does not generalize beyond the shell, and its pressure-gap content is recorded as a consequence defect rather than the strictness witness itself.

\subsubsection{Name and Host}

The Cantor realization name is \(\mathsf{CantorShell}_{\mathrm{aud}}\). Its host is the audited Cantor shell context \(\mathsf H_{\mathrm{CantorShell}}\): a scoped audited application host built over a fixed audited shell-stable comparison domain \(\mathcal S_{\mathrm{aud}}\), recording the base and target packages, the bridge data, the pressure-gap defect record, and the audit/nonclaim register. The host is not the base finite theorem host of Sections 3 and 4, but a scoped audited application host in the sense of Subsection 11.1; the host instrument is \(I_{\mathrm{Cantor}}\), an admissible instrument under which the realization claim of Subsubsection 13.3.8 is recorded.

The realization object is the finite tuple

\[
\mathsf{CantorShell}_{\mathrm{aud}}
\;=\;
(\,\mathcal S_{\mathrm{aud}},\;T_0,\;T_1,\;\pi_0,\;\pi_1,\;\mathcal F,\;\Delta,\;A,\;\mathcal N\,),
\]

where \(T_0\) is the base cocycle-pressure package, \(T_1\) is the completion-packaged target package, \(\mathcal F\) is the persistent packaged-strata family, \(\Delta\) is the pressure-gap consequence record, \(A\) is the finite audit/theoremlet support record, and \(\mathcal N\) is the Cantor nonclaim register of Subsubsection 13.3.9.

\subsubsection{\texorpdfstring{Base Package \(T_0\)}{Base Package T\_0}}

The base package is

\[
\mathcal T_0^{\mathrm{Cantor}}\;=\;(\,\mathcal S_{\mathrm{aud}},\;\pi_0,\;\Sigma_{\pi_0}^{\mathrm{cocycle}},\;E_0^{\mathrm{cocycle}},\;A_{T_0}\,),
\]

a finite typed package in the schema of Subsection 3.3, with:

\begin{itemize}
\tightlist
\item
  \(\mathcal S_{\mathrm{aud}}\), the audited shell-stable comparison domain;
\item
  \(\pi_0:\mathcal S_{\mathrm{aud}}\to\mathcal O_0\), the old cocycle-level object map onto a finite collection \(\mathcal O_0\) of cocycle-level objects;
\item
  \(\Sigma_{\pi_0}^{\mathrm{cocycle}}\), the visible-content interface of the base package, recording the cocycle-level claims that \(T_0\) exposes;
\item
  \(E_0^{\mathrm{cocycle}}=(E_{\pi_0},\operatorname{Cl}_{P_{T_0}})\), the base closure field combining the object-map closure \(E_{\pi_0}\) with the closure of the base pressure object;
\item
  \(A_{T_0}\), the finite audit record on \(\mathcal T_0^{\mathrm{Cantor}}\) under \(I_{\mathrm{Cantor}}\).
\end{itemize}

The base pressure object is

\[
P_{T_0}(s)\;=\;\lim_{n\to\infty}\frac{1}{n}\,\log\!\left(\sup_{x\in\mathcal S_{\mathrm{aud}}}a_n(s;x)\right),
\]

a finite-shell scalar functional of a parameter \(s\), treated as a recorded object of \(T_0\) rather than as an external object.

\subsubsection{\texorpdfstring{Target Package \(T_1\)}{Target Package T\_1}}

The target package is

\[
\mathcal T_1^{\mathrm{Cantor}}\;=\;(\,\mathcal S_{\mathrm{aud}},\;\pi_1,\;\Sigma_{\pi_1}^{\mathrm{pkg}},\;E_1^{\mathrm{pkg}},\;A_{T_1}\,),
\]

with \(\pi_1:\mathcal S_{\mathrm{aud}}\to\mathcal O_1\) the new packaged-object map onto a finite collection \(\mathcal O_1\) of completion-packaged objects, \(\Sigma_{\pi_1}^{\mathrm{pkg}}\) the packaged-object visible-content interface, \(A_{T_1}\) the audit record on \(\mathcal T_1^{\mathrm{Cantor}}\), and packaged completion field

\[
E_1^{\mathrm{pkg}}\;=\;(\,E_{\tau,\ell},\;\mathrm{Fix}(E_{\tau,\ell}),\;\mathcal F,\;\mathrm{Sat},\;\mathrm{Force}_{4\leftarrow 5}\,),
\]

with \(E_{\tau,\ell}(\mu)=U_{\ell}\bigl(Q_{\ell}(\mu K^{\tau})\bigr)\), \(\mathcal F\) the persistent packaged-strata family, \(\mathrm{Sat}\) the saturation operator, and \(\mathrm{Force}_{4\leftarrow 5}\) the forcing operator that records the \(P_4\leftarrow P_5\) completion-packaging relation.

\subsubsection{Cantor Strict Bridge}

The Cantor strict bridge is

\[
B_{0\to 1}^{\mathrm{Cantor}}\;=\;(\,\mathcal T_0^{\mathrm{Cantor}},\;\mathcal T_1^{\mathrm{Cantor}},\;\pi_0,\;\pi_1,\;L_{\mathrm{Cantor}},\;V_{\mathrm{Cantor}},\;\Theta_{\mathrm{Cantor}},\;A_{\mathrm{Cantor}},\;\delta_{\mathrm{Cantor}},\;\mathcal N_{\mathrm{Cantor}}\,),
\]

a directed strict bridge in the promotion-bridge schema of Subsection 10.1, with lift data

\[
L_{\mathrm{Cantor}}\;=\;(\,E_{\tau,\ell},\;Q_{\ell},\;U_{\ell},\;\mathrm{Sat},\;\mathrm{Force}_{4\leftarrow 5}\,),
\]

and bridge-defect bundle

\[
\delta_{\mathrm{Cantor}}\;=\;(\,\delta_{\mathrm{fact}},\;\delta_{\mathrm{sat}},\;\delta_{4\leftarrow 5},\;\delta_{\mathrm{macro}},\;\delta_{\mathrm{press}}\,),
\]

aggregating the factorization defect of Subsubsection 13.3.5, the saturation defect, the \(4\leftarrow 5\) forcing defect, the macro defect, and the pressure-gap consequence defect of Subsubsection 13.3.6.

\subsubsection{Strictness Witness}

The strictness witness is the nonfactorization

\[
\pi_1\not\factor\pi_0\quad\text{on}\quad\mathcal S_{\mathrm{aud}},
\]

equivalently,

\[
\exists\,x,x'\in\mathcal S_{\mathrm{aud}}:\;\pi_0(x)=\pi_0(x')\;\land\;\pi_1(x)\neq\pi_1(x').
\]

The factorization defect record is

\[
\Delta_{\mathrm{fact}}^{\mathrm{Cantor}}(\pi_0,\pi_1)\;=\;\bigl\{\,(x,x')\in\mathcal S_{\mathrm{aud}}\times\mathcal S_{\mathrm{aud}}\;:\;\pi_0(x)=\pi_0(x')\;\land\;\pi_1(x)\neq\pi_1(x')\,\bigr\},
\]

a finite witness set on the audited shell. The strictness gate of Subsection 10.2 returns

\[
G_{\mathrm{strict}}\;=\;\texttt{strict\_pass}\quad\Longleftrightarrow\quad\Delta_{\mathrm{fact}}^{\mathrm{Cantor}}(\pi_0,\pi_1)\;\neq\;\varnothing,
\]

so the strictness witness is the nonempty factorization defect on the audited shell. The witness is host-bound: it is a property of \(\pi_0\), \(\pi_1\), and \(\mathcal S_{\mathrm{aud}}\), recorded under \(I_{\mathrm{Cantor}}\), and not a property of \(\pi_0\) and \(\pi_1\) outside the shell.

\subsubsection{Pressure-Gap Consequence Defect}

The Cantor strict bridge carries a pressure-gap consequence defect

\[
\Delta_{T_0\to T_1}(s;x)\;=\;P_{T_0}(s)\;-\;\sum_{\Sigma\in\mathcal F(x)}w_{\Sigma}(x)\,P_{\Sigma}(s),
\]

with \(\mathcal F(x)\) the persistent packaged-strata family at \(x\), \(w_{\Sigma}(x)\) the per-stratum weight, and \(P_{\Sigma}(s)\) the per-stratum pressure object. The defect record is

\[
\mathsf{Def}^{\mathrm{press}}_{0\to 1}(s;x)\;=\;(\,\delta_{\mathrm{press}},\;\Omega_{\mathrm{press}},\;\Theta_{\mathrm{press}},\;\mathrm{polarity}_{\mathrm{press}},\;W_{\mathrm{press}},\;A_{\mathrm{press}},\;V_{\mathrm{press}},\;\sigma_{\mathrm{press}}\,),
\]

with \(\delta_{\mathrm{press}}=\Delta_{T_0\to T_1}(s;x)\) and the remaining fields drawn from the defect-record schema of Subsection 5.2.

The pressure-gap defect is a \emph{consequence} defect, not the strictness witness itself: it is activated when

\[
\Delta_{T_0\to T_1}(s;x)\;\not\equiv\;0
\]

holds on the audited shell or on a recorded audited witness set, and it records a quantitative pressure deficit associated with the strict bridge, without itself supplying the nonfactorization data. The fragment discipline of Subsection 13.1 forbids reading the pressure-gap defect as the strictness witness, and Subsubsection 13.3.9 records the corresponding nonclaim.

\subsubsection{Gate Results}

For the Cantor strict bridge, the promotion gates of Subsection 10.2 return the following verdicts under \(I_{\mathrm{Cantor}}\):

\begin{longtable}[]{@{}ll@{}}
\toprule
gate & result\tabularnewline
\midrule
\endhead
\(G_{\mathrm{suff}}\) & \(\texttt{pass}\)\tabularnewline
\(G_{\mathrm{desc}}\) & \(\texttt{not\_required}\)\tabularnewline
\(G_{\mathrm{stab}}\) & \(\texttt{pass}\)\tabularnewline
\(G_{\mathrm{ctrl}}\) & \(\texttt{pass}\)\tabularnewline
\(G_{\mathrm{nosmuggle}}\) & \(\texttt{pass}\)\tabularnewline
\(G_{\mathrm{vis}}\) & \(\texttt{pass}\)\tabularnewline
\(G_{\mathrm{audit}}\) & \(\texttt{pass}\)\tabularnewline
\(G_{\mathrm{strict}}\) & \(\texttt{strict\_pass}\)\tabularnewline
\(G_{\mathrm{locglob}}\) & \(\texttt{not\_required}\)\tabularnewline
\bottomrule
\end{longtable}

The descent gate \(G_{\mathrm{desc}}\) is \(\texttt{not\_required}\) because the bridge claims object-map strictness on the audited shell, not closed macro dynamics; the local-global gate \(G_{\mathrm{locglob}}\) is \(\texttt{not\_required}\) because the realization is bound to the shell. The remaining gates pass, and the strictness gate \(G_{\mathrm{strict}}\) returns \(\texttt{strict\_pass}\) from the nonempty factorization defect of Subsubsection 13.3.5.

\subsubsection{Model Theorem 34: Cantor Strict-Bridge Realization}

\textbf{Model Theorem 34 (Cantor Strict-Bridge Realization; model-grade).} Suppose the Cantor realization object is admissible:

\[
\mathsf{AdmCantorShell}\bigl(\mathsf{CantorShell}_{\mathrm{aud}}\bigr)
\]

holds, the strictness witness of Subsubsection 13.3.5 records a nonempty factorization defect on \(\mathcal S_{\mathrm{aud}}\), and the gate record of Subsubsection 13.3.7 holds. Then

\[
\mathsf{CantorShell}_{\mathrm{aud}}\;\models\;\mathbf{StrictBridge}_{0\to 1}^{\mathrm{audited\ shell}},
\]

and the instrument-indexed form is

\[
\Gamma\,;\;\mathsf H_{\mathrm{CantorShell}}\,;\;I_{\mathrm{Cantor}}\;\vdash\;\mathsf{CantorShell}_{\mathrm{aud}}\;\models\;\mathbf{StrictBridge}_{0\to 1}^{\mathrm{audited\ shell}}\;:\;\texttt{model\_realization}.
\]

Equivalently, the bridge promotion under \(\mathcal T_0^{\mathrm{Cantor}}\) classifies as \(\texttt{strict}\):

\[
\Gamma\,;\;\mathcal T_0^{\mathrm{Cantor}}\,;\;I_{\mathrm{Cantor}}\;\vdash\;\mathrm{Promote}\bigl(B_{0\to 1}^{\mathrm{Cantor}}\bigr)\;:\;\texttt{strict}.
\]

\emph{Proof.} By Subsubsection 13.3.4, \(B_{0\to 1}^{\mathrm{Cantor}}\) is a finite typed strict-bridge record in the promotion-bridge schema of Subsection 10.1. Subsubsection 13.3.5 records that the factorization defect \(\Delta_{\mathrm{fact}}^{\mathrm{Cantor}}\) is nonempty on \(\mathcal S_{\mathrm{aud}}\), so \(G_{\mathrm{strict}}=\texttt{strict\_pass}\). Subsubsection 13.3.7 records that the remaining gates of Subsection 10.2 either pass or are \(\texttt{not\_required}\) for the audited shell. The promotion classifier of Subsection 5.6.4, applied with the gate record of Subsubsection 13.3.7, returns \(\mathrm{Promote}(B_{0\to 1}^{\mathrm{Cantor}}):\texttt{strict}\). By the realization-map convention of Subsection 13.1.1, the model-realization claim therefore receives the model-grade verdict \(\texttt{model\_realization}\) on the declared structure \(\mathbf{StrictBridge}_{0\to 1}^{\mathrm{audited\ shell}}\). \(\square\)

\subsubsection{Cantor Nonclaims}

The Cantor nonclaim register \(\mathcal N_{\mathrm{Cantor}}\) contains:

\begin{enumerate}
\def\labelenumi{\arabic{enumi}.}
\tightlist
\item
  The realization is scoped to the audited Cantor shell \(\mathcal S_{\mathrm{aud}}\); it is not a shell-general theorem and does not assert a strict-bridge passage outside the shell.
\item
  There is no shell-general strict-bridge theorem; extending the realization beyond \(\mathcal S_{\mathrm{aud}}\) requires a separately admissible bridge.
\item
  There is no broader external-class theorem; the realization does not assert structural facts about external classes of theories unless an admissible inter-theory bridge is supplied.
\item
  Strictness does not imply macro closure: the strictness witness of Subsubsection 13.3.5 is object-map nonfactorization, not closure of macro dynamics.
\item
  Strictness does not imply drive: the realization does not supply a \(\mathrm{P6}_{\mathrm{drive}}\) assertion in the sense of Section 8 without a separately admissible drive bridge.
\item
  The pressure-gap consequence defect of Subsubsection 13.3.6 is not a Kullback--Leibler closure deficit unless an admissible bridge to the KL closure-deficit framework is recorded.
\item
  The pressure-gap defect is not the strictness witness; it is a consequence of the witness, recorded as a separate defect record.
\item
  The pressure-gap functional \(\Delta_{T_0\to T_1}\) is a conditional pressure disintegration on the persistent strata family \(\mathcal F\), and is not a direct stratumwise root separation.
\item
  Cantor is a model realization/application, not a universal strict-extension theory.
\item
  The target completion \(T_1\) is not merely repeated application of one fixed idempotent completion; it is the packaged completion of Subsubsection 13.3.3.
\end{enumerate}

The register \(\mathcal N_{\mathrm{Cantor}}\) is finite, recorded in \(\mathsf{CantorShell}_{\mathrm{aud}}\), and is part of the admissibility predicate \(\mathsf{AdmCantorShell}\).

\subsection{Graph/Cohomology Realization}

This subsection records the graph/cohomology realization sheet \(\mathsf{CohGraph}_{\mathrm{fin}}\) and connects the cohomology theorems of Section 8 to the model-realization convention of Subsection 13.1. The realization is host-bound to the finite graph/cohomology host and fragment-scoped to the drive-support fragment \(P_2/P_6^{H^1}\text{-drive}/P_3\text{-separation}\).

\subsubsection{Host}

The graph/cohomology realization name is \(\mathsf{CohGraph}_{\mathrm{fin}}\). Its host is the finite graph/cohomology context \(\mathsf H_{\mathrm{graph}}\), with realization object

\[
\mathsf{CohGraph}_{\mathrm{fin}}\;=\;(\,G,\;C^{0}(G;k),\;C^{1}(G;k),\;d,\;H^{1}(G;k),\;\mathsf{Cycles}(G),\;a,\;\mathsf{Gate},\;A,\;V,\;\mathcal N\,),
\]

with components:

\begin{itemize}
\tightlist
\item
  \(G=(V,E)\), a finite support graph (finite vertex set \(V\) and finite edge set \(E\));
\item
  \(C^{0}(G;k)\) and \(C^{1}(G;k)\), the finite \(0\)- and \(1\)-cochain spaces over a fixed coefficient field \(k\);
\item
  \(d:C^{0}(G;k)\to C^{1}(G;k)\), the coboundary operator;
\item
  \(H^{1}(G;k)\), the first cohomology group of \(G\) with coefficients in \(k\);
\item
  \(\mathsf{Cycles}(G)\), the finite cycle space of \(G\);
\item
  \(a\in C^{1}(G;k)\), the affinity/drive cochain assigning a finite scalar to each oriented edge of \(G\);
\item
  \(\mathsf{Gate}\), a finite typed record of admissible edge or support deletions/restrictions \(G\rightsquigarrow G'\);
\item
  \(A\) and \(V\), the host audit and visibility records;
\item
  \(\mathcal N\), the cohomology nonclaim register.
\end{itemize}

The host instrument is \(I_{\mathrm{graph}}\), an admissible instrument under which all cohomology and gating data are audited. The realization object is admissible, written \(\mathsf{AdmCohGraph}(\mathsf{CohGraph}_{\mathrm{fin}})\), when each component is finite and \(A\) records the verifications of the cochain spaces, the coboundary, the cycle space, the cohomology, and the affinity cochain.

\subsubsection{Realized Fragment}

The graph/cohomology realization realizes the \(\mathbf{BirdInt}^{\mathrm{aud}}_{\mathrm{fin}}\) fragment

\[
\mathbf{BirdInt}^{\mathrm{aud}}_{\mathrm{fin}}{}^{\,P_{2}/P_{6}^{H^{1}}\text{-drive}/P_{3}\text{-separation}},
\]

with realization map

\begin{longtable}[]{@{}ll@{}}
\toprule
\begin{minipage}[b]{0.26\columnwidth}\raggedright
primitive\strut
\end{minipage} & \begin{minipage}[b]{0.68\columnwidth}\raggedright
cohomology realization\strut
\end{minipage}\tabularnewline
\midrule
\endhead
\begin{minipage}[t]{0.26\columnwidth}\raggedright
\(P_{2}\)\strut
\end{minipage} & \begin{minipage}[t]{0.68\columnwidth}\raggedright
support gating / edge deletion / cycle-rank capacity\strut
\end{minipage}\tabularnewline
\begin{minipage}[t]{0.26\columnwidth}\raggedright
\(\mathrm{P6}_{\mathrm{drive}}^{H^1}\)\strut
\end{minipage} & \begin{minipage}[t]{0.68\columnwidth}\raggedright
cohomological drive certificate \([a]\neq 0\)\strut
\end{minipage}\tabularnewline
\begin{minipage}[t]{0.26\columnwidth}\raggedright
\(P_{6}\) audit face\strut
\end{minipage} & \begin{minipage}[t]{0.68\columnwidth}\raggedright
audit of cochains, cycles, source-of-truth, and visibility\strut
\end{minipage}\tabularnewline
\begin{minipage}[t]{0.26\columnwidth}\raggedright
\(P_{3}\)\strut
\end{minipage} & \begin{minipage}[t]{0.68\columnwidth}\raggedright
optional route/holonomy witness, kept separate from drive\strut
\end{minipage}\tabularnewline
\begin{minipage}[t]{0.26\columnwidth}\raggedright
\(P_{4}\)\strut
\end{minipage} & \begin{minipage}[t]{0.68\columnwidth}\raggedright
optional local/global cover or gluing of cochains (out of core scope)\strut
\end{minipage}\tabularnewline
\begin{minipage}[t]{0.26\columnwidth}\raggedright
\(P_{5}\)\strut
\end{minipage} & \begin{minipage}[t]{0.68\columnwidth}\raggedright
optional quotient/support package of graph state (out of core scope)\strut
\end{minipage}\tabularnewline
\bottomrule
\end{longtable}

The core sheet of the realization focuses on \(P_{2}\) and \(\mathrm{P6}_{\mathrm{drive}}^{H^1}\); the \(P_{4}\) and \(P_{5}\) assignments above are optional extensions, recorded for completeness but not part of the core fragment, and the \(P_{3}\) separation is enforced as a nonclaim of Subsubsection 13.4.4.

\subsubsection{Drive and No-Drive Certificates}

The cohomological drive certificate is the finite typed record

\[
\kappa_{6}^{\mathrm{drive}}(a)\;:\;[a]\neq 0\;\in\;H^{1}(G;k),
\]

equivalent on the finite cycle space to the existence of a cycle integral witness:

\[
[a]\;\neq\;0\quad\Longleftrightarrow\quad\exists\,\gamma\in\mathsf{Cycles}(G):\;\oint_{\gamma}a\;\neq\;0.
\]

The drive certificate is a \(\mathrm{P6}_{\mathrm{drive}}^{H^1}\) record in the sense of Section 8, distinct from a generic \(P_{6}\) audit record: it records that the affinity cochain \(a\) has a nonzero cohomology class, equivalently a nonzero cycle integral on at least one cycle.

No-drive certificates are the dual finite records

\[
[a]\;=\;0,\qquad a\;=\;d\phi\;\;\text{for some}\;\phi\in C^{0}(G;k),\qquad H^{1}(G;k)\;=\;0,\qquad\beta_{1}(G)\;=\;0.
\]

The four no-drive certificates are connected by finite implications recorded in Section 8: \(a=d\phi\Rightarrow\oint_{\gamma}a=0\) for every cycle \(\gamma\); \(\beta_{1}(G)=0\Rightarrow H^{1}(G;k)=0\); and \(H^{1}(G;k)=0\Rightarrow[a]=0\) for every \(a\in C^{1}(G;k)\). The gating realization

\[
G\;=\;(V,E)\;\rightsquigarrow\;G'\;=\;(V',E')
\]

records a \(P_{2}\) support restriction with cycle-rank loss defect

\[
\delta_{2}^{\mathrm{cycle}}(G,G')\;=\;\beta_{1}(G)\;-\;\beta_{1}(G')
\]

and forest threshold \(\Theta_{2}^{\mathrm{forest}}:\beta_{1}(G')=0\). When the gate sends \(G\) to a forest \(G'\), the consequence \(H^{1}(G';k)=0\) records the absence of any cohomological drive certificate on \(G'\), in the sense that \(\mathrm{P6}_{\mathrm{drive}}^{H^1}\) is absent with record on \(G'\).

The cohomology theorems of Section 8 --- forest no-drive (\(\beta_{1}(G)=0\Rightarrow H^{1}(G;k)=0\)), exact-form null-drive (\(a=d\phi\Rightarrow\oint_{\gamma}a=0\)), nonzero-affinity equivalence (\([a]\neq 0\iff\exists\gamma:\oint_{\gamma}a\neq 0\)), and gating-affinity suppression (gating to a forest kills \(\mathrm{P6}_{\mathrm{drive}}^{H^1}\)) --- are recorded under the realization-map convention as the host-bound model-grade record of the drive-support fragment. Their theorem-grade status is preserved from Section 8; the model-realization convention adds the host-bound record under \(\mathsf H_{\mathrm{graph}}\) and \(I_{\mathrm{graph}}\), with claim-strength verdict

\[
\Gamma\,;\;\mathsf H_{\mathrm{graph}}\,;\;I_{\mathrm{graph}}\;\vdash\;\mathsf{CohGraph}_{\mathrm{fin}}\;\models\;\mathbf{BirdInt}^{\mathrm{aud}}_{\mathrm{fin}}{}^{\,P_{2}/P_{6}^{H^{1}}\text{-drive}/P_{3}\text{-separation}}\;:\;\texttt{model\_realization}.
\]

The realization claim is admissible whenever \(\mathsf{AdmCohGraph}\) holds and the cohomology theorems of Section 8 are imported as their theorem-grade \(\texttt{accepted}\) verdicts.

\subsubsection{Nonclaims}

The graph/cohomology nonclaim register \(\mathcal N_{\mathrm{coh}}\) contains:

\begin{enumerate}
\def\labelenumi{\arabic{enumi}.}
\tightlist
\item
  The drive certificate \(\mathrm{P6}_{\mathrm{drive}}^{H^1}\) with \([a]\neq 0\) in \(H^{1}(G;k)\) is one drive certificate family, not every possible drive notion.
\item
  A \(P_{3}\)-route/holonomy witness is not a \(\mathrm{P6}_{\mathrm{drive}}^{H^1}\) assertion in the sense of Section 8 in the absence of an admissible bridge \(B_{3\to 6}^{\mathrm{drive}}\).
\item
  Gating to a forest kills cycle-supported drive, not all forms of activity on \(G\).
\item
  Exact \(1\)-cochains \(a=d\phi\) have zero cycle drive: \(\oint_{\gamma}a=0\) for every cycle \(\gamma\).
\item
  The condition \(H^{1}(G;k)=0\) is a no-drive certificate for cohomological drive on \(G\), and not a no-drive certificate for any other notion of activity.
\item
  Nonzero route holonomy on \(G\) can coexist with a zero drive cochain \(a\); the two are distinct records.
\item
  The graph/cohomology host \(\mathsf H_{\mathrm{graph}}\) is finite and scoped; it is not a universal cohomological framework.
\item
  Cohomology does not replace the PICA, Cantor, or abstract-interpretation realizations of Subsections 13.2, 13.3, and Section 12.
\item
  No edge-support claim about \(G\) is accepted under \(I_{\mathrm{graph}}\) without the visibility and audit conditions of \(V\) and \(A\).
\item
  Positive holonomy-to-drive claims --- that a positive route holonomy implies a positive drive --- require an admissible drive bridge \(B_{3\to 6}^{\mathrm{drive}}\) in the sense of Section 8.
\end{enumerate}

The register \(\mathcal N_{\mathrm{coh}}\) is finite, recorded in \(\mathsf{CohGraph}_{\mathrm{fin}}\), and is part of \(\mathsf{AdmCohGraph}\).

\subsection{Combined Model-Realization Matrix}

This subsection collects the four realizations of Sections 12 and 13 into a single matrix, with one row per realization and one column per scoped attribute. The matrix records what each realization realizes, under what host, with what claim strength, and where the supporting record is located. It does not assert any cross-model fact beyond what the per-row realizations already record; in particular, it does not state a model-coverage theorem, in keeping with Subsection 13.6.

\subsubsection{Rows}

The rows of the combined matrix are the four model realizations of the first paper:

\begin{itemize}
\tightlist
\item
  \emph{PICA}: \(\mathsf{PICAReal}_{\mathrm{fin}}\) of Subsection 13.2;
\item
  \emph{Cantor}: \(\mathsf{CantorShell}_{\mathrm{aud}}\) of Subsection 13.3;
\item
  \emph{Abstract interpretation}: \(\mathsf{AIFam}_{\mathrm{fin}}\) of Section 12;
\item
  \emph{Graph/cohomology}: \(\mathsf{CohGraph}_{\mathrm{fin}}\) of Subsection 13.4.
\end{itemize}

\subsubsection{Columns}

The columns of the combined matrix are the per-realization attributes:

\begin{itemize}
\tightlist
\item
  \emph{Host}: the host context \(\mathsf H\) and instrument \(I\) under which the realization claim is recorded.
\item
  \emph{Realized fragment}: the declared structure \(\mathcal S\) realized by the model.
\item
  \emph{Main theorem}: the theorem-grade or model-grade theorem of the realization, with its number in this paper.
\item
  \emph{Blocked overclaims}: the principal nonclaims of the realization that the fragment discipline of Subsection 13.1.2 enforces.
\item
  \emph{Appendix support}: the appendix sheet that records the long-form realization data for the model.
\end{itemize}

The combined matrix is

\begin{longtable}[]{@{}llllll@{}}
\toprule
\begin{minipage}[b]{0.08\columnwidth}\raggedright
model\strut
\end{minipage} & \begin{minipage}[b]{0.12\columnwidth}\raggedright
host / instrument\strut
\end{minipage} & \begin{minipage}[b]{0.22\columnwidth}\raggedright
realized fragment\strut
\end{minipage} & \begin{minipage}[b]{0.12\columnwidth}\raggedright
main theorem\strut
\end{minipage} & \begin{minipage}[b]{0.22\columnwidth}\raggedright
blocked overclaims\strut
\end{minipage} & \begin{minipage}[b]{0.08\columnwidth}\raggedright
appendix support\strut
\end{minipage}\tabularnewline
\midrule
\endhead
\begin{minipage}[t]{0.08\columnwidth}\raggedright
\(\mathsf{PICAReal}_{\mathrm{fin}}\)\strut
\end{minipage} & \begin{minipage}[t]{0.12\columnwidth}\raggedright
\(\mathsf H_{\mathrm{PICA}}\) / \(I_{\mathrm{PICA}}\)\strut
\end{minipage} & \begin{minipage}[t]{0.22\columnwidth}\raggedright
\(\mathbf{BirdInt}^{\mathrm{aud}}_{\mathrm{fin\text{-}stoch}}{}^{\,\mathrm{cell}/\mathrm{pair}/\mathrm{prov}}\)\strut
\end{minipage} & \begin{minipage}[t]{0.12\columnwidth}\raggedright
Model Theorem 33 (model-grade)\strut
\end{minipage} & \begin{minipage}[t]{0.22\columnwidth}\raggedright
no PICA universality; no \(6\times 6\) universal law; raw status is not formal status\strut
\end{minipage} & \begin{minipage}[t]{0.08\columnwidth}\raggedright
Appendix D.1\strut
\end{minipage}\tabularnewline
\begin{minipage}[t]{0.08\columnwidth}\raggedright
\(\mathsf{CantorShell}_{\mathrm{aud}}\)\strut
\end{minipage} & \begin{minipage}[t]{0.12\columnwidth}\raggedright
\(\mathsf H_{\mathrm{CantorShell}}\) / \(I_{\mathrm{Cantor}}\)\strut
\end{minipage} & \begin{minipage}[t]{0.22\columnwidth}\raggedright
\(\mathbf{StrictBridge}_{0\to 1}^{\mathrm{audited\ shell}}\)\strut
\end{minipage} & \begin{minipage}[t]{0.12\columnwidth}\raggedright
Model Theorem 34 (model-grade)\strut
\end{minipage} & \begin{minipage}[t]{0.22\columnwidth}\raggedright
no shell-general theorem; pressure-gap is consequence, not strictness witness; strictness does not imply drive\strut
\end{minipage} & \begin{minipage}[t]{0.08\columnwidth}\raggedright
Appendix D.2\strut
\end{minipage}\tabularnewline
\begin{minipage}[t]{0.08\columnwidth}\raggedright
\(\mathsf{AIFam}_{\mathrm{fin}}\)\strut
\end{minipage} & \begin{minipage}[t]{0.12\columnwidth}\raggedright
\(\mathsf H_{\mathrm{AI}}\) / \(I_{\mathrm{AI}}\)\strut
\end{minipage} & \begin{minipage}[t]{0.22\columnwidth}\raggedright
\(\mathbf{BirdInt}^{\mathrm{aud}}_{\mathrm{fin}}{}^{\,P_{1}/P_{2}/P_{4}/P_{5}/P_{6}}\)\strut
\end{minipage} & \begin{minipage}[t]{0.12\columnwidth}\raggedright
Model Theorem 29 (model-grade)\strut
\end{minipage} & \begin{minipage}[t]{0.22\columnwidth}\raggedright
no AI totality; no soundness-as-completeness; no closure-implies-soundness\strut
\end{minipage} & \begin{minipage}[t]{0.08\columnwidth}\raggedright
Appendix D.3\strut
\end{minipage}\tabularnewline
\begin{minipage}[t]{0.08\columnwidth}\raggedright
\(\mathsf{CohGraph}_{\mathrm{fin}}\)\strut
\end{minipage} & \begin{minipage}[t]{0.12\columnwidth}\raggedright
\(\mathsf H_{\mathrm{graph}}\) / \(I_{\mathrm{graph}}\)\strut
\end{minipage} & \begin{minipage}[t]{0.22\columnwidth}\raggedright
\(\mathbf{BirdInt}^{\mathrm{aud}}_{\mathrm{fin}}{}^{\,P_{2}/P_{6}^{H^{1}}\text{-drive}/P_{3}\text{-separation}}\)\strut
\end{minipage} & \begin{minipage}[t]{0.12\columnwidth}\raggedright
Section 8 theorem group plus CohGraph model-realization statement (model-grade)\strut
\end{minipage} & \begin{minipage}[t]{0.22\columnwidth}\raggedright
one drive certificate family, not all drive; \(P_{3}\) holonomy separated from drive\strut
\end{minipage} & \begin{minipage}[t]{0.08\columnwidth}\raggedright
Appendix D.4\strut
\end{minipage}\tabularnewline
\bottomrule
\end{longtable}

The four rows record the four scoped realizations recognized by the calculus in this paper. Each row is read horizontally as a host-bound, fragment-scoped record; the matrix is not read vertically as a coverage assertion, in keeping with the cross-model role-coverage observation of Subsection 13.6.

The instrument-indexed claim form for each row is the \(\models\,:\,\texttt{model\_realization}\) statement recorded in the corresponding subsection: Subsubsection 13.2.9 for PICA, Subsubsection 13.3.8 for Cantor, Subsubsection 12.5.1 for abstract interpretation, and Subsubsection 13.4.3 for graph/cohomology. The matrix collects these four claims in one place; it does not introduce a new claim. The blocked-overclaims column is a summary view of the per-realization nonclaim registers \(\mathcal N_{\mathrm{PICA}}\), \(\mathcal N_{\mathrm{Cantor}}\), \(\mathcal N_{\mathrm{AI}}\), and \(\mathcal N_{\mathrm{coh}}\), and is not a separate, free-standing nonclaim register.

\subsection{Cross-Model Role Coverage}

This subsection records the cross-model role-coverage observation: which primitives \(P_{1},\ldots,P_{6}\) each realization realizes within its declared fragment, and which it does not. The observation is not a coverage theorem; the calculus does not assert that the union of the four realizations covers the six-role calculus, and Subsubsection 13.6.2 records the corresponding coverage nonclaim.

\subsubsection{Coverage Table}

The role-coverage table records, per primitive, what each realization realizes within its declared fragment:

\begin{longtable}[]{@{}lllll@{}}
\toprule
\begin{minipage}[b]{0.10\columnwidth}\raggedright
role\strut
\end{minipage} & \begin{minipage}[b]{0.17\columnwidth}\raggedright
\(\mathsf{PICAReal}_{\mathrm{fin}}\)\strut
\end{minipage} & \begin{minipage}[b]{0.19\columnwidth}\raggedright
\(\mathsf{CantorShell}_{\mathrm{aud}}\)\strut
\end{minipage} & \begin{minipage}[b]{0.23\columnwidth}\raggedright
\(\mathsf{AIFam}_{\mathrm{fin}}\)\strut
\end{minipage} & \begin{minipage}[b]{0.17\columnwidth}\raggedright
\(\mathsf{CohGraph}_{\mathrm{fin}}\)\strut
\end{minipage}\tabularnewline
\midrule
\endhead
\begin{minipage}[t]{0.10\columnwidth}\raggedright
\(P_{1}\)\strut
\end{minipage} & \begin{minipage}[t]{0.17\columnwidth}\raggedright
descent / closure update\strut
\end{minipage} & \begin{minipage}[t]{0.19\columnwidth}\raggedright
macro obstruction; not primary\strut
\end{minipage} & \begin{minipage}[t]{0.23\columnwidth}\raggedright
sound descent \(\alpha F\le_{A}F^{\sharp}\alpha\)\strut
\end{minipage} & \begin{minipage}[t]{0.17\columnwidth}\raggedright
not primary\strut
\end{minipage}\tabularnewline
\begin{minipage}[t]{0.10\columnwidth}\raggedright
\(P_{2}\)\strut
\end{minipage} & \begin{minipage}[t]{0.17\columnwidth}\raggedright
gating / support / admissibility\strut
\end{minipage} & \begin{minipage}[t]{0.19\columnwidth}\raggedright
admissibility obstruction\strut
\end{minipage} & \begin{minipage}[t]{0.23\columnwidth}\raggedright
representability \(\rho(p)=p\)\strut
\end{minipage} & \begin{minipage}[t]{0.17\columnwidth}\raggedright
support gating\strut
\end{minipage}\tabularnewline
\begin{minipage}[t]{0.10\columnwidth}\raggedright
\(P_{3}\)\strut
\end{minipage} & \begin{minipage}[t]{0.17\columnwidth}\raggedright
route mismatch\strut
\end{minipage} & \begin{minipage}[t]{0.19\columnwidth}\raggedright
not central\strut
\end{minipage} & \begin{minipage}[t]{0.23\columnwidth}\raggedright
optional noncommuting closures (deferred extension)\strut
\end{minipage} & \begin{minipage}[t]{0.17\columnwidth}\raggedright
separated from drive\strut
\end{minipage}\tabularnewline
\begin{minipage}[t]{0.10\columnwidth}\raggedright
\(P_{4}\)\strut
\end{minipage} & \begin{minipage}[t]{0.17\columnwidth}\raggedright
refinement / staging\strut
\end{minipage} & \begin{minipage}[t]{0.19\columnwidth}\raggedright
\(P_{4}\!\leftarrow\!P_{5}\) completion-packaging\strut
\end{minipage} & \begin{minipage}[t]{0.23\columnwidth}\raggedright
refinement \(A_{k}\rightsquigarrow A_{\ell}\)\strut
\end{minipage} & \begin{minipage}[t]{0.17\columnwidth}\raggedright
optional local/global cover\strut
\end{minipage}\tabularnewline
\begin{minipage}[t]{0.10\columnwidth}\raggedright
\(P_{5}\)\strut
\end{minipage} & \begin{minipage}[t]{0.17\columnwidth}\raggedright
package / completion\strut
\end{minipage} & \begin{minipage}[t]{0.19\columnwidth}\raggedright
completion-packaged objects\strut
\end{minipage} & \begin{minipage}[t]{0.23\columnwidth}\raggedright
closure / packaging \(\rho=\gamma\alpha\)\strut
\end{minipage} & \begin{minipage}[t]{0.17\columnwidth}\raggedright
optional graph package\strut
\end{minipage}\tabularnewline
\begin{minipage}[t]{0.10\columnwidth}\raggedright
\(P_{6}\)\strut
\end{minipage} & \begin{minipage}[t]{0.17\columnwidth}\raggedright
audit / provenance\strut
\end{minipage} & \begin{minipage}[t]{0.19\columnwidth}\raggedright
audit of shell bridge\strut
\end{minipage} & \begin{minipage}[t]{0.23\columnwidth}\raggedright
audit of maps and soundness\strut
\end{minipage} & \begin{minipage}[t]{0.17\columnwidth}\raggedright
drive certificate \([a]\neq 0\) and audit\strut
\end{minipage}\tabularnewline
\bottomrule
\end{longtable}

A cell of the table records what the column's realization carries for the row's primitive \emph{within its declared fragment}. The phrase ``not primary'' records that the primitive is not part of the realization's core declared fragment; the phrase ``deferred extension'' records that the primitive is realized only by an optional extension that is not part of the core sheet.

The table reads horizontally and vertically only as a record of per-cell facts; it does not assert any cross-row equivalence of realizations of the same primitive, and it does not assert any vertical aggregation across columns.

\subsubsection{Coverage Nonclaim}

The cross-model role-coverage observation does not, by itself, support a model-coverage theorem for \(\mathbf{BirdInt}^{\mathrm{aud}}_{\mathrm{fin}}\). The framework records the following coverage nonclaim:

\begin{itemize}
\item
  \emph{No model-coverage theorem.} For every primitive \(P_{i}\) realized by some realization in the table, the calculus does not assert that the four realizations together realize the whole calculus. Each cell of the table is a per-row, per-fragment record; the union of the cells is not lifted to a coverage statement.
\item
  \emph{No cross-model identity.} Two cells in the same row carrying entries for the same primitive \(P_{i}\) are not identified: the PICA realization of \(P_{6}\) as audit/provenance and the cohomology realization of \(P_{6}\) as a drive certificate are distinct realizations of distinct facets of the audit role, and the matrix does not assert their identity. Cross-model identification of any per-primitive realization requires an admissible instrument-transfer bridge of Subsection 11.9.
\item
  \emph{Realization is not theorem grade.} The coverage table is a model-grade summary view. The \(\texttt{accepted}\) verdicts on theorem-grade claims of the calculus are recorded in Sections 6, 7, 8, 10, and 11; they are not displaced or summarized by the coverage table.
\end{itemize}

The coverage nonclaim is an instance of the fragment discipline of Subsection 13.1.2 and the model-realization nonclaim of Subsection 2.4: the calculus records the per-row \(\texttt{model\_realization}\) verdicts of Subsection 13.5, and does not promote them, jointly or severally, to a coverage statement about \(\mathbf{BirdInt}^{\mathrm{aud}}_{\mathrm{fin}}\).
 \section{Secondary High-Structure Semantics}

This section records a finite decorated-square semantics for \(\mathbf{BirdInt}^{\mathrm{aud}}_{\mathrm{fin}}\) and proves three high-structure recovery theorems: descent data appear as decorated squares (Theorem 30), noncommuting completions appear as pasting defects (Theorem 31), and strict extensions appear as no-filler conditions (Theorem 32). The semantics is \emph{secondary}: it is an organizational layer over the core calculus, recording how finite categorical decorations recover descent, completion, and strict-extension data, and is not a replacement for the core calculus or its status families. Subsection 14.1 fixes the finite double-category fragment, Subsections 14.2--14.4 prove the three theorems, and Subsection 14.5 records the high-structure nonclaims.

The discipline of the section is the secondary-organization discipline: the decorated-square fragment supplies a finite, audit-checkable representation of selected core constructions, and does not assert that the core calculus is reducible to the decorated-square fragment. Status families remain separate, in keeping with the no-total-algebra and status-family-separation results of Section 6.

\subsection{Finite Double-Category Fragment}

This subsection records the finite double-category data used by the section: package objects as objects, refinements/lenses as horizontal arrows, transformations as vertical arrows, and the decorated squares carrying descent, quotient, completion, promotion, visibility, audit, and factorization data. The fragment is finite throughout: the object class, horizontal class, vertical class, and square class are all finite typed records audited by the host instrument \(I_{\mathrm{DC}}\).

\subsubsection{Objects, Horizontal Arrows, Vertical Arrows}

The finite double-category fragment \(\mathsf H_{\mathrm{DC}}\) is the finite typed record

\[
\mathsf H_{\mathrm{DC}}\;=\;(\,\mathrm{Ob},\;\mathrm{Hor},\;\mathrm{Ver},\;\mathrm{Sq},\;A_{\mathrm{DC}},\;V_{\mathrm{DC}},\;\mathcal N_{\mathrm{DC}}\,),
\]

with:

\begin{itemize}
\tightlist
\item
  \(\mathrm{Ob}\), the finite class of \emph{package objects}: each object is a finite typed package \(\mathcal T\) in the schema of Subsection 3.3, recording carrier, object map, visible interface, closure, and audit record;
\item
  \(\mathrm{Hor}\), the finite class of \emph{horizontal arrows}: each horizontal arrow is a typed refinement, lens, or package morphism in the sense of Subsection 3.3, or a promotion bridge when the square type requires one;
\item
  \(\mathrm{Ver}\), the finite class of \emph{vertical arrows}: each vertical arrow is a typed transformation between packages or within a package, recording the actor-update content \(U\) and decorations such as audit, source, and visibility, drawn from the witness/update discipline of Subsection 3.7 and the directed-cell schema of Subsection 4.3;
\item
  \(\mathrm{Sq}\), the finite class of \emph{decorated squares}, defined in Subsubsection 14.1.2;
\item
  \(A_{\mathrm{DC}}\), the host audit record, recording verifications that \(\mathrm{Ob}\), \(\mathrm{Hor}\), and \(\mathrm{Ver}\) are finite and that their boundary and composition data are admissible;
\item
  \(V_{\mathrm{DC}}\), the host visibility record;
\item
  \(\mathcal N_{\mathrm{DC}}\), the host nonclaim register of Subsection 14.5.
\end{itemize}

Composition of horizontal arrows is along shared targets/sources, and composition of vertical arrows is along shared targets/sources, with both compositions admitted only when the resulting boundary records are finite and audited under \(I_{\mathrm{DC}}\).

\subsubsection{Decorated Squares}

A decorated square \(\Xi\in\mathrm{Sq}\) is a finite typed record

\[
\Xi\;=\;(\,X,\;F,\;q,\;F^{\sharp},\;\mathrm{src}_{\Xi},\;\mathrm{tgt}_{\Xi},\;\mathrm{type}_{\Xi},\;\mathsf{Decor}_{\Xi},\;A_{\Xi},\;V_{\Xi},\;\delta_{\Xi},\;\mathcal N_{\Xi}\,),
\]

with:

\begin{itemize}
\tightlist
\item
  \(X\), the source object of the square (an element of \(\mathrm{Ob}\));
\item
  \(F\) and \(q\), the horizontal and vertical arrow components of the square's boundary, drawn from \(\mathrm{Hor}\) and \(\mathrm{Ver}\);
\item
  \(F^{\sharp}\), the second horizontal arrow forming the bottom edge, drawn from \(\mathrm{Hor}\);
\item
  \(\mathrm{src}_{\Xi}\) and \(\mathrm{tgt}_{\Xi}\), the source and target boundary tuples;
\item
  \(\mathrm{type}_{\Xi}\in\{\,\texttt{descent},\;\texttt{quotient},\;\texttt{completion},\;\texttt{promotion},\;\texttt{visibility},\;\texttt{audit},\;\texttt{factorization}\,\}\), the square type;
\item
  \(\mathsf{Decor}_{\Xi}\), the decoration record carrying the type-specific data of the square: descent data for \(\texttt{descent}\), quotient data for \(\texttt{quotient}\), completion data for \(\texttt{completion}\), promotion data for \(\texttt{promotion}\), visibility-bridge data for \(\texttt{visibility}\), audit-bridge data for \(\texttt{audit}\), and factorization data for \(\texttt{factorization}\);
\item
  \(A_{\Xi}\) and \(V_{\Xi}\), the per-square audit and visibility records;
\item
  \(\delta_{\Xi}\), the per-square defect record drawn from the general defect-record schema of Subsection 5.2 and classified in the square-status family of Subsection 5.1;
\item
  \(\mathcal N_{\Xi}\), the per-square nonclaim register.
\end{itemize}

A decorated square is \emph{exact} under \(I_{\mathrm{DC}}\) when the per-type exactness condition recorded in \(\mathsf{Decor}_{\Xi}\) holds and \(\delta_{\Xi}=\varnothing\). The exactness condition is type-specific: for \(\texttt{descent}\), it is the descent commutation equality of Subsubsection 14.2.1; for \(\texttt{completion}\), it is the completion-commutation equality of Subsubsection 14.3.1; for \(\texttt{factorization}\), it is the factorization-filler condition of Subsubsection 14.4.1; for \(\texttt{visibility}\) and \(\texttt{audit}\), it is the bridge-admissibility condition of Subsections 11.3 and 11.7.

\subsubsection{Pasting}

The decorated-square fragment supports finite \emph{strict pasting}: given two decorated squares \(\Xi_{1}\) and \(\Xi_{2}\) that share a horizontal or vertical edge, their pasted square \(\Xi_{1}\cdot\Xi_{2}\) is recorded as a single decorated square with concatenated boundary, concatenated decoration, and aggregated defect

\[
\delta_{\Xi_{1}\cdot\Xi_{2}}\;=\;\delta_{\Xi_{1}}\;\cup\;\delta_{\Xi_{2}}\;\cup\;\delta_{\mathrm{paste}}(\Xi_{1},\Xi_{2}),
\]

where \(\delta_{\mathrm{paste}}(\Xi_{1},\Xi_{2})\) is the finite paste-defect record for the boundary mismatch of \(\Xi_{1}\) and \(\Xi_{2}\). Pasting is admissible when the shared edge data agree on \(\mathrm{Hor}\) or \(\mathrm{Ver}\) and the per-square admissibility conditions hold for both factors. Strict pasting is the only pasting form admitted in this section: pseudo-pasting and weak pasting are out of scope, and the paste-defect record is the locus where any pasting failure is recorded.

The fragment \(\mathsf H_{\mathrm{DC}}\) together with the strict-pasting law is the finite decorated double-category fragment used by Theorems 30, 31, and 32 of Subsections 14.2, 14.3, and 14.4.

\subsection{Theorem 30: Descent-Square Recovery}

This subsection proves Theorem 30: in the finite decorated double-category fragment of Subsection 14.1, descent data are recovered as decorated squares of type \(\texttt{descent}\), with exactness recorded as the commutation equality between the concrete map \(F\) and the quotient map \(q\) followed by the descended map \(F^{\sharp}\).

\subsubsection{Statement}

Let \(X\in\mathrm{Ob}\) be a finite package object, \(F\in\mathrm{Hor}\) a horizontal map \(F:X\to X\), \(q\in\mathrm{Ver}\) a quotient map \(q:X\to\mathrm{im}(q)\), and \(F^{\sharp}\in\mathrm{Hor}\) the candidate descended horizontal map \(F^{\sharp}:\mathrm{im}(q)\to\mathrm{im}(q)\). Form the decorated square

\[
\Xi_{\mathrm{desc}}\;:\;
\begin{array}{ccc}
X & \xrightarrow{F} & X\\
{\scriptstyle q}\downarrow & \Downarrow\,\Xi_{\mathrm{desc}} & \downarrow{\scriptstyle q}\\
\mathrm{im}(q) & \xrightarrow{F^{\sharp}} & \mathrm{im}(q)
\end{array}
\]

with \(\mathrm{type}_{\Xi}=\texttt{descent}\) and \(\mathsf{Decor}_{\Xi}\) recording the descent data \((F,q,F^{\sharp})\).

\textbf{Theorem 30 (DescentSquareRecovery; theorem-grade).} The decorated square \(\Xi_{\mathrm{desc}}\) is exact iff the descent commutation equality

\[
F^{\sharp}\circ q_{\mathrm{im}}\;=\;q_{\mathrm{im}}\circ F
\]

holds on \(X\), where \(q_{\mathrm{im}}:X\to\mathrm{im}(q)\) is the corestriction of \(q\) to its image. Equivalently, a \emph{filler} exists for \(\Xi_{\mathrm{desc}}\) iff the well-definedness condition

\[
q(x)\;=\;q(x')\quad\Longrightarrow\quad qF(x)\;=\;qF(x')
\]

holds for all \(x,x'\in X\).

The instrument-indexed form is

\[
\Gamma\,;\;\mathsf H_{\mathrm{DC}}\,;\;I_{\mathrm{DC}}\;\vdash\;\textsc{DescentSquareRecovery}\;:\;\texttt{accepted}.
\]

\subsubsection{Proof}

The proof has two parts: the existence of a filler is equivalent to the well-definedness condition, and the existence of a filler is equivalent to the commutation equality. Both parts use only the finite set-theoretic structure of \(X\) and \(\mathrm{im}(q)\).

\emph{Filler from well-definedness.} Suppose \(q(x)=q(x')\Rightarrow qF(x)=qF(x')\). Define \(F^{\sharp}:\mathrm{im}(q)\to\mathrm{im}(q)\) by \(F^{\sharp}(y)=qF(x)\) for any \(x\in X\) with \(q(x)=y\). The well-definedness condition guarantees that the value does not depend on the choice of representative \(x\), so \(F^{\sharp}\) is a well-defined function. By construction \(F^{\sharp}\circ q_{\mathrm{im}}=q_{\mathrm{im}}\circ F\), so the descent commutation equality holds and \(\Xi_{\mathrm{desc}}\) is exact.

\emph{Well-definedness from filler.} Suppose \(\Xi_{\mathrm{desc}}\) is exact, so \(F^{\sharp}\circ q_{\mathrm{im}}=q_{\mathrm{im}}\circ F\). Suppose \(q(x)=q(x')\). Then \(q_{\mathrm{im}}(x)=q_{\mathrm{im}}(x')\), so \(F^{\sharp}(q_{\mathrm{im}}(x))=F^{\sharp}(q_{\mathrm{im}}(x'))\). Applying the commutation equality on both sides gives \(q_{\mathrm{im}}F(x)=q_{\mathrm{im}}F(x')\), that is, \(qF(x)=qF(x')\).

The two parts together yield the equivalence of the descent commutation equality and the well-definedness condition. Each step is a finite check on the finite host data of \(\mathsf H_{\mathrm{DC}}\), recorded in \(A_{\mathrm{DC}}\) and \(A_{\Xi}\). Under the admissibility assumptions for \(\mathsf H_{\mathrm{DC}}\) and the instrument \(I_{\mathrm{DC}}\), the in-scope, visibility, evidence, threshold, audit, and nonclaim conditions for the headline claim pass, and the instrument-indexed verdict is \(\texttt{accepted}\). \(\square\)

The descent square thus recovers descent data inside the decorated double-category fragment: a descent record \((F,q,F^{\sharp})\) is exact in \(\mathsf H_{\mathrm{DC}}\) iff \(F^{\sharp}\) is the well-defined descended map of \(F\) along \(q\). The recovery is a representational result; it does not assert that every descent fact about the core calculus is reducible to a descent square, in keeping with the secondary-organization discipline of the section.

\subsection{Theorem 31: Completion-Pasting Recovery}

This subsection proves Theorem 31: in the finite decorated double-category fragment of Subsection 14.1, the noncommutation of two finite completions appears as a finite pasting defect of type \(\texttt{completion}\), and the pasting square is exact iff the two completions commute.

\subsubsection{Statement}

Let \(C\in\mathrm{Ob}\) be a finite package object and let \(E_{1},E_{2}:C\to C\) be two finite completions in the sense of Subsection 7.4, with \(E_{2}\) used as the horizontal edge and \(E_{1}\) used as the vertical edge of the square. Form the decorated square

\[
\Xi_{\mathrm{comp}}\;:\;
\begin{array}{ccc}
C & \xrightarrow{E_{2}} & C\\
{\scriptstyle E_{1}}\downarrow & \Downarrow\,\Xi_{\mathrm{comp}} & \downarrow{\scriptstyle E_{1}}\\
C & \xrightarrow{E_{2}} & C
\end{array}
\]

with \(\mathrm{type}_{\Xi}=\texttt{completion}\) and \(\mathsf{Decor}_{\Xi}\) recording the completion pair \((E_{1},E_{2})\).

\textbf{Theorem 31 (CompletionPastingRecovery; theorem-grade).} The decorated square \(\Xi_{\mathrm{comp}}\) is exact iff the two completions commute, that is,

\[
E_{1}\,E_{2}\;=\;E_{2}\,E_{1}\;\;\text{on}\;\;C.
\]

The square is \emph{obstructed} iff there exists \(c\in C\) with \(E_{1}E_{2}(c)\neq E_{2}E_{1}(c)\). In that case the noncommutation defect

\[
\delta_{\mathrm{comp}}\bigl(E_{1},E_{2}\bigr)\;=\;\bigl\{\,c\in C\;:\;E_{1}E_{2}(c)\neq E_{2}E_{1}(c)\,\bigr\}
\]

is a finite typed defect record under \(I_{\mathrm{DC}}\), and the noncommutation \(E_{1}E_{2}\neq E_{2}E_{1}\) is recovered as a \(P_{3}\)-pasting defect of the calculus.

The instrument-indexed form is

\[
\Gamma\,;\;\mathsf H_{\mathrm{DC}}\,;\;I_{\mathrm{DC}}\;\vdash\;\textsc{CompletionPastingRecovery}\;:\;\texttt{accepted}.
\]

\subsubsection{Proof}

The decorated square \(\Xi_{\mathrm{comp}}\) has both horizontal edges equal to \(E_{2}\) and both vertical edges equal to \(E_{1}\). Exactness in the sense of Subsubsection 14.1.2 is the commutation equality of the boundary composition: the upper-right path \(E_{1}\circ E_{2}:C\to C\) equals the lower-left path \(E_{2}\circ E_{1}:C\to C\). By pointwise equality of maps on the finite carrier \(C\), this is the equation \(E_{1}E_{2}=E_{2}E_{1}\) on \(C\).

\emph{Forward direction.} Suppose \(\Xi_{\mathrm{comp}}\) is exact. Then by definition \(E_{1}E_{2}(c)=E_{2}E_{1}(c)\) for every \(c\in C\), hence \(E_{1}E_{2}=E_{2}E_{1}\) on \(C\).

\emph{Reverse direction.} Suppose \(E_{1}E_{2}=E_{2}E_{1}\) on \(C\). Then the boundary composition of \(\Xi_{\mathrm{comp}}\) commutes pointwise on the finite carrier \(C\), so by the exactness criterion of Subsubsection 14.1.2 the square is exact.

For the obstruction direction, suppose there exists \(c_{0}\in C\) with \(E_{1}E_{2}(c_{0})\neq E_{2}E_{1}(c_{0})\). The set

\[
\delta_{\mathrm{comp}}(E_{1},E_{2})\;=\;\bigl\{\,c\in C:E_{1}E_{2}(c)\neq E_{2}E_{1}(c)\,\bigr\}
\]

is a finite subset of \(C\) containing \(c_{0}\), and is recorded as the finite typed defect record \(\delta_{\Xi}\) of the square. Since \(\delta_{\Xi}\neq\varnothing\), the completion square is obstructed rather than exact. As in the completion-pasting defect of Subsection 7.5, this noncommutation of two finite completions is the route/pasting witness for \(P_{3}\). Hence \(E_{1}E_{2}\neq E_{2}E_{1}\) is recovered as a \(P_{3}\)-pasting defect.

Each step is a finite check on the finite host data of \(\mathsf H_{\mathrm{DC}}\), recorded in \(A_{\mathrm{DC}}\) and \(A_{\Xi}\). Under the admissibility assumptions for \(\mathsf H_{\mathrm{DC}}\) and the instrument \(I_{\mathrm{DC}}\), the headline claim receives instrument-indexed verdict \(\texttt{accepted}\). \(\square\)

The completion-pasting square thus records the noncommutation of two finite completions inside the decorated double-category fragment: commuting completions yield exact squares, and noncommuting completions yield finite pasting-defect records that recover the \(P_{3}\) route-mismatch content of the calculus.

\subsection{Theorem 32: Strict Extension as No-Filler}

This subsection proves Theorem 32: in the finite decorated double-category fragment of Subsection 14.1, a strict extension between two object maps \(\pi_{0},\pi_{1}\) on a common carrier \(S\) is recovered as the absence of a factorization filler in a decorated square of type \(\texttt{factorization}\). The theorem connects the strict-bridge nonfactorization witness of Section 7 (and its Cantor instance in Subsection 13.3) to the no-filler condition in \(\mathsf H_{\mathrm{DC}}\).

\subsubsection{Statement}

Let \(S\) be a finite carrier object and let \(\pi_{0}\in\mathrm{Hor}\) and \(\pi_{1}\in\mathrm{Hor}\) be two horizontal object maps \(\pi_{0}:S\to O_{0}\) and \(\pi_{1}:S\to O_{1}\) with finite codomains. Write \(\mathrm{im}(\pi_{0})\) for the effective old interface, and let \(\phi:\mathrm{im}(\pi_{0})\to O_{1}\) be a candidate vertical filler map. Form the decorated factorization square

\[
\Xi_{\mathrm{fact}}(\phi)\;:\;
\begin{array}{ccc}
S & \xrightarrow{\pi_{0}} & \mathrm{im}(\pi_{0})\\
{\scriptstyle\mathrm{id}_{S}}\downarrow & \Downarrow\,\Xi_{\mathrm{fact}}(\phi) & \downarrow{\scriptstyle\phi}\\
S & \xrightarrow{\pi_{1}} & O_{1}
\end{array}
\]

with \(\mathrm{type}_{\Xi}=\texttt{factorization}\) and \(\mathsf{Decor}_{\Xi}\) recording the factorization candidate \((\pi_{0},\pi_{1},\phi)\).

\textbf{Theorem 32 (StrictExtensionAsNoFiller; theorem-grade).} An exact filler \(\phi\) for the factorization square \(\Xi_{\mathrm{fact}}(\phi)\) exists iff

\[
\pi_{1}\;=\;\phi\circ\pi_{0}\;\;\text{on}\;\;S.
\]

Hence

\[
\pi_{1}\not\factor\pi_{0}\quad\Longleftrightarrow\quad\text{no exact factorization filler exists for}\;(\pi_{0},\pi_{1}),
\]

with \(\mathrm{Fill}(\pi_{0},\pi_{1})\;=\;\bigl\{\,\phi:\mathrm{im}(\pi_{0})\to O_{1}\;:\;\pi_{1}=\phi\circ\pi_{0}\,\bigr\}\) and

\[
\mathrm{Fill}(\pi_{0},\pi_{1})\;=\;\varnothing\quad\Longleftrightarrow\quad\exists\,s,s'\in S:\;\pi_{0}(s)=\pi_{0}(s')\;\land\;\pi_{1}(s)\neq\pi_{1}(s').
\]

The instrument-indexed form is

\[
\Gamma\,;\;\mathsf H_{\mathrm{DC}}\,;\;I_{\mathrm{DC}}\;\vdash\;\textsc{StrictExtensionAsNoFiller}\;:\;\texttt{accepted}.
\]

\subsubsection{Proof}

The proof translates the nonfactorization witness \(\pi_{1}\not\factor\pi_{0}\) of Subsection 7.5 into the no-filler condition for \(\Xi_{\mathrm{fact}}\).

\emph{Filler iff factorization.} Suppose \(\phi:\mathrm{im}(\pi_{0})\to O_{1}\) satisfies \(\pi_{1}=\phi\circ\pi_{0}\) on \(S\). Then the boundary composition of \(\Xi_{\mathrm{fact}}(\phi)\) commutes pointwise: for every \(s\in S\), the upper-right path returns \(\phi(\pi_{0}(s))=\pi_{1}(s)\), and the lower-left path returns \(\pi_{1}(\mathrm{id}_{S}(s))=\pi_{1}(s)\). By the exactness criterion of Subsubsection 14.1.2, \(\Xi_{\mathrm{fact}}(\phi)\) is exact, that is, \(\phi\) is a filler. Conversely, if \(\phi\) is a filler, then by definition the boundary composition commutes pointwise, so \(\pi_{1}(s)=\phi(\pi_{0}(s))\) for every \(s\in S\), that is, \(\pi_{1}=\phi\circ\pi_{0}\) on \(S\).

\emph{No filler iff nonfactorization witness.} Suppose \(\mathrm{Fill}(\pi_{0},\pi_{1})=\varnothing\). By the equivalence above, no \(\phi\) satisfies \(\pi_{1}=\phi\circ\pi_{0}\). Equivalently, the assignment \(\pi_{0}(s)\mapsto\pi_{1}(s)\) is not a well-defined function on \(\mathrm{im}(\pi_{0})\), so there exist \(s,s'\in S\) with \(\pi_{0}(s)=\pi_{0}(s')\) and \(\pi_{1}(s)\neq\pi_{1}(s')\). Conversely, suppose there exist such \(s,s'\in S\). Then any candidate \(\phi\) with \(\pi_{1}=\phi\circ\pi_{0}\) would require \(\phi(\pi_{0}(s))=\pi_{1}(s)\) and \(\phi(\pi_{0}(s))=\phi(\pi_{0}(s'))=\pi_{1}(s')\), contradicting \(\pi_{1}(s)\neq\pi_{1}(s')\). Hence \(\mathrm{Fill}(\pi_{0},\pi_{1})=\varnothing\).

The two equivalences together yield the theorem: \(\pi_{1}\not\factor\pi_{0}\) on \(S\) in the sense of Subsection 7.5 holds iff no exact filler exists for \(\Xi_{\mathrm{fact}}\), iff there exist \(s,s'\in S\) with \(\pi_{0}(s)=\pi_{0}(s')\) and \(\pi_{1}(s)\neq\pi_{1}(s')\). Each step is a finite check on the finite host data of \(\mathsf H_{\mathrm{DC}}\), recorded in \(A_{\mathrm{DC}}\) and \(A_{\Xi}\). Under the admissibility assumptions for \(\mathsf H_{\mathrm{DC}}\) and the instrument \(I_{\mathrm{DC}}\), the headline claim receives instrument-indexed verdict \(\texttt{accepted}\). \(\square\)

The factorization square thus records the strict-extension nonfactorization witness inside the decorated double-category fragment: a strict extension is exactly the absence of an exact filler for the factorization square. The Cantor strict-bridge realization of Subsection 13.3 supplies a concrete instance: the audited shell carries \(\pi_{1}\not\factor\pi_{0}\) on \(\mathcal S_{\mathrm{aud}}\), so \(\mathrm{Fill}(\pi_{0},\pi_{1})=\varnothing\) on the shell, and \(\Xi_{\mathrm{fact}}\) has no filler on the shell.

\subsection{High-Structure Nonclaims}

This subsection records the high-structure nonclaims that fix the boundaries of Section 14: the decorated-square fragment is secondary, the status family discipline is preserved, and the general propagation of pasting defects is future work.

\subsubsection{No Double-Category Completeness}

The finite decorated double-category fragment \(\mathsf H_{\mathrm{DC}}\) of Subsection 14.1 does not replace the core calculus \(\mathbf{BirdInt}^{\mathrm{aud}}_{\mathrm{fin}}\). Theorems 30, 31, and 32 record three specific recovery results --- descent data as descent squares, completion noncommutation as completion squares, and strict-extension nonfactorization as factorization squares --- and not a general claim that every theorem-grade fact about the calculus is reducible to a decorated-square fact.

In particular, the section does not claim:

\begin{itemize}
\tightlist
\item
  that \(\mathbf{BirdInt}^{\mathrm{aud}}_{\mathrm{fin}}\) embeds into \(\mathsf H_{\mathrm{DC}}\) as a sub-double-category;
\item
  that every directed-cell judgment of Subsection 4.3 is recovered by a decorated square;
\item
  that the priority order of the classifier of Subsection 11.2.1 is recovered by the per-square exactness conditions of Subsubsection 14.1.2.
\end{itemize}

The decorated-square fragment is a secondary organizational layer over the core calculus; it supplies finite, audit-checkable representations of selected core constructions where those representations are useful, and does not promote those representations to a completeness statement about the calculus.

\subsubsection{No Status Erasure}

The decorated-square fragment does not collapse the status families of Subsection 5.1 into a single square-status family. The role-status, cell-status, pair-status, promotion-status, claim-status, gate-status, and square-status families remain separate judgment-indexed families, in keeping with the no-total-algebra and status-family-separation results of Section 6.

In particular, the per-square exactness conditions of Subsubsection 14.1.2 do not assert that:

\begin{itemize}
\tightlist
\item
  every directed-cell status of Subsection 5.6.2 is determined by the exactness of a decorated square that contains the cell;
\item
  every pair-observable status of Subsection 5.6.3 is determined by the exactness of a decorated square that contains the pair;
\item
  every claim status of Subsection 11.2.1 is determined by the exactness of a decorated square.
\end{itemize}

A decorated square carries its own square-status verdict only as part of a full decorated/audited record: witness and update data, language, visibility, threshold, audit, defect, square status, and nonclaim data must be present when the square is used as a \(\mathbf{BirdInt}^{\mathrm{aud}}_{\mathrm{fin}}\) claim. A bare commutative diagram is not an admissible claim. The square-status verdict does not export to the other status families, and the status-separation discipline of the calculus is preserved by Section 14.

\subsubsection{No General Obstruction Propagation Theorem}

The pasting law of Subsubsection 14.1.3 records the aggregation rule for defects under finite strict pasting: \(\delta_{\Xi_{1}\cdot\Xi_{2}}=\delta_{\Xi_{1}}\cup\delta_{\Xi_{2}}\cup\delta_{\mathrm{paste}}(\Xi_{1},\Xi_{2})\). This aggregation is a syntactic record on the finite paste, not a propagation theorem.

In particular, Section 14 does not claim a general theorem of the form

\begin{quote}
a defect at one square in a finite pasting diagram propagates to a defect at every square downstream of it,
\end{quote}

nor does it claim a general theorem of the form

\begin{quote}
the absence of a defect at one square in a finite pasting diagram suffices for the absence of a defect at every square in the diagram.
\end{quote}

Both directions of general defect propagation are future work. Specific propagation facts under specific pasting diagrams may be admissible under the per-diagram audit data of \(A_{\Xi_{1}\cdot\Xi_{2}}\), but no general propagation theorem is asserted by the calculus in this paper.

The three nonclaims above together fix the local boundary of the secondary high-structure semantics within \(\mathbf{BirdInt}^{\mathrm{aud}}_{\mathrm{fin}}\): the decorated-square fragment supplies finite recoveries of descent, completion, and strict-extension data, preserves the separation of the status families, and records pasted defects as finite aggregated records without a general propagation theorem.
 \section{Lean 4 Mechanization and Validation}

This section records the completed Lean 4 validation artifact for the finite audited scope of \(\mathbf{BirdInt}^{\mathrm{aud}}_{\mathrm{fin}}\) and Proposition 35, the audit-grade validation statement for that artifact. The Lean development is a repository artifact, not merely a future target: it is the Lake project in \texttt{lean/full}, rooted at the module \texttt{SixBirdsIII}, checked with the toolchain \texttt{leanprover/lean4:v4.28.0}, and audited by the manifest-driven validator \texttt{scripts/validate\_lean.sh}.

The validation scope is finite and explicit. The required paper definitions and numbered claims are listed in \texttt{lean/manifest.toml}, with the paper inventory recorded in \texttt{lean/paper\_inventory.toml}. The allowed non-project axiom trust base for theorem audits is recorded in \texttt{lean/trust\_base.txt}.

\subsection{Lean Artifact}

The Lean project is \texttt{lean/full}, with:

\begin{itemize}
\tightlist
\item
  root module \texttt{SixBirdsIII};
\item
  Lake configuration \texttt{lean/full/lakefile.toml};
\item
  toolchain \texttt{lean/full/lean-toolchain}, containing \texttt{leanprover/lean4:v4.28.0};
\item
  module tree \texttt{lean/full/SixBirdsIII/*.lean};
\item
  semantic validation probes in \texttt{lean/full/SixBirdsIII/SemanticTests.lean}.
\end{itemize}

The root import module \texttt{lean/full/SixBirdsIII.lean} imports the full mechanized stack: primitive labels, levels, hosts, status families, finite record schemas, the core interaction theorems, finite-map and completion theorems, graph/cohomology theorems, promotion and instrument theorems, abstract-interpretation theorems, high-structure recovery theorems, model-realization records, and the mechanization-feasibility record.

\subsection{Validator}

The strict validation command is \texttt{./scripts/validate\_lean.sh\ -\/-tier\ full\ -\/-require-full}.

The validator checks:

\begin{enumerate}
\def\labelenumi{\arabic{enumi}.}
\tightlist
\item
  the pinned toolchain;
\item
  the Lake build of \texttt{lean/full};
\item
  coverage of every manifest entry;
\item
  existence probes for the listed Lean declarations;
\item
  expected declaration-kind checks for definitions and theorems;
\item
  theorem axiom dependency audits against \texttt{lean/trust\_base.txt};
\item
  placeholder scans;
\item
  semantic-audit probes.
\end{enumerate}

With the completed track, strict validation reports 21 full schema definitions, 35 theorem-backed required claims, no unformalized definitions or claims, 56/56 declaration probes passing, 56/56 declaration-kind checks passing, 35/35 theorem axiom audits passing, no placeholders, no semantic-audit findings, and the overall result \texttt{FULL\ MECHANIZATION\ PASS}.

\subsection{Proposition 35: Mechanization Target Feasibility}

\textbf{Proposition 35 (MechanizationTargetFeasibility).} The finite audited manifest scope of \(\mathbf{BirdInt}^{\mathrm{aud}}_{\mathrm{fin}}\) is represented by a Lean 4 validation artifact. Specifically:

\begin{enumerate}
\def\labelenumi{\arabic{enumi}.}
\tightlist
\item
  every required schema definition listed in \texttt{lean/manifest.toml} is represented by a Lean declaration in the \texttt{SixBirdsIII} namespace;
\item
  every required numbered claim listed in \texttt{lean/manifest.toml} is represented by a theorem-backed Lean declaration;
\item
  the declarations build under \texttt{leanprover/lean4:v4.28.0};
\item
  theorem declarations use only the allowed non-project axiom trust base recorded in \texttt{lean/trust\_base.txt};
\item
  the strict validator passes with no unformalized manifest entries, no placeholder findings, and no semantic-audit findings.
\end{enumerate}

The instrument-indexed form is

\[
\Gamma\,;\;\mathsf H_{\mathrm{fin}}\,;\;I_{\mathrm{meta}}\;\vdash\;\text{the manifest-scoped Lean artifact validates the finite audited paper declarations}\;:\;\texttt{audit\_report}.
\]

The verdict \(\texttt{audit\_report}\) records a repository-level validation audit. It does not change theorem-grade results into model-realization results, and it does not promote model-grade statements to universal theorem-grade claims.

\subsection{Mechanization Scope Nonclaim}

The Lean artifact validates the finite audited scope represented by \texttt{lean/manifest.toml}. It does not assert:

\begin{itemize}
\tightlist
\item
  any theorem, definition, model realization, or audit record not listed in \texttt{lean/manifest.toml};
\item
  empirical correctness of any concrete PICA, Cantor, abstract-interpretation, graph, or implementation system;
\item
  a universal exact-six theorem;
\item
  a universal six-symbol algebra;
\item
  a unique decomposition theorem;
\item
  complete model coverage;
\item
  complete same-level self-certification;
\item
  proof-assistant certification beyond the finite declarations audited by the strict validator.
\end{itemize}

This is the mechanization boundary used by Appendix F and Section 16. The Lean track is completed for the manifest-scoped finite audited artifact; claims outside that manifest require their own declarations, manifest entries, and validation audit before they can be reported as mechanized.
 \section{Limitations, Nonclaims, and Future Work}

This section collects the final scope discipline of \(\mathbf{BirdInt}^{\mathrm{aud}}_{\mathrm{fin}}\) and the future-work boundaries of the paper. Subsection 16.1 restates the positive contribution of the paper across six layers --- finite typed interaction calculus, statused profile-relative interaction, the theorem-grade finite core, the countermodel/no-go atlas, the instrument and self-reference scope, and the model-realization scope. Subsection 16.2 records the required nonclaims, organized by the six topical groupings of the calculus. Subsection 16.3 records the six classes of future work that the paper does not undertake. Subsection 16.4 records the final paper-inclusion nonclaim block.

The discipline of the section is uniform: the paper records exactly what it asserts and exactly what it does not assert, and treats the nonclaims and future-work entries as part of the paper's content rather than as exclusions of content. The section does not introduce any new theorem or model-realization claim; it collects, restates, and frames what the earlier sections have already recorded.

\subsection{Positive Scope Summary}

This subsection restates the positive contribution of the paper across the six structural layers of \(\mathbf{BirdInt}^{\mathrm{aud}}_{\mathrm{fin}}\).

\subsubsection{Finite Typed Interaction Calculus}

The first paper constructs \(\mathbf{BirdInt}^{\mathrm{aud}}_{\mathrm{fin}}\), a finite typed interaction calculus over six primitive labels \(\mathbb P=\{P_{1},\ldots,P_{6}\}\), with a finite typed schema of directed-cell records, pair-observable records, promotion-bridge records, claim records, instrument records, source-of-truth records, audit records, and visibility records. The calculus is defined in Sections 3 and 4 as a finite typed package with finite carriers, finite stores, and finite admissibility predicates on each component. Every datum in the calculus is finite, every carrier is finite, and every record is interpreted under an explicit host instrument and audit record.

\subsubsection{Statused Profile-Relative Interaction}

The second positive layer is the statused profile-relative interaction discipline of Section 5. Each judgment carries a status family: the role status of Subsection 5.1.1, the directed-cell status of Subsection 5.1.2, the pair-observable status of Subsection 5.1.3, the promotion status of Subsection 5.1.4, the claim status of Subsection 5.1.5, the gate status of Subsection 5.1.6, and the square status of Subsection 5.1.7. The status families are typed-separately judgment-indexed: a directed-cell status is not a pair status, a pair status is not a promotion status, and the priority orders within each family are frozen by the classifier of Subsection 5.6.

The same primitive pair \((P_{i},P_{j})\) receives different statuses on different profiles: profile relativity is enforced by the per-cell profile \(\lambda\) of Subsection 4.3 and reflected in Theorem 1's status non-factorization through the primitive pair (Section 6).

\subsubsection{Theorem-Grade Finite Core}

The third positive layer is the numbered theorem and proposition spine of the paper. The theorem-grade finite core is concentrated in Sections 6, 7, 8, 10, 11, 12, and 14; the model-grade and audit-report statements are recorded separately.

\begin{itemize}
\tightlist
\item
  Section 6: the no-total-algebra theorem (Theorem 1) and the four core finite-interaction theorems (Theorems 2--5).
\item
  Section 7: nine finite-domain theorems on descent, refinement, packaging, completion, strictness, and audit (Theorems 6--14).
\item
  Section 8: four cohomological theorems on graph-supported drive (Theorems 15--18).
\item
  Section 10: four promotion-and-stacking theorems (Theorems 19--22).
\item
  Section 11: three self-reference and visibility theorems (Theorems 23--25).
\item
  Section 12: the abstract-interpretation theorem trio (Theorems 26--28) and the abstract-interpretation model-family theorem (Model Theorem 29).
\item
  Section 14: the high-structure recovery theorems (Theorems 30--32).
\item
  Section 13: the PICA and Cantor model-realization theorems (Model Theorems 33--34), together with the graph/cohomology model-realization statement of Subsection 13.4.
\item
  Section 15: the Lean 4 validation audit report (Proposition 35).
\end{itemize}

Each numbered statement is recorded with its appropriate instrument-indexed claim form and verdict: \(\texttt{accepted}\) for theorem-grade results, \(\texttt{model\_realization}\) for scoped model realizations, and \(\texttt{audit\_report}\) for the manifest-scoped Lean validation artifact.

\subsubsection{Countermodel/No-Go Scope}

The fourth positive layer is the countermodel atlas of Section 9 and Appendix C. The paper records twelve finite countermodels, each blocking a specific overclaim about the calculus: active package without macro closure (CM\_1), active route/holonomy without drive (CM\_2), local packages failing globally (CM\_3), directed cell active while the pair observable is blocked (CM\_4), refinement worsening closure (CM\_5), same-level self-audit circularity (CM\_6), same primitive pair with different statuses (CM\_7), idempotent completions that do not commute (CM\_8), strict extension without macro closure (CM\_9), strict extension without drive (CM\_10), fallback source unable to support a real pair observable (CM\_11), and gating that destroys cycle-supported drive (CM\_12).

Each countermodel is a finite typed structure under an instrument whose visible content blocks the corresponding overclaim under \(I\). The atlas is finite, audited, and host-bound, in keeping with the no-go scope discipline of Subsection 9.1.

\subsubsection{Instrument/Self-Reference Scope}

The fifth positive layer is the instrument and self-reference scope of Section 11. The paper records the instrument-indexed claim semantics of Subsection 11.1, the claim classifier priority of Subsection 11.2, the visibility and suppression discipline of Subsection 11.3, the no-overreading suppression theorem (Theorem 23), the same-level self-audit failure theorem (Theorem 24), the \(P_{6}\leftarrow P_{6}\) status separation of Subsection 11.6, the rotating-audit chain and bridge of Subsection 11.7, the finite rotating-audit theorem (Theorem 25), and the instrument-transfer rule of Subsection 11.9.

The discipline records that no claim under any instrument \(I\) overreads its visible-or-bridged content, that no instrument certifies its own complete self-soundness without a level-shifted audit, that finite rotating-audit chains are admissible but non-final, and that instrument transfer is bridge-mediated and non-implication.

\subsubsection{Model-Realization Scope}

The sixth positive layer is the model-realization scope of Section 13 and Section 12. The paper records four scoped model realizations: the PICA realization \(\mathsf{PICAReal}_{\mathrm{fin}}\) of Subsection 13.2 with Model Theorem 33, the Cantor strict-bridge realization \(\mathsf{CantorShell}_{\mathrm{aud}}\) of Subsection 13.3 with Model Theorem 34, the abstract-interpretation family \(\mathsf{AIFam}_{\mathrm{fin}}\) of Section 12 with Model Theorem 29, and the graph/cohomology realization \(\mathsf{CohGraph}_{\mathrm{fin}}\) of Subsection 13.4. Each realization is host-bound, fragment-scoped, and certified at claim-strength \(\texttt{model\_realization}\), in keeping with the model-realization convention of Subsection 13.1.

\subsection{Required Nonclaims}

This subsection records the required nonclaims of the paper, organized into six topical groupings: algebra and universality, status, package and strictness, route/refinement/locality, audit, and model. Each nonclaim is a scoped negative statement; together they fix the boundary of \(\mathbf{BirdInt}^{\mathrm{aud}}_{\mathrm{fin}}\) as a finite typed audited interaction calculus and prevent the positive results of Subsection 16.1 from being read as universal foundations.

\subsubsection{Algebra and Universality Nonclaims}

\begin{itemize}
\tightlist
\item
  \emph{No universal six-symbol algebra.} The paper does not assert a total operation \(\ast:\mathbb P^{2}\to\mathbb P\) that produces a faithful universal interaction algebra over the six primitive labels. The non-factorization of cell status through the primitive pair (Theorem 1, Section 6) blocks any such operation as a status-decoding map.
\item
  \emph{No all-mathematics reduction.} The paper does not assert that all mathematics is six birds, or that every mathematical object is a \(\mathbf{BirdInt}^{\mathrm{aud}}_{\mathrm{fin}}\) object under any reading.
\item
  \emph{No unrestricted exact-six theorem.} The paper does not assert that every typed interaction is exactly representable as one of the thirty-six directed cells \(P_{i}\leftarrow P_{j}\). The exact-six structure is recorded as a finite typed schema with admissibility conditions, not as an unrestricted decomposition theorem.
\item
  \emph{No unique decomposition.} The paper does not assert that a typed interaction has a unique decomposition into primitive directed cells. Decomposition records are admissible only under a declared profile and audit, and different profiles can yield different admissible decompositions.
\end{itemize}

\subsubsection{Status Nonclaims}

\begin{itemize}
\tightlist
\item
  \emph{No primitive-pair status determination.} The paper does not assert that the primitive pair \((P_{i},P_{j})\) alone determines the directed-cell status. The status non-factorization theorem (Theorem 1, Section 6) records the formal counterexample.
\item
  \emph{No status-family collapse.} The paper does not assert that the seven status families of Subsection 5.1 collapse into a single status family. Role status, cell status, pair status, promotion status, claim status, gate status, and square status remain typed-separately judgment-indexed; Theorem 5 of Section 6 records representative separations.
\item
  \emph{No directed-cell-to-pair collapse.} The paper does not assert that a directed-cell status determines a pair-observable status. CM\_4 of Section 9 records a directed cell with status \(\texttt{action}\) whose corresponding pair observable is \(\texttt{blocked}\).
\item
  \emph{No role-channel-to-cell collapse.} The paper does not assert that the role status of a primitive determines the directed-cell status of any cell whose actor or informant is that primitive. Role status is a separate judgment family from cell status; Theorem 5 of Section 6 includes the channel-not-cell separation.
\end{itemize}

\subsubsection{Package and Strictness Nonclaims}

\begin{itemize}
\tightlist
\item
  \emph{No package-implies-closure.} The paper does not assert that a \(P_{5}\)-packaging operator on a finite carrier implies macro-dynamics closure of any associated transformer. CM\_1 of Section 9 records an active package whose macro dynamics are not closed under \(I\).
\item
  \emph{No idempotence-implies-commutation.} The paper does not assert that two idempotent finite completions \(E_{1},E_{2}:C\to C\) commute. CM\_8 of Appendix C, named in Subsection 9.8, records two idempotent completions with \(E_{1}E_{2}\neq E_{2}E_{1}\).
\item
  \emph{No repeated fixed completion as theory growth.} The paper does not assert that repeated application of a single fixed idempotent completion is the same as the packaged completion of a target package. The Cantor target package \(T_{1}\) of Subsubsection 13.3.3 uses a packaged completion \(E_{\tau,\ell}=U_{\ell}\circ Q_{\ell}\circ K^{\tau}\), not a repeated single completion.
\item
  \emph{No strictness-implies-closure.} The paper does not assert that a strict-bridge passage \(\mathrm{Promote}(B):\texttt{strict}\) implies macro-dynamics closure of either side of the bridge. CM\_9 of Appendix C, named in Subsection 9.8, records a strict bridge without closure.
\item
  \emph{No strictness-implies-drive.} The paper does not assert that a strict-bridge passage implies a \(\mathrm{P6}_{\mathrm{drive}}\) certificate on either side. CM\_10 of Appendix C, named in Subsection 9.8, records a strict bridge without drive.
\end{itemize}

\subsubsection{Route, Refinement, and Locality Nonclaims}

\begin{itemize}
\tightlist
\item
  \emph{No route-mismatch-implies-drive.} The paper does not assert that a \(P_{3}\) route-mismatch witness implies a \(\mathrm{P6}_{\mathrm{drive}}\) certificate. CM\_2 of Section 9 records a route mismatch without drive, and the cohomology nonclaims of Subsubsection 13.4.4 record the corresponding model-level form.
\item
  \emph{No refinement-implies-improvement.} The paper does not assert that a refinement \(A_{k}\rightsquigarrow A_{\ell}\) improves every defect record. CM\_5 of Section 9 records a refinement that worsens closure under a chosen defect measure.
\item
  \emph{No local-implies-global.} The paper does not assert that local admissibility of a package implies global admissibility. CM\_3 of Section 9 records local packages that fail globally.
\end{itemize}

\subsubsection{Audit Nonclaims}

\begin{itemize}
\tightlist
\item
  \emph{No audit-is-truth.} The paper does not assert that a claim status \(\texttt{accepted}\) under instrument \(I\) is instrument-free truth. The instrument-indexed claim form of Subsection 11.1.3 fixes acceptance as relative to \(I\) and the surrounding context \(\Gamma;\mathcal T\).
\item
  \emph{No same-level complete self-certification.} The paper does not assert that a single instrument \(I_{0}\) certifies its own complete soundness. Theorem 24 of Section 11 records the formal failure: same-level complete self-soundness is either \(\texttt{outside\_scope}\) or \(\texttt{undefined\_circular}\) under \(I_{0}\).
\item
  \emph{No rotating-audit finality.} The paper does not assert that a finite rotating-audit chain is final. Theorem 25 of Section 11 and the non-finality clause of Subsubsection 11.8.3 record that a finite rotating-audit chain audits the lower stack but does not certify its top.
\item
  \emph{No accepted claim overreads hidden content.} The paper does not assert that an accepted claim under \(I\) implies any fact about \(I\)-suppressed content. Theorem 23 of Section 11 records the formal no-overreading result.
\end{itemize}

\subsubsection{Model Nonclaims}

\begin{itemize}
\tightlist
\item
  \emph{No model-realization-as-universal-theorem.} The paper does not assert that a \(\texttt{model\_realization}\) verdict promotes to a theorem-grade \(\texttt{accepted}\) verdict, or that it proves the declared structure universally, without an admissible bridge. The fragment discipline of Subsection 13.1.2 records the corresponding rule.
\item
  \emph{No PICA universality.} The paper does not assert that PICA realizes the entire \(\mathbf{BirdInt}^{\mathrm{aud}}_{\mathrm{fin}}\) calculus. The PICA realization is host-bound and fragment-scoped to the finite stochastic cell/pair/provenance fragment.
\item
  \emph{No PICA \(6\times 6\) universal algebra.} The paper does not assert that the recovered \(6\times 6\) multiplicity count of Subsubsection 13.2.5 is a universal \(6\times 6\) law of \(\mathbf{BirdInt}^{\mathrm{aud}}_{\mathrm{fin}}\). It is a PICA-profile model fact under the cell profile \(\lambda_{\mathrm{PICA}}\).
\item
  \emph{No raw implementation status as formal status.} The paper does not assert that a raw PICA implementation status is a formal directed-cell status without recovery and classifier check; Subsubsection 13.2.6 records the recovery step.
\item
  \emph{No fallback-as-real-source.} The paper does not assert that a fallback record supports a real pair observable. The fallback-default cap of Subsubsection 13.2.8 records that fallback supports at most \(\texttt{real\_provisional}\) unless an explicit upgrade bridge is recorded, and CM\_11 of Appendix C gives the corresponding finite source-policy countermodel.
\item
  \emph{No dirty-state realness.} The paper does not assert that a dirty record can support a real pair observable under any host configuration of \(\mathsf H_{\mathrm{PICA}}\); Subsubsection 13.2.8 records the dirty-state exclusion.
\item
  \emph{No Cantor generalization beyond audited shell.} The paper does not assert that the Cantor strict-bridge realization extends beyond the audited shell \(\mathcal S_{\mathrm{aud}}\). Subsubsection 13.3.9 records the corresponding nonclaim register.
\item
  \emph{No AI totality.} The paper does not assert that the abstract-interpretation host realizes the entire calculus. The realized fragment is \(P_{1}/P_{2}/P_{4}/P_{5}/P_{6}\); a full six-role AI realization would require additional \(P_{3}\) data, and no such extension is asserted here.
\item
  \emph{No graph cohomology as all drive.} The paper does not assert that the cohomological drive certificate \([a]\neq 0\) exhausts every notion of drive. The cohomology nonclaims of Subsubsection 13.4.4 record \(\mathrm{P6}_{\mathrm{drive}}^{H^1}\) as one finite certificate family, not all drive.
\item
  \emph{No double-category completeness.} The paper does not assert that the finite decorated-square fragment of Section 14 replaces or supersedes the core calculus. Subsubsection 14.5.1 records the corresponding nonclaim.
\item
  \emph{No mechanization beyond manifest scope.} The paper does not assert proof-assistant certification for any theorem, definition, model realization, or audit record outside \texttt{lean/manifest.toml}. The audit-grade verdict of Proposition 35 records the completed Lean validation artifact for the finite manifest scope, not a certificate for unlisted statements, empirical implementations, or unrestricted host extensions.
\end{itemize}

\subsection{Future Work}

This subsection records the six classes of future work that the present paper does not undertake. Each class is a scoped pointer to admissible extension targets for \(\mathbf{BirdInt}^{\mathrm{aud}}_{\mathrm{fin}}\); none of the entries below is a claim about results in this paper.

\subsubsection{Classifier and Dependency Completeness}

Future work on the claim classifier of Subsection 11.2 includes:

\begin{itemize}
\tightlist
\item
  a complete status-classifier theorem, characterizing the priority order over \(\mathrm{ClaimStatus}\) under stated soundness and admissibility conditions for the claim classifier of Subsection 5.6.5;
\item
  a defect-to-status soundness theorem, characterizing the routing from the defect families of Subsections 5.2--5.5 through the classifier rules of Subsection 5.6;
\item
  a dependency-lattice theorem, organizing the per-component admissibility predicates of Section 4 and Section 5 into a lattice of dependent admissibility checks with a finite reduction theorem.
\end{itemize}

The present paper records the priority order of Subsection 11.2.1 and the per-status routing rules of Subsection 5.6, but does not establish completeness or soundness in these stronger forms.

\subsubsection{Promotion and Stacking Theory}

Future work on promotion and stacking includes:

\begin{itemize}
\tightlist
\item
  a complete promotion classifier, extending the eight-gate machinery of Subsection 10.2 to a full first-match procedure with stated soundness and completeness conditions over \(\mathrm{PromotionStatus}\);
\item
  a promotion well-formedness decidability theorem, certifying that admissibility of a promotion bridge is decidable on finite records under the host instrument;
\item
  a general stacking composition theorem, characterizing the composition of admissible promotion bridges across host stacks beyond the per-bridge results of Section 10 and the explicit boundary in Subsection 10.7;
\item
  a local-to-global promotion theorem, identifying the conditions under which a finite family of locally admissible promotion bridges supports a global promotion verdict.
\end{itemize}

The present paper records Theorems 19--22 of Section 10, the eight-gate machinery, and the per-bridge classifier; the four extension targets above are recorded as future work.

\subsubsection{Verified Model Realizations}

Future work on the model-realization sheets includes:

\begin{itemize}
\tightlist
\item
  a deeper PICA realization audit, extending Model Theorem 33 of Subsection 13.2 toward implementation-correctness records for the actor-update map, informant-witness map, status recovery, and pair-observable realness predicate;
\item
  a deeper Cantor strict-bridge realization audit, extending Model Theorem 34 of Subsection 13.3 toward implementation-correctness records for the shell-bound strictness witness and the gate-result table;
\item
  a deeper abstract-interpretation realization audit, extending Model Theorem 29 of Section 12 beyond the manifest-scoped Lean support toward richer concrete domains and transformers;
\item
  an implementation-correctness layer for any concrete PICA implementation, audited under a separate proof-assistant-bound instrument.
\end{itemize}

The present paper records the four realization sheets at model-grade verdict \(\texttt{model\_realization}\) and records manifest-scoped Lean declarations for the required model claims. Implementation correctness and richer external model audits remain future work.

\subsubsection{Cantor Shell Expansion}

Future work on the Cantor strict-bridge realization includes:

\begin{itemize}
\tightlist
\item
  a full Cantor strictness proof, extending the audited-shell strict-bridge realization beyond \(\mathcal S_{\mathrm{aud}}\) to a wider scope under an admissible expansion bridge;
\item
  a pressure-gap-to-KL bridge, recording an admissible bridge from the pressure-gap consequence defect of Subsubsection 13.3.6 to the Kullback--Leibler closure-deficit framework;
\item
  a stratumwise root-separation theorem, characterizing the persistent strata family \(\mathcal F\) via a separation property that supports stronger consequences than the conditional pressure disintegration of Subsubsection 13.3.6.
\end{itemize}

The present paper records the audited-shell realization with verdict \(\texttt{model\_realization}\) on \(\mathbf{StrictBridge}_{0\to 1}^{\mathrm{audited\ shell}}\); the three extensions above are recorded as future work.

\subsubsection{Drive and High-Structure Semantics}

Future work on the drive and high-structure layers includes:

\begin{itemize}
\tightlist
\item
  a general drive classification, extending the cohomological drive certificate \([a]\neq 0\) of Subsubsection 13.4.3 to a finite family of drive certificates over multiple drive notions;
\item
  a holonomy-to-drive bridge \(B_{3\to 6}^{\mathrm{drive}}\), recording an admissible bridge from a \(P_{3}\) route/holonomy witness to a \(\mathrm{P6}_{\mathrm{drive}}\) certificate, in the sense of Subsubsection 13.4.4 nonclaim 2;
\item
  gating beyond forests, characterizing \(P_{2}\)-gating effects on drive on graphs with non-trivial cycle structure under additional gating policies;
\item
  weak double-category coherence, extending the strict pasting law of Subsubsection 14.1.3 to weak coherence relations on finite decorated squares;
\item
  a general obstruction propagation theorem, characterizing the propagation of square defects under finite strict pasting in the sense of Subsubsection 14.5.3.
\end{itemize}

The present paper records the cohomological theorems of Section 8 and the high-structure recovery theorems of Section 14; the extensions above are recorded as future work.

\subsubsection{Self-Reference and Instrument Transfer}

Future work on self-reference and instrument transfer includes:

\begin{itemize}
\tightlist
\item
  a full instrument-transfer classifier, fixing the per-claim-type mapping from source-side acceptance and bridge defect to target-side \(\chi\in\{\,\texttt{accepted},\;\texttt{provisional},\;\texttt{blocked}\,\}\) under \(J\), in the sense of Subsubsection 11.9.2;
\item
  an infinite rotating-audit semantics, extending the finite rotating-audit chain of Subsection 11.7 to an inverse limit or coalgebraic semantics admissible under a stronger meta-instrument;
\item
  a complete self-reference classification, extending Theorem 24 of Subsection 11.5 to a comprehensive classification of same-level, level-shifted, and rotating self-reference patterns.
\end{itemize}

The present paper records the instrument-transfer rule of Subsection 11.9.1, the rotating-audit chain of Subsection 11.7, and Theorems 24 and 25; the extensions above are recorded as future work.

\subsection{Final Paper-Inclusion Nonclaim Block}

This subsection records the final paper-inclusion nonclaim block of the paper: a manuscript-prose adaptation of the thirty-six numbered nonclaims \(\mathrm{NC}\text{-}1\) through \(\mathrm{NC}\text{-}36\) of the final scope and nonclaims packet, organized into the six topical groupings of Subsection 16.2 and reflected in the finite nonclaim records \(\mathcal N\) carried by instruments, judgments, bridges, and model-realization sheets.

The block is the paper's consolidated roster of negative statements: every nonclaim recorded in the body of the paper (per-section nonclaim subsections, per-realization nonclaim registers, per-bridge nonclaim records) is covered by, or refines, one of the thirty-six entries below. The block is finite and records the paper's commitment to the boundary it has fixed.

\subsubsection{Algebra and Universality}

\begin{itemize}
\tightlist
\item
  \(\mathrm{NC}\text{-}1\). No universal six-symbol algebra: there is no faithful total operation \(\ast:\mathbb P^{2}\to\mathbb P\) whose composition with a status-decoding map recovers typed, statused, profile-relative directed-cell status across all profiles.
\item
  \(\mathrm{NC}\text{-}2\). No ``all mathematics is six birds'' claim: \(\mathbf{BirdInt}^{\mathrm{aud}}_{\mathrm{fin}}\) is a finite typed audited interaction calculus, not a foundation for all mathematics.
\end{itemize}

\subsubsection{Status}

\begin{itemize}
\tightlist
\item
  \(\mathrm{NC}\text{-}3\). No primitive-pair status determination: the directed-cell status does not factor through the primitive pair \((P_{i},P_{j})\) alone.
\item
  \(\mathrm{NC}\text{-}4\). No status-family collapse: \(\mathrm{RoleStatus}\), \(\mathrm{CellStatus}\), \(\mathrm{PairStatus}\), \(\mathrm{PromotionStatus}\), \(\mathrm{ClaimStatus}\), \(\mathrm{GateStatus}\), and \(\mathrm{SqStatus}\) are typed-separately judgment-indexed.
\item
  \(\mathrm{NC}\text{-}5\). No directed-cell-to-pair collapse: a directed-cell status does not determine the pair-observable status.
\item
  \(\mathrm{NC}\text{-}6\). No role-channel-to-cell collapse: the role status of a primitive does not determine the directed-cell status of any cell carrying that primitive.
\end{itemize}

\subsubsection{Package and Strictness}

\begin{itemize}
\tightlist
\item
  \(\mathrm{NC}\text{-}7\). No package-implies-closure: a \(P_{5}\)-packaging operator does not imply macro-dynamics closure of an associated transformer.
\item
  \(\mathrm{NC}\text{-}8\). No idempotence-implies-commutation: two idempotent finite completions \(E_{1},E_{2}:C\to C\) may fail to commute.
\item
  \(\mathrm{NC}\text{-}9\). No repeated fixed completion as theory growth: a packaged completion is not a single fixed idempotent applied repeatedly.
\item
  \(\mathrm{NC}\text{-}10\). No strictness-implies-closure: a strict-bridge passage \(\mathrm{Promote}(B):\texttt{strict}\) does not imply macro-dynamics closure.
\item
  \(\mathrm{NC}\text{-}11\). No strictness-implies-drive: a strict-bridge passage does not imply a \(\mathrm{P6}_{\mathrm{drive}}\) certificate.
\end{itemize}

\subsubsection{Route, Refinement, and Locality}

\begin{itemize}
\tightlist
\item
  \(\mathrm{NC}\text{-}12\). No route-mismatch-implies-drive: a \(P_{3}\) route-mismatch witness does not imply a \(\mathrm{P6}_{\mathrm{drive}}\) certificate.
\item
  \(\mathrm{NC}\text{-}13\). No refinement-implies-improvement: a refinement \(A_{k}\rightsquigarrow A_{\ell}\) does not improve every defect record.
\item
  \(\mathrm{NC}\text{-}14\). No local-implies-global: local admissibility of a finite family of packages does not imply global admissibility.
\item
  \(\mathrm{NC}\text{-}\mathrm{TD}\). No structural-downward-implies-channel claim: a completion map, feasibility gate, macro profile, phase record, promoted object, or context-indexed update supplies at most a structural downward path or constraint. It certifies top-down causal difference-making only when a finite \(\mathsf{TopDownChannelRecord}\) supplies admissible macro interventions, matched controls, source-of-truth response records, visibility, audit, no-smuggling, and a passed host-specific effect threshold.
\end{itemize}

\subsubsection{Audit}

\begin{itemize}
\tightlist
\item
  \(\mathrm{NC}\text{-}15\). No audit-is-truth: a passing audit record or claim status \(\texttt{accepted}\) under an instrument is not instrument-free truth.
\item
  \(\mathrm{NC}\text{-}16\). No same-level complete self-certification: an instrument \(I_{0}\) does not certify its own complete soundness at the same level as an accepted global self-certification.
\item
  \(\mathrm{NC}\text{-}17\). No rotating-audit finality: a finite rotating-audit chain is non-final.
\item
  \(\mathrm{NC}\text{-}18\). No accepted claim overreads hidden content: an \(I\)-accepted claim does not imply any fact about \(I\)-suppressed unbridged content.
\item
  \(\mathrm{NC}\text{-}19\). No model-realization-as-theorem overclaim: a \(\texttt{model\_realization}\) verdict does not promote to theorem-grade \(\texttt{accepted}\) without an admissible bridge.
\end{itemize}

\subsubsection{Model}

\begin{itemize}
\tightlist
\item
  \(\mathrm{NC}\text{-}20\). No PICA universality: \(\mathsf{PICAReal}_{\mathrm{fin}}\) realizes a finite stochastic fragment of the calculus, not its entirety.
\item
  \(\mathrm{NC}\text{-}21\). No PICA \(6\times 6\) universal algebra: the recovered PICA \(6\times 6\) multiplicity count is a profile fact, not a universal law.
\item
  \(\mathrm{NC}\text{-}22\). No raw implementation status as formal status: raw PICA status is not formal directed-cell status without recovery and classifier check.
\item
  \(\mathrm{NC}\text{-}23\). No fallback-as-real-source: a fallback record supports at most \(\texttt{real\_provisional}\) unless an admissible upgrade bridge is recorded.
\item
  \(\mathrm{NC}\text{-}24\). No dirty-state realness: a dirty record cannot support a real pair observable.
\item
  \(\mathrm{NC}\text{-}25\). No Cantor generalization beyond audited shell: the Cantor strict-bridge realization is shell-bound.
\item
  \(\mathrm{NC}\text{-}26\). No Cantor pressure gap as KL closure deficit: the pressure-gap consequence defect is not a Kullback--Leibler closure-deficit without a separately admissible bridge.
\item
  \(\mathrm{NC}\text{-}27\). No Cantor pressure gap as strictness witness: the pressure-gap defect is a consequence record, not the strictness witness.
\item
  \(\mathrm{NC}\text{-}28\). No direct autonomous stratum theorem: the persistent strata family \(\mathcal F\) does not support a direct stratumwise root-separation theorem in the present paper.
\item
  \(\mathrm{NC}\text{-}29\). No AI totality: the abstract-interpretation host realizes the fragment \(P_{1}/P_{2}/P_{4}/P_{5}/P_{6}\), not the whole calculus.
\item
  \(\mathrm{NC}\text{-}30\). No AI soundness-as-completeness: sound abstraction does not imply complete representation.
\item
  \(\mathrm{NC}\text{-}31\). No AI closure-implies-soundness: the closure operator \(\rho=\gamma\alpha\) does not imply transformer soundness.
\item
  \(\mathrm{NC}\text{-}32\). No graph cohomology as all drive: the cohomological drive certificate \([a]\neq 0\) is one drive certificate family, not all drive.
\item
  \(\mathrm{NC}\text{-}33\). No gating-preserves-drive: \(P_{2}\)-gating may suppress cycle-supported drive on \(G\).
\item
  \(\mathrm{NC}\text{-}34\). No double-category completeness: the finite decorated-square fragment of Section 14 does not replace the core calculus.
\item
  \(\mathrm{NC}\text{-}35\). No high-structure status erasure: the decorated-square fragment does not replace statuses, defects, audits, or instruments, and does not collapse the seven status families into a single square-status family.
\item
  \(\mathrm{NC}\text{-}36\). No mechanization beyond manifest scope: the Lean 4 artifact validates the finite audited declarations listed in \texttt{lean/manifest.toml}, and does not certify unlisted statements, empirical implementation correctness, a universal exact-six theorem, or a universal six-symbol algebra.
\end{itemize}

The thirty-six nonclaims above are the paper's final manuscript-prose record. Each is supported by either a theorem, a countermodel, a per-realization nonclaim register, or a scope-discipline statement recorded earlier in the paper, and each is carried by the relevant finite nonclaim records \(\mathcal N\) of the calculus. The block closes the boundary of \(\mathbf{BirdInt}^{\mathrm{aud}}_{\mathrm{fin}}\): the paper records what it asserts, what it does not assert, and what it leaves to future work.
 \section{Conclusion}

This section closes the paper by restating the exact earned contribution: the constructed object \(\mathbf{BirdInt}^{\mathrm{aud}}_{\mathrm{fin}}\), the layered theorem stack it admits, the scope discipline it imposes, and the final boundary it does not cross.

\subsection{The Constructed Object}

The paper constructs \(\mathbf{BirdInt}^{\mathrm{aud}}_{\mathrm{fin}}\), a finite audited typed interaction calculus for six-bird theory. The object is finite throughout: every carrier is finite, every store is finite, every admissibility predicate is decidable on finite records, and every claim is recorded under an explicit host instrument. Its components are six primitive labels \(\mathbb P=\{P_{1},\ldots,P_{6}\}\), a finite typed schema of directed-cell, pair-observable, promotion-bridge, claim, instrument, source-of-truth, audit, and visibility records, seven finite status families with their classifier or per-value rules where applicable, a finite typed family of defect records that route judgments and claims to their status consequences, and a finite nonclaim register that records the negative boundary of every admissible claim.

The calculus is not a universal algebra over six symbols, an all-mathematics decomposition, or a foundation for empirical phenomena. It is a finite typed audited interaction calculus, defined by Sections 3, 4, and 5, and admitting the theorem-grade, model-grade, and audit-grade results recorded in Sections 6 through 15.

\subsection{The Theorem Stack}

The paper records its earned content as a layered theorem stack:

\begin{itemize}
\tightlist
\item
  \emph{Finite typed calculus and core theorems.} Sections 6 and 7 record the no-total-algebra theorem (Theorem 1), the four core finite-interaction theorems (Theorems 2--5), and the nine descent/packaging/completion/strictness theorems (Theorems 6--14), all theorem-grade with verdict \(\texttt{accepted}\).
\item
  \emph{No-go and cohomological theorems.} Section 8 records the four cohomological no-drive theorems (Theorems 15--18) on graph-supported drive.
\item
  \emph{Countermodel/no-go atlas.} Section 9 and Appendix C record twelve countermodels (CM\_1--CM\_12), each blocking a specific overclaim about the calculus under a finite typed structure.
\item
  \emph{Promotion and instrument-indexed semantics.} Section 10 records the promotion gate and strictness theorems (Theorems 19--21) together with the Toy22 strict-bridge example (Theorem 22), and Section 11 records the no-overreading suppression theorem (Theorem 23), the same-level self-audit failure theorem (Theorem 24), and the finite rotating-audit theorem (Theorem 25), each accompanied by the finite typed instrument-indexed claim form.
\item
  \emph{Scoped model realizations.} Section 12 records theorem-grade abstract-interpretation host results (Theorems 26--28) and the model-grade abstract-interpretation family result (Model Theorem 29). Section 13 records the PICA realization (Model Theorem 33), the Cantor strict-bridge realization (Model Theorem 34), and the graph/cohomology realization, all model-grade with verdict \(\texttt{model\_realization}\) on declared fragments.
\item
  \emph{High-structure semantics as secondary organization.} Section 14 records three decorated-square recovery theorems (Theorems 30--32) over a finite double-category fragment, presented as secondary organization rather than as a replacement for the core calculus.
\item
  \emph{Lean validation audit.} Section 15 records the manifest-scoped Lean 4 validation artifact and Proposition 35, with the audit-grade verdict \(\texttt{audit\_report}\).
\end{itemize}

The seven layers together record the paper's earned content. The theorem and model-realization claims carry their appropriate claim-strength verdicts, while the countermodels and Lean validation proposition are recorded with their corresponding finite dependency or audit records.

\subsection{The Scope Discipline}

The paper replaces a naive total algebra over six primitive symbols with a typed, finite, statused, audited interaction calculus. The replacement is structural: cell status does not factor through the primitive pair (Theorem 1), the judgment-attached status families do not collapse onto a single family and the gate and square families remain schema-separate (Theorem 5 and Section 5), descent is well-definedness on a finite quotient (Theorem 6), strict extension is fiber mismatch on a finite carrier (Theorem 12), accepted claims do not overread suppressed content (Theorem 23), no instrument certifies its own complete soundness at the same level (Theorem 24), and finite rotating audits are admissible but non-final (Theorem 25). Each replacement is recorded as a finite typed result under explicit host, instrument, admissibility, and nonclaim assumptions, with defect schemas explicit where they drive status.

The scope discipline is uniform: every theorem-grade claim, model-realization claim, promotion judgment, and audit-report statement is accompanied by an instrument-indexed or host-indexed form, a finite admissibility predicate, and a nonclaim record that fixes the boundary of the claim. The paper records exactly what it asserts and exactly what it does not assert; the negative scope is recorded in Section 16 and across the per-section nonclaim subsections, not as a defensive afterthought but as part of the calculus itself.

\subsection{Final Boundary}

The final boundary of the paper is the finite audited scope. The calculus \(\mathbf{BirdInt}^{\mathrm{aud}}_{\mathrm{fin}}\) is a finite typed audited interaction calculus, admitting the theorem-grade, model-grade, and audit-grade results recorded in this paper, and not admitting an unrestricted exact-six theorem, a universal six-symbol algebra, a unique decomposition theorem, complete model-coverage, complete instrument-family completeness, or final self-certification.

The paper closes with this boundary in hand: the constructed object, the layered theorem stack, the scope discipline, and the explicit non-extension of any of the above into a universal foundation claim. The earned contribution is the calculus and what it admits --- finite, typed, statused, audited, and bounded by an explicit nonclaim register.
 
\appendix
\section{Full Definitions}

This appendix records the long-form definitions and tables of \(\mathbf{BirdInt}^{\mathrm{aud}}_{\mathrm{fin}}\), supplementing the body of the paper. Subsection A.1 records the full domain tuple. Subsection A.2 records the full primitive witness table. Subsection A.3 records the full primitive update table. Subsection A.4 records the full profile grammar and admissibility table. Subsection A.5 records the full status-family tables. Subsection A.6 records the full primitive and promotion defect tables.

\subsection{Full Domain Tuple}

The full domain tuple of \(\mathbf{BirdInt}^{\mathrm{aud}}_{\mathrm{fin}}\) is the sixteen-component finite typed record

\[
\mathbf{BirdInt}^{\mathrm{aud}}_{\mathrm{fin}}\;=\;(\,\mathsf{Host},\;\mathsf{Cont},\;\mathbb P,\;\mathsf{Lev},\;\mathsf{Instr},\;\mathsf{Pkg},\;\mathsf{Prof},\;\mathsf{Wit},\;\mathsf{Upd},\;\mathsf{Def},\;\mathsf{Judg},\;\mathsf{Status},\;\mathsf{Gate},\;\mathsf{Dep},\;\mathsf{Audit},\;\mathsf{Nonclaim}\,),
\]

with components:

\begin{itemize}
\tightlist
\item
  \(\mathsf{Host}\), the finite and scoped host inventory, recording finite carriers, finite maps, finite records, finite relations, finite diagrams, finite graphs, and finite audit records (Subsection 2.1);
\item
  \(\mathsf{Cont}\), the finite content set whose records may be visible, suppressed, outside-scope, or unknown under an instrument (Subsections 3.2 and 3.5);
\item
  \(\mathbb P=\{P_{1},P_{2},P_{3},P_{4},P_{5},P_{6}\}\), the finite primitive labels with stable role glosses (Subsection 2.2);
\item
  \(\mathsf{Lev}\), the finite level set \(\{\mathsf{Beh},\mathsf{Ver},\mathsf{Str}\}\) used by the calculus (Subsection 2.2.3);
\item
  \(\mathsf{Instr}\), the finite class of instrument records, with components introduced in Subsection 3.4 and used in the claim semantics of Subsection 11.1;
\item
  \(\mathsf{Pkg}\), the finite class of theory packages with components \((Z,f,\Sigma_{f},E,\mathcal A)\) (Subsection 3.3);
\item
  \(\mathsf{Prof}\), the finite typed profile family \(\lambda\) with admissibility predicates (Subsection 3.6);
\item
  \(\mathsf{Wit}\), the finite typed witness family for primitive informants (Subsection 3.7 and Table A.2);
\item
  \(\mathsf{Upd}\), the finite typed update family for primitive actors (Subsection 3.7 and Table A.3);
\item
  \(\mathsf{Def}\), the finite typed defect family with per-component schemas (Subsections 5.2 through 5.5 and Table A.6);
\item
  \(\mathsf{Judg}\), the finite typed judgment family --- role-channel, directed-cell, pair-observable, promotion, claim, gate, square --- with per-judgment admissibility predicates (Subsections 4.2 through 4.7, 10.2, 11.1, and 14.1);
\item
  \(\mathsf{Status}\), the seven finite status families, with their family-specific classifier and priority rules where defined (Subsection 5.1 and Table A.5);
\item
  \(\mathsf{Gate}\), the finite typed gate family --- eight promotion gates, the strict-pass gate, and bridge-admissibility gates (Subsection 10.2);
\item
  \(\mathsf{Dep}\), the finite dependency record, recording producer/consumer relations among store entries and the dependency-graph schema (Subsection 4.6);
\item
  \(\mathsf{Audit}\), the finite typed audit family --- host audit, claim audit, bridge audit, rotating-audit (Subsection 4.8 and Section 11);
\item
  \(\mathsf{Nonclaim}\), the finite nonclaim register \(\mathcal N\) (Subsection 16.4 and per-record nonclaim subfields throughout).
\end{itemize}

Every strong claim under the calculus is instrument-indexed:

\[
\Gamma\,;\;\mathcal T\,;\;I\;\vdash\;\varphi\;:\;\chi,
\]

with \(\chi\in\mathrm{ClaimStatus}\) as recorded in Subsection 5.1.5. The shorthand \(I\vdash\varphi(\mathcal T)\) abbreviates \(\Gamma;\mathcal T;I\vdash\varphi:\texttt{accepted}\) under the acceptance rule of Subsection 11.1.3 and never abbreviates instrument-free truth.

The domain is \emph{finite}: every component is finite, every admissibility predicate is decidable on finite records, and every claim under the domain is recorded with a finite typed instrument-indexed claim form. The design principle of the domain is that interaction is typed, statused, profile-relative, instrument-indexed, and audited --- and that the primitive-pair collapse \(P_{i}\ast P_{j}=P_{k}\) is rejected as a universal representation of interaction in keeping with Theorem 1 of Section 6.

\subsection{Full Primitive Witness Table}

The witness families are typed by the informant primitive: \(\mathsf{Wit}:\mathbb P\to\mathrm{Type}\), with \(W_{j}\in\mathsf{Wit}(P_{j})\) for the informant \(P_{j}\) of any directed cell \(P_{i}\leftarrow P_{j}\). The full witness families per primitive are:

\begin{longtable}[]{@{}ll@{}}
\toprule
\begin{minipage}[b]{0.08\columnwidth}\raggedright
primitive\strut
\end{minipage} & \begin{minipage}[b]{0.86\columnwidth}\raggedright
witness family\strut
\end{minipage}\tabularnewline
\midrule
\endhead
\begin{minipage}[t]{0.08\columnwidth}\raggedright
\(P_{1}\)\strut
\end{minipage} & \begin{minipage}[t]{0.86\columnwidth}\raggedright
descent success/failure, rewrite obstruction, closure deficit\strut
\end{minipage}\tabularnewline
\begin{minipage}[t]{0.08\columnwidth}\raggedright
\(P_{2}\)\strut
\end{minipage} & \begin{minipage}[t]{0.86\columnwidth}\raggedright
representability, nonrepresentability, gate, support, capacity\strut
\end{minipage}\tabularnewline
\begin{minipage}[t]{0.08\columnwidth}\raggedright
\(P_{3}\)\strut
\end{minipage} & \begin{minipage}[t]{0.86\columnwidth}\raggedright
route mismatch, holonomy, commutator, critical pair, pasting mismatch\strut
\end{minipage}\tabularnewline
\begin{minipage}[t]{0.08\columnwidth}\raggedright
\(P_{4}\)\strut
\end{minipage} & \begin{minipage}[t]{0.86\columnwidth}\raggedright
stage, refinement, filtration, cover, local/global obstruction\strut
\end{minipage}\tabularnewline
\begin{minipage}[t]{0.08\columnwidth}\raggedright
\(P_{5}\)\strut
\end{minipage} & \begin{minipage}[t]{0.86\columnwidth}\raggedright
package map, quotient, completion, idempotence, fixed point\strut
\end{minipage}\tabularnewline
\begin{minipage}[t]{0.08\columnwidth}\raggedright
\(P_{6}\)\strut
\end{minipage} & \begin{minipage}[t]{0.86\columnwidth}\raggedright
audit, provenance, source-of-truth, threshold check, nonclaim, drive certificate\strut
\end{minipage}\tabularnewline
\bottomrule
\end{longtable}

The typed actor/informant rule of Subsection 4.3 fixes the admissibility of any directed-cell judgment: a cell record \(P_{i}\leftarrow P_{j}\) is admissible only when its witness \(W_{j}\in\mathsf{Wit}(P_{j})\) has type drawn from the row of \(P_{j}\) above and its actor update has type drawn from the row of \(P_{i}\) in Subsection A.3 below. The witness families are finite for each primitive, and admissibility of the typed witness assignment is a decidable check on the finite records of the cell.

\subsection{Full Primitive Update Table}

The update families are typed by the actor primitive: \(\mathsf{Upd}:\mathbb P\to\mathrm{Type}\), with \(U_{i}\in\mathsf{Upd}(P_{i})\) for the actor \(P_{i}\) of any directed cell \(P_{i}\leftarrow P_{j}\). The full update/effect families per primitive are:

\begin{longtable}[]{@{}ll@{}}
\toprule
\begin{minipage}[b]{0.08\columnwidth}\raggedright
primitive\strut
\end{minipage} & \begin{minipage}[b]{0.86\columnwidth}\raggedright
update/effect family\strut
\end{minipage}\tabularnewline
\midrule
\endhead
\begin{minipage}[t]{0.08\columnwidth}\raggedright
\(P_{1}\)\strut
\end{minipage} & \begin{minipage}[t]{0.86\columnwidth}\raggedright
mark descent failure, install descended update, rewrite, repair\strut
\end{minipage}\tabularnewline
\begin{minipage}[t]{0.08\columnwidth}\raggedright
\(P_{2}\)\strut
\end{minipage} & \begin{minipage}[t]{0.86\columnwidth}\raggedright
mark representable/nonrepresentable, gate support, update admissibility\strut
\end{minipage}\tabularnewline
\begin{minipage}[t]{0.08\columnwidth}\raggedright
\(P_{3}\)\strut
\end{minipage} & \begin{minipage}[t]{0.86\columnwidth}\raggedright
record route mismatch, holonomy, commutator, critical pair\strut
\end{minipage}\tabularnewline
\begin{minipage}[t]{0.08\columnwidth}\raggedright
\(P_{4}\)\strut
\end{minipage} & \begin{minipage}[t]{0.86\columnwidth}\raggedright
refine lens, advance stage, record local/global obstruction\strut
\end{minipage}\tabularnewline
\begin{minipage}[t]{0.08\columnwidth}\raggedright
\(P_{5}\)\strut
\end{minipage} & \begin{minipage}[t]{0.86\columnwidth}\raggedright
install package, apply completion, record fixed points, saturate\strut
\end{minipage}\tabularnewline
\begin{minipage}[t]{0.08\columnwidth}\raggedright
\(P_{6}\)\strut
\end{minipage} & \begin{minipage}[t]{0.86\columnwidth}\raggedright
add audit, check provenance, commit source, record nonclaim, certify drive/no-drive\strut
\end{minipage}\tabularnewline
\bottomrule
\end{longtable}

The actor update is the second half of the typed actor/informant rule of Subsection 4.3: a cell record \(P_{i}\leftarrow P_{j}\) is admissible only when both \(W_{j}\in\mathsf{Wit}(P_{j})\) and \(U_{i}\in\mathsf{Upd}(P_{i})\) hold for the witness and update assigned to the cell. The update families are finite for each primitive, and admissibility of the typed update assignment is a decidable check on the finite records of the cell.

\subsection{Full Profile Grammar}

The profile grammar of \(\mathbf{BirdInt}^{\mathrm{aud}}_{\mathrm{fin}}\) is the finite typed record schema

\[
\lambda\;=\;(\,\ell_{i},\;\ell_{j},\;\mathsf{Scale},\;\mathsf{Regime},\;\mathsf{BridgeType},\;\mathsf{InstrumentMode},\;\mathsf{Locality}\,),
\]

with components:

\begin{longtable}[]{@{}ll@{}}
\toprule
\begin{minipage}[b]{0.47\columnwidth}\raggedright
component\strut
\end{minipage} & \begin{minipage}[b]{0.47\columnwidth}\raggedright
admissible values / role\strut
\end{minipage}\tabularnewline
\midrule
\endhead
\begin{minipage}[t]{0.47\columnwidth}\raggedright
\(\ell_i\)\strut
\end{minipage} & \begin{minipage}[t]{0.47\columnwidth}\raggedright
actor-side level tag in \(\mathsf{Lev}=\{\mathsf{Beh},\mathsf{Ver},\mathsf{Str}\}\)\strut
\end{minipage}\tabularnewline
\begin{minipage}[t]{0.47\columnwidth}\raggedright
\(\ell_j\)\strut
\end{minipage} & \begin{minipage}[t]{0.47\columnwidth}\raggedright
informant-side level tag in \(\mathsf{Lev}=\{\mathsf{Beh},\mathsf{Ver},\mathsf{Str}\}\)\strut
\end{minipage}\tabularnewline
\begin{minipage}[t]{0.47\columnwidth}\raggedright
\(\mathsf{Scale}\)\strut
\end{minipage} & \begin{minipage}[t]{0.47\columnwidth}\raggedright
finite scale label for the carrier, partition, kernel, graph, or other hosted data used by the judgment\strut
\end{minipage}\tabularnewline
\begin{minipage}[t]{0.47\columnwidth}\raggedright
\(\mathsf{Regime}\)\strut
\end{minipage} & \begin{minipage}[t]{0.47\columnwidth}\raggedright
finite regime label, such as dynamical, structural, or audit, as fixed by the surrounding judgment family\strut
\end{minipage}\tabularnewline
\begin{minipage}[t]{0.47\columnwidth}\raggedright
\(\mathsf{BridgeType}\)\strut
\end{minipage} & \begin{minipage}[t]{0.47\columnwidth}\raggedright
bridge type admitted by \(\mathsf{BridgePolicy}_I\), such as visibility, level, instrument-transfer, drive, or model-realization\strut
\end{minipage}\tabularnewline
\begin{minipage}[t]{0.47\columnwidth}\raggedright
\(\mathsf{InstrumentMode}\)\strut
\end{minipage} & \begin{minipage}[t]{0.47\columnwidth}\raggedright
finite instrument-mode label, such as default, audit-only, transfer, or rotating-audit\strut
\end{minipage}\tabularnewline
\begin{minipage}[t]{0.47\columnwidth}\raggedright
\(\mathsf{Locality}\)\strut
\end{minipage} & \begin{minipage}[t]{0.47\columnwidth}\raggedright
finite locality label, such as local, semi-global, or global use of the judgment data\strut
\end{minipage}\tabularnewline
\bottomrule
\end{longtable}

When individual fields of a profile must be named, the paper writes them as \(\mathsf{Scale}_{\lambda}\), \(\mathsf{Regime}_{\lambda}\), \(\mathsf{BridgeType}_{\lambda}\), \(\mathsf{InstrumentMode}_{\lambda}\), and \(\mathsf{Locality}_{\lambda}\).

Named profiles include

\[
\lambda_{\mathrm{cell}},\quad\lambda_{\mathrm{prom}},\quad\lambda_{\mathrm{AI}},\quad\lambda_{\mathrm{PICA}},\quad\lambda_{\mathrm{Cantor}},
\]

corresponding to directed-cell judgments, promotion-bridge records, abstract-interpretation host claims, PICA-host claims, and Cantor-shell claims.

Profile admissibility is the predicate

\[
\mathsf{AdmProfile}(\lambda,I,\mathcal T)
\]

on a profile \(\lambda\), an instrument \(I\), and a target package \(\mathcal T\). The admissibility predicate is finite and decidable on the finite components of \(\lambda\), \(I\), and \(\mathcal T\).

\begin{longtable}[]{@{}ll@{}}
\toprule
\begin{minipage}[b]{0.47\columnwidth}\raggedright
check\strut
\end{minipage} & \begin{minipage}[b]{0.47\columnwidth}\raggedright
admissibility condition\strut
\end{minipage}\tabularnewline
\midrule
\endhead
\begin{minipage}[t]{0.47\columnwidth}\raggedright
level\strut
\end{minipage} & \begin{minipage}[t]{0.47\columnwidth}\raggedright
\(\ell_i,\ell_j\in\mathsf{Lev}\), with cross-level use requiring an admissible level bridge\strut
\end{minipage}\tabularnewline
\begin{minipage}[t]{0.47\columnwidth}\raggedright
scale/regime\strut
\end{minipage} & \begin{minipage}[t]{0.47\columnwidth}\raggedright
\(\mathsf{Scale}\) and \(\mathsf{Regime}\) are consistent with the host of \(\mathcal T\) and with the surrounding judgment family\strut
\end{minipage}\tabularnewline
\begin{minipage}[t]{0.47\columnwidth}\raggedright
bridge type\strut
\end{minipage} & \begin{minipage}[t]{0.47\columnwidth}\raggedright
\(\mathsf{BridgeType}\) is admitted by \(\mathsf{BridgePolicy}_I\)\strut
\end{minipage}\tabularnewline
\begin{minipage}[t]{0.47\columnwidth}\raggedright
instrument mode\strut
\end{minipage} & \begin{minipage}[t]{0.47\columnwidth}\raggedright
\(\mathsf{InstrumentMode}\) is consistent with \(I\)'s declared mode for the surrounding judgment\strut
\end{minipage}\tabularnewline
\begin{minipage}[t]{0.47\columnwidth}\raggedright
locality\strut
\end{minipage} & \begin{minipage}[t]{0.47\columnwidth}\raggedright
\(\mathsf{Locality}\) is consistent with the carrier, cover, staging, or support data of \(\mathcal T\)\strut
\end{minipage}\tabularnewline
\bottomrule
\end{longtable}

Profile mutation is not allowed without an explicit profile update record: a judgment under profile \(\lambda\) does not become a judgment under a different profile \(\lambda'\) unless a finite typed update record for that profile change is present in \(\mathsf{Cont}\) and passes the relevant visibility, source, and audit checks. This rule records the profile-relative discipline used later by the no-total-algebra theorem; it does not prove that theorem here.

\subsection{Full Status Tables}

The seven status families and their values are recorded below. The family-specific classifier and priority rules, where applicable, are fixed in Section 5 and Section 11; this subsection records only the value sets and the brief gloss per value.

\subsubsection{Role Status}

\[
\mathrm{RoleStatus}\;=\;\{\,\texttt{active\_projection},\;\texttt{below\_threshold},\;\texttt{collapsed\_boundary\_case},\;\texttt{absent\_with\_record},\;\texttt{inapplicable\_outside\_scope}\,\}.
\]

\begin{longtable}[]{@{}ll@{}}
\toprule
\begin{minipage}[b]{0.33\columnwidth}\raggedright
value\strut
\end{minipage} & \begin{minipage}[b]{0.61\columnwidth}\raggedright
meaning\strut
\end{minipage}\tabularnewline
\midrule
\endhead
\begin{minipage}[t]{0.33\columnwidth}\raggedright
\(\texttt{active\_projection}\)\strut
\end{minipage} & \begin{minipage}[t]{0.61\columnwidth}\raggedright
role \(P_{i}\) is active and visible in the package/profile\strut
\end{minipage}\tabularnewline
\begin{minipage}[t]{0.33\columnwidth}\raggedright
\(\texttt{below\_threshold}\)\strut
\end{minipage} & \begin{minipage}[t]{0.61\columnwidth}\raggedright
role evidence exists but does not meet activation threshold\strut
\end{minipage}\tabularnewline
\begin{minipage}[t]{0.33\columnwidth}\raggedright
\(\texttt{collapsed\_boundary\_case}\)\strut
\end{minipage} & \begin{minipage}[t]{0.61\columnwidth}\raggedright
role degenerates to identity, singleton, total collapse, or boundary\strut
\end{minipage}\tabularnewline
\begin{minipage}[t]{0.33\columnwidth}\raggedright
\(\texttt{absent\_with\_record}\)\strut
\end{minipage} & \begin{minipage}[t]{0.61\columnwidth}\raggedright
absence of the role/witness is audited\strut
\end{minipage}\tabularnewline
\begin{minipage}[t]{0.33\columnwidth}\raggedright
\(\texttt{inapplicable\_outside\_scope}\)\strut
\end{minipage} & \begin{minipage}[t]{0.61\columnwidth}\raggedright
instrument or host does not adjudicate this role\strut
\end{minipage}\tabularnewline
\bottomrule
\end{longtable}

The role/cell separation rule of Section 5 records that an active role channel does not imply \(\texttt{action}\) on every cell carrying that role.

\subsubsection{Cell Status}

\[
\mathrm{CellStatus}\;=\;\{\,\texttt{action},\;\texttt{implicit},\;\texttt{trivial},\;\texttt{blocked},\;\texttt{undefined\_circular},\;\texttt{below\_threshold},\;\texttt{collapsed},\;\texttt{absent},\;\texttt{outside\_scope}\,\}.
\]

\begin{longtable}[]{@{}ll@{}}
\toprule
\begin{minipage}[b]{0.22\columnwidth}\raggedright
value\strut
\end{minipage} & \begin{minipage}[b]{0.72\columnwidth}\raggedright
meaning\strut
\end{minipage}\tabularnewline
\midrule
\endhead
\begin{minipage}[t]{0.22\columnwidth}\raggedright
\(\texttt{action}\)\strut
\end{minipage} & \begin{minipage}[t]{0.72\columnwidth}\raggedright
witness materially and directly parameterizes actor update\strut
\end{minipage}\tabularnewline
\begin{minipage}[t]{0.22\columnwidth}\raggedright
\(\texttt{implicit}\)\strut
\end{minipage} & \begin{minipage}[t]{0.72\columnwidth}\raggedright
effect is real but mediated by refresh, recomputation, or dependency closure\strut
\end{minipage}\tabularnewline
\begin{minipage}[t]{0.22\columnwidth}\raggedright
\(\texttt{trivial}\)\strut
\end{minipage} & \begin{minipage}[t]{0.72\columnwidth}\raggedright
well-formed cell exists but is identity-like, saturated, or no-op\strut
\end{minipage}\tabularnewline
\begin{minipage}[t]{0.22\columnwidth}\raggedright
\(\texttt{blocked}\)\strut
\end{minipage} & \begin{minipage}[t]{0.72\columnwidth}\raggedright
in scope, but missing bridge, wrong type, hidden source, failed visibility, or failed audit prevents use\strut
\end{minipage}\tabularnewline
\begin{minipage}[t]{0.22\columnwidth}\raggedright
\(\texttt{undefined\_circular}\)\strut
\end{minipage} & \begin{minipage}[t]{0.72\columnwidth}\raggedright
cell depends on an unstratified circular audit/support loop\strut
\end{minipage}\tabularnewline
\begin{minipage}[t]{0.22\columnwidth}\raggedright
\(\texttt{below\_threshold}\)\strut
\end{minipage} & \begin{minipage}[t]{0.72\columnwidth}\raggedright
evidence exists but activation threshold fails\strut
\end{minipage}\tabularnewline
\begin{minipage}[t]{0.22\columnwidth}\raggedright
\(\texttt{collapsed}\)\strut
\end{minipage} & \begin{minipage}[t]{0.72\columnwidth}\raggedright
candidate degenerates and no longer instantiates the intended cell structure\strut
\end{minipage}\tabularnewline
\begin{minipage}[t]{0.22\columnwidth}\raggedright
\(\texttt{absent}\)\strut
\end{minipage} & \begin{minipage}[t]{0.72\columnwidth}\raggedright
required witness/update is absent with an absence record\strut
\end{minipage}\tabularnewline
\begin{minipage}[t]{0.22\columnwidth}\raggedright
\(\texttt{outside\_scope}\)\strut
\end{minipage} & \begin{minipage}[t]{0.72\columnwidth}\raggedright
host/instrument/profile does not adjudicate the cell\strut
\end{minipage}\tabularnewline
\bottomrule
\end{longtable}

\subsubsection{Pair Status}

\[
\mathrm{PairStatus}\;=\;\{\,\texttt{real},\;\texttt{real\_provisional},\;\texttt{real\_duplicated},\;\texttt{blocked}\,\}.
\]

\begin{longtable}[]{@{}ll@{}}
\toprule
\begin{minipage}[b]{0.19\columnwidth}\raggedright
value\strut
\end{minipage} & \begin{minipage}[b]{0.75\columnwidth}\raggedright
meaning\strut
\end{minipage}\tabularnewline
\midrule
\endhead
\begin{minipage}[t]{0.19\columnwidth}\raggedright
\(\texttt{real}\)\strut
\end{minipage} & \begin{minipage}[t]{0.75\columnwidth}\raggedright
pair observable has source-of-truth, compatibility, visibility, threshold, and audit\strut
\end{minipage}\tabularnewline
\begin{minipage}[t]{0.19\columnwidth}\raggedright
\(\texttt{real\_provisional}\)\strut
\end{minipage} & \begin{minipage}[t]{0.75\columnwidth}\raggedright
partial/model/fallback/ablation evidence supports provisional pair status\strut
\end{minipage}\tabularnewline
\begin{minipage}[t]{0.19\columnwidth}\raggedright
\(\texttt{real\_duplicated}\)\strut
\end{minipage} & \begin{minipage}[t]{0.75\columnwidth}\raggedright
pair is real but duplicates an already-accepted equivalent pair observable\strut
\end{minipage}\tabularnewline
\begin{minipage}[t]{0.19\columnwidth}\raggedright
\(\texttt{blocked}\)\strut
\end{minipage} & \begin{minipage}[t]{0.75\columnwidth}\raggedright
missing source, missing bridge, incompatible branches, circularity, visibility failure, or profile exclusion\strut
\end{minipage}\tabularnewline
\bottomrule
\end{longtable}

The pair-realness rule of Section 5 records that \(\texttt{real}\) requires a \(\mathsf{SourceOfTruthOK}\) record on the underlying source datum.

\subsubsection{Promotion Status}

\[
\mathrm{PromotionStatus}\;=\;\{\,\texttt{candidate},\;\texttt{accepted},\;\texttt{strict},\;\texttt{non\_strict},\;\texttt{failed\_descent},\;\texttt{failed\_stability},\;\texttt{failed\_audit},\;\texttt{failed\_no\_smuggling},\;\texttt{local\_only},\;\texttt{globally\_obstructed},\;\texttt{outside\_scope}\,\}.
\]

The values \(\texttt{strict}\) and \(\texttt{non\_strict}\) are accepted-status refinements: a \(\texttt{strict}\) verdict implies the accepted core and \(\pi_{j+1}\not\factor\pi_{j}\), and a \(\texttt{non\_strict}\) verdict implies the accepted core and \(\pi_{j+1}\factor\pi_{j}\), in keeping with Theorems 19, 20, 21 of Section 10.

\subsubsection{Claim Status}

\[
\mathrm{ClaimStatus}\;=\;\{\,\texttt{accepted},\;\texttt{rejected},\;\texttt{provisional},\;\texttt{blocked},\;\texttt{outside\_scope},\;\texttt{absent\_with\_record},\;\texttt{below\_threshold},\;\texttt{failed\_audit},\;\texttt{undefined\_circular}\,\}.
\]

\begin{longtable}[]{@{}ll@{}}
\toprule
\begin{minipage}[b]{0.26\columnwidth}\raggedright
value\strut
\end{minipage} & \begin{minipage}[b]{0.68\columnwidth}\raggedright
meaning\strut
\end{minipage}\tabularnewline
\midrule
\endhead
\begin{minipage}[t]{0.26\columnwidth}\raggedright
\(\texttt{accepted}\)\strut
\end{minipage} & \begin{minipage}[t]{0.68\columnwidth}\raggedright
all required checks pass\strut
\end{minipage}\tabularnewline
\begin{minipage}[t]{0.26\columnwidth}\raggedright
\(\texttt{rejected}\)\strut
\end{minipage} & \begin{minipage}[t]{0.68\columnwidth}\raggedright
instrument has an admissible refuting witness\strut
\end{minipage}\tabularnewline
\begin{minipage}[t]{0.26\columnwidth}\raggedright
\(\texttt{provisional}\)\strut
\end{minipage} & \begin{minipage}[t]{0.68\columnwidth}\raggedright
partial evidence exists and profile permits provisional status\strut
\end{minipage}\tabularnewline
\begin{minipage}[t]{0.26\columnwidth}\raggedright
\(\texttt{blocked}\)\strut
\end{minipage} & \begin{minipage}[t]{0.68\columnwidth}\raggedright
in scope but unadjudicable due to missing evidence, bridge, source, or visibility\strut
\end{minipage}\tabularnewline
\begin{minipage}[t]{0.26\columnwidth}\raggedright
\(\texttt{outside\_scope}\)\strut
\end{minipage} & \begin{minipage}[t]{0.68\columnwidth}\raggedright
instrument or host does not adjudicate claim\strut
\end{minipage}\tabularnewline
\begin{minipage}[t]{0.26\columnwidth}\raggedright
\(\texttt{absent\_with\_record}\)\strut
\end{minipage} & \begin{minipage}[t]{0.68\columnwidth}\raggedright
required evidence absent and absence is audited\strut
\end{minipage}\tabularnewline
\begin{minipage}[t]{0.26\columnwidth}\raggedright
\(\texttt{below\_threshold}\)\strut
\end{minipage} & \begin{minipage}[t]{0.68\columnwidth}\raggedright
evidence exists but threshold fails\strut
\end{minipage}\tabularnewline
\begin{minipage}[t]{0.26\columnwidth}\raggedright
\(\texttt{failed\_audit}\)\strut
\end{minipage} & \begin{minipage}[t]{0.68\columnwidth}\raggedright
required audit fails\strut
\end{minipage}\tabularnewline
\begin{minipage}[t]{0.26\columnwidth}\raggedright
\(\texttt{undefined\_circular}\)\strut
\end{minipage} & \begin{minipage}[t]{0.68\columnwidth}\raggedright
claim depends on an unstratified circular support/audit loop\strut
\end{minipage}\tabularnewline
\bottomrule
\end{longtable}

The full priority order of \(\mathrm{ClaimStatus}\) is fixed in Subsection 11.2.1.

\subsubsection{Gate Status}

\[
\mathrm{GateStatus}\;=\;\{\,\texttt{pass},\;\texttt{fail},\;\texttt{not\_required},\;\texttt{not\_checked},\;\texttt{outside\_scope}\,\}.
\]

The strictness gate refines \(\mathrm{GateStatus}\):

\[
\mathrm{StrictGateStatus}\;=\;\{\,\texttt{strict\_pass},\;\texttt{non\_strict\_pass},\;\texttt{not\_required},\;\texttt{not\_checked},\;\texttt{bad\_audit},\;\texttt{outside\_scope}\,\}.
\]

The core promotion gates are \(\mathsf{CoreGates}=\{G_{\mathrm{suff}},G_{\mathrm{desc}},G_{\mathrm{stab}},G_{\mathrm{ctrl}},G_{\mathrm{nosmuggle}},G_{\mathrm{vis}},G_{\mathrm{audit}}\}\), with the additional gate \(G_{\mathrm{strict}}\) and the optional gate \(G_{\mathrm{locglob}}\) (Subsection 10.2).

\subsubsection{Square Status}

\[
\mathrm{SqStatus}\;=\;\{\,\texttt{exact},\;\texttt{lax},\;\texttt{obstructed},\;\texttt{blocked},\;\texttt{outside\_scope},\;\texttt{failed\_audit},\;\texttt{undefined\_circular},\;\texttt{provisional}\,\}.
\]

\begin{longtable}[]{@{}ll@{}}
\toprule
\begin{minipage}[b]{0.28\columnwidth}\raggedright
value\strut
\end{minipage} & \begin{minipage}[b]{0.66\columnwidth}\raggedright
meaning\strut
\end{minipage}\tabularnewline
\midrule
\endhead
\begin{minipage}[t]{0.28\columnwidth}\raggedright
\(\texttt{exact}\)\strut
\end{minipage} & \begin{minipage}[t]{0.66\columnwidth}\raggedright
square defect is empty and exact comparison holds\strut
\end{minipage}\tabularnewline
\begin{minipage}[t]{0.28\columnwidth}\raggedright
\(\texttt{lax}\)\strut
\end{minipage} & \begin{minipage}[t]{0.66\columnwidth}\raggedright
ordered/approximate comparison holds (e.g.~\(\alpha F\le F^{\sharp}\alpha\))\strut
\end{minipage}\tabularnewline
\begin{minipage}[t]{0.28\columnwidth}\raggedright
\(\texttt{obstructed}\)\strut
\end{minipage} & \begin{minipage}[t]{0.66\columnwidth}\raggedright
square is in scope and audited but defect is nonempty\strut
\end{minipage}\tabularnewline
\begin{minipage}[t]{0.28\columnwidth}\raggedright
\(\texttt{blocked}\)\strut
\end{minipage} & \begin{minipage}[t]{0.66\columnwidth}\raggedright
square data missing, hidden, unbridged, or unaudited\strut
\end{minipage}\tabularnewline
\begin{minipage}[t]{0.28\columnwidth}\raggedright
\(\texttt{outside\_scope}\)\strut
\end{minipage} & \begin{minipage}[t]{0.66\columnwidth}\raggedright
host/instrument does not adjudicate square\strut
\end{minipage}\tabularnewline
\begin{minipage}[t]{0.28\columnwidth}\raggedright
\(\texttt{failed\_audit}\)\strut
\end{minipage} & \begin{minipage}[t]{0.66\columnwidth}\raggedright
square audit fails\strut
\end{minipage}\tabularnewline
\begin{minipage}[t]{0.28\columnwidth}\raggedright
\(\texttt{undefined\_circular}\)\strut
\end{minipage} & \begin{minipage}[t]{0.66\columnwidth}\raggedright
square support/audit is circular\strut
\end{minipage}\tabularnewline
\begin{minipage}[t]{0.28\columnwidth}\raggedright
\(\texttt{provisional}\)\strut
\end{minipage} & \begin{minipage}[t]{0.66\columnwidth}\raggedright
partial/model evidence supports provisional square\strut
\end{minipage}\tabularnewline
\bottomrule
\end{longtable}

The seven status families are typed-separately judgment-indexed and remain non-collapsing, in keeping with Theorem 5 of Section 6 and the status-discipline of Section 5.

\subsection{Full Defect Tables}

The general defect record schema and the per-primitive and promotion defect tables are recorded below.

\subsubsection{General Defect Record Schema}

A typed defect record under instrument \(I\) on a target package \(\mathcal T\) has the schema

\[
\mathsf{Def}_{i}^{\lambda}(\mathcal T,I)\;=\;(\,\delta_{i},\;\Omega_{i},\;\Theta_{i},\;\mathrm{polarity}_{i},\;W_{i},\;A_{i},\;V_{i},\;\sigma_{i}\,),
\]

with components:

\begin{longtable}[]{@{}ll@{}}
\toprule
\begin{minipage}[b]{0.21\columnwidth}\raggedright
field\strut
\end{minipage} & \begin{minipage}[b]{0.73\columnwidth}\raggedright
meaning\strut
\end{minipage}\tabularnewline
\midrule
\endhead
\begin{minipage}[t]{0.21\columnwidth}\raggedright
\(\delta_{i}\)\strut
\end{minipage} & \begin{minipage}[t]{0.73\columnwidth}\raggedright
defect value (set, function, scalar, or other finite typed value)\strut
\end{minipage}\tabularnewline
\begin{minipage}[t]{0.21\columnwidth}\raggedright
\(\Omega_{i}\)\strut
\end{minipage} & \begin{minipage}[t]{0.73\columnwidth}\raggedright
value object / defect type\strut
\end{minipage}\tabularnewline
\begin{minipage}[t]{0.21\columnwidth}\raggedright
\(\Theta_{i}\)\strut
\end{minipage} & \begin{minipage}[t]{0.73\columnwidth}\raggedright
threshold or pass condition\strut
\end{minipage}\tabularnewline
\begin{minipage}[t]{0.21\columnwidth}\raggedright
\(\mathrm{polarity}_{i}\)\strut
\end{minipage} & \begin{minipage}[t]{0.73\columnwidth}\raggedright
obstruction, certificate, gap, residual, capacity loss, or other polarity tag\strut
\end{minipage}\tabularnewline
\begin{minipage}[t]{0.21\columnwidth}\raggedright
\(W_{i}\)\strut
\end{minipage} & \begin{minipage}[t]{0.73\columnwidth}\raggedright
witness, certificate, or counterexample for \(\delta_{i}\)\strut
\end{minipage}\tabularnewline
\begin{minipage}[t]{0.21\columnwidth}\raggedright
\(A_{i}\)\strut
\end{minipage} & \begin{minipage}[t]{0.73\columnwidth}\raggedright
audit record for \(\delta_{i}\)\strut
\end{minipage}\tabularnewline
\begin{minipage}[t]{0.21\columnwidth}\raggedright
\(V_{i}\)\strut
\end{minipage} & \begin{minipage}[t]{0.73\columnwidth}\raggedright
visibility record for \(\delta_{i}\)\strut
\end{minipage}\tabularnewline
\begin{minipage}[t]{0.21\columnwidth}\raggedright
\(\sigma_{i}\)\strut
\end{minipage} & \begin{minipage}[t]{0.73\columnwidth}\raggedright
induced status or status annotation\strut
\end{minipage}\tabularnewline
\bottomrule
\end{longtable}

The general threshold rule on a defect record is

\[
\mathsf{Eval}_{\Theta_{i}}(\delta_{i})\;\in\;\{\,\texttt{pass},\;\texttt{fail},\;\texttt{not\_checked},\;\texttt{outside\_scope}\,\}.
\]

\subsubsection{Threshold Types}

The admissible thresholds form the finite typed family

\[
\Theta\;\in\;\{\,\Theta_{=0},\;\Theta_{\neq0},\;\Theta_{\le\varepsilon},\;\Theta_{\ge\varepsilon},\;\Theta_{\in S},\;\Theta_{\notin S},\;\Theta_{\mathrm{pass}},\;\Theta_{\mathrm{source}},\;\Theta_{\mathrm{bridge}},\;\Theta_{\mathrm{audit}},\;\Theta_{\mathrm{vis}}\,\},
\]

with intended readings recorded in Subsection 5.2.

\subsubsection{Primitive Defect Families}

The per-primitive defect records are as follows. Each row is a finite typed instance of the general defect schema of Subsubsection A.6.1.

\begin{longtable}[]{@{}lll@{}}
\toprule
\begin{minipage}[b]{0.30\columnwidth}\raggedright
primitive\strut
\end{minipage} & \begin{minipage}[b]{0.30\columnwidth}\raggedright
defect record\strut
\end{minipage} & \begin{minipage}[b]{0.30\columnwidth}\raggedright
threshold / status use\strut
\end{minipage}\tabularnewline
\midrule
\endhead
\begin{minipage}[t]{0.30\columnwidth}\raggedright
\(P_1\)\strut
\end{minipage} & \begin{minipage}[t]{0.30\columnwidth}\raggedright
\(\delta_1^{\mathrm{split}}(q,F)=\{\,(x,x'):\ q(x)=q(x')\wedge qF(x)\neq qF(x')\,\}\)\strut
\end{minipage} & \begin{minipage}[t]{0.30\columnwidth}\raggedright
\(\Theta_1^{\mathrm{desc}}:\delta_1^{\mathrm{split}}=\varnothing\); nonempty records an active \(P_1\) obstruction\strut
\end{minipage}\tabularnewline
\begin{minipage}[t]{0.30\columnwidth}\raggedright
\(P_1\)\strut
\end{minipage} & \begin{minipage}[t]{0.30\columnwidth}\raggedright
\(\delta_1^{\mathrm{cand}}(q,F,F^\sharp)=\{\,x:\ F^\sharp(q_{\mathrm{im}}(x))\neq q_{\mathrm{im}}F(x)\,\}\)\strut
\end{minipage} & \begin{minipage}[t]{0.30\columnwidth}\raggedright
candidate descended-map defect\strut
\end{minipage}\tabularnewline
\begin{minipage}[t]{0.30\columnwidth}\raggedright
\(P_1\)\strut
\end{minipage} & \begin{minipage}[t]{0.30\columnwidth}\raggedright
\(\mathrm{CD}_{\tau}^{\mu}(\Pi)=I(X_t;\Pi(X_{t+\tau})\mid\Pi(X_t))\)\strut
\end{minipage} & \begin{minipage}[t]{0.30\columnwidth}\raggedright
\(\Theta_1^{\mathrm{closure}}:\mathrm{CD}_{\tau}^{\mu}(\Pi)=0\) for exact closure, or \(\mathrm{CD}_{\tau}^{\mu}(\Pi)\le\varepsilon\) for approximate closure\strut
\end{minipage}\tabularnewline
\begin{minipage}[t]{0.30\columnwidth}\raggedright
\(P_2\)\strut
\end{minipage} & \begin{minipage}[t]{0.30\columnwidth}\raggedright
\(\delta_2^{\mathrm{rep}}(s;f)=0\) if \(s\in\Sigma_f\), and \(1\) if \(s\notin\Sigma_f\)\strut
\end{minipage} & \begin{minipage}[t]{0.30\columnwidth}\raggedright
in full finite-lens mode, \(s\in\Sigma_f\) iff \(s=f^{-1}(B)\) for some \(B\subseteq\mathrm{im}(f)\); nonrepresentability is witnessed by \(z,z'\) with \(f(z)=f(z')\) and \(\mathbf 1_s(z)\neq\mathbf 1_s(z')\)\strut
\end{minipage}\tabularnewline
\begin{minipage}[t]{0.30\columnwidth}\raggedright
\(P_2\)\strut
\end{minipage} & \begin{minipage}[t]{0.30\columnwidth}\raggedright
\(\delta_2^{\mathrm{cycle}}(G,G')=\beta_1(G)-\beta_1(G')\)\strut
\end{minipage} & \begin{minipage}[t]{0.30\columnwidth}\raggedright
drive-suppression threshold \(\Theta_2^{\mathrm{forest}}:\beta_1(G')=0\)\strut
\end{minipage}\tabularnewline
\begin{minipage}[t]{0.30\columnwidth}\raggedright
\(P_3\)\strut
\end{minipage} & \begin{minipage}[t]{0.30\columnwidth}\raggedright
\(\delta_3^{\mathrm{comp}}(E_1,E_2)=\{\,c:\ E_1E_2(c)\neq E_2E_1(c)\,\}\)\strut
\end{minipage} & \begin{minipage}[t]{0.30\columnwidth}\raggedright
route mismatch is active when \(\delta_3^{\mathrm{comp}}\neq\varnothing\); route coherence holds when \(\delta_3^{\mathrm{comp}}=\varnothing\)\strut
\end{minipage}\tabularnewline
\begin{minipage}[t]{0.30\columnwidth}\raggedright
\(P_3\)\strut
\end{minipage} & \begin{minipage}[t]{0.30\columnwidth}\raggedright
\(\delta_3^{\mathrm{hol}}(\gamma)=\operatorname{Hol}(\gamma)-\mathrm{id}\)\strut
\end{minipage} & \begin{minipage}[t]{0.30\columnwidth}\raggedright
active when \(\operatorname{Hol}(\gamma)\neq\mathrm{id}\); no \(P6_{\mathrm{drive}}\) consequence follows without a drive bridge\strut
\end{minipage}\tabularnewline
\begin{minipage}[t]{0.30\columnwidth}\raggedright
\(P_4\)\strut
\end{minipage} & \begin{minipage}[t]{0.30\columnwidth}\raggedright
\(\delta_4^{\mathrm{ref}}(f_0,f_1,\rho)=\{\,z:\ f_0(z)\neq\rho f_1(z)\,\}\)\strut
\end{minipage} & \begin{minipage}[t]{0.30\columnwidth}\raggedright
refinement passes iff \(\delta_4^{\mathrm{ref}}=\varnothing\)\strut
\end{minipage}\tabularnewline
\begin{minipage}[t]{0.30\columnwidth}\raggedright
\(P_4\)\strut
\end{minipage} & \begin{minipage}[t]{0.30\columnwidth}\raggedright
\(\delta_4^{\mathrm{glue}}=\{\,x\in U\cap V:\mathrm{res}_{U,U\cap V}(x_U)\neq\mathrm{res}_{V,U\cap V}(x_V)\,\}\)\strut
\end{minipage} & \begin{minipage}[t]{0.30\columnwidth}\raggedright
gluing passes iff \(\delta_4^{\mathrm{glue}}=\varnothing\); failed gluing may induce \(\texttt{globally\_obstructed}\)\strut
\end{minipage}\tabularnewline
\begin{minipage}[t]{0.30\columnwidth}\raggedright
\(P_5\)\strut
\end{minipage} & \begin{minipage}[t]{0.30\columnwidth}\raggedright
\(\delta_5^{\mathrm{idem}}(E)=\{\,x:\ E^2(x)\neq E(x)\,\}\)\strut
\end{minipage} & \begin{minipage}[t]{0.30\columnwidth}\raggedright
\(\Theta_5^{\mathrm{idem}}:\delta_5^{\mathrm{idem}}=\varnothing\); under a fixed exact completion profile, idempotence yields \(P_5\leftarrow P_5:\texttt{trivial}\)\strut
\end{minipage}\tabularnewline
\begin{minipage}[t]{0.30\columnwidth}\raggedright
\(P_5\)\strut
\end{minipage} & \begin{minipage}[t]{0.30\columnwidth}\raggedright
\(\delta_5^{\mathrm{fix}}(o,E)=E(o)-o\), or set-valued \(\{\,o:\ E(o)\neq o\,\}\)\strut
\end{minipage} & \begin{minipage}[t]{0.30\columnwidth}\raggedright
fixed-point residual\strut
\end{minipage}\tabularnewline
\begin{minipage}[t]{0.30\columnwidth}\raggedright
\(P_5\)\strut
\end{minipage} & \begin{minipage}[t]{0.30\columnwidth}\raggedright
package collapse values such as \(\texttt{identity\_package}\), \(\texttt{total\_collapse}\), and \(\texttt{singleton\_collapse}\)\strut
\end{minipage} & \begin{minipage}[t]{0.30\columnwidth}\raggedright
may induce \(\texttt{collapsed}\) or \(\texttt{trivial}\) depending on profile\strut
\end{minipage}\tabularnewline
\begin{minipage}[t]{0.30\columnwidth}\raggedright
\(P_6\)\strut
\end{minipage} & \begin{minipage}[t]{0.30\columnwidth}\raggedright
\(\delta_6^{\mathrm{audit}}\in\{\texttt{audit\_passes},\texttt{missing\_audit},\texttt{failed\_audit},\texttt{source\_failure},\texttt{visibility\_failure},\texttt{threshold\_failure},\texttt{nonclaim\_failure},\texttt{circular\_audit}\}\)\strut
\end{minipage} & \begin{minipage}[t]{0.30\columnwidth}\raggedright
\(\Theta_6^{\mathrm{audit}}:\delta_6^{\mathrm{audit}}=\texttt{audit\_passes}\); circular audit induces \(\texttt{undefined\_circular}\)\strut
\end{minipage}\tabularnewline
\begin{minipage}[t]{0.30\columnwidth}\raggedright
\(P_6\)\strut
\end{minipage} & \begin{minipage}[t]{0.30\columnwidth}\raggedright
\(\delta_6^{\mathrm{source}}\in\{\texttt{committed},\texttt{dirty},\texttt{fallback\_forbidden},\texttt{unknown},\texttt{contradictory},\texttt{missing}\}\)\strut
\end{minipage} & \begin{minipage}[t]{0.30\columnwidth}\raggedright
real pair claims require \(\delta_6^{\mathrm{source}}=\texttt{committed}\)\strut
\end{minipage}\tabularnewline
\begin{minipage}[t]{0.30\columnwidth}\raggedright
\(P_6\)\strut
\end{minipage} & \begin{minipage}[t]{0.30\columnwidth}\raggedright
drive certificate \(\kappa_6^{\mathrm{drive}}\), kept separate from \(\delta_6^{\mathrm{audit}}\)\strut
\end{minipage} & \begin{minipage}[t]{0.30\columnwidth}\raggedright
cohomological drive uses \([a]\neq0\in H^1(G)\), equivalently \(\exists\gamma:\oint_{\gamma}a\neq0\); no-drive certificates include \([a]=0\), \(a=d\phi\), and \(H^1(G)=0\)\strut
\end{minipage}\tabularnewline
\bottomrule
\end{longtable}

\subsubsection{Promotion-Specific Defects}

The promotion bridge \(B_{j\to j+1}\) carries the promotion-specific defects recorded in Subsection 5.4, in addition to any per-primitive defect contributions of its fields.

\begin{longtable}[]{@{}ll@{}}
\toprule
\begin{minipage}[b]{0.47\columnwidth}\raggedright
defect\strut
\end{minipage} & \begin{minipage}[b]{0.47\columnwidth}\raggedright
threshold / status use\strut
\end{minipage}\tabularnewline
\midrule
\endhead
\begin{minipage}[t]{0.47\columnwidth}\raggedright
\(\Delta_{\mathrm{fact}}(\pi_j,\pi_{j+1})=\{\,(s,s'):\pi_j(s)=\pi_j(s')\wedge\pi_{j+1}(s)\neq\pi_{j+1}(s')\,\}\)\strut
\end{minipage} & \begin{minipage}[t]{0.47\columnwidth}\raggedright
\(\Theta_{\mathrm{strict}}:\Delta_{\mathrm{fact}}\neq\varnothing\) gives \(G_{\mathrm{strict}}=\texttt{strict\_pass}\); \(\Theta_{\mathrm{nonstrict}}:\Delta_{\mathrm{fact}}=\varnothing\) gives \(G_{\mathrm{strict}}=\texttt{non\_strict\_pass}\)\strut
\end{minipage}\tabularnewline
\begin{minipage}[t]{0.47\columnwidth}\raggedright
\(\delta_{\mathrm{nosmuggle}}(B)\)\strut
\end{minipage} & \begin{minipage}[t]{0.47\columnwidth}\raggedright
contains violations: target-layer data used to prove target novelty; suppressed content used without bridge; simulation used as theorem; local evidence used as global without gluing; \(P_3\)-witness used as \(P6_{\mathrm{drive}}\) without drive bridge; lossy summary used as exact data. Threshold \(\Theta_{\mathrm{nosmuggle}}:\delta_{\mathrm{nosmuggle}}=\varnothing\); failure induces \(\mathrm{Promote}(B):\texttt{failed\_no\_smuggling}\)\strut
\end{minipage}\tabularnewline
\begin{minipage}[t]{0.47\columnwidth}\raggedright
\(\delta_{\mathrm{stab}}=\{\,o:\ E(o)\neq o\,\}\), or its approximate residual analogue\strut
\end{minipage} & \begin{minipage}[t]{0.47\columnwidth}\raggedright
if required and failed, induces \(\mathrm{Promote}(B):\texttt{failed\_stability}\)\strut
\end{minipage}\tabularnewline
\begin{minipage}[t]{0.47\columnwidth}\raggedright
descent defect inside promotion\strut
\end{minipage} & \begin{minipage}[t]{0.47\columnwidth}\raggedright
if macro closure is claimed, the promotion requires \(G_{\mathrm{desc}}=\texttt{pass}\); failure induces \(\mathrm{Promote}(B):\texttt{failed\_descent}\); if no dynamics or closure claim is made, \(G_{\mathrm{desc}}=\texttt{not\_required}\)\strut
\end{minipage}\tabularnewline
\bottomrule
\end{longtable}

Each defect record above conforms to the general schema of Subsubsection A.6.1, with the polarity tag, threshold, and induced promotion status drawn from the schemas of Subsection 5.4 and the promotion classifier of Subsection 5.6.4.
 \section{Long Proofs}

This appendix records long-form proofs that are referenced from the main text. The proofs of Sections 6, 7, 8, 10, 11, 12, and 14 are written in full where feasible inside the main text; this appendix supplements them with the longer technical content of the closure-deficit theorem, the cohomological theorems of Section 8, the completion-pasting theorems of Sections 7 and 14, and the fixed-interface definability bound of Section 7.

Each subsection records the long-form proof under the same instrument-indexed claim form as its main-text statement, and adds technical lemmas only where the main-text proof would otherwise become unreadable. The full statements remain at their main-text locations; this appendix is a pointer-and-extension document, not a re-statement of theorems.

\subsection{Closure-Deficit Proof}

The closure-deficit theorem of Subsection 7.3 (Theorem 8) is the finite Markov/KL result on the host \(\mathsf H_{\mathrm{prob}}\). The main-text proof gives the decomposition; this subsection records the same calculation in long form and identifies the associated \(P_1\) closure-deficit record. It does not introduce a new theorem.

\emph{Statement (recall).} Fix a finite state space \(X\), a finite stochastic kernel on \(X\), a fixed distribution, a finite abstraction \(\Pi:X\to Y\), a time lag \(\tau\), the declared KL predictive closure loss \(\mathcal L_\tau\), and an admissible macro-kernel class \(\mathcal K_{\mathrm{adm}}\) containing the conditional macro kernel

\[
K^\ast(\cdot\mid y)=P(\Pi(X_{t+\tau})\mid\Pi(X_t)=y)
\]

on the support of \(\Pi(X_t)\). Then

\[
\min_{K\in\mathcal K_{\mathrm{adm}}}\mathcal L_\tau(K)
\;=\;
I\bigl(X_t;\Pi(X_{t+\tau})\mid\Pi(X_t)\bigr).
\]

\emph{Proof.} Write

\[
Y_t=\Pi(X_t),\qquad Y'=\Pi(X_{t+\tau}).
\]

For a macro kernel \(K:Y\to\Delta(Y)\), the declared loss is

\[
\mathcal L_\tau(K)
\;=\;
\mathbb E\!\left[
D_{\mathrm{KL}}\bigl(P(Y'\mid X_t)\,\|\,K(\cdot\mid Y_t)\bigr)
\right].
\]

Because the state space is finite, the expectation decomposes over the finite support of \(Y_t\). For any \(y\) with \(P(Y_t=y)>0\), expand the conditional loss:

\[
\mathbb E\!\left[
D_{\mathrm{KL}}\bigl(P(Y'\mid X_t)\,\|\,K(\cdot\mid y)\bigr)
\mid Y_t=y
\right].
\]

Insert and subtract the conditional macro law \(P(Y'\mid Y_t=y)\). The finite KL projection identity gives

\[
\begin{aligned}
&\mathbb E\!\left[
D_{\mathrm{KL}}\bigl(P(Y'\mid X_t)\,\|\,K(\cdot\mid y)\bigr)
\mid Y_t=y
\right] \\
&\quad =
\mathbb E\!\left[
D_{\mathrm{KL}}\bigl(P(Y'\mid X_t)\,\|\,P(Y'\mid Y_t=y)\bigr)
\mid Y_t=y
\right]
+D_{\mathrm{KL}}\bigl(P(Y'\mid Y_t=y)\,\|\,K(\cdot\mid y)\bigr).
\end{aligned}
\]

The first term is independent of \(K\). The second term is non-negative and is minimized, with value zero, by

\[
K^\ast(\cdot\mid y)=P(Y'\mid Y_t=y),
\]

which belongs to \(\mathcal K_{\mathrm{adm}}\) by hypothesis. Therefore the minimizing loss is the expectation of the first term:

\[
\min_{K\in\mathcal K_{\mathrm{adm}}}\mathcal L_\tau(K)
\;=\;
\mathbb E\!\left[
D_{\mathrm{KL}}\bigl(P(Y'\mid X_t)\,\|\,P(Y'\mid Y_t)\bigr)
\right].
\]

For finite variables this expectation is the conditional mutual information \(I(X_t;Y'\mid Y_t)\). Substituting \(Y_t=\Pi(X_t)\) and \(Y'=\Pi(X_{t+\tau})\) proves the equality.

Support-zero cases are ignored by restricting the calculation to the finite support of \(Y_t\), equivalently by the standard convention that zero-probability conditioning events do not contribute to the expectation. The finite-state hypothesis makes every sum finite and the minimizer \(K^\ast\) a finite macro kernel.

The associated \(P_1\) closure-deficit value is

\[
\mathrm{CD}_{\tau}^{\mu}(\Pi)
\;=\;
I\bigl(X_t;\Pi(X_{t+\tau})\mid\Pi(X_t)\bigr),
\]

with exact-closure threshold \(\Theta_1^{\mathrm{closure}}:\mathrm{CD}_{\tau}^{\mu}(\Pi)=0\) and approximate threshold \(\mathrm{CD}_{\tau}^{\mu}(\Pi)\le\varepsilon\), as in the defect schema of Subsection 5.3.1. Written as a defect record,

\[
\delta_{\mathrm{cl}}^{\mathrm{prob}}
\;=\;
(\,\mathrm{CD}_{\tau}^{\mu}(\Pi),\;\Omega_{\mathrm{cl}},\;\Theta_1^{\mathrm{closure}},\;\mathrm{polarity}_{\mathrm{gap}},\;W_{\mathrm{cl}},\;A_{\mathrm{prob}},\;V_{\mathrm{prob}},\;\sigma_{\mathrm{cl}}\,),
\]

where \(W_{\mathrm{cl}}\) is the finite collection of conditional-distribution entries used in the KL calculation, \(A_{\mathrm{prob}}\) is the host audit record, \(V_{\mathrm{prob}}\) is the visibility record, and \(\sigma_{\mathrm{cl}}\) is the status annotation assigned by the later classifier. This record is a supporting defect record, not an additional theorem statement.

The instrument-indexed verdict remains the main-text theorem line

\[
\Gamma\,;\;\mathsf H_{\mathrm{prob}}\,;\;I_{\mathrm{prob}}\;\vdash\;\textsc{FiniteMarkovClosureDeficit}\;:\;\texttt{accepted}.
\]

\subsection{Graph/Cohomology Proofs}

This subsection records long-form proofs of the four graph/cohomology theorems of Section 8: forest no-drive (Theorem 15), exact-form null-drive (Theorem 16), nonzero-affinity equivalence (Theorem 17), and gating-affinity suppression (Theorem 18). The main-text proofs in Section 8 are sufficient for the headline statements; this subsection records the long-form arguments and the associated \(P_2\) and \(\mathrm{P6}_{\mathrm{drive}}^{H^{1}}\) defect bookkeeping.

\subsubsection{Forest No-Drive (Theorem 15)}

\emph{Statement (recall).} Let \(G=(V,E)\) be a finite graph with first Betti number \(\beta_{1}(G)=0\), that is, \(G\) is a forest. Let \(C^{0}(G;k)\), \(C^{1}(G;k)\), and \(d:C^{0}(G;k)\to C^{1}(G;k)\) be the finite cochain complex over the declared coefficient field or free coefficient host \(k\). Then \(H^{1}(G;k)=0\).

\emph{Proof.} For a finite graph over the declared coefficient field, or over a declared free coefficient host with the corresponding free-rank interpretation, finite graph cohomology satisfies

\[
H^{1}(G;k)\cong\mathrm{Hom}(H_{1}(G;k),k).
\]

The first Betti number is the finite dimension, or free rank, of \(H_1(G;k)\):

\[
\dim H_1(G;k)=\beta_1(G).
\]

If \(\beta_1(G)=0\), then \(H_1(G;k)=0\). Therefore \(\mathrm{Hom}(H_1(G;k),k)=0\), and hence \(H^1(G;k)=0\). Equivalently, a forest has no nontrivial cycle space, so there is no nonzero cycle-supported first cohomology.

The associated no-drive defect record is

\[
\delta_{6}^{\mathrm{forest}}\;=\;(\,H^{1}(G;k),\;\Omega_{\mathrm{drive}},\;\Theta_{=0},\;\mathrm{polarity}_{\mathrm{no\text{-}drive}},\;W_{\mathrm{forest}},\;A_{\mathrm{graph}},\;V_{\mathrm{graph}},\;\sigma_{\mathrm{no\text{-}drive}}\,),
\]

with \(W_{\mathrm{forest}}\) the finite forest-witness data (the absence of cycles in \(G\)).

The instrument-indexed verdict is

\[
\Gamma\,;\;\mathsf H_{\mathrm{graph}}\,;\;I_{\mathrm{graph}}\;\vdash\;\textsc{ForestNoDrive}\;:\;\texttt{accepted}.\quad\square
\]

\subsubsection{Exact-Form Null-Drive (Theorem 16)}

\emph{Statement (recall).} Let \(a\in C^{1}(G;k)\) be a \(1\)-cochain with \(a=d\phi\) for some \(\phi\in C^{0}(G;k)\). Then for every cycle \(\gamma\in\mathsf{Cycles}(G)\), the cycle integral satisfies \(\oint_{\gamma}a=0\).

\emph{Proof.} By definition, the cycle integral is the pairing \(\oint_{\gamma}a=\langle a,\gamma\rangle=\sum_{e\in\gamma}\epsilon(e)\,a(e)\), where \(\epsilon(e)\in\{+1,-1\}\) records the orientation of edge \(e\) in the cycle \(\gamma\). Substituting \(a=d\phi\), we have \(a(e)=\phi(e^{+})-\phi(e^{-})\) for edge \(e\) with head \(e^{+}\) and tail \(e^{-}\). The pairing becomes a telescoping sum

\[
\oint_{\gamma}a\;=\;\sum_{e\in\gamma}\epsilon(e)\bigl(\phi(e^{+})-\phi(e^{-})\bigr)\;=\;0,
\]

because every vertex of \(\gamma\) is visited the same number of times as a head (with \(+1\) orientation) and as a tail (with \(-1\) orientation), so each vertex contribution cancels.

The associated no-drive defect record is

\[
\delta_{6}^{\mathrm{exact}}\;=\;(\,\{\oint_{\gamma}a:\gamma\in\mathsf{Cycles}(G)\},\;\Omega_{\mathrm{drive}},\;\Theta_{=0},\;\mathrm{polarity}_{\mathrm{no\text{-}drive}},\;W_{\mathrm{exact}}=\phi,\;A_{\mathrm{graph}},\;V_{\mathrm{graph}},\;\sigma_{\mathrm{no\text{-}drive}}\,).
\]

The instrument-indexed verdict is \(\Gamma;\mathsf H_{\mathrm{graph}};I_{\mathrm{graph}}\vdash\textsc{ExactFormNullDrive}:\texttt{accepted}\). \(\square\)

\subsubsection{Nonzero-Affinity Equivalence (Theorem 17)}

\emph{Statement (recall).} For \(a\in C^{1}(G;k)\), the cohomology class \([a]\neq 0\) in \(H^{1}(G;k)\) iff there exists a cycle \(\gamma\in\mathsf{Cycles}(G)\) with \(\oint_{\gamma}a\neq 0\).

\emph{Proof.} Forward direction. Suppose \([a]\neq 0\). Then \(a\notin\operatorname{im} d^{0}\), so \(a\) is not exact. By the universal coefficient theorem on the finite truncated complex, the pairing \(H^{1}(G;k)\times H_{1}(G;k)\to k\) is non-degenerate. Since \([a]\neq 0\), there exists a homology class \([\gamma]\in H_{1}(G;k)\) with \(\langle[a],[\gamma]\rangle\neq 0\), that is, \(\oint_{\gamma}a\neq 0\).

Reverse direction. Suppose there exists \(\gamma\in\mathsf{Cycles}(G)\) with \(\oint_{\gamma}a\neq 0\). By the exact-form null-drive theorem (Subsubsection B.2.2), an exact \(1\)-cochain has zero cycle integral on every cycle. Hence \(a\) is not exact, so \([a]\neq 0\) in \(H^{1}(G;k)\).

The associated drive certificate record is

\[
\kappa_{6}^{\mathrm{drive}}(a)\;=\;(\,[a]\neq 0,\;\Omega_{\mathrm{drive}},\;\Theta_{\neq0},\;\mathrm{polarity}_{\mathrm{drive}},\;W=\gamma,\;A_{\mathrm{graph}},\;V_{\mathrm{graph}},\;\sigma_{\mathrm{drive}}\,).
\]

The instrument-indexed verdict is \(\Gamma;\mathsf H_{\mathrm{graph}};I_{\mathrm{graph}}\vdash\textsc{NonzeroAffinityEquivalence}:\texttt{accepted}\). \(\square\)

\subsubsection{Gating-Affinity Suppression (Theorem 18)}

\emph{Statement (recall).} Let \(G\rightsquigarrow G'=(V',E')\) be a \(P_{2}\)-gating that sends \(G\) to a forest, that is, \(\beta_{1}(G')=0\). Then \(H^{1}(G';k)=0\), and the cycle-supported cohomological drive certificate \(\mathrm{P6}_{\mathrm{drive}}^{H^1}\) is absent-with-record on \(G'\).

\emph{Proof.} By the forest no-drive theorem (Subsubsection B.2.1), \(\beta_{1}(G')=0\) implies \(H^{1}(G';k)=0\). Hence every cohomology class on \(G'\) is zero; in particular, the restriction \([a|_{G'}]\) is the zero class in \(H^{1}(G';k)\), so the drive certificate \(\kappa_{6}^{\mathrm{drive}}(a)\) does not survive the gating \(G\rightsquigarrow G'\).

The cycle-rank loss defect \(\delta_{2}^{\mathrm{cycle}}(G,G')=\beta_{1}(G)-\beta_{1}(G')\) records the number of cycles destroyed by the gating; the forest threshold \(\Theta_{2}^{\mathrm{forest}}:\beta_{1}(G')=0\) records the gating-to-forest condition.

The instrument-indexed verdict is \(\Gamma;\mathsf H_{\mathrm{graph}};I_{\mathrm{graph}}\vdash\textsc{GatingAffinitySuppression}:\texttt{accepted}\). \(\square\)

The four graph/cohomology theorems together record the no-drive boundary of the graph/cohomology realization of Subsection 13.4: forest hosts have no cycle-supported cohomological drive, exact \(1\)-cochains have no cycle drive, nonzero cohomology is equivalent to a nonzero cycle integral, and gating to a forest kills cycle-supported drive.

\subsection{Completion and Pasting Proofs}

This subsection records long-form proofs of the noncommuting-completions theorem (Theorem 9) and the completion-pasting recovery theorem (Theorem 31) in their finite double-category form. The main-text proofs are in Subsections 7.4 and 14.3; this subsection records the long-form arguments together with the explicit pasting-defect record.

\subsubsection{Noncommuting Completions (Theorem 9)}

\emph{Statement (recall).} There exist finite idempotent completions \(E_{1},E_{2}:C\to C\) on a finite poset \((C,\le)\) that fail to commute, that is, \(E_{1}E_{2}\neq E_{2}E_{1}\).

\emph{Proof.} It suffices to exhibit the finite witness used in the main text. Let

\[
U=\{a,b,c\},\qquad C=\mathcal P(U).
\]

Define \(E_{1},E_{2}:C\to C\) by

\[
E_1(S)=
\begin{cases}
S\cup\{b\}, & a\in S,\\
S, & a\notin S,
\end{cases}
\qquad
E_2(S)=
\begin{cases}
S\cup\{c\}, & b\in S,\\
S, & b\notin S.
\end{cases}
\]

Both maps are finite completion endomaps of \(C\). They are idempotent. For \(E_1\), if \(a\in S\), then applying \(E_1\) once adds \(b\) and applying it again adds no further element:

\[
E_1(E_1(S))=(S\cup\{b\})\cup\{b\}=S\cup\{b\}=E_1(S).
\]

If \(a\notin S\), then \(E_1(S)=S\), so \(E_1(E_1(S))=E_1(S)\). Thus \(E_1^2=E_1\). The same argument gives \(E_2^2=E_2\): if \(b\in S\), one application adds \(c\) and a second application adds no further element; if \(b\notin S\), the set is fixed.

Now take the witness set \(S_0=\{a\}\). Since \(b\notin S_0\), \(E_2(S_0)=S_0\), and therefore

\[
E_1E_2(S_0)=E_1(\{a\})=\{a,b\}.
\]

On the other route, \(E_1(S_0)=\{a,b\}\), and since \(b\in\{a,b\}\),

\[
E_2E_1(S_0)=E_2(\{a,b\})=\{a,b,c\}.
\]

Hence \(E_1E_2(S_0)\neq E_2E_1(S_0)\), so \(E_1E_2\neq E_2E_1\).

The completion-pasting defect record is

\[
\Delta_3^{\mathrm{comp}}\bigl(E_{1},E_{2}\bigr)\;=\;\{\,S\in C:E_1E_2(S)\neq E_2E_1(S)\,\},
\]

a nonempty finite subset of the carrier on which the two completions disagree; in particular \(S_0=\{a\}\) belongs to it. Theorem 31 of Subsection 14.3 records that this nonempty defect is recovered as a \(P_{3}\)-pasting defect of the calculus.

The instrument-indexed verdict is

\[
\Gamma\,;\;\mathsf H_{\mathrm{fin}}\,;\;I_{\mathrm{fin}}\;\vdash\;\textsc{NoncommutingCompletions}\;:\;\texttt{accepted}.\quad\square
\]

\subsubsection{Completion-Pasting Recovery (Theorem 31)}

\emph{Statement (recall).} In the finite decorated double-category fragment of Subsection 14.1, the decorated square \(\Xi_{\mathrm{comp}}\) for a pair of finite completions \(E_{1},E_{2}:C\to C\) is exact iff \(E_{1}E_{2}=E_{2}E_{1}\) on \(C\). The square is obstructed iff there exists \(c\in C\) with \(E_{1}E_{2}(c)\neq E_{2}E_{1}(c)\), and the noncommutation \(E_{1}E_{2}\neq E_{2}E_{1}\) is recovered as a \(P_{3}\)-pasting defect of the calculus.

\emph{Proof.} The decorated square \(\Xi_{\mathrm{comp}}\) has both horizontal edges equal to \(E_{2}\) and both vertical edges equal to \(E_{1}\). Exactness in the sense of Subsubsection 14.1.2 is the commutation equality of the boundary composition: the upper-right path \(E_{1}\circ E_{2}:C\to C\) equals the lower-left path \(E_{2}\circ E_{1}:C\to C\). By definition of composition on the finite poset \(C\), this is the equation \(E_{1}E_{2}=E_{2}E_{1}\) on \(C\).

Forward direction: \(\Xi_{\mathrm{comp}}\) exact \(\Rightarrow\) \(E_{1}E_{2}=E_{2}E_{1}\) pointwise on \(C\), since exactness is the pointwise commutation. Reverse direction: \(E_{1}E_{2}=E_{2}E_{1}\) on \(C\) \(\Rightarrow\) the boundary composition commutes pointwise, so the square is exact.

For the obstruction direction, the existence of \(c_{0}\in C\) with \(E_{1}E_{2}(c_{0})\neq E_{2}E_{1}(c_{0})\) produces a nonempty finite subset

\[
\delta_{\mathrm{comp}}(E_{1},E_{2})\;=\;\{\,c\in C:E_{1}E_{2}(c)\neq E_{2}E_{1}(c)\,\}\;\subseteq\;C,
\]

which is the per-square defect record \(\delta_{\Xi}\). A nonempty \(\delta_{\Xi}\) obstructs exactness of the square.

The recovery as \(P_{3}\)-pasting defect: the route-mismatch primitive \(P_{3}\) records typed route and pasting mismatch data (Subsection 2.2.2 and Table A.2). The two routes here are \(E_{1}\circ E_{2}\) and \(E_{2}\circ E_{1}\), differing on \(\delta_{\mathrm{comp}}(E_{1},E_{2})\), and the corresponding \(P_{3}\)-pasting defect is the per-element witness collection of this disagreement.

The instrument-indexed verdict is

\[
\Gamma\,;\;\mathsf H_{\mathrm{DC}}\,;\;I_{\mathrm{DC}}\;\vdash\;\textsc{CompletionPastingRecovery}\;:\;\texttt{accepted}.\quad\square
\]

The two theorems together record the completion-noncommutation phenomenon at the level of the finite double-category fragment: noncommuting completions exist (Theorem 9), and their failure to commute is recovered as a finite typed \(P_{3}\)-pasting defect (Theorem 31).

\subsection{Definability Bound Proof}

This subsection records the long-form proof of the fixed-interface definability bound (Theorem 14) of Section 7. The main-text proof is in Subsection 7.9; this subsection records the explicit counting argument and the definability set on a fixed visible interface.

\emph{Statement (recall).} Let \(f:X\to Y\) be a finite map on \(\mathsf H_{\mathrm{fin}}\). The full finite-lens definability of \(f\) is

\[
\mathrm{Def}(f)\;=\;\{\,f^{-1}(B):B\subseteq\mathrm{im}(f)\,\}.
\]

Then \(\mathrm{Def}(f)\) has cardinality

\[
|\mathrm{Def}(f)|\;=\;2^{|\mathrm{im}(f)|}.
\]

\emph{Proof.} The set \(\mathrm{Def}(f)\) is in bijection with the powerset of \(\mathrm{im}(f)\). Define

\[
\Phi:\mathcal P(\mathrm{im}(f))\longrightarrow\mathrm{Def}(f),
\qquad
\Phi(B)=f^{-1}(B).
\]

The map is well-defined because each \(B\subseteq\mathrm{im}(f)\) gives a subset \(f^{-1}(B)\subseteq X\), and by definition this subset lies in \(\mathrm{Def}(f)\).

\emph{Injectivity.} Suppose \(B,B'\subseteq\mathrm{im}(f)\) and \(B\neq B'\). Choose \(y\in B\triangle B'\). Without loss of generality, let \(y\in B\setminus B'\). Since \(y\in\mathrm{im}(f)\), there exists \(x\in X\) with \(f(x)=y\). Then \(x\in f^{-1}(B)\) but \(x\notin f^{-1}(B')\), so \(f^{-1}(B)\neq f^{-1}(B')\). Hence \(\Phi(B)\neq\Phi(B')\).

\emph{Surjectivity.} By definition, every element of \(\mathrm{Def}(f)\) is of the form \(f^{-1}(B)\) for some \(B\subseteq\mathrm{im}(f)\). Hence every definable predicate lies in the image of \(\Phi\).

The map \(\Phi\) is therefore a bijection, and

\[
|\mathrm{Def}(f)|
\;=\;
|\mathcal P(\mathrm{im}(f))|
\;=\;
2^{|\mathrm{im}(f)|}.
\]

The fixed-interface definability bound records that the number of package distinctions definable through the fixed lens \(f\) depends only on the finite interface \(\mathrm{im}(f)\). It does not count open-ended definability after changing the interface, the package, the instrument, or the host.

The instrument-indexed verdict is

\[
\Gamma\,;\;\mathsf H_{\mathrm{fin}}\,;\;I_{\mathrm{fin}}\;\vdash\;\textsc{FixedInterfaceDefinabilityBound}\;:\;\texttt{accepted}.\quad\square
\]

The bound is finite and decidable on the finite records. It supports the fixed-interface discipline used in Section 10: strict promotion is recorded as an interface change rather than as theory growth at the same fixed interface.
 \section{Full Countermodel Atlas}

This appendix records the full twelve-countermodel atlas. Sheets \(\mathsf{CM}_1\) through \(\mathsf{CM}_6\) summarize the constructions of Section 9 with their packet-grade dependency judgments and point back to the corresponding subsections of Section 9. Sheets \(\mathsf{CM}_7\) through \(\mathsf{CM}_{12}\) supply the full constructions for the appendix-only countermodels referenced in Subsection 9.8.

\subsection{Atlas Convention}

Each sheet uses the form fixed in Subsection 9.1.1: construction, supported no-go (with dependency judgment), what it blocks, what it does not block, and required positive bridge or repair where applicable. The internal packet form \((\mathrm{id},H,\mathsf{Data},\mathsf{ActiveClaim},\mathsf{FailedClaim},W,A,V,\mathcal N)\) of Subsection 9.1.1 is the underlying representation of each sheet; the prose fields below unpack that representation for the reader.

Where a sheet repeats a Section 9 entry, the sheet is brief and points back to the corresponding main-text subsection. Where the sheet is appendix-only, the construction is given in full with the data of the calculus made explicit.

\subsection{\texorpdfstring{\(\mathsf{CM}_1\): Active Package but No Macro Closure}{C.2 \textbackslash mathsf\{CM\}\_1: Active Package but No Macro Closure}}

Construction, blocked overclaim, and required bridge are recorded in Subsection 9.2: a finite package \(q:\{a,b,c\}\to\{A,B\}\) with \(q(a)=q(b)=A\), \(q(c)=B\), together with \(F:X\to X\) given by \(F(a)=a\), \(F(b)=c\), \(F(c)=c\). The split-pair \((a,b)\) witnesses \(\delta_1^{\mathrm{split}}(q,F)\neq\varnothing\), so no descended update through \(q\) exists.

The dependency judgment is

\[
\Gamma\,;\;\mathsf H_{\mathrm{fin}}\,;\;I_{\mathrm{fin}}\;\vdash\;P_5(q)\;\mathrel{\mathsf{dep}_{\texttt{insufficient}}}\;P_1(qF=F^\sharp q),
\]

supporting the no-go \(P_5\not\Rightarrow P_1\text{ macro closure}\) (NC-7). The required positive bridge is the descent record \(\delta_1^{\mathrm{split}}(q,F)=\varnothing\) together with the audit and visibility data of the surrounding judgment, with Theorem 7 supplying the equivalence between the empty descent defect and the existence of \(F^\sharp\).

\subsection{\texorpdfstring{\(\mathsf{CM}_2\): Active Route / Holonomy but No Drive}{C.3 \textbackslash mathsf\{CM\}\_2: Active Route / Holonomy but No Drive}}

Construction and blocked overclaim are recorded in Subsection 9.3: a finite cycle graph \(G\) on \(\mathsf H_{\mathrm{graph}}\) (for instance the triangle \(0\to 1\to 2\to 0\)) with active \(P_3\)-route data \(\operatorname{Hol}(\gamma)\neq\mathrm{id}\) and zero drive cochain \(a=0\). Every cycle integral vanishes, so \([a]=0\in H^1(G)\) and no positive cycle-supported cohomological drive certificate is admissible on \(G\) for this \(a\).

The dependency judgment is

\[
\Gamma\,;\;\mathsf H_{\mathrm{graph}}\,;\;I_{\mathrm{graph}}\;\vdash\;P_3^{\mathrm{holonomy}}\;\mathrel{\mathsf{dep}_{\texttt{requires\_bridge}}}\;\mathrm{P6}_{\mathrm{drive}}^{H^1},
\]

supporting NC-12 (no route-mismatch-implies-drive). The required positive bridge is the drive bridge \(B_{3\to 6}^{\mathrm{drive}}\) in \(\mathsf{BridgePolicy}_I\) that connects the route/holonomy data to a specific cochain \(a\) with \([a]\neq 0\).

\subsection{\texorpdfstring{\(\mathsf{CM}_3\): Local Packages Fail Globally}{C.4 \textbackslash mathsf\{CM\}\_3: Local Packages Fail Globally}}

Construction and blocked overclaim are recorded in Subsection 9.4: a finite carrier \(Z=\{a,b,c\}\) with cover \(U=\{a,b\}\), \(V=\{b,c\}\) and overlap \(U\cap V=\{b\}\), with local packages \(q_U,q_V\) that assign incompatible values \(0,1\) on the overlap \(\{b\}\). The gluing defect \(\delta_4^{\mathrm{glue}}\) is non-empty, with \(b\) as a witness, and no global package extends both local data.

The dependency judgment is

\[
\Gamma\,;\;\mathsf H_{\mathrm{fin}}\,;\;I_{\mathrm{fin}}\;\vdash\;\bigl\{P_5(q_U),\,P_5(q_V)\bigr\}\;\mathrel{\mathsf{dep}_{\texttt{insufficient}}}\;P_5(q_{\mathrm{global}}),
\]

supporting NC-14 (no local-implies-global). The required positive bridge is a \(P_4\)-gluing bridge \(B_{\mathrm{glue}}\) with \(\delta_4^{\mathrm{glue}}=\varnothing\) together with a \(P_6\)-audit compatibility record on the overlaps.

\subsection{\texorpdfstring{\(\mathsf{CM}_4\): Directed Cell Active but Pair Observable Blocked}{C.5 \textbackslash mathsf\{CM\}\_4: Directed Cell Active but Pair Observable Blocked}}

Construction and blocked overclaim are recorded in Subsection 9.5: a directed cell \(\mathsf{Cell}_{63}^{\lambda_{\mathrm{audit}}}\) with \(W_3=\mathsf{RouteMismatchWit}\), \(U_6=\mathsf{AddAuditRecord}\) and audit-update profile, classified \(\texttt{action}\) by the cell classifier of Subsection 5.6.2; together with a pair-observable record \(\mathsf{PairObs}_{\{3,6\}}^{\mu_{\mathrm{drive}}}\) on the same primitives with profile \(\mu_{\mathrm{drive}}\) asserting drive content, but with no admissible drive bridge and no admissible drive source-of-truth, classified \(\texttt{blocked}\) by the pair classifier of Subsection 5.6.3.

The dependency judgment is

\[
\Gamma\,;\;\mathbf{BirdInt}^{\mathrm{aud}}_{\mathrm{fin}}\,;\;I_{\mathrm{fin}}\;\vdash\;\mathsf{Cell}_{63}^{\lambda_{\mathrm{audit}}}:\texttt{action}\;\mathrel{\mathsf{dep}_{\texttt{insufficient}}}\;\mathsf{PairObs}_{\{3,6\}}^{\mu_{\mathrm{drive}}}:\texttt{real},
\]

supporting NC-5 (no cell-to-pair collapse). The required positive bridge is a real or fallback source under \(\mathsf{SourcePolicy}_I\) together with the drive bridge \(B_{3\to 6}^{\mathrm{drive}}\) when the pair claims drive content; Theorem 5 of Section 6 records the status-family separation in theorem-grade form.

\subsection{\texorpdfstring{\(\mathsf{CM}_5\): Refinement Worsens Closure}{C.6 \textbackslash mathsf\{CM\}\_5: Refinement Worsens Closure}}

Construction and blocked overclaim are recorded in Subsection 9.6: a finite carrier \(X=\{a,b,c,d\}\) with coarse abstraction \(q_0:X\to\{*\}\) and refinement \(q_1:X\to\{A,B,C\}\) given by \(q_1(a)=q_1(b)=A\), \(q_1(c)=B\), \(q_1(d)=C\), with \(q_0\) factoring through \(q_1\) via the forgetting map \(r:\{A,B,C\}\to\{*\}\). The finite update \(F(a)=c\), \(F(b)=d\), \(F(c)=c\), \(F(d)=d\) descends through \(q_0\) trivially but creates a \(q_1\)-split at \((a,b)\), so \(\delta_1^{\mathrm{split}}(q_1,F)\neq\varnothing\).

The dependency judgment is

\[
\Gamma\,;\;\mathsf H_{\mathrm{fin}}\,;\;I_{\mathrm{fin}}\;\vdash\;P_4(q_0\leadsto q_1)\;\mathrel{\mathsf{dep}_{\texttt{insufficient}}}\;\bigl(\delta_1^{\mathrm{split}}(q_1,F)\;\le\;\delta_1^{\mathrm{split}}(q_0,F)\bigr),
\]

supporting NC-13 (no refinement-implies-improvement). A positive refinement-improvement theorem requires a \(P_4\)-monotonicity bridge with explicit hypotheses on the refinement, the update, and the closure functional, recorded as future work in Subsection 16.3.1.

\subsection{\texorpdfstring{\(\mathsf{CM}_6\): Same-Level Self-Audit Circularity}{C.7 \textbackslash mathsf\{CM\}\_6: Same-Level Self-Audit Circularity}}

Construction and blocked overclaim are recorded in Subsection 9.7: a finite admissible instrument \(I_0\) under package \(\mathcal T_0\) asked to certify \(\mathrm{Sound}(I_0)\) at the same level. Two cases arise --- outside-scope when \(\mathrm{Sound}(I_0)\notin\mathsf{ClaimTypes}(I_0)\), and undefined-circular when the soundness claim is syntactically in \(\mathsf{ClaimTypes}(I_0)\) but its acceptance under \(I_0\) depends on \(I_0\)'s own soundness verdict --- and in neither case is \(\mathrm{Sound}(I_0)\) accepted under \(I_0\) at the same level.

The dependency judgment is

\[
\Gamma\,;\;\mathsf H_{\mathrm{instr}}\,;\;I_{\mathrm{instr}}\;\vdash\;I_0\text{ audits }\mathcal T_0\;\mathrel{\mathsf{dep}_{\texttt{insufficient}}}\;\bigl(I_0\vdash\mathrm{Sound}(I_0):\texttt{accepted}\bigr),
\]

supporting NC-16 (no same-level complete self-certification). The required positive repair is a level shift via a higher instrument \(I_1\) that audits \(I_0\) through an admissible audit bridge \(B^{\mathrm{audit}}_{0\to 1}\); Theorem 24 of Section 11 records the same-level self-audit failure at theorem-grade, and Theorem 25 records the rotating-audit theorem with NC-17 (no rotating-audit finality).

\subsection{\texorpdfstring{\(\mathsf{CM}_7\): Same Primitive Pair, Different Statuses}{C.8 \textbackslash mathsf\{CM\}\_7: Same Primitive Pair, Different Statuses}}

\subsubsection{Construction}

Take the primitive pair \(P_6\leftarrow P_3\) in \(\mathbf{BirdInt}^{\mathrm{aud}}_{\mathrm{fin}}\).

\textbf{Cell 1 --- audit-update profile.} Take witness \(W_3=\mathsf{RouteMismatchWit}\in\mathsf{Wit}(P_3)\) and update \(U_6=\mathsf{AddAuditRecord}\in\mathsf{Upd}(P_6)\), with the witness visible and audited under the surrounding instrument \(I\) and a profile \(\lambda_{\mathrm{audit}}\) that records audit-update mode. The cell

\[
\mathsf{Cell}_{63}^{\lambda_{\mathrm{audit}}}\;=\;(P_6\leftarrow P_3,\,W_3,\,U_6,\,L_{\mathrm{audit}},\,V_{\mathrm{audit}},\,\Theta_{\mathrm{audit}},\,A_{\mathrm{audit}},\,\delta=\varnothing)
\]

is well-formed, and the classifier of Subsection 5.6.2 assigns

\[
\mathsf{Cell}_{63}^{\lambda_{\mathrm{audit}}}\;:\;\texttt{action}.
\]

\textbf{Cell 2 --- drive-certification profile.} Use the same primitive pair \(P_6\leftarrow P_3\), but with update \(U_6=\mathsf{CertifyDrive}\in\mathsf{Upd}(P_6)\) and a profile \(\lambda_{\mathrm{drive}}\) that requires a drive bridge. Suppose the drive bridge \(B_{3\to 6}^{\mathrm{drive}}\) is absent from \(\mathcal L\) and the cell's bridge data fails \(\mathsf{BridgePolicy}_I\). The cell record

\[
\mathsf{Cell}_{63}^{\lambda_{\mathrm{drive}}}
\]

fails the bridge check; its classifier verdict is

\[
\mathsf{Cell}_{63}^{\lambda_{\mathrm{drive}}}\;:\;\texttt{blocked}.
\]

The same primitive pair \((P_6,P_3)\) thus carries two cell records with different profile, update, and bridge requirements, and these records receive distinct cell statuses.

\subsubsection{Supported No-Go}

The construction supports the schema-grade observation that cell status does not factor through the primitive pair: there is no map \(g:\mathbb P\times\mathbb P\to\mathrm{CellStatus}\) recovering both verdicts from the same pair. Equivalently, no total operation \(*:\mathbb P\times\mathbb P\to\mathbb P\) together with a decoder \(d:\mathbb P\to\mathrm{CellStatus}\) can faithfully encode typed, statused, profile-relative interaction. The construction is the working content of NC-1 and NC-3 of Subsection 2.4.3, and Theorem 1 of Section 6 makes the non-factorization formal.

\subsubsection{What It Blocks}

The overclaim ``the \(6\times 6\) primitive-pair table is a universal interaction algebra'' is rejected. NC-1 and NC-3 are the corresponding nonclaims.

\subsubsection{What It Does Not Block}

The construction does not say cell records on \((P_6,P_3)\) are unattainable. Both are finite covered records for the directed-cell classifier; what fails is the inference from the primitive pair alone to a cell status.

\subsubsection{Required Positive Bridge}

A positive cell-status claim on \((P_6,P_3)\) requires the surrounding profile, witness, update, threshold, audit, visibility, and bridge data of the cell record; the primitive pair alone is not sufficient input.

\subsection{\texorpdfstring{\(\mathsf{CM}_8\): Idempotent Completions Do Not Commute}{C.9 \textbackslash mathsf\{CM\}\_8: Idempotent Completions Do Not Commute}}

\subsubsection{Construction}

Take a finite three-element set \(U=\{a,b,c\}\) and let \(C=\mathcal P(U)\) be its finite power set on \(\mathsf H_{\mathrm{fin}}\). Define two finite endomaps \(E_1,E_2:C\to C\) by

\[
E_1(S)\;=\;\begin{cases}\,S\cup\{b\}, & a\in S,\\[2pt] S, & a\notin S,\end{cases}\qquad E_2(S)\;=\;\begin{cases}\,S\cup\{c\}, & b\in S,\\[2pt] S, & b\notin S.\end{cases}
\]

\textbf{Idempotence.} For \(E_1\): if \(a\in S\), then \(E_1(S)=S\cup\{b\}\) still contains \(a\), so \(E_1(E_1(S))=(S\cup\{b\})\cup\{b\}=S\cup\{b\}=E_1(S)\). If \(a\notin S\), then \(E_1(S)=S\), hence \(E_1(E_1(S))=E_1(S)\). So \(E_1^2=E_1\). The argument for \(E_2\) is symmetric: \(E_2^2=E_2\).

\textbf{Noncommutation.} Take \(S_0=\{a\}\). Then

\[
E_1\,E_2(\{a\})\;=\;E_1(\{a\})\;=\;\{a,b\},
\]

since \(E_2\) leaves \(\{a\}\) unchanged (\(b\notin\{a\}\)). On the other hand,

\[
E_2\,E_1(\{a\})\;=\;E_2(\{a,b\})\;=\;\{a,b,c\},
\]

since \(E_1\) adds \(b\) (because \(a\in\{a\}\)) and then \(E_2\) adds \(c\) (because \(b\in\{a,b\}\)). Hence \(E_1E_2(\{a\})\neq E_2E_1(\{a\})\), and \(\Delta_3^{\mathrm{comp}}(E_1,E_2)\neq\varnothing\) by Theorem 10.

\subsubsection{Supported No-Go}

The construction supports the dependency judgment

\[
\Gamma\,;\;\mathsf H_{\mathrm{fin}}\,;\;I_{\mathrm{fin}}\;\vdash\;\bigl(E_1^2=E_1\,\wedge\,E_2^2=E_2\bigr)\;\mathrel{\mathsf{dep}_{\texttt{insufficient}}}\;\bigl(E_1E_2=E_2E_1\bigr),
\]

recording that individual idempotence does not imply commutation, equivalently that \(P_5\)-individual exactness does not imply \(P_3\)-route coherence.

\subsubsection{What It Blocks}

The overclaim ``valid completions always commute'' is rejected. NC-8 of Subsection 2.4.3 records the corresponding nonclaim, and Theorem 9 of Section 7 makes the noncommuting-completions result theorem-grade.

\subsubsection{What It Does Not Block}

The construction does not say completions are always non-commuting, or that route coherence is unattainable. Specific pairs of idempotent completions on a finite carrier may commute --- for instance, when the two maps are identical or when one is the identity and commutes with the other --- but commutation is not implied by individual idempotence.

\subsubsection{Required Positive Bridge}

A positive commutation claim \(E_1E_2=E_2E_1\) requires an explicit route-coherence record: the empty defect \(\Delta_3^{\mathrm{comp}}(E_1,E_2)=\varnothing\) together with the audit and visibility data of the surrounding judgment. Theorem 10 makes the equivalence between the empty defect and commutation precise.

\subsection{\texorpdfstring{\(\mathsf{CM}_9\): Strict Extension Without Macro Closure}{C.10 \textbackslash mathsf\{CM\}\_9: Strict Extension Without Macro Closure}}

\subsubsection{Construction}

Take a finite carrier \(S=\{a,b,c\}\) on \(\mathsf H_{\mathrm{fin}}\). Define an old object map \(\pi_0:S\to O_0=\{*\}\) collapsing every element to a single value:

\[
\pi_0(a)\;=\;\pi_0(b)\;=\;\pi_0(c)\;=\;*.
\]

Define a new object map \(\pi_1:S\to O_1=\{U,V\}\) with

\[
\pi_1(a)\;=\;U,\qquad \pi_1(b)\;=\;U,\qquad \pi_1(c)\;=\;V.
\]

Since \(\pi_0\) collapses everything but \(\pi_1\) separates \(c\), there exist \(s,s'\in S\) --- for example \(s=a\) and \(s'=c\) --- with \(\pi_0(s)=\pi_0(s')=*\) but \(\pi_1(s)=U\neq V=\pi_1(s')\). By Theorem 12, \(\pi_1\not\factor\pi_0\), so the bridge is strict at the object-map level: the strictness gate \(G_{\mathrm{strict}}\) records \(\texttt{strict\_pass}\) from \(\Delta_{\mathrm{fact}}(\pi_0,\pi_1)\neq\varnothing\).

Now define a finite update \(F:S\to S\) by

\[
F(a)\;=\;a,\qquad F(b)\;=\;c,\qquad F(c)\;=\;c.
\]

Then \(\pi_1(a)=\pi_1(b)=U\), but \(\pi_1F(a)=\pi_1(a)=U\) and \(\pi_1F(b)=\pi_1(c)=V\), so \(\pi_1F(a)\neq\pi_1F(b)\). The descent defect \(\delta_1^{\mathrm{split}}(\pi_1,F)\) is non-empty, with \((a,b)\) as a witness, and by Theorem 7 no descended dynamics exists for \(F\) through \(\pi_1\). The bridge is strict, but macro closure on the target package fails.

\subsubsection{Supported No-Go}

The construction supports the dependency judgment

\[
\Gamma\,;\;\mathsf H_{\mathrm{fin}}\,;\;I_{\mathrm{prom}}\;\vdash\;\mathrm{Promote}(B):\texttt{strict}\;\mathrel{\mathsf{dep}_{\texttt{insufficient}}}\;\bigl(\pi_1F=F^\sharp\pi_1\bigr),
\]

recording that strict promotion does not by itself imply macro closure on the target.

\subsubsection{What It Blocks}

The overclaim ``strict extension automatically gives closed dynamics'' is rejected. NC-10 of Subsection 2.4.3 records the corresponding nonclaim.

\subsubsection{What It Does Not Block}

The construction does not say strict promotion is incompatible with macro closure. A different finite update \(F'\) --- one whose action is consistent with the fibers of \(\pi_1\) --- would descend through \(\pi_1\) and produce an admissible descended dynamics on the target package; the construction shows only that descent is independent data from strictness.

\subsubsection{Required Positive Bridge}

A positive macro-closure claim on a strict promotion requires an explicit descent record: the empty defect \(\delta_1^{\mathrm{split}}(\pi_1,F)=\varnothing\) recorded by the descent gate \(G_{\mathrm{desc}}\) of Subsection 5.4.4, together with the other audit, visibility, and required gate data of the bridge.

\subsection{\texorpdfstring{\(\mathsf{CM}_{10}\): Strict Extension Without Drive}{C.11 \textbackslash mathsf\{CM\}\_\{10\}: Strict Extension Without Drive}}

\subsubsection{Construction}

Use the same strict object maps from \(\mathsf{CM}_9\): \(S=\{a,b,c\}\), \(\pi_0:S\to\{*\}\) collapsing all elements, and \(\pi_1:S\to\{U,V\}\) separating \(c\). By Theorem 12, \(\pi_1\not\factor\pi_0\), so the bridge is strict at the object-map level.

Now place the model on a finite support graph \(G\) that is a forest, for example the path

\[
a\;\text{---}\;b\;\text{---}\;c
\]

with vertex set \(V(G)=\{a,b,c\}\) and edge set \(E(G)=\{(a,b),(b,c)\}\). Then \(|E(G)|=2\), \(|V(G)|=3\), \(c(G)=1\), so \(\beta_1(G)=2-3+1=0\). By Theorem 15,

\[
H^1(G)\;=\;0,
\]

and no nonzero cohomology class \([a]\in H^1(G)\) exists, so no admissible cycle-supported cohomological drive certificate \(\mathrm{P6}_{\mathrm{drive}}^{H^1}\) is available on \(G\). The bridge is strict at the object-map level but the drive face is absent on the support graph.

\subsubsection{Supported No-Go}

The construction supports the dependency judgment

\[
\Gamma\,;\;\mathsf H_{\mathrm{graph}}\,;\;I_{\mathrm{graph}}\;\vdash\;\mathrm{Promote}(B):\texttt{strict}\;\mathrel{\mathsf{dep}_{\texttt{insufficient}}}\;\mathrm{P6}_{\mathrm{drive}}^{H^1},
\]

recording that strict promotion does not by itself imply a cycle-supported cohomological drive certificate.

\subsubsection{What It Blocks}

The overclaim ``strictness implies drive'' is rejected. NC-11 of Subsection 2.4.3 records the corresponding nonclaim.

\subsubsection{What It Does Not Block}

The construction does not say strict promotion is incompatible with drive. A different support graph with \(\beta_1\ge 1\) --- for example, a triangle --- could carry a positive drive certificate; the construction shows only that the choice of support graph is independent data from the strictness of the promotion bridge.

\subsubsection{Required Positive Bridge}

A positive drive claim on a strict promotion requires an explicit drive bridge in \(\mathsf{BridgePolicy}_I\) together with a support graph \(G'\) on which \(H^1(G')\neq 0\) and an admissible cochain \(a\) with \([a]\neq 0\in H^1(G')\). Theorem 17 makes the equivalence between \([a]\neq 0\) and existence of a cycle with \(\oint_\gamma a\neq 0\) precise.

\subsection{\texorpdfstring{\(\mathsf{CM}_{11}\): Fallback Source Cannot Support Real Pair}{C.12 \textbackslash mathsf\{CM\}\_\{11\}: Fallback Source Cannot Support Real Pair}}

\subsubsection{Construction}

Take a pair-observable record \(Q=\mathsf{PairObs}_{\{i,j\}}^{\mu}\) on a primitive pair \(\{i,j\}\), with profile \(\mu\) that requires the source-of-truth record to have type \(\texttt{committed\_state}\) or \(\texttt{independent\_pair\_witness}\) --- that is, a real source under \(\mathsf{SourcePolicy}_I\) in the schema of Subsection 4.8.1.

Suppose the only available source record has type \(\texttt{fallback}\):

\[
\mathrm{src}\;=\;(\mathrm{id}_{\mathrm{src}},\;\texttt{fallback},\;\mathrm{trace}_{\mathrm{src}},\;\mathrm{valid}_{\mathrm{src}}),
\]

and that the profile \(\mu\) does not include a fallback-upgrade record that would license \(\texttt{fallback}\) as a real source for this pair. The pair classifier of Subsection 5.6.3 then assigns at most

\[
Q\;:\;\texttt{real\_provisional}
\]

(via the fallback-source case of the classifier), and never \(Q:\texttt{real}\). If \(\mu\) explicitly forbids fallback sources, the classifier instead routes \(Q\) to \(\texttt{blocked}\).

\subsubsection{Supported No-Go}

The construction supports the dependency judgment

\[
\Gamma\,;\;\mathsf H_{\mathrm{PICA}}\,;\;I_{\mathrm{PICA}}\;\vdash\;\mathrm{src}:\texttt{fallback}\;\mathrel{\mathsf{dep}_{\texttt{insufficient}}}\;Q:\texttt{real},
\]

recording that fallback evidence does not by itself support a real pair-observable verdict.

\subsubsection{What It Blocks}

The overclaim ``placeholder or fallback source counts as a real pair source-of-truth'' is rejected. The pair-observable source-of-truth discipline of Subsection 4.4.3 records the schema-level statement; this countermodel records the failure case at the dependency level.

\subsubsection{What It Does Not Block}

The construction does not say fallback sources are useless. A fallback source under a profile that admits fallback-upgrade can support \(\texttt{real\_provisional}\), which is a legitimate pair-observable verdict. The construction shows only that a fallback source does not by itself license the strongest verdict \(\texttt{real}\).

\subsubsection{Required Positive Bridge}

A positive \(Q:\texttt{real}\) verdict requires a real source under \(\mathsf{SourcePolicy}_I\) --- one of the five real-source classes from Subsection 4.8.1: \(\texttt{committed\_state}\), \(\texttt{audited\_cell\_records}\), \(\texttt{independent\_pair\_witness}\), \(\texttt{simulation\_trace}\), or \(\texttt{ablation\_record}\), admitted by \(\mu\). A fallback source can support \(\texttt{real}\) only when the profile explicitly upgrades fallback through an audited bridge or upgrade record admitted by the surrounding source and bridge policies; absent that upgrade, the fallback source supports at most \(\texttt{real\_provisional}\), or \(\texttt{blocked}\) if fallback is forbidden.

\subsection{\texorpdfstring{\(\mathsf{CM}_{12}\): Gating Destroys Cycle-Supported Drive}{C.13 \textbackslash mathsf\{CM\}\_\{12\}: Gating Destroys Cycle-Supported Drive}}

\subsubsection{Construction}

Take the triangle graph \(G\) on \(\mathsf H_{\mathrm{graph}}\) with vertex set \(V(G)=\{0,1,2\}\) and oriented edge set \(E(G)=\{(0,1),(1,2),(2,0)\}\). Then \(|E(G)|=3\), \(|V(G)|=3\), \(c(G)=1\), so

\[
\beta_1(G)\;=\;3\;-\;3\;+\;1\;=\;1,
\]

and the cycle space \(Z_1(G)\) is one-dimensional, spanned by the triangle cycle \(\gamma=(0,1)(1,2)(2,0)\).

Choose a finite \(1\)-cochain \(a\in C^1(G)\) with nonzero cycle integral, for example \(a((0,1))=1\), \(a((1,2))=0\), \(a((2,0))=0\) over a free coefficient host. Then

\[
\oint_\gamma a\;=\;a((0,1))\;+\;a((1,2))\;+\;a((2,0))\;=\;1\;\neq\;0,
\]

so \([a]\neq 0\in H^1(G)\) and \(\mathrm{P6}_{\mathrm{drive}}^{H^1}\) is active on \(G\) for this \(a\).

Now apply a \(P_2\)-style gate that deletes the edge \((2,0)\): the resulting graph \(G'\) has vertex set \(V(G')=\{0,1,2\}\) and edge set \(E(G')=\{(0,1),(1,2)\}\), that is, the path \(0\to 1\to 2\). Then \(|E(G')|=2\), \(|V(G')|=3\), \(c(G')=1\), so

\[
\beta_1(G')\;=\;2\;-\;3\;+\;1\;=\;0.
\]

By Theorem 15, \(H^1(G')=0\), and by Theorem 18 the cohomological drive certificate \(\mathrm{P6}_{\mathrm{drive}}^{H^1}\) is absent-with-record on \(G'\). The gate has destroyed the cycle-supported drive that was active on \(G\).

\subsubsection{Supported No-Go}

The construction supports the dependency judgment

\[
\Gamma\,;\;\mathsf H_{\mathrm{graph}}\,;\;I_{\mathrm{graph}}\;\vdash\;P_2(\text{edge deletion to forest})\;\mathrel{\mathsf{dep}_{\texttt{destroys}}}\;\mathrm{P6}_{\mathrm{drive}}^{H^1},
\]

recording that \(P_2\)-style gating that produces a forest destroys the cycle-supported cohomological drive certificate.

\subsubsection{What It Blocks}

The overclaim ``gating only removes irrelevant support'' is rejected. The gating-suppression discipline of Section 8 records the corresponding theorem (Theorem 18), and this sheet gives the finite countermodel instance: gating to a forest can destroy a previously active drive certificate.

\subsubsection{What It Does Not Block}

The construction does not say all gates destroy drive. A \(P_2\)-style gate whose output retains nontrivial cycles (\(\beta_1(G')\ge 1\)) may leave drive certificates intact; the construction shows only that gates producing forests are sufficient to suppress cycle-supported cohomological drive.

\subsubsection{Required Positive Repair}

A positive drive-preservation claim under gating requires that the output graph \(G'\) have \(\beta_1(G')\ge 1\) and that the cohomology class \([a]\) remain nonzero in \(H^1(G')\) --- equivalently, by Theorem 17, that some cycle \(\gamma'\in Z_1(G')\) have \(\oint_{\gamma'}a\neq 0\). Without these conditions, the drive certificate is absent on \(G'\).
 \section{Model-Realization Sheets}

This appendix records the four scoped model-realization sheets in long form: PICA (D.1), Cantor strict-bridge (D.2), abstract interpretation (D.3), and graph/cohomology (D.4). Subsection D.5 records the combined coverage matrix across the four realizations. The sheets supplement the main-text records of Section 13 and Section 12 with per-realization reference data drawn from the model-realization packet of the source material.

Each sheet declares the host, the realization object, the realization map, the realized fragment, the actor-update map and informant-witness map (where applicable), the per-judgment realizations, the source-of-truth and provenance discipline (where applicable), the realization theorem, and the realization-specific nonclaim register. Each sheet is host-bound and fragment-scoped: the main-text claim that the corresponding model realizes a declared fragment of \(\mathbf{BirdInt}^{\mathrm{aud}}_{\mathrm{fin}}\) is recorded in Section 13 (or Section 12 for AI), and this appendix does not introduce a new realization claim.

\subsection{PICA Sheet}

The PICA realization sheet \(\mathsf{PICAReal}_{\mathrm{fin}}\) records the realization of the finite stochastic cell/pair/provenance fragment of \(\mathbf{BirdInt}^{\mathrm{aud}}_{\mathrm{fin}}\) on the finite stochastic implementation host \(\mathsf H_{\mathrm{PICA}}\). The main-text record is in Subsection 13.2; this subsection records the long-form sheet.

\subsubsection{Host and Realization Object}

The host \(\mathsf H_{\mathrm{PICA}}\) has components

\[
\mathsf H_{\mathrm{PICA}}\;=\;(\,\Omega,\;K,\;\mathbb P,\;\mathsf{Bus},\;\mathsf{CellRules},\;\mathsf{WitnessStore},\;\mathsf{UpdateStore},\;\mathsf{PairStore},\;\mathsf{DepGraph},\;A,\;V\,),
\]

with \(\Omega\) a finite stochastic state space, \(K\) a finite Markov kernel on \(\Omega\), \(\mathbb P=\{P_{1},\ldots,P_{6}\}\), and the remaining fields finite stores. The realization object \(\mathcal P\) extends the host with \(\mathsf{StatusStore}\), \(\mathsf{Abl}\), and \(\mathcal N\):

\[
\mathcal P\;=\;(\,\Omega,\;K,\;\mathbb P,\;\mathsf{Bus},\;\mathsf{CellRules},\;\mathsf{WitnessStore},\;\mathsf{UpdateStore},\;\mathsf{StatusStore},\;\mathsf{PairStore},\;\mathsf{DepGraph},\;\mathsf{Abl},\;A,\;V,\;\mathcal N\,).
\]

The host instrument is \(I_{\mathrm{PICA}}\).

\subsubsection{Realization Map and Induced Package}

The realization map \(\mathcal R_{\mathrm{PICA}}:\mathsf{PICAReal}_{\mathrm{fin}}\to\mathbf{BirdInt}^{\mathrm{aud}}_{\mathrm{fin\text{-}stoch}}{}^{\mathrm{cell}/\mathrm{pair}/\mathrm{prov}}\) sends the realization object \(\mathcal P\) to the induced package

\[
\mathcal T_{\mathrm{PICA}}\;=\;(\,\Omega,\;f_{\mathrm{PICA}},\;\Sigma_{f_{\mathrm{PICA}}},\;E_{\mathrm{PICA}},\;\mathcal A_{\mathrm{PICA}}\,),
\]

with PICA-field interpretation: \(\Omega\) is the finite stochastic carrier, \(K\) provides finite Markov dynamics, the bus/lens projection provides \(f_{\mathrm{PICA}}\), the visible PICA claims provide \(\Sigma_{f_{\mathrm{PICA}}}\), the commit/update closure provides \(E_{\mathrm{PICA}}\), and the finite audit/provenance ledger \(A\) assembles the inherited package-level audit functional \(\mathcal A_{\mathrm{PICA}}\). The audit record \(A\) remains distinct from the audit functional \(\mathcal A_{\mathrm{PICA}}\).

\subsubsection{Actor-Update and Informant-Witness Maps}

The actor-update map sends each primitive to its PICA update sort, and the informant-witness map sends each primitive to its PICA witness sort:

\begin{longtable}[]{@{}lll@{}}
\toprule
\begin{minipage}[b]{0.08\columnwidth}\raggedright
primitive\strut
\end{minipage} & \begin{minipage}[b]{0.46\columnwidth}\raggedright
actor update sort\strut
\end{minipage} & \begin{minipage}[b]{0.37\columnwidth}\raggedright
informant witness sort\strut
\end{minipage}\tabularnewline
\midrule
\endhead
\begin{minipage}[t]{0.08\columnwidth}\raggedright
\(P_{1}\)\strut
\end{minipage} & \begin{minipage}[t]{0.46\columnwidth}\raggedright
descent repair / closure update\strut
\end{minipage} & \begin{minipage}[t]{0.37\columnwidth}\raggedright
descent/closure defect\strut
\end{minipage}\tabularnewline
\begin{minipage}[t]{0.08\columnwidth}\raggedright
\(P_{2}\)\strut
\end{minipage} & \begin{minipage}[t]{0.46\columnwidth}\raggedright
support gate / admissibility update\strut
\end{minipage} & \begin{minipage}[t]{0.37\columnwidth}\raggedright
gate/support/admissibility witness\strut
\end{minipage}\tabularnewline
\begin{minipage}[t]{0.08\columnwidth}\raggedright
\(P_{3}\)\strut
\end{minipage} & \begin{minipage}[t]{0.46\columnwidth}\raggedright
route-memory / mismatch update\strut
\end{minipage} & \begin{minipage}[t]{0.37\columnwidth}\raggedright
route mismatch / route-memory witness\strut
\end{minipage}\tabularnewline
\begin{minipage}[t]{0.08\columnwidth}\raggedright
\(P_{4}\)\strut
\end{minipage} & \begin{minipage}[t]{0.46\columnwidth}\raggedright
refinement / staging update\strut
\end{minipage} & \begin{minipage}[t]{0.37\columnwidth}\raggedright
refinement/stage witness\strut
\end{minipage}\tabularnewline
\begin{minipage}[t]{0.08\columnwidth}\raggedright
\(P_{5}\)\strut
\end{minipage} & \begin{minipage}[t]{0.46\columnwidth}\raggedright
package / completion / closure-summary update\strut
\end{minipage} & \begin{minipage}[t]{0.37\columnwidth}\raggedright
package/closure/completion witness\strut
\end{minipage}\tabularnewline
\begin{minipage}[t]{0.08\columnwidth}\raggedright
\(P_{6}\)\strut
\end{minipage} & \begin{minipage}[t]{0.46\columnwidth}\raggedright
audit / provenance / source update\strut
\end{minipage} & \begin{minipage}[t]{0.37\columnwidth}\raggedright
audit/provenance/source witness\strut
\end{minipage}\tabularnewline
\bottomrule
\end{longtable}

\subsubsection{Directed-Cell and Status Recovery}

Each PICA cell is realized as \(\Gamma;\mathcal T_{\mathrm{PICA}};I_{\mathrm{PICA}}\vdash\mathsf{Cell}_{ij}^{\lambda_{\mathrm{PICA}}}:\sigma_{ij}\) with cell record \((P_{i}\leftarrow P_{j},W_{j},U_{i},L,V,\Theta,A,\delta)\) and PICA cell profile \(\lambda_{\mathrm{PICA}}\). The recovered \(6\times 6\) multiplicity count under the formal classifier is \(25\;\texttt{action},\;8\;\texttt{implicit},\;2\;\texttt{trivial},\;1\;\texttt{undefined\_circular}\), recorded as a PICA-profile model fact rather than a universal \(6\times 6\) law (PICA nonclaim 2 of Subsubsection 13.2.10).

The raw-to-formal status recovery map is

\[
\texttt{Action}\;\to\;\texttt{action},\quad\texttt{Implicit}\;\to\;\texttt{implicit},\quad\texttt{Trivial}\;\to\;\texttt{trivial},
\]

with \(\texttt{Undefined}\to\text{reason-dependent formal status}\in\{\,\texttt{blocked},\;\texttt{undefined\_circular},\;\texttt{outside\_scope},\;\texttt{absent},\;\texttt{below\_threshold},\;\texttt{collapsed}\,\}\), decided by the formal classifier of Subsection 5.6. Raw PICA status never overrides the formal verdict.

\subsubsection{Pair-Observable Realization and Provenance}

PICA pair-status recovery: \(\texttt{real}\to\texttt{real}\), \(\texttt{provisional}\to\texttt{real\_provisional}\), \(\texttt{duplicated}\to\texttt{real\_duplicated}\), \(\texttt{blocked}\to\texttt{blocked}\).

A PICA pair observable record is

\[
\mathsf{PairObs}_{\{i,j\}}^{\mu}\;=\;(\,\{P_{i},P_{j}\},\;\mu,\;\mathsf{ObservableType},\;\mathsf{BranchData},\;\mathsf{SourceOfTruth},\;L,\;V,\;\Theta,\;A,\;\delta,\;\mathcal N\,),
\]

with realness condition \(\mathsf{SourceOfTruthOK}\land\mathsf{VisibilityPasses}\land\mathsf{AuditPasses}\land\mathsf{ThresholdPasses}\). The source-of-truth predicate is

\[
\mathsf{SourceOfTruthOK}(R,Q)\;\Leftrightarrow\;\mathsf{Committed}(R)\land\mathsf{CorrectProducer}(R,Q)\land\mathsf{CorrectProvenance}(R,Q)\land\mathsf{NotFallbackUnlessAllowed}(R,Q)\land\mathsf{Visible}(R)\land\mathsf{Audited}(R).
\]

Dirty exclusion: \(R.\mathsf{dirty\_since\_step}\neq\varnothing\Rightarrow\neg\mathsf{SourceOfTruthOK}(R,Q)\). Fallback default: a fallback record supports at most \(\texttt{real\_provisional}\) unless an explicit upgrade bridge is recorded.

\subsubsection{Realization Theorem and Nonclaims}

\textbf{Model Theorem 33 (PICA Model-Realization; model-grade).} If \(\mathcal P\in\mathsf{PICAReal}_{\mathrm{fin}}\) is admissible, then

\[
\Gamma;\mathsf H_{\mathrm{PICA}};I_{\mathrm{PICA}}\vdash\mathsf{PICAReal}_{\mathrm{fin}}\models\mathbf{BirdInt}^{\mathrm{aud}}_{\mathrm{fin\text{-}stoch}}{}^{\mathrm{cell}/\mathrm{pair}/\mathrm{prov}}:\texttt{model\_realization}.
\]

The PICA nonclaim register \(\mathcal N_{\mathrm{PICA}}\) contains the ten nonclaims of Subsubsection 13.2.10: no PICA universality, no \(6\times 6\) universal law, raw status not formal status, \(\texttt{action}\) not pair-realness, fallback not real source, ablation evidence is model-grade, \(P_{3}\) not \(\mathrm{P6}_{\mathrm{drive}}\) without bridge, dirty records not real, finite stochastic fragment only, implementation correctness future work.

\subsection{Cantor Sheet}

The Cantor strict-bridge realization sheet \(\mathsf{CantorShell}_{\mathrm{aud}}\) records the realization of \(\mathbf{StrictBridge}_{0\to 1}^{\mathrm{audited\ shell}}\) on the audited Cantor shell host \(\mathsf H_{\mathrm{CantorShell}}\). The main-text record is in Subsection 13.3.

\subsubsection{Host and Realization Object}

The host \(\mathsf H_{\mathrm{CantorShell}}\) is a scoped audited application host built over a fixed audited shell-stable comparison domain \(\mathcal S_{\mathrm{aud}}\). The realization object is

\[
\mathsf{CantorShell}_{\mathrm{aud}}\;=\;(\,\mathcal S_{\mathrm{aud}},\;T_{0},\;T_{1},\;\pi_{0},\;\pi_{1},\;\mathcal F,\;\Delta,\;A,\;\mathcal N\,),
\]

with \(\mathcal S_{\mathrm{aud}}\) the audited shell-stable comparison domain, \(T_{0}\) the base cocycle-pressure theory, \(T_{1}\) the completion-packaged theory, \(\pi_{0}\) the old cocycle-level object map, \(\pi_{1}\) the new packaged-object map, \(\mathcal F\) the persistent packaged-strata family, \(\Delta\) the pressure-gap consequence, \(A\) the audit/theoremlet support record, and \(\mathcal N\) the Cantor nonclaim register. The host instrument is \(I_{\mathrm{Cantor}}\).

\subsubsection{\texorpdfstring{Base Package \(T_{0}\) and Target Package \(T_{1}\)}{D.2.2 Base Package T\_\{0\} and Target Package T\_\{1\}}}

The base package is \(\mathcal T_{0}^{\mathrm{Cantor}}=(\mathcal S_{\mathrm{aud}},\pi_{0},\Sigma_{\pi_{0}}^{\mathrm{cocycle}},E_{0}^{\mathrm{cocycle}},A_{T_{0}})\), with \(\pi_{0}:\mathcal S_{\mathrm{aud}}\to\mathcal O_{0}\), \(E_{0}^{\mathrm{cocycle}}=(E_{\pi_{0}},\operatorname{Cl}_{P_{T_{0}}})\), and base pressure object

\[
P_{T_{0}}(s)\;=\;\lim_{n\to\infty}\frac{1}{n}\,\log\!\left(\sup_{x\in\mathcal S_{\mathrm{aud}}}a_{n}(s;x)\right).
\]

The target package is \(\mathcal T_{1}^{\mathrm{Cantor}}=(\mathcal S_{\mathrm{aud}},\pi_{1},\Sigma_{\pi_{1}}^{\mathrm{pkg}},E_{1}^{\mathrm{pkg}},A_{T_{1}})\), with \(\pi_{1}:\mathcal S_{\mathrm{aud}}\to\mathcal O_{1}\) and packaged completion field \(E_{1}^{\mathrm{pkg}}=(E_{\tau,\ell},\mathrm{Fix}(E_{\tau,\ell}),\mathcal F,\mathrm{Sat},\mathrm{Force}_{4\leftarrow 5})\) with \(E_{\tau,\ell}(\mu)=U_{\ell}(Q_{\ell}(\mu K^{\tau}))\).

\subsubsection{Cantor Strict Bridge}

The Cantor strict bridge is

\[
B_{0\to 1}^{\mathrm{Cantor}}\;=\;(\,\mathcal T_{0}^{\mathrm{Cantor}},\;\mathcal T_{1}^{\mathrm{Cantor}},\;\pi_{0},\;\pi_{1},\;L_{\mathrm{Cantor}},\;V_{\mathrm{Cantor}},\;\Theta_{\mathrm{Cantor}},\;A_{\mathrm{Cantor}},\;\delta_{\mathrm{Cantor}},\;\mathcal N_{\mathrm{Cantor}}\,),
\]

with lift data \(L_{\mathrm{Cantor}}=(E_{\tau,\ell},Q_{\ell},U_{\ell},\mathrm{Sat},\mathrm{Force}_{4\leftarrow 5})\) and bridge-defect bundle \(\delta_{\mathrm{Cantor}}=(\delta_{\mathrm{fact}},\delta_{\mathrm{sat}},\delta_{4\leftarrow 5},\delta_{\mathrm{macro}},\delta_{\mathrm{press}})\).

\subsubsection{Strictness Witness and Pressure-Gap Defect}

The strictness witness is the nonfactorization \(\pi_{1}\not\factor\pi_{0}\) on \(\mathcal S_{\mathrm{aud}}\), equivalently \(\exists\,x,x'\in\mathcal S_{\mathrm{aud}}:\pi_{0}(x)=\pi_{0}(x')\land\pi_{1}(x)\neq\pi_{1}(x')\). The factorization defect is

\[
\Delta_{\mathrm{fact}}^{\mathrm{Cantor}}(\pi_{0},\pi_{1})\;=\;\{\,(x,x'):\pi_{0}(x)=\pi_{0}(x')\land\pi_{1}(x)\neq\pi_{1}(x')\,\}.
\]

The strictness gate \(G_{\mathrm{strict}}=\texttt{strict\_pass}\Leftrightarrow\Delta_{\mathrm{fact}}^{\mathrm{Cantor}}\neq\varnothing\).

The pressure-gap consequence defect is

\[
\Delta_{T_{0}\to T_{1}}(s;x)\;=\;P_{T_{0}}(s)\;-\;\sum_{\Sigma\in\mathcal F(x)}w_{\Sigma}(x)\,P_{\Sigma}(s).
\]

It is a \emph{consequence} defect, not the strictness witness itself.

\subsubsection{Gate Results}

\begin{longtable}[]{@{}ll@{}}
\toprule
gate & result\tabularnewline
\midrule
\endhead
\(G_{\mathrm{suff}}\) & \(\texttt{pass}\)\tabularnewline
\(G_{\mathrm{desc}}\) & \(\texttt{not\_required}\)\tabularnewline
\(G_{\mathrm{stab}}\) & \(\texttt{pass}\)\tabularnewline
\(G_{\mathrm{ctrl}}\) & \(\texttt{pass}\)\tabularnewline
\(G_{\mathrm{nosmuggle}}\) & \(\texttt{pass}\)\tabularnewline
\(G_{\mathrm{vis}}\) & \(\texttt{pass}\)\tabularnewline
\(G_{\mathrm{audit}}\) & \(\texttt{pass}\)\tabularnewline
\(G_{\mathrm{strict}}\) & \(\texttt{strict\_pass}\)\tabularnewline
\(G_{\mathrm{locglob}}\) & \(\texttt{not\_required}\)\tabularnewline
\bottomrule
\end{longtable}

\subsubsection{Realization Theorem and Nonclaims}

\textbf{Model Theorem 34 (Cantor Strict-Bridge Realization; model-grade).} If \(\mathsf{AdmCantorShell}(\mathsf{CantorShell}_{\mathrm{aud}})\) holds, the strictness witness records a nonempty \(\Delta_{\mathrm{fact}}^{\mathrm{Cantor}}\) on \(\mathcal S_{\mathrm{aud}}\), and the gate record holds, then

\[
\Gamma;\mathsf H_{\mathrm{CantorShell}};I_{\mathrm{Cantor}}\vdash\mathsf{CantorShell}_{\mathrm{aud}}\models\mathbf{StrictBridge}_{0\to 1}^{\mathrm{audited\ shell}}:\texttt{model\_realization},
\]

and the associated promotion judgment is \(\Gamma;\mathcal T_{0}^{\mathrm{Cantor}};I_{\mathrm{Cantor}}\vdash\mathrm{Promote}(B_{0\to 1}^{\mathrm{Cantor}}):\texttt{strict}\).

The Cantor nonclaim register \(\mathcal N_{\mathrm{Cantor}}\) contains the ten nonclaims of Subsubsection 13.3.9: shell-bound only, no shell-general theorem, no broader external-class theorem, strictness not macro closure, strictness not drive, pressure gap not KL closure deficit without bridge, pressure gap not strictness witness, \(\Delta_{T_{0}\to T_{1}}\) is conditional disintegration not stratumwise root separation, model realization not universal strict-extension theory, \(T_{1}\) not single repeated idempotent.

\subsection{Abstract-Interpretation Sheet}

The abstract-interpretation realization sheet \(\mathsf{AIFam}_{\mathrm{fin}}\) records the realization of the role fragment \(P_{1}/P_{2}/P_{4}/P_{5}/P_{6}\) on the finite abstract-interpretation host \(\mathsf H_{\mathrm{AI}}\), supported by the theorem-grade theorems 26--28 of Section 12 and the model-grade Model Theorem 29.

\subsubsection{Host}

A finite abstract-interpretation host is the eight-component finite typed record

\[
H\;=\;(\,C,\;A,\;\alpha,\;\gamma,\;\rho,\;F,\;F^{\sharp},\;A_{\mathrm{AI}}\,),
\]

with \((C,\le_{C})\) a finite concrete ordered domain, \((A,\le_{A})\) a finite abstract ordered domain, \(\alpha\dashv\gamma\) a Galois connection between \(C\) and \(A\), \(\rho=\gamma\alpha:C\to C\) the host closure operator, \(F:C\to C\) and \(F^{\sharp}:A\to A\) monotone transformers, and \(A_{\mathrm{AI}}\) the host audit record. The host instrument is \(I_{\mathrm{AI}}\).

\subsubsection{Single-Host Record and Field Glosses}

\begin{longtable}[]{@{}ll@{}}
\toprule
field & meaning\tabularnewline
\midrule
\endhead
\(C\) & finite concrete ordered domain\tabularnewline
\(A\) & finite abstract ordered domain\tabularnewline
\(\alpha:C\to A\) & abstraction map\tabularnewline
\(\gamma:A\to C\) & concretization map\tabularnewline
\(\rho=\gamma\alpha\) & host closure operator on \(C\)\tabularnewline
\(F:C\to C\) & monotone concrete transformer\tabularnewline
\(F^{\sharp}:A\to A\) & monotone abstract transformer\tabularnewline
\(A_{\mathrm{AI}}\) & host audit record\tabularnewline
\bottomrule
\end{longtable}

For the model-family realization, the finite family is recorded as

\[
\mathsf{AIFam}_{\mathrm{fin}}\;=\;(\,\{H_k\}_{k\in K},\;\mathsf{Refine},\;\mathsf{Audit},\;\mathcal N_{\mathrm{AI}}\,),
\]

with finite index set \(K\) and

\[
H_k=(\,C,\;A_k,\;\alpha_k,\;\gamma_k,\;\rho_k,\;F,\;F_k^{\sharp},\;A_{\mathrm{AI},k}\,).
\]

A refinement edge \(A_k\leadsto A_\ell\) is recorded by a finite audited refinement record, for example by a forgetting map \(r_{\ell k}:A_\ell\to A_k\) with \(\alpha_k=r_{\ell k}\alpha_\ell\), or by a recorded closure-refinement comparison \(\rho_\ell\le\rho_k\). The canonical AI package for a host is \(\mathcal T_{\mathrm{AI}}=(C,\alpha,\Sigma_\alpha,\rho,\mathcal A_{\mathrm{AI}})\), with \(\Sigma_\alpha=\{\alpha^{-1}(B):B\subseteq A\}\) in full finite-lens mode; the audit functional \(\mathcal A_{\mathrm{AI}}\) is assembled from the finite host audit record \(A_{\mathrm{AI}}\).

\subsubsection{Realized Roles and Theorem Trio}

The realization map sends each realized primitive to its AI-host witness:

\begin{longtable}[]{@{}lll@{}}
\toprule
primitive & AI realization & core condition\tabularnewline
\midrule
\endhead
\(P_{1}\) & descent / sound transformer & \(\alpha F\le_{A}F^{\sharp}\alpha\)\tabularnewline
\(P_{2}\) & representability & \(p\;\text{representable}\iff\rho(p)=p\)\tabularnewline
\(P_{4}\) & refinement of abstract domain & \(A_{k}\rightsquigarrow A_{\ell}\)\tabularnewline
\(P_{5}\) & closure / packaging & \(\rho=\gamma\alpha,\;\rho^{2}=\rho\)\tabularnewline
\(P_{6}\) & audit / provenance & \(A_{\mathrm{AI}}\) checks maps and defects\tabularnewline
\bottomrule
\end{longtable}

The supporting theorems are:

\begin{itemize}
\tightlist
\item
  \textbf{Theorem 26 (AIClosureAsPackaging; theorem-grade).} \(\rho=\gamma\alpha\) is extensive, monotone, and idempotent.
\item
  \textbf{Theorem 27 (AIDescentAsSoundTransformer; theorem-grade).} \(\alpha F\le_{A}F^{\sharp}\alpha\iff F\gamma\le_{C}\gamma F^{\sharp}\).
\item
  \textbf{Theorem 28 (AIRepresentabilityAsFixedness; theorem-grade).} \(p\,\text{representable}\iff\rho(p)=p\), with \(\mathrm{Rep}_{\alpha}(C)=\mathrm{Fix}(\rho)=\mathrm{im}(\gamma)\).
\end{itemize}

\subsubsection{\texorpdfstring{Optional \(P_{3}\) Extension}{D.3.4 Optional P\_\{3\} Extension}}

A \(P_{3}\) extension may be added by recording two AI host closures \(\rho_{1}=\gamma_{1}\alpha_{1}\) and \(\rho_{2}=\gamma_{2}\alpha_{2}\) and the noncommutation witness \(\rho_{1}\rho_{2}\neq\rho_{2}\rho_{1}\), in which case AI realizes the route-mismatch primitive at a directed cell \(P_{3}\leftarrow P_{5}:\texttt{action}\). This extension is not part of the core five-role realization of Subsection 12.5.

\subsubsection{Realization Theorem and Nonclaims}

\textbf{Model Theorem 29 (AIModelFamily; model-grade).} If \(\mathsf{AIFam}_{\mathrm{fin}}\) is an admissible finite family of AI hosts, then

\[
\Gamma;\mathsf H_{\mathrm{AI}};I_{\mathrm{AI}}\vdash\mathsf{AIFam}_{\mathrm{fin}}\models\mathbf{BirdInt}^{\mathrm{aud}}_{\mathrm{fin}}{}^{\,P_{1}/P_{2}/P_{4}/P_{5}/P_{6}}:\texttt{model\_realization}.
\]

The AI nonclaim register \(\mathcal N_{\mathrm{AI}}\) contains: AI is not the whole calculus, soundness is not completeness, closure does not imply descent, representability does not imply sound transformer, refinement does not automatically improve every defect, AI five-role realization does not include full \(P_{3}\) unless an extension is added, abstract claims cannot overread concrete distinctions invisible to \(\alpha\), audit of soundness is not instrument-free truth, AI model-family claims are finite-host scoped, and \(P_{5}\)-closure does not imply drive.

\subsection{Graph/Cohomology Sheet}

The graph/cohomology realization sheet \(\mathsf{CohGraph}_{\mathrm{fin}}\) records the realization of the drive-support fragment \(P_{2}/\mathrm{P6}_{\mathrm{drive}}^{H^{1}}/P_{3}\text{-separation}\) on the finite graph/cohomology host \(\mathsf H_{\mathrm{graph}}\). The main-text record is in Subsection 13.4.

\subsubsection{Host and Realization Object}

The realization object is

\[
\mathsf{CohGraph}_{\mathrm{fin}}\;=\;(\,G,\;C^{0}(G;k),\;C^{1}(G;k),\;d,\;H^{1}(G;k),\;\mathsf{Cycles}(G),\;a,\;\mathsf{Gate},\;A,\;V,\;\mathcal N\,),
\]

with \(G=(V,E)\) a finite support graph, \(C^{0}(G;k)\) and \(C^{1}(G;k)\) the finite cochain spaces, \(d:C^{0}\to C^{1}\) the coboundary, \(H^{1}(G;k)\) the first cohomology, \(\mathsf{Cycles}(G)\) the finite cycle space, \(a\in C^{1}(G;k)\) the affinity/drive cochain, \(\mathsf{Gate}\) the gating record, \(A\) and \(V\) host audit and visibility records, and \(\mathcal N\) the cohomology nonclaim register. The host instrument is \(I_{\mathrm{graph}}\).

\subsubsection{Realized Fragment and Realization Map}

The realization fragment is \(\mathbf{BirdInt}^{\mathrm{aud}}_{\mathrm{fin}}{}^{\,P_{2}/\mathrm{P6}_{\mathrm{drive}}^{H^{1}}/P_{3}\text{-separation}}\), with realization map

\begin{longtable}[]{@{}ll@{}}
\toprule
\begin{minipage}[b]{0.34\columnwidth}\raggedright
primitive\strut
\end{minipage} & \begin{minipage}[b]{0.60\columnwidth}\raggedright
cohomology realization\strut
\end{minipage}\tabularnewline
\midrule
\endhead
\begin{minipage}[t]{0.34\columnwidth}\raggedright
\(P_{2}\)\strut
\end{minipage} & \begin{minipage}[t]{0.60\columnwidth}\raggedright
support gating / edge deletion / cycle-rank capacity\strut
\end{minipage}\tabularnewline
\begin{minipage}[t]{0.34\columnwidth}\raggedright
\(\mathrm{P6}_{\mathrm{drive}}^{H^{1}}\)\strut
\end{minipage} & \begin{minipage}[t]{0.60\columnwidth}\raggedright
cohomological drive certificate \([a]\neq 0\)\strut
\end{minipage}\tabularnewline
\begin{minipage}[t]{0.34\columnwidth}\raggedright
\(P_{6}^{\mathrm{audit}}\)\strut
\end{minipage} & \begin{minipage}[t]{0.60\columnwidth}\raggedright
audit of cochains, cycles, source-of-truth, and visibility\strut
\end{minipage}\tabularnewline
\begin{minipage}[t]{0.34\columnwidth}\raggedright
\(P_{3}\)\strut
\end{minipage} & \begin{minipage}[t]{0.60\columnwidth}\raggedright
optional route/holonomy witness, kept separate from drive\strut
\end{minipage}\tabularnewline
\begin{minipage}[t]{0.34\columnwidth}\raggedright
\(P_{4}\)\strut
\end{minipage} & \begin{minipage}[t]{0.60\columnwidth}\raggedright
optional local/global cover or gluing of cochains\strut
\end{minipage}\tabularnewline
\begin{minipage}[t]{0.34\columnwidth}\raggedright
\(P_{5}\)\strut
\end{minipage} & \begin{minipage}[t]{0.60\columnwidth}\raggedright
optional quotient/support package of graph state\strut
\end{minipage}\tabularnewline
\bottomrule
\end{longtable}

The core sheet focuses on \(P_{2}\) and \(\mathrm{P6}_{\mathrm{drive}}^{H^{1}}\).

\subsubsection{Drive and No-Drive Certificates}

The drive certificate is \(\kappa_{6}^{\mathrm{drive}}(a):[a]\neq 0\in H^{1}(G;k)\), equivalent to \(\exists\,\gamma\in\mathsf{Cycles}(G):\oint_{\gamma}a\neq 0\). No-drive certificates: \([a]=0\), \(a=d\phi\), \(H^{1}(G;k)=0\), \(\beta_{1}(G)=0\).

\subsubsection{Gating Realization}

A \(P_{2}\)-gate \(G\rightsquigarrow G'\) has cycle-rank loss defect \(\delta_{2}^{\mathrm{cycle}}(G,G')=\beta_{1}(G)-\beta_{1}(G')\) and forest threshold \(\Theta_{2}^{\mathrm{forest}}:\beta_{1}(G')=0\). When \(G'\) is a forest, \(H^{1}(G';k)=0\) and \(\mathrm{P6}_{\mathrm{drive}}^{H^{1}}\) is absent with record on \(G'\).

\subsubsection{Core Theorems}

The core graph/cohomology theorems of Section 8 (with long-form proofs in Appendix B.2) are: forest no-drive (Theorem 15: \(\beta_{1}(G)=0\Rightarrow H^{1}(G;k)=0\)), exact-form null-drive (Theorem 16: \(a=d\phi\Rightarrow\oint_{\gamma}a=0\)), nonzero-affinity equivalence (Theorem 17: \([a]\neq 0\iff\exists\gamma:\oint_{\gamma}a\neq 0\)), and gating-affinity suppression (Theorem 18: gating to a forest kills cycle-supported drive).

\subsubsection{Realization Theorem and Nonclaims}

\textbf{Cohomology Support Realization (model-grade).} Every admissible \(\mathsf{CohGraph}_{\mathrm{fin}}\) realizes the fragment containing \(P_{2}\)-support gating and cycle-rank loss, \(\mathrm{P6}_{\mathrm{drive}}^{H^{1}}\) certificates, no-drive certificates, and \(P_{3}\)/drive separation:

\[
\Gamma;\mathsf H_{\mathrm{graph}};I_{\mathrm{graph}}\vdash\mathsf{CohGraph}_{\mathrm{fin}}\models\mathbf{BirdInt}^{\mathrm{aud}}_{\mathrm{fin}}{}^{\,P_{2}/\mathrm{P6}_{\mathrm{drive}}^{H^{1}}/P_{3}\text{-separation}}:\texttt{model\_realization}.
\]

The cohomology nonclaim register \(\mathcal N_{\mathrm{coh}}\) contains the ten nonclaims of Subsubsection 13.4.4: \(\mathrm{P6}_{\mathrm{drive}}^{H^{1}}\) is one drive certificate family not all drive, \(P_{3}\) holonomy not \(\mathrm{P6}_{\mathrm{drive}}\) without a bridge, gating to a forest kills cycle-supported drive only, exact \(1\)-cochains have zero cycle drive, \(H^{1}(G;k)=0\) is a no-drive certificate for cohomological drive only, nonzero holonomy can coexist with zero drive cochain, host is finite and scoped, cohomology does not replace PICA/Cantor/AI, no edge-support claim accepted without visibility/audit, positive holonomy-to-drive claims require \(B_{3\to 6}^{\mathrm{drive}}\).

\subsection{Combined Coverage Matrix}

This subsection records the combined coverage matrix across the four model realizations. The matrix has one row per realization and one column per scoped attribute (host/instrument, realized fragment, claim strength); Subsubsection D.5.2 separately records per-primitive role coverage. The matrix collects the per-realization data of Subsections D.1--D.4 in one place; it is the long-form analogue of the matrix recorded in Subsection 13.5.

\subsubsection{Realization Matrix}

\begin{longtable}[]{@{}llll@{}}
\toprule
\begin{minipage}[b]{0.14\columnwidth}\raggedright
model\strut
\end{minipage} & \begin{minipage}[b]{0.22\columnwidth}\raggedright
host / instrument\strut
\end{minipage} & \begin{minipage}[b]{0.42\columnwidth}\raggedright
realized fragment\strut
\end{minipage} & \begin{minipage}[b]{0.11\columnwidth}\raggedright
claim strength\strut
\end{minipage}\tabularnewline
\midrule
\endhead
\begin{minipage}[t]{0.14\columnwidth}\raggedright
\(\mathsf{PICAReal}_{\mathrm{fin}}\)\strut
\end{minipage} & \begin{minipage}[t]{0.22\columnwidth}\raggedright
\(\mathsf H_{\mathrm{PICA}}\) / \(I_{\mathrm{PICA}}\)\strut
\end{minipage} & \begin{minipage}[t]{0.42\columnwidth}\raggedright
\(\mathbf{BirdInt}^{\mathrm{aud}}_{\mathrm{fin\text{-}stoch}}{}^{\mathrm{cell}/\mathrm{pair}/\mathrm{prov}}\)\strut
\end{minipage} & \begin{minipage}[t]{0.11\columnwidth}\raggedright
\(\texttt{model\_realization}\)\strut
\end{minipage}\tabularnewline
\begin{minipage}[t]{0.14\columnwidth}\raggedright
\(\mathsf{CantorShell}_{\mathrm{aud}}\)\strut
\end{minipage} & \begin{minipage}[t]{0.22\columnwidth}\raggedright
\(\mathsf H_{\mathrm{CantorShell}}\) / \(I_{\mathrm{Cantor}}\)\strut
\end{minipage} & \begin{minipage}[t]{0.42\columnwidth}\raggedright
\(\mathbf{StrictBridge}_{0\to 1}^{\mathrm{audited\ shell}}\)\strut
\end{minipage} & \begin{minipage}[t]{0.11\columnwidth}\raggedright
\(\texttt{model\_realization}\) (application-scoped)\strut
\end{minipage}\tabularnewline
\begin{minipage}[t]{0.14\columnwidth}\raggedright
\(\mathsf{AIFam}_{\mathrm{fin}}\)\strut
\end{minipage} & \begin{minipage}[t]{0.22\columnwidth}\raggedright
\(\mathsf H_{\mathrm{AI}}\) / \(I_{\mathrm{AI}}\)\strut
\end{minipage} & \begin{minipage}[t]{0.42\columnwidth}\raggedright
\(\mathbf{BirdInt}^{\mathrm{aud}}_{\mathrm{fin}}{}^{\,P_{1}/P_{2}/P_{4}/P_{5}/P_{6}}\)\strut
\end{minipage} & \begin{minipage}[t]{0.11\columnwidth}\raggedright
\(\texttt{model\_realization}\)\strut
\end{minipage}\tabularnewline
\begin{minipage}[t]{0.14\columnwidth}\raggedright
\(\mathsf{CohGraph}_{\mathrm{fin}}\)\strut
\end{minipage} & \begin{minipage}[t]{0.22\columnwidth}\raggedright
\(\mathsf H_{\mathrm{graph}}\) / \(I_{\mathrm{graph}}\)\strut
\end{minipage} & \begin{minipage}[t]{0.42\columnwidth}\raggedright
\(\mathbf{BirdInt}^{\mathrm{aud}}_{\mathrm{fin}}{}^{\,P_{2}/\mathrm{P6}_{\mathrm{drive}}^{H^{1}}/P_{3}\text{-separation}}\)\strut
\end{minipage} & \begin{minipage}[t]{0.11\columnwidth}\raggedright
\(\texttt{model\_realization}\)\strut
\end{minipage}\tabularnewline
\bottomrule
\end{longtable}

\subsubsection{Cross-Model Role Coverage}

The cross-model role-coverage table records, per primitive, what each realization carries within its declared fragment:

\begin{longtable}[]{@{}lllll@{}}
\toprule
\begin{minipage}[b]{0.05\columnwidth}\raggedright
role\strut
\end{minipage} & \begin{minipage}[b]{0.18\columnwidth}\raggedright
\(\mathsf{PICA}\)\strut
\end{minipage} & \begin{minipage}[b]{0.20\columnwidth}\raggedright
\(\mathsf{Cantor}\)\strut
\end{minipage} & \begin{minipage}[b]{0.25\columnwidth}\raggedright
\(\mathsf{AI}\)\strut
\end{minipage} & \begin{minipage}[b]{0.19\columnwidth}\raggedright
\(\mathsf{CohGraph}\)\strut
\end{minipage}\tabularnewline
\midrule
\endhead
\begin{minipage}[t]{0.05\columnwidth}\raggedright
\(P_{1}\)\strut
\end{minipage} & \begin{minipage}[t]{0.18\columnwidth}\raggedright
descent / closure update\strut
\end{minipage} & \begin{minipage}[t]{0.20\columnwidth}\raggedright
macro obstruction; not primary\strut
\end{minipage} & \begin{minipage}[t]{0.25\columnwidth}\raggedright
sound descent \(\alpha F\le_{A}F^{\sharp}\alpha\)\strut
\end{minipage} & \begin{minipage}[t]{0.19\columnwidth}\raggedright
not primary\strut
\end{minipage}\tabularnewline
\begin{minipage}[t]{0.05\columnwidth}\raggedright
\(P_{2}\)\strut
\end{minipage} & \begin{minipage}[t]{0.18\columnwidth}\raggedright
gating / support / admissibility\strut
\end{minipage} & \begin{minipage}[t]{0.20\columnwidth}\raggedright
admissibility obstruction\strut
\end{minipage} & \begin{minipage}[t]{0.25\columnwidth}\raggedright
representability \(\rho(p)=p\)\strut
\end{minipage} & \begin{minipage}[t]{0.19\columnwidth}\raggedright
support gating\strut
\end{minipage}\tabularnewline
\begin{minipage}[t]{0.05\columnwidth}\raggedright
\(P_{3}\)\strut
\end{minipage} & \begin{minipage}[t]{0.18\columnwidth}\raggedright
route mismatch\strut
\end{minipage} & \begin{minipage}[t]{0.20\columnwidth}\raggedright
not central\strut
\end{minipage} & \begin{minipage}[t]{0.25\columnwidth}\raggedright
optional noncommuting closures (deferred extension)\strut
\end{minipage} & \begin{minipage}[t]{0.19\columnwidth}\raggedright
separated from drive\strut
\end{minipage}\tabularnewline
\begin{minipage}[t]{0.05\columnwidth}\raggedright
\(P_{4}\)\strut
\end{minipage} & \begin{minipage}[t]{0.18\columnwidth}\raggedright
refinement / staging\strut
\end{minipage} & \begin{minipage}[t]{0.20\columnwidth}\raggedright
\(P_{4}\!\leftarrow\!P_{5}\) completion-packaging\strut
\end{minipage} & \begin{minipage}[t]{0.25\columnwidth}\raggedright
refinement \(A_{k}\rightsquigarrow A_{\ell}\)\strut
\end{minipage} & \begin{minipage}[t]{0.19\columnwidth}\raggedright
optional local/global cover\strut
\end{minipage}\tabularnewline
\begin{minipage}[t]{0.05\columnwidth}\raggedright
\(P_{5}\)\strut
\end{minipage} & \begin{minipage}[t]{0.18\columnwidth}\raggedright
package / completion\strut
\end{minipage} & \begin{minipage}[t]{0.20\columnwidth}\raggedright
completion-packaged objects\strut
\end{minipage} & \begin{minipage}[t]{0.25\columnwidth}\raggedright
closure / packaging \(\rho=\gamma\alpha\)\strut
\end{minipage} & \begin{minipage}[t]{0.19\columnwidth}\raggedright
optional graph package\strut
\end{minipage}\tabularnewline
\begin{minipage}[t]{0.05\columnwidth}\raggedright
\(P_{6}\)\strut
\end{minipage} & \begin{minipage}[t]{0.18\columnwidth}\raggedright
audit / provenance\strut
\end{minipage} & \begin{minipage}[t]{0.20\columnwidth}\raggedright
audit of shell bridge\strut
\end{minipage} & \begin{minipage}[t]{0.25\columnwidth}\raggedright
audit of maps and soundness\strut
\end{minipage} & \begin{minipage}[t]{0.19\columnwidth}\raggedright
drive certificate \([a]\neq 0\) and audit\strut
\end{minipage}\tabularnewline
\bottomrule
\end{longtable}

\subsubsection{Combined Nonclaims}

The combined nonclaim register records ten cross-model nonclaims that hold uniformly across the four realizations:

\begin{enumerate}
\def\labelenumi{\arabic{enumi}.}
\tightlist
\item
  No model realization is the whole of \(\mathbf{BirdInt}^{\mathrm{aud}}_{\mathrm{fin}}\).
\item
  PICA is not the universal six-symbol algebra.
\item
  Cantor is audited-shell scoped.
\item
  The abstract-interpretation host realizes a five-role fragment unless an explicit \(P_{3}\) extension is added.
\item
  The graph/cohomology host realizes cohomological drive, not all drive notions.
\item
  Model-realization evidence is not theorem-grade unless an admissible bridge is supplied.
\item
  Ablations, simulations, and fallbacks require source-of-truth contracts.
\item
  Strictness does not imply macro closure or drive.
\item
  Audit is not final self-certification.
\item
  No model may overread suppressed content without an admissible bridge.
\end{enumerate}

The combined matrix records four scoped realizations and the per-primitive role coverage they admit. The matrix is not read vertically as a coverage assertion: the framework does not assert that the union of the four realizations covers the whole calculus, in keeping with the model-coverage nonclaim of Subsubsection 13.6.2.
 \section{High-Structure Semantics}

This appendix records the long-form definitions and recoveries of the finite decorated-square fragment of Section 14. Subsection E.1 records the finite decorated-square fragment in full. Subsection E.2 records the descent-square recovery in detail. Subsection E.3 records the completion-pasting recovery in detail. Subsection E.4 records the strict-extension as no-filler recovery in detail. Subsection E.5 records the high-structure nonclaims.

The appendix is host-bound to the finite decorated double-category fragment \(\mathsf H_{\mathrm{DC}}\) and supplements Section 14; it does not introduce new theorems beyond Theorems 30, 31, and 32.

\subsection{Finite Decorated-Square Fragment}

This subsection expands the schema of Subsection 14.1. The decorated-square fragment is a secondary organizational layer over the core calculus: it does not replace the directed-cell, pair-observable, promotion, claim, gate, or square status families, and it does not assert a completeness theorem for the calculus.

\subsubsection{Objects, Horizontal Arrows, Vertical Arrows}

The finite decorated double-category fragment \(\mathsf H_{\mathrm{DC}}\) is the finite typed record

\[
\mathsf H_{\mathrm{DC}}\;=\;(\,\mathrm{Ob},\;\mathrm{Hor},\;\mathrm{Ver},\;\mathrm{Sq},\;A_{\mathrm{DC}},\;V_{\mathrm{DC}},\;\mathcal N_{\mathrm{DC}}\,),
\]

with components fixed in Subsubsection 14.1.1: \(\mathrm{Ob}\) the finite class of package objects \(\mathcal T\) of Subsection 3.3; \(\mathrm{Hor}\) the finite class of typed refinements, lenses, and package morphisms; \(\mathrm{Ver}\) the finite class of typed transformations carrying actor-update content with audit, source, and visibility decorations; \(\mathrm{Sq}\) the finite class of decorated squares; \(A_{\mathrm{DC}}\) the host audit record; \(V_{\mathrm{DC}}\) the host visibility record; and \(\mathcal N_{\mathrm{DC}}\) the host nonclaim register.

Compositions of horizontal and vertical arrows are admitted only when the resulting boundary records are finite and audited under \(I_{\mathrm{DC}}\).

\subsubsection{Decorated Square Schema}

A decorated square \(\Xi\in\mathrm{Sq}\) is a finite typed record

\[
\Xi\;=\;(\,X,\;F,\;q,\;F^{\sharp},\;\mathrm{src}_{\Xi},\;\mathrm{tgt}_{\Xi},\;\mathrm{type}_{\Xi},\;\mathsf{Decor}_{\Xi},\;A_{\Xi},\;V_{\Xi},\;\delta_{\Xi},\;\mathcal N_{\Xi}\,),
\]

with \(\mathrm{type}_{\Xi}\in\{\,\texttt{descent},\;\texttt{quotient},\;\texttt{completion},\;\texttt{promotion},\;\texttt{visibility},\;\texttt{audit},\;\texttt{factorization}\,\}\) and \(\mathsf{Decor}_{\Xi}\) recording the type-specific decoration data.

\begin{longtable}[]{@{}ll@{}}
\toprule
\begin{minipage}[b]{0.47\columnwidth}\raggedright
field\strut
\end{minipage} & \begin{minipage}[b]{0.47\columnwidth}\raggedright
meaning\strut
\end{minipage}\tabularnewline
\midrule
\endhead
\begin{minipage}[t]{0.47\columnwidth}\raggedright
\(X\)\strut
\end{minipage} & \begin{minipage}[t]{0.47\columnwidth}\raggedright
source object of the square, an element of \(\mathrm{Ob}\)\strut
\end{minipage}\tabularnewline
\begin{minipage}[t]{0.47\columnwidth}\raggedright
\(F\)\strut
\end{minipage} & \begin{minipage}[t]{0.47\columnwidth}\raggedright
horizontal arrow component of the boundary\strut
\end{minipage}\tabularnewline
\begin{minipage}[t]{0.47\columnwidth}\raggedright
\(q\)\strut
\end{minipage} & \begin{minipage}[t]{0.47\columnwidth}\raggedright
vertical arrow component of the boundary\strut
\end{minipage}\tabularnewline
\begin{minipage}[t]{0.47\columnwidth}\raggedright
\(F^{\sharp}\)\strut
\end{minipage} & \begin{minipage}[t]{0.47\columnwidth}\raggedright
bottom horizontal arrow or candidate descended/completed arrow\strut
\end{minipage}\tabularnewline
\begin{minipage}[t]{0.47\columnwidth}\raggedright
\(\mathrm{src}_{\Xi},\mathrm{tgt}_{\Xi}\)\strut
\end{minipage} & \begin{minipage}[t]{0.47\columnwidth}\raggedright
source and target boundary tuples\strut
\end{minipage}\tabularnewline
\begin{minipage}[t]{0.47\columnwidth}\raggedright
\(\mathrm{type}_{\Xi}\)\strut
\end{minipage} & \begin{minipage}[t]{0.47\columnwidth}\raggedright
one of the seven square types listed above\strut
\end{minipage}\tabularnewline
\begin{minipage}[t]{0.47\columnwidth}\raggedright
\(\mathsf{Decor}_{\Xi}\)\strut
\end{minipage} & \begin{minipage}[t]{0.47\columnwidth}\raggedright
type-specific decoration: descent, quotient, completion, promotion, visibility, audit, or factorization data\strut
\end{minipage}\tabularnewline
\begin{minipage}[t]{0.47\columnwidth}\raggedright
\(A_{\Xi},V_{\Xi}\)\strut
\end{minipage} & \begin{minipage}[t]{0.47\columnwidth}\raggedright
per-square audit and visibility records\strut
\end{minipage}\tabularnewline
\begin{minipage}[t]{0.47\columnwidth}\raggedright
\(\delta_{\Xi}\)\strut
\end{minipage} & \begin{minipage}[t]{0.47\columnwidth}\raggedright
per-square defect record, classified in the square-status family\strut
\end{minipage}\tabularnewline
\begin{minipage}[t]{0.47\columnwidth}\raggedright
\(\mathcal N_{\Xi}\)\strut
\end{minipage} & \begin{minipage}[t]{0.47\columnwidth}\raggedright
per-square nonclaim register\strut
\end{minipage}\tabularnewline
\bottomrule
\end{longtable}

The exactness condition for each square type is recorded in Subsubsection 14.1.2: descent commutation equality for \(\texttt{descent}\), completion-commutation equality for \(\texttt{completion}\), factorization-filler condition for \(\texttt{factorization}\), and bridge-admissibility conditions for \(\texttt{visibility}\) and \(\texttt{audit}\).

\subsubsection{Strict Pasting}

The strict-pasting law of Subsubsection 14.1.3 records the aggregation rule for finite paste of two decorated squares \(\Xi_{1}\) and \(\Xi_{2}\) along a shared horizontal or vertical edge:

\[
\delta_{\Xi_{1}\cdot\Xi_{2}}\;=\;\delta_{\Xi_{1}}\;\cup\;\delta_{\Xi_{2}}\;\cup\;\delta_{\mathrm{paste}}(\Xi_{1},\Xi_{2}),
\]

where \(\delta_{\mathrm{paste}}(\Xi_{1},\Xi_{2})\) is the finite paste-defect record on the boundary mismatch. Pasting is admissible when the shared edge data agree on \(\mathrm{Hor}\) or \(\mathrm{Ver}\) and the per-square admissibility conditions hold for both factors. Pseudo-pasting and weak pasting are out of scope.

\subsection{Descent Squares}

This subsection records the long-form descent-square recovery (Theorem 30 of Subsection 14.2) on the finite decorated double-category fragment \(\mathsf H_{\mathrm{DC}}\).

For a finite package object \(X\in\mathrm{Ob}\), a horizontal arrow \(F:X\to X\), a quotient vertical arrow \(q:X\to\mathrm{im}(q)\), and a candidate descended horizontal arrow \(F^{\sharp}:\mathrm{im}(q)\to\mathrm{im}(q)\), the descent square is

\[
\Xi_{\mathrm{desc}}\;:\;
\begin{array}{ccc}
X & \xrightarrow{F} & X\\
{\scriptstyle q}\downarrow & \Downarrow\,\Xi_{\mathrm{desc}} & \downarrow{\scriptstyle q}\\
\mathrm{im}(q) & \xrightarrow{F^{\sharp}} & \mathrm{im}(q)
\end{array}
\]

with \(\mathrm{type}_{\Xi}=\texttt{descent}\) and \(\mathsf{Decor}_{\Xi}\) recording the descent triple \((F,q,F^{\sharp})\).

The square is exact under \(I_{\mathrm{DC}}\) iff the descent commutation equality \(F^{\sharp}\circ q_{\mathrm{im}}=q_{\mathrm{im}}\circ F\) holds on \(X\), equivalently iff the well-definedness condition

\[
q(x)=q(x')\Rightarrow qF(x)=qF(x')
\]

holds for every \(x,x'\in X\). The two conditions are equivalent on the finite carrier by Theorem 30 (Subsubsection 14.2.2).

The recovery records descent data in the decorated double-category fragment: a finite descent record \((F,q,F^{\sharp})\) corresponds to an exact decorated square of type \(\texttt{descent}\). The descent defect record on the fragment is

\[
\delta_{\mathrm{desc}}(F,q)\;=\;\{\,(x,x')\in X\times X\;:\;q(x)=q(x')\;\land\;qF(x)\neq qF(x')\,\},
\]

a finite typed defect on the finite carrier \(X\). A nonempty \(\delta_{\mathrm{desc}}(F,q)\) obstructs the existence of a well-defined descended map \(F^{\sharp}\) satisfying \(F^{\sharp}\circ q_{\mathrm{im}}=q_{\mathrm{im}}\circ F\), equivalently obstructs the exactness of \(\Xi_{\mathrm{desc}}\). The instrument-indexed verdict is recorded in Subsubsection 14.2.1: \(\Gamma;\mathsf H_{\mathrm{DC}};I_{\mathrm{DC}}\vdash\textsc{DescentSquareRecovery}:\texttt{accepted}\).

This recovery is representational: it records the descent data handled by Theorem 30 inside \(\mathsf H_{\mathrm{DC}}\) and does not assert that every descent fact in the core calculus reduces to a decorated square.

\subsection{Completion Pasting}

This subsection records the long-form completion-pasting recovery (Theorem 31 of Subsection 14.3).

For a finite package object \(C\in\mathrm{Ob}\) and two finite completions \(E_{1},E_{2}\in\mathrm{Hor}\), \(E_{1},E_{2}:C\to C\), the completion-pasting square is

\[
\Xi_{\mathrm{comp}}\;:\;
\begin{array}{ccc}
C & \xrightarrow{E_{2}} & C\\
{\scriptstyle E_{1}}\downarrow & \Downarrow\,\Xi_{\mathrm{comp}} & \downarrow{\scriptstyle E_{1}}\\
C & \xrightarrow{E_{2}} & C
\end{array}
\]

with \(\mathrm{type}_{\Xi}=\texttt{completion}\) and \(\mathsf{Decor}_{\Xi}\) recording the completion pair \((E_{1},E_{2})\).

The square is exact iff \(E_{1}E_{2}=E_{2}E_{1}\) on \(C\), and is obstructed iff there exists \(c\in C\) with \(E_{1}E_{2}(c)\neq E_{2}E_{1}(c)\). The completion-pasting defect record is

\[
\delta_{\mathrm{comp}}(E_{1},E_{2})\;=\;\{\,c\in C:E_{1}E_{2}(c)\neq E_{2}E_{1}(c)\,\}\;\subseteq\;C,
\]

a finite typed defect on the finite carrier \(C\). A nonempty \(\delta_{\mathrm{comp}}(E_{1},E_{2})\) obstructs the exactness of \(\Xi_{\mathrm{comp}}\), and the noncommutation \(E_{1}E_{2}\neq E_{2}E_{1}\) is recovered as a \(P_{3}\)-pasting defect of the calculus, in keeping with Theorem 31 (Subsubsection 14.3.2).

The instrument-indexed verdict is \(\Gamma;\mathsf H_{\mathrm{DC}};I_{\mathrm{DC}}\vdash\textsc{CompletionPastingRecovery}:\texttt{accepted}\). The recovery is representational; it records the completion-noncommutation phenomenon at the level of \(\mathsf H_{\mathrm{DC}}\) and does not assert that every completion-related fact reduces to a decorated square.

\subsection{Strict Extension as No-Filler}

This subsection records the long-form strict-extension as no-filler recovery (Theorem 32 of Subsection 14.4).

For a finite package object \(S\in\mathrm{Ob}\), two horizontal object maps \(\pi_{0}:S\to O_{0}\) and \(\pi_{1}:S\to O_{1}\) with finite codomains, and a candidate filler \(\phi:\mathrm{im}(\pi_{0})\to O_{1}\), the factorization square is

\[
\Xi_{\mathrm{fact}}(\phi)\;:\;
\begin{array}{ccc}
S & \xrightarrow{\pi_{0}} & \mathrm{im}(\pi_{0})\\
{\scriptstyle\mathrm{id}_{S}}\downarrow & \Downarrow\,\Xi_{\mathrm{fact}}(\phi) & \downarrow{\scriptstyle\phi}\\
S & \xrightarrow{\pi_{1}} & O_{1}
\end{array}
\]

with \(\mathrm{type}_{\Xi}=\texttt{factorization}\) and \(\mathsf{Decor}_{\Xi}\) recording the factorization candidate \((\pi_{0},\pi_{1},\phi)\).

An exact filler \(\phi\) exists iff \(\pi_{1}=\phi\circ\pi_{0}\) on \(S\). Hence

\[
\pi_{1}\not\factor\pi_{0}\quad\Longleftrightarrow\quad\mathrm{Fill}(\pi_{0},\pi_{1})=\varnothing\quad\Longleftrightarrow\quad\exists\,s,s'\in S:\;\pi_{0}(s)=\pi_{0}(s')\;\land\;\pi_{1}(s)\neq\pi_{1}(s'),
\]

with \(\mathrm{Fill}(\pi_{0},\pi_{1})=\{\phi:\mathrm{im}(\pi_{0})\to O_{1}\;:\;\pi_{1}=\phi\circ\pi_{0}\}\). The fiber-mismatch witness \(\{(s,s')\in S\times S:\pi_{0}(s)=\pi_{0}(s')\land\pi_{1}(s)\neq\pi_{1}(s')\}\) is a finite typed defect on \(S\times S\).

The instrument-indexed verdict is \(\Gamma;\mathsf H_{\mathrm{DC}};I_{\mathrm{DC}}\vdash\textsc{StrictExtensionAsNoFiller}:\texttt{accepted}\). The Cantor strict-bridge realization of Subsection 13.3 supplies a concrete instance: \(\pi_{1}\not\factor\pi_{0}\) on \(\mathcal S_{\mathrm{aud}}\), so \(\mathrm{Fill}(\pi_{0},\pi_{1})=\varnothing\) on the shell, and \(\Xi_{\mathrm{fact}}\) has no filler on the shell. The recovery is representational and does not assert that every strict-extension claim reduces to a decorated square.

\subsection{High-Structure Nonclaims}

This subsection records the long-form high-structure nonclaims of Subsection 14.5.

\subsubsection{No Double-Category Completeness}

The finite decorated double-category fragment \(\mathsf H_{\mathrm{DC}}\) does not replace the core calculus \(\mathbf{BirdInt}^{\mathrm{aud}}_{\mathrm{fin}}\). Theorems 30, 31, and 32 record three specific recoveries --- descent data as descent squares, completion noncommutation as completion squares, strict-extension nonfactorization as factorization squares --- and not a general claim that every theorem-grade fact about the calculus is reducible to a decorated-square fact.

In particular, the appendix does not claim:

\begin{itemize}
\tightlist
\item
  that \(\mathbf{BirdInt}^{\mathrm{aud}}_{\mathrm{fin}}\) embeds into \(\mathsf H_{\mathrm{DC}}\) as a sub-double-category;
\item
  that every directed-cell judgment of Subsection 4.3 is recovered by a decorated square;
\item
  that the priority order of the classifier of Subsection 11.2.1 is recovered by the per-square exactness conditions.
\end{itemize}

\subsubsection{No Status Erasure}

The decorated-square fragment does not collapse the status families of Subsection 5.1 into a single square-status family. The role-status, cell-status, pair-status, promotion-status, claim-status, gate-status, and square-status families remain separate judgment-indexed families. A decorated square carries its own square-status verdict only as part of a full decorated/audited record

\[
(W,U,L,V,\Theta,A,\delta,\sigma,\mathcal N),
\]

including witness data, update data, language, visibility, threshold, audit record, defect record, square status, and nonclaim data. A bare commutative diagram without this record is not an admissible \(\mathbf{BirdInt}^{\mathrm{aud}}_{\mathrm{fin}}\) claim.

\subsubsection{No General Obstruction Propagation Theorem}

The pasting law of Subsubsection 14.1.3 records a syntactic aggregation rule for defects under finite strict pasting. The appendix does not claim a general theorem of the form ``a defect at one square in a finite pasting diagram propagates to a defect at every square downstream of it'' or ``the absence of a defect at one square suffices for the absence of a defect at every square in the diagram.'' Both directions of general defect propagation are future work, in keeping with Subsubsection 14.5.3.

The three nonclaims above fix the local boundary of the secondary high-structure semantics within \(\mathbf{BirdInt}^{\mathrm{aud}}_{\mathrm{fin}}\).
 \section{Lean 4 Mechanization Appendix}

This appendix records the repository-level Lean 4 artifact that supports Proposition 35. It supplements Section 15 with the module map, validator command, toolchain, manifest, trust base, and claim-strength boundary for the completed validation track.

\subsection{Reproducibility}

From the repository root, the Lean project builds with \texttt{cd\ lean/full\ \&\&\ lake\ build}.

The strict validation command is \texttt{./scripts/validate\_lean.sh\ -\/-tier\ full\ -\/-require-full}.

The pinned toolchain is \texttt{leanprover/lean4:v4.28.0}.

The validator can also emit a JSON report with \texttt{./scripts/validate\_lean.sh\ -\/-tier\ full\ -\/-require-full\ -\/-json\ lean-full-validation.json}.

The JSON report is a reproducibility artifact, not an additional paper claim.

\subsection{Files Audited by the Validator}

The validator uses the following repository records:

\begin{itemize}
\tightlist
\item
  \texttt{lean/full/lean-toolchain}, the exact Lean toolchain;
\item
  \texttt{lean/full/lakefile.toml}, the Lake project configuration;
\item
  \texttt{lean/full/SixBirdsIII.lean}, the root Lean module;
\item
  \texttt{lean/manifest.toml}, the required paper-to-Lean declaration manifest;
\item
  \texttt{lean/paper\_inventory.toml}, the broader paper inventory from which the manifest is curated;
\item
  \texttt{lean/trust\_base.txt}, the allowed non-project axiom trust base;
\item
  \texttt{lean/full/SixBirdsIII/SemanticTests.lean}, the semantic probe module.
\end{itemize}

The allowed non-project axioms are \texttt{Classical.choice}, \texttt{propext}, \texttt{Quot.sound}, \texttt{choice}, and \texttt{funext}.

All theorem-backed manifest entries are audited against this list.

\subsection{Module Map}

The root module \texttt{SixBirdsIII} imports the following module stack:

\begin{longtable}[]{@{}ll@{}}
\toprule
\begin{minipage}[b]{0.47\columnwidth}\raggedright
Lean module\strut
\end{minipage} & \begin{minipage}[b]{0.47\columnwidth}\raggedright
Paper role\strut
\end{minipage}\tabularnewline
\midrule
\endhead
\begin{minipage}[t]{0.47\columnwidth}\raggedright
\texttt{Basic.lean}\strut
\end{minipage} & \begin{minipage}[t]{0.47\columnwidth}\raggedright
primitive labels, levels, hosts, and finite tags\strut
\end{minipage}\tabularnewline
\begin{minipage}[t]{0.47\columnwidth}\raggedright
\texttt{Status.lean}\strut
\end{minipage} & \begin{minipage}[t]{0.47\columnwidth}\raggedright
seven status families and classifier-facing tags\strut
\end{minipage}\tabularnewline
\begin{minipage}[t]{0.47\columnwidth}\raggedright
\texttt{Domain.lean}\strut
\end{minipage} & \begin{minipage}[t]{0.47\columnwidth}\raggedright
finite audited domain and core package data\strut
\end{minipage}\tabularnewline
\begin{minipage}[t]{0.47\columnwidth}\raggedright
\texttt{Records.lean}\strut
\end{minipage} & \begin{minipage}[t]{0.47\columnwidth}\raggedright
profiles, instruments, directed cells, pair observables, bridges, defects, and claim records\strut
\end{minipage}\tabularnewline
\begin{minipage}[t]{0.47\columnwidth}\raggedright
\texttt{Mech0.lean}\strut
\end{minipage} & \begin{minipage}[t]{0.47\columnwidth}\raggedright
small core interaction package and no-total/descent/saturation/nonfactorization support\strut
\end{minipage}\tabularnewline
\begin{minipage}[t]{0.47\columnwidth}\raggedright
\texttt{CoreInteraction.lean}\strut
\end{minipage} & \begin{minipage}[t]{0.47\columnwidth}\raggedright
Section 6 core finite-interaction theorems\strut
\end{minipage}\tabularnewline
\begin{minipage}[t]{0.47\columnwidth}\raggedright
\texttt{FiniteMaps.lean}\strut
\end{minipage} & \begin{minipage}[t]{0.47\columnwidth}\raggedright
finite map, quotient, image, and fiber support\strut
\end{minipage}\tabularnewline
\begin{minipage}[t]{0.47\columnwidth}\raggedright
\texttt{Completion.lean}\strut
\end{minipage} & \begin{minipage}[t]{0.47\columnwidth}\raggedright
completion and idempotent-operator support\strut
\end{minipage}\tabularnewline
\begin{minipage}[t]{0.47\columnwidth}\raggedright
\texttt{Strictness.lean}\strut
\end{minipage} & \begin{minipage}[t]{0.47\columnwidth}\raggedright
strict extension and factorization support\strut
\end{minipage}\tabularnewline
\begin{minipage}[t]{0.47\columnwidth}\raggedright
\texttt{Defects.lean}\strut
\end{minipage} & \begin{minipage}[t]{0.47\columnwidth}\raggedright
defect equivalence records\strut
\end{minipage}\tabularnewline
\begin{minipage}[t]{0.47\columnwidth}\raggedright
\texttt{CompletionExamples.lean}\strut
\end{minipage} & \begin{minipage}[t]{0.47\columnwidth}\raggedright
noncommuting completion and pasting examples\strut
\end{minipage}\tabularnewline
\begin{minipage}[t]{0.47\columnwidth}\raggedright
\texttt{Definability.lean}\strut
\end{minipage} & \begin{minipage}[t]{0.47\columnwidth}\raggedright
fixed-interface definability bound\strut
\end{minipage}\tabularnewline
\begin{minipage}[t]{0.47\columnwidth}\raggedright
\texttt{FiniteProbability.lean}\strut
\end{minipage} & \begin{minipage}[t]{0.47\columnwidth}\raggedright
finite stochastic and closure-deficit support\strut
\end{minipage}\tabularnewline
\begin{minipage}[t]{0.47\columnwidth}\raggedright
\texttt{GraphCohomology.lean}\strut
\end{minipage} & \begin{minipage}[t]{0.47\columnwidth}\raggedright
Section 8 graph/cohomology drive theorems\strut
\end{minipage}\tabularnewline
\begin{minipage}[t]{0.47\columnwidth}\raggedright
\texttt{Promotion.lean}\strut
\end{minipage} & \begin{minipage}[t]{0.47\columnwidth}\raggedright
promotion gate, strictness, and non-strict factorization theorems\strut
\end{minipage}\tabularnewline
\begin{minipage}[t]{0.47\columnwidth}\raggedright
\texttt{InstrumentClaims.lean}\strut
\end{minipage} & \begin{minipage}[t]{0.47\columnwidth}\raggedright
no-overreading and instrument-indexed claim support\strut
\end{minipage}\tabularnewline
\begin{minipage}[t]{0.47\columnwidth}\raggedright
\texttt{RotatingAudit.lean}\strut
\end{minipage} & \begin{minipage}[t]{0.47\columnwidth}\raggedright
same-level self-audit and finite rotating-audit theorems\strut
\end{minipage}\tabularnewline
\begin{minipage}[t]{0.47\columnwidth}\raggedright
\texttt{AbstractInterpretation.lean}\strut
\end{minipage} & \begin{minipage}[t]{0.47\columnwidth}\raggedright
Section 12 abstract-interpretation theorems and model-family record\strut
\end{minipage}\tabularnewline
\begin{minipage}[t]{0.47\columnwidth}\raggedright
\texttt{HighStructure.lean}\strut
\end{minipage} & \begin{minipage}[t]{0.47\columnwidth}\raggedright
Section 14 decorated-square recovery theorems\strut
\end{minipage}\tabularnewline
\begin{minipage}[t]{0.47\columnwidth}\raggedright
\texttt{ModelRealizations.lean}\strut
\end{minipage} & \begin{minipage}[t]{0.47\columnwidth}\raggedright
PICA and Cantor model-realization records\strut
\end{minipage}\tabularnewline
\begin{minipage}[t]{0.47\columnwidth}\raggedright
\texttt{MechanizationFeasibility.lean}\strut
\end{minipage} & \begin{minipage}[t]{0.47\columnwidth}\raggedright
Proposition 35 and manifest-level feasibility record\strut
\end{minipage}\tabularnewline
\begin{minipage}[t]{0.47\columnwidth}\raggedright
\texttt{SemanticTests.lean}\strut
\end{minipage} & \begin{minipage}[t]{0.47\columnwidth}\raggedright
validator semantic probes\strut
\end{minipage}\tabularnewline
\bottomrule
\end{longtable}

The module map is an audit map: it records where the finite support for each manifest entry lives. It is not a claim that the Lean artifact certifies statements outside the manifest.

\subsection{What Strict Validation Checks}

Strict validation checks that:

\begin{enumerate}
\def\labelenumi{\arabic{enumi}.}
\tightlist
\item
  the Lean project builds under the pinned toolchain;
\item
  every \texttt{required\ =\ true} manifest entry is present;
\item
  every required definition is marked \texttt{status\ =\ "full"};
\item
  every required numbered claim is marked \texttt{status\ =\ "theorem"};
\item
  every listed Lean declaration exists;
\item
  every listed declaration has the expected Lean declaration kind;
\item
  every theorem-backed declaration has only allowed non-project axiom dependencies;
\item
  the Lean sources contain no blocked placeholder markers;
\item
  the semantic probe module exposes the required semantic checks.
\end{enumerate}

In the completed artifact, strict validation reports 21 full schema definitions and 35 theorem-backed required claims, with all probes, declaration-kind checks, theorem axiom audits, placeholder scans, and semantic audits passing.

\subsection{Claim-Strength Boundary}

The paper distinguishes theorem-grade, model-grade, and audit-grade claims. The Lean validator uses \texttt{status\ =\ "theorem"} for every required numbered claim because the validation question is whether there is a Lean theorem declaration backing the manifest entry. That does not erase the paper-facing claim-strength distinction:

\begin{itemize}
\tightlist
\item
  theorem-grade paper claims remain theorem-grade finite claims;
\item
  model-grade paper claims remain scoped model-realization claims;
\item
  Proposition 35 remains an audit-grade validation claim.
\end{itemize}

The strict validator certifies the existence, declaration kind, and audited dependencies of the Lean declarations listed in \texttt{lean/manifest.toml}. It does not certify unlisted claims, empirical implementation correctness, universal exact-six structure, a universal six-symbol algebra, or any host extension outside the finite audited manifest scope.
 \section{Dependency Graph}

This appendix records the long-form dependency graph of \(\mathbf{BirdInt}^{\mathrm{aud}}_{\mathrm{fin}}\). Subsection G.1 records the definition nodes. Subsection G.2 records the theorem build order in topological form. Subsection G.3 records the critical proof paths.

The dependency graph is a finite typed record of the producer/consumer relations among definitions, lemmas, theorems, countermodels, model-realization records, and high-structure recovery theorems in the paper. It is a bookkeeping appendix, not an additional theorem.

\subsection{Definition Nodes}

The definition nodes of \(\mathbf{BirdInt}^{\mathrm{aud}}_{\mathrm{fin}}\) are the finite typed schemas that the calculus introduces and uses throughout the paper. Each node is a finite typed record schema; the dependency graph records which schemas consume the records of which other schemas as part of their admissibility predicate.

The eleven definition-node families are:

\begin{enumerate}
\def\labelenumi{\arabic{enumi}.}
\tightlist
\item
  \textbf{Host} --- finite host schema (Subsection 2.1).
\item
  \textbf{Content} --- finite content universe with visibility and suppression records (Subsections 3.2 and 3.5).
\item
  \textbf{Theory packages} --- \(\mathcal T=(Z,f,\Sigma_{f},E,\mathcal A)\) (Subsection 3.3).
\item
  \textbf{Instruments} --- instrument records \(I\) (Subsections 3.4 and 11.1.1).
\item
  \textbf{Profiles} --- typed profile \(\lambda\) (Subsection 3.6 and Appendix A.4).
\item
  \textbf{Witness/update maps} --- \(\mathsf{Wit}:\mathbb P\to\mathrm{Type}\) and \(\mathsf{Upd}:\mathbb P\to\mathrm{Type}\) (Subsection 3.7 and Tables A.2, A.3).
\item
  \textbf{Statuses} --- seven typed status families (Subsection 5.1 and Table A.5).
\item
  \textbf{Defects} --- typed defect families per primitive and per promotion (Subsections 5.2--5.5 and Table A.6).
\item
  \textbf{Gates} --- promotion gates and bridge-admissibility gates (Subsection 10.2).
\item
  \textbf{Audits} --- host audit, claim audit, bridge audit, rotating-audit (Subsection 4.8 and Section 11).
\item
  \textbf{Model-realization maps} --- realization maps \(\mathrm{real}_{\mathcal M\to\mathcal S}\) (Subsection 13.1.1).
\end{enumerate}

The producer/consumer relations among these eleven node families are recorded as a directed dependency graph: hosts support content and packages; packages, instruments, profiles, witness/update maps, visibility records, thresholds, audits, and defects feed the judgment schemas; classifiers consume judgment and defect data to assign statuses; gates feed promotion-status verdicts; audits feed visibility, source, and instrument-bound claim verdicts; model-realization maps consume the host, package, judgment, status, and audit nodes to produce model-grade \(\texttt{model\_realization}\) verdicts on declared fragments.

The dependency graph is a finite typed record on the eleven node families; it is admissible as a bookkeeping object when every node is finite and every dependency edge is recorded in the topological build order of Subsection G.2.

\subsection{Theorem Build Order}

The theorem build order is the topological order in which the paper's theorems and supporting records are introduced and proved. It is the order under which every consumer of a record is preceded by its producer.

The nine-stage build order is:

\begin{enumerate}
\def\labelenumi{\arabic{enumi}.}
\tightlist
\item
  \textbf{Formal host and records.} Sections 2--5: finite host mathematics, primitive labels, claim strengths, status discipline, theory packages, instruments, source-of-truth records, audit records, visibility/suppression records, directed-cell records, pair-observable records, promotion bridges, profiles, status families, and defect schemas.
\item
  \textbf{Finite map/completion lemmas.} Sections 6 and 7: finite-cell admissibility predicates, effective-image representative lemmas, idempotent-iteration lemmas, and finite fiber-mismatch arguments.
\item
  \textbf{Core finite theorem spine.} Section 6 (Theorems 1--5) and Section 7 (Theorems 6--14): no-total-algebra, status-family separation, well-formedness decidability, covered-cell status uniqueness, descent through quotient, descent defect equivalence, Markov closure deficit, noncommuting completions, completion-pasting defect, idempotent saturation, strict-extension nonfactorization, factorization defect equivalence, fixed-interface definability bound.
\item
  \textbf{No-go and countermodels.} Section 8 (Theorems 15--18) and Section 9 plus Appendix C (CM\_1--CM\_12): forest no-drive, exact-form null-drive, nonzero-affinity equivalence, gating-affinity suppression, plus the twelve finite countermodels.
\item
  \textbf{Promotion.} Section 10 (Theorems 19--22): promotion gate soundness, strict promotion implies nonfactorization, non-strict promotion implies factorization, Toy22 strict-bridge example.
\item
  \textbf{Instruments and self-reference.} Section 11 (Theorems 23--25): no-overreading suppression, same-level self-audit failure, finite rotating-audit theorem.
\item
  \textbf{Model families.} Section 12 (Theorems 26--28 and Model Theorem 29) and Section 13 (Model Theorems 33, 34, plus the graph/cohomology realization): closure as packaging, descent as sound transformer, representability as fixedness, AI model-family theorem, PICA realization, Cantor strict-bridge realization, graph/cohomology realization. The realization statements in this stage are model-grade, not theorem-grade.
\item
  \textbf{High-structure semantics.} Section 14 (Theorems 30--32): descent-square recovery, completion-pasting recovery, strict extension as no-filler.
\item
  \textbf{Lean validation appendix.} Section 15 (Proposition 35) and Appendix F: Lean module map, \texttt{lean/manifest.toml}, validator command, toolchain, trust base, and manifest-scoped mechanization nonclaims.
\end{enumerate}

Each stage consumes only the records and theorems of earlier stages. The order is a topological sort on the dependency graph of Subsection G.1, with the host and definition nodes at the bottom and the model-realization and Lean-validation nodes at the top. The build order does not claim closure under any larger dependency closure; it records only the per-stage producer/consumer ordering used by the paper.

\subsection{Critical Proof Paths}

This subsection records the five critical proof paths in the dependency graph. Each path is the chain of definitions and theorems that supports a specific theorem-grade or model-grade conclusion.

\subsubsection{Strict Promotion Path}

The strict promotion path is the chain of records and theorems supporting the strict-promotion verdict \(\mathrm{Promote}(B):\texttt{strict}\):

theory packages \(\to\) promotion bridges \(\to\) promotion gates (Subsection 10.2) \(\to\) promotion classifier (Subsection 5.6.4) \(\to\) Theorem 19 (gate soundness) \(\to\) Theorem 20 (strict \(\Rightarrow\) nonfactorization) \(\to\) Theorem 12 (strict-extension nonfactorization).

The path is consumed by the Cantor strict-bridge realization (Model Theorem 34) and by the strict-extension recovery in the high-structure fragment (Theorem 32).

\subsubsection{Package versus Closure Path}

The package-versus-closure path is the chain supporting the \(P_{5}\)-packaging vs macro-closure separation:

theory packages \(\to\) quotient/package map \(q\) and update \(F\) \(\to\) effective-image descent setup \(\to\) Theorem 6 (descent through quotient) \(\to\) Theorem 7 (descent defect equivalence) \(\to\) CM\_1 (active package without macro closure).

The path is consumed by package and closure nonclaims throughout the paper, including the AI host realization (Theorem 26 closure as packaging, paired with Theorem 27 descent as sound transformer) and the Cantor strict-bridge realization (Model Theorem 34), where macro-dynamics closure is \emph{not} asserted without the required descent data.

\subsubsection{Completion versus Route Coherence Path}

The completion-versus-route path is the chain supporting the completion-noncommutation vs \(P_{3}\)-pasting recovery:

finite completions \(E_{1},E_{2}\) \(\to\) Theorem 9 (noncommuting completions) \(\to\) CM\_8 (idempotent completions do not commute) \(\to\) Theorem 10 (completion-pasting defect) \(\to\) Theorem 31 (completion-pasting recovery in \(\mathsf H_{\mathrm{DC}}\)).

The path is consumed by Theorem 31 of Section 14 and the corresponding decorated-square recovery in Appendix E.3.

\subsubsection{Self-Reference Honesty Path}

The self-reference honesty path is the chain supporting the no-overreading and same-level self-audit failure theorems:

instrument records \(\to\) visibility/suppression discipline (Subsection 11.3) \(\to\) classifier priority (Subsection 11.2) \(\to\) Theorem 23 (no-overreading suppression) \(\to\) Theorem 24 (same-level self-audit failure) \(\to\) rotating-audit chain (Subsection 11.7) \(\to\) Theorem 25 (finite rotating-audit) and its non-finality clause.

The path is consumed by every model-realization sheet of Section 13 (whose claim-grade verdicts are scoped under instruments) and by the manifest-scope mechanization boundary of Section 15 and Appendix F.

\subsubsection{Drive Separation Path}

The drive separation path is the chain supporting the cohomological-drive vs holonomy separation:

finite graph \(G\) and cycle space \(\mathsf{Cycles}(G)\) \(\to\) first cohomology \(H^{1}(G;k)\) \(\to\) Theorem 15 (forest no-drive), Theorem 16 (exact-form null-drive), and Theorem 17 (nonzero-affinity equivalence) \(\to\) CM\_2 (active route/holonomy without drive) \(\to\) Theorem 18 (gating-affinity suppression) \(\to\) CM\_12 (gating destroys cycle-supported drive).

The path is consumed by the graph/cohomology realization (Subsection 13.4) and by the cohomology nonclaim register \(\mathcal N_{\mathrm{coh}}\) of Subsubsection 13.4.4.

The five critical paths are the principal supporting chains for the paper's theorem-grade and model-grade conclusions. They do not exhaust the dependency graph; they record the load-bearing paths that the paper relies on most heavily.
 \section{Deferred and Dropped Claims}

This appendix records the deferred and dropped claims of \(\mathbf{BirdInt}^{\mathrm{aud}}_{\mathrm{fin}}\). Subsection H.1 records the deferred claims, organized by future-paper group. Subsection H.2 records the dropped claims, organized as outright exclusions from the calculus's positive scope.

The appendix is a finite typed register that consolidates the future-work items of Subsection 16.3 and the negative-scope items of Subsection 16.4 (and the per-section nonclaim subsections) in one place. It does not introduce new theorems; it records what is \emph{not} in this paper and \emph{not} a target of this paper.

\subsection{Deferred Claims}

This subsection records the deferred claims by future-paper group, mirroring the six classes of future work in Subsection 16.3.

\subsubsection{Classifier and Dependency Completeness}

Deferred items:

\begin{itemize}
\tightlist
\item
  a complete status-classifier theorem characterizing the priority order over \(\mathrm{ClaimStatus}\);
\item
  a defect-to-status soundness theorem characterizing the routing from the defect families of Subsections 5.2--5.5 through the classifier rules of Subsection 5.6;
\item
  a dependency-lattice theorem organizing the per-component admissibility predicates of Sections 4 and 5 into a lattice with a finite reduction theorem.
\end{itemize}

\subsubsection{Promotion and Stacking Theory}

Deferred items:

\begin{itemize}
\tightlist
\item
  a complete promotion classifier extending the eight-gate machinery of Subsection 10.2 with stated soundness and completeness conditions over \(\mathrm{PromotionStatus}\);
\item
  a promotion well-formedness decidability theorem certifying that admissibility of a promotion bridge is decidable on finite records;
\item
  a general stacking composition theorem extending Subsection 10.7 to admissible promotion-bridge composition across host stacks;
\item
  a local-to-global promotion theorem identifying conditions under which a finite family of locally admissible promotion bridges supports a global promotion verdict.
\end{itemize}

\subsubsection{Verified Model Realizations}

Deferred items:

\begin{itemize}
\tightlist
\item
  a machine-checked PICA realization extending Model Theorem 33 with a Lean 4 / Mathlib development of the actor-update map, informant-witness map, status recovery, and pair-observable realness predicate;
\item
  a machine-checked Cantor strict-bridge realization extending Model Theorem 34 with a verified shell-bound development of the strictness witness and gate-result table;
\item
  a machine-checked abstract-interpretation realization extending Model Theorem 29 with a verified Galois-pair, closure-operator, and soundness-inequality development;
\item
  an implementation-correctness layer for any concrete PICA implementation, audited under a separate proof-assistant-bound instrument.
\end{itemize}

\subsubsection{Cantor Shell Expansion}

Deferred items:

\begin{itemize}
\tightlist
\item
  a full Cantor strictness proof extending the audited-shell strict-bridge realization beyond \(\mathcal S_{\mathrm{aud}}\) via an admissible expansion bridge;
\item
  a pressure-gap-to-KL bridge recording an admissible bridge from the pressure-gap consequence defect of Subsubsection 13.3.6 to the Kullback--Leibler closure-deficit framework;
\item
  a stratumwise root-separation theorem characterizing the persistent strata family \(\mathcal F\) via a separation property stronger than the conditional pressure disintegration of Subsubsection 13.3.6.
\end{itemize}

\subsubsection{Drive and High-Structure Semantics}

Deferred items:

\begin{itemize}
\tightlist
\item
  a general drive classification extending the cohomological drive certificate \([a]\neq 0\) of Subsubsection 13.4.3 to a finite family of drive certificates over multiple drive notions;
\item
  a holonomy-to-drive bridge \(B_{3\to 6}^{\mathrm{drive}}\) from a \(P_{3}\) route/holonomy witness to a \(\mathrm{P6}_{\mathrm{drive}}\) certificate;
\item
  gating beyond forests, characterizing \(P_{2}\)-gating effects on drive on graphs with non-trivial cycle structure under additional gating policies;
\item
  weak double-category coherence extending the strict pasting law of Subsubsection 14.1.3 to weak coherence relations on finite decorated squares;
\item
  a general obstruction propagation theorem characterizing the propagation of square defects under finite strict pasting, in the sense of Subsubsection 14.5.3.
\end{itemize}

\subsubsection{Self-Reference and Instrument Transfer}

Deferred items:

\begin{itemize}
\tightlist
\item
  a full instrument-transfer classifier fixing the per-claim-type mapping from source-side acceptance and bridge defect to target-side \(\chi\in\{\,\texttt{accepted},\;\texttt{provisional},\;\texttt{blocked}\,\}\) under \(J\);
\item
  an infinite rotating-audit semantics extending the finite rotating-audit chain of Subsection 11.7 to an inverse limit or coalgebraic semantics admissible under a stronger meta-instrument;
\item
  a complete self-reference classification extending Theorem 24 of Subsection 11.5 to a comprehensive classification of same-level, level-shifted, and rotating self-reference patterns.
\end{itemize}

The six groups above record the principal future-work items of the paper. Each item is a scoped pointer to an admissible extension target, in keeping with the future-work boundary of Subsection 16.3.

\subsection{Dropped Claims}

This subsection records claims that are dropped outright from the calculus's positive scope. Each item below is \emph{not} deferred to future work but explicitly excluded as a positive claim of \(\mathbf{BirdInt}^{\mathrm{aud}}_{\mathrm{fin}}\).

The seven dropped claims are:

\begin{enumerate}
\def\labelenumi{\arabic{enumi}.}
\tightlist
\item
  \textbf{Universal six-symbol algebra.} No total operation \(\ast:\mathbb P^{2}\to\mathbb P\) recovers cell status across profiles. Theorem 1 of Section 6 records the formal counterexample, and \(\mathrm{NC}\text{-}1\) of Subsection 16.4 records the corresponding nonclaim.
\item
  \textbf{All mathematics is six birds.} The claim that every mathematical structure decomposes into the six primitives is dropped; \(\mathbf{BirdInt}^{\mathrm{aud}}_{\mathrm{fin}}\) is a finite typed audited interaction calculus, not a foundation for all mathematics. \(\mathrm{NC}\text{-}2\) of Subsection 16.4 records the corresponding nonclaim.
\item
  \textbf{Unrestricted exact-six.} The unrestricted decomposition claim, that every typed interaction is exactly representable as one of the thirty-six directed cells, is dropped. The exact-six structure is a finite typed schema with admissibility conditions, not an unrestricted decomposition theorem. The dropped claim is recorded in Subsubsection 1.3.2 and the corresponding nonclaim in \(\mathrm{NC}\text{-}1\) and the algebra/universality grouping of Subsubsection 16.2.1.
\item
  \textbf{PICA as canonical algebra.} The claim that PICA defines the canonical \(6\times 6\) algebra of the calculus is dropped. PICA is a finite stochastic model realization, scoped to the cell/pair/provenance fragment, and the recovered \(6\times 6\) multiplicity count is a profile fact, not a universal law. PICA nonclaims 1 and 2 of Subsubsection 13.2.10 and \(\mathrm{NC}\text{-}20\), \(\mathrm{NC}\text{-}21\) of Subsection 16.4 record the corresponding nonclaims.
\item
  \textbf{Cantor pressure gap as unbridged strictness witness.} The claim that the pressure-gap defect \(\Delta_{T_{0}\to T_{1}}\) is itself the strictness witness, rather than a consequence of the strictness witness, is dropped. The strictness witness is the factorization defect \(\Delta_{\mathrm{fact}}^{\mathrm{Cantor}}\); the pressure-gap defect is recorded as a consequence record. Cantor nonclaim 7 of Subsubsection 13.3.9 and \(\mathrm{NC}\text{-}27\) of Subsection 16.4 record the corresponding nonclaim.
\item
  \textbf{Provenance as proof.} The claim that provenance records constitute proof of any underlying mathematical fact is dropped. Provenance records are finite typed records of producer, source, audit, and visibility data; they support pair-realness and source-of-truth predicates but are not theorem-grade proofs of mathematical claims. The audit nonclaim \(\mathrm{NC}\text{-}15\) of Subsection 16.4 (no audit-is-truth) records the corresponding boundary.
\item
  \textbf{High-structure semantics as replacement calculus.} The claim that the finite decorated double-category fragment of Section 14 replaces or supersedes the core calculus is dropped. The decorated-square fragment is a secondary organizational layer; it records three specific recoveries (Theorems 30--32) and not a general claim that every theorem-grade fact reduces to a decorated-square fact. \(\mathrm{NC}\text{-}34\), \(\mathrm{NC}\text{-}35\) of Subsection 16.4 record the corresponding nonclaims.
\end{enumerate}

The seven dropped claims above are explicitly outside the positive scope of the paper. Each is recorded as a negative scope item in Section 16, and each is supported by a theorem, a countermodel, or a per-realization nonclaim register elsewhere in the paper. The appendix records them in one place to make the paper's negative scope register explicit.

The seven dropped claims and the six groups of deferred claims in Subsection H.1 together record the boundary of \(\mathbf{BirdInt}^{\mathrm{aud}}_{\mathrm{fin}}\): the calculus admits the theorem-grade, model-grade, and audit-grade results recorded in this paper; defers the items of Subsection H.1 to future work; and excludes the items of Subsection H.2 from any positive claim.
 
\clearpage
\nocite{*}
\bibliographystyle{plainnat}
\bibliography{references}

\end{document}
